% 本文件由 LaTeX 排版 Agent 根据有效 translation.json 自动生成。
% author=Origen; work_id=1016; title=玛窦福音注释

\chapter{玛窦福音注释}

% structure: p0001 source=h1 target=omitted reason=page-title-already-rendered-by-parent

\Emph{第一卷} 残篇\newline{}\Emph{第二卷} 残篇\newline{}\Emph{第十卷}\newline{}\Emph{第十一卷}\newline{}\Emph{第十二卷}\newline{}\Emph{第十三卷}\newline{}\Emph{第十四卷}\par

\SourceRecord{Translated by John Patrick.；Ante-Nicene Fathers；Vol. 9.；Edited by Allan Menzies.；Buffalo, NY: Christian Literature Publishing Co.,；1896.；<http://www.newadvent.org/fathers/1016.htm>.}

% structure: p0001 source=h1 target=section reason=page-title

\section{\WorkTitle{玛窦福音注释（卷一）}}

关于天下天主教会中无可争议的四部福音，我由传统得知：\BibleBook{玛窦}{福音}——他原是一个税吏，后来成为耶稣基督的宗徒——是首先写成的；他用希伯来语撰写，并为从犹太教皈依的人而发表。第二部写成的，是\BibleBook{马尔谷}{福音}，他按照伯多禄的指示撰写；伯多禄在他的公函中称他为儿子，说：\BibleQuote{在巴比伦与你们一同蒙选的教会向你们问安，我儿子马尔谷也向你们问安。} \BibleRef{伯前 5:13} 第三部，是\BibleBook{路加}{福音}，是保禄所称赞的，他为从外邦人中皈依的人而撰写。最后一部，是\BibleBook{若望}{福音}。\par

\SourceRecord{Translated by John Patrick.；Ante-Nicene Fathers；Vol. 9.；Edited by Allan Menzies.；Buffalo, NY: Christian Literature Publishing Co.,；1896.；<http://www.newadvent.org/fathers/101601.htm>.}

% structure: p0001 source=h1 target=section reason=page-title

\section{玛窦福音释义（卷二）}

% structure: p0002 source=h2 target=subsection reason=relative-heading-rank; base=h2

\subsection{圣经的统一与和谐}

\BibleQuote{\Emph{缔造和平的人是有福的}……}\BibleRef{玛 5:9}对那在任一意义上成为缔造和平的人来说，在圣经中没有弯曲或乖僻的，因为对理解的人而言，它们全是平直的。\BibleRef{箴 8:8}因为对这样的人没有弯曲或乖僻，所以他看见所有圣经中的丰盛和平，即使在那些看似冲突和彼此矛盾的部分。同样，他成为第三种缔造和平的人，因为他证明那在别人看来是圣经中的冲突其实并非冲突，并显明它们的和谐与和平，无论是旧约与新约之间，还是法律与先知之间，或是福音与宗徒著作之间，或是宗徒著作彼此之间。因为，按照训道者，全部圣经都是\BibleQuote{智者的话就像刺棒，又像钉牢的钉子，都是一个牧者所赐。}\BibleRef{训 12:11}其中没有多余的东西。但圣言是理性万物的唯一牧者，这些事物对于没有耳朵可听的人可能显得不和谐，但事实上是完全和谐的。因为正如竖琴或里拉琴的不同琴弦，每一根发出某种自己的声音，似乎与其他琴弦的声音不同，一个不谙音乐、不懂音乐和谐原理的人会认为它们不和谐，因为声音的差异；同样，那些不熟练于在圣经中聆听天主和谐的人，认为旧约与新约不和谐，或先知与法律不和谐，或福音彼此不和谐，或宗徒与福音不和谐，或宗徒与自己不和谐，或与其他宗徒不和谐。但那个来到并受教于天主的音乐的人，是一个在言语和行为上智慧的人，因此，像另一位达味——按解释为\Quote{善于用手}——他将奏出天主的音乐之声，学会了适时弹拨琴弦，时而法律的琴弦，时而与它们和谐的福音琴弦，时而又先知的琴弦，并且当理性要求时，与先知和谐的宗徒琴弦，同样宗徒的琴弦与福音的琴弦和谐。因为他知道全部圣经是天主的一个完美而和谐的工具，它从不同的声音发出一个拯救的声音给那些愿意学习的人，它阻止并抑制一切恶灵的活动，正如达味的音乐平息了撒乌耳身上的恶灵，那恶灵也曾使他窒息。\BibleRef{撒上 16:14}\par

\SourceRecord{Translated by John Patrick.；Ante-Nicene Fathers；Vol. 9.；Edited by Allan Menzies.；Buffalo, NY: Christian Literature Publishing Co.,；1896.；<http://www.newadvent.org/fathers/101602.htm>.}

% structure: p0001 source=h1 target=section reason=page-title

\section{\WorkTitle{玛窦福音释义（第十卷）}}

% structure: p0002 source=h2 target=subsection reason=relative-heading-rank; base=h2

\subsection{一、莠子比喻：耶稣的家。}

\BibleQuote{\Emph{他离开群众，进了他的家，门徒来到他跟前说：请给我们讲解田里莠子的比喻。}}\BibleRef{玛 13:36} 当\Person{耶稣}与群众在一起时，他不在他的家中，因为群众在屋外，出于他对人的爱，他离开家，走向那些不能到他这里来的人。现在，他充分向群众讲了比喻，便打发他们走，回到自己的家，在那里，他的门徒——他们没有与那些被他打发走的人一起留下——来到他跟前。凡是\Person{耶稣}更真诚的听众，先跟随他，然后询问他的住处，被允许看见，来到，看见并与他同住，一整天，甚至有些人更长。依我见，在\BibleBook{若望}{福音}中以下的话指示了这样的事：\BibleQuote{次日，\Person{若翰}又站在那里，他的两个门徒。}\BibleRef{若 1:35} 为了解释那些被许可与\Person{耶稣}同去并看见他住处的人中，更卓越的那一位也成为宗徒的事实，补充了这句话：\BibleQuote{那两个听见\Person{若翰}说话而跟随他的人中，有一个是\Person{西满伯多禄}的兄弟\Person{安德肋}。}\BibleRef{若 1:40} 那么，如果我们不像被他打发走的群众，而是愿意听\Person{耶稣}，去他家，获得比群众更好的东西，让我们成为\Person{耶稣}的朋友，以便作为他的门徒，在他进入家时来到他面前，来到后询问比喻的解释，无论是莠子比喻还是其他。为了更准确地理解\Person{耶稣}的家所代表的意义，让人从福音书中收集一切关于\Person{耶稣}的家所说的话，以及他在其中所说或所行的事；所有收集在一起的段落将说服任何致力于阅读的人：福音的文字并非如有些人认为的那样绝对简单，而是由于天主的恩赐对简单的人变得简单；但对于那些有意志和力量更敏锐地听的人，其中有隐藏的智慧，与天主圣言相称。\par

% structure: p0004 source=h2 target=subsection reason=relative-heading-rank; base=h2

\subsection{二、比喻的讲解。}

\BibleQuote{\Emph{在这些事之后，他回答他们说：那撒好种子的，就是人子。}}\BibleRef{玛 13:37} 虽然我们在前几章中已经按我们的能力讨论了这些事，但我们仍要说与它们一致的话，即使有合理的理由另作解释。现在思考，除了我们已经讲述的之外，你们是否还可以将好种子理解为王国之子，因为凡在人的灵魂中所撒的好事物，都是天主王国的产物，并由那起初与天主同在的天主圣言所撒，\BibleRef{若 1:2} 因此关于任何事物的健康话语都是王国之子。但当人们睡觉——他们不按照\Person{耶稣}的命令行动，\BibleQuote{你们要醒寤祈祷，以免陷于诱惑}\BibleRef{玛 26:41}——魔鬼伺机撒下所谓的莠子，即邪恶的意见，散播在某些人所谓的自然概念中，即来自圣言的好种子之上和之中。据此，整个世界可被称为田地，而不仅是天主的教会，因为在整个世界，人子撒了好种子，但恶者撒了莠子——即邪恶的话语——这些出自邪恶，是恶者的子女。在事物的终结，被称为\Quote{世代终结}的时候，必然会有收割，好让天主为这工作所派遣的天使们收集灵魂上所滋生的邪恶意见，并将它们倾覆，投入那据说焚烧的火中，使它们被销毁。因此，圣言的天使和仆役将从基督的整个王国中收集一切使人灵魂绊倒的事物和制造邪恶的推理，他们将散播并投入燃烧的火炉中。那时，那些意识到自己因为睡觉而接受了恶者的种子的人，将哀号，并仿佛对自己发怒；因为这就是\BibleQuote{切齿}\BibleRef{玛 13:42}。因此，在圣咏中也说：\Quote{他们向我切齿。} 然后尤其是\BibleQuote{义人将发光}，不再像起初那样有差异，而是全部\BibleQuote{如同太阳，在他们父的王国里。}\BibleRef{玛 13:43} 然后，仿佛为了指出在一切与比喻的讲解相关的事中，也许特别是在这句话\BibleQuote{那时义人将像太阳一样在他们父的王国里发光}中，确有隐藏的意义，救主补充说：\BibleQuote{有耳听的，让他听。}\BibleRef{玛 13:43} 从而教导那些认为在讲解中比喻已被如此完美清晰地陈述以至于平民也能理解的人，即使与比喻的解释相关的事也需要说明。\par

% structure: p0006 source=h2 target=subsection reason=relative-heading-rank; base=h2

\subsection{三、义人的发光。其解释。}

但如同我们上文论及\Quote{那时，义人将要发光如同太阳}这句话时所说的，义人将不再像从前那样不同地发光，而是如同一个太阳，我们有必要阐述我们对此的看法。达尼尔知道明智者是世界的光，而义人的群众在荣耀上各有不同，他似乎说了这话：\BibleQuote{明智者将发光像穹苍的光辉，义人中的群众像星辰，直到永远。}\BibleRef{达 12:3}而在那段落中，\BibleQuote{太阳有一种荣耀，月亮有另一种荣耀，星辰又有另一种荣耀；因为星辰与星辰在荣耀上各有区别：死人复活也是这样，}\BibleRef{格前 15:41}宗徒说了与达尼尔同样的话，从后者的预言中取用了这思想。有人或许会问，为何有人谈论义人之间光的差别，而救主却说：\Quote{他们要像一个太阳那样发光。}那么，我认为，在那正在得救者享受真福的开端（因为不是这种人尚未净化），得救者之光所关联的差别是存在的；但是，正如我们已指出的，当他聚集了基督整个王国中一切使人跌倒的事，并将行不义的思念投入火炉，而更恶劣的元素完全烧尽；当这事发生时，那些接受了恶者之子的话的人便恢复了自我意识；那时，义人已成了一个太阳的光，要在他们父的国里发光。他们将为何人发光？为他们下面的人，这些人将享受他们的光，正如现在太阳为地上的人发光一样；因为，当然，他们不会为自己发光。但也许那句\BibleQuote{你们的光当在人前照耀}\BibleRef{玛 5:16}可以写在\Quote{心里的版上}，正如撒罗满所说，以三重方式；因此，即使是现在，耶稣门徒的光在人面前照耀；死后在复活前照耀；复活后\BibleQuote{直到众人达到成熟的人}\BibleRef{弗 4:13}，并且众人都成为一个太阳。那时，他们要像太阳一样在他们父的国里发光。\par

% structure: p0008 source=h2 target=subsection reason=relative-heading-rank; base=h2

\subsection{四、论藏于田地之宝藏的比喻。比喻与相似之区别。}

\BibleQuote{\Emph{天国又好像藏在地里的宝藏，一个人找到了，就把它藏起来。}}\BibleRef{玛 13:44}以前的比喻是对群众讲的；但这一个和接下来的两个，不是比喻，而是与天国有关的相似，他似乎在屋里对门徒讲的。关于这比喻和接下来的两个，让那\BibleQuote{注意阅读}\BibleRef{弟前 4:13}的人去查考它们究竟是不是比喻。在后者的情形中，圣经毫不迟疑地在每种情况下加上比喻的名称；但在此处，它并没有这样做；这是很自然的。因为如果他从前用比喻对群众讲话，并且\BibleQuote{他用比喻对群众讲了这一切，若不用比喻，就不给他们讲什么，}\BibleRef{玛 13:34}但是，当他进了屋，他不是对群众讲话，而是对来到他那里的门徒讲话；显然，在屋里所讲的不是比喻：因为，对外面的人，甚至对那些\BibleQuote{天国的奥秘没有赐给他们知道}\BibleRef{玛 13:11}的人，他用比喻讲话。有人会说，如果他们确实不是比喻，那是什么？我们要根据圣经的措辞说它们是相似（比较）吗？相似不同于比喻；因为马尔谷记载：\BibleQuote{我们可把天主的国比作什么？或用哪个比喻来形容它？}\BibleRef{谷 4:30}由此可知，相似与比喻之间有区别。相似似乎是属，而比喻是种。或许也如同相似，作为比喻的最高属，包含比喻为它的一个种，同样它也包含那一特定的相似形式，该形式与属同名。这与其他词语的情况相同，正如那些精于多名称呼的人所观察到的；他们说\Quote{冲动}是许多品种的最高属，例如\Quote{厌恶}和\Quote{倾向}，并说在那种与属同名的种的情况下，\Quote{倾向}是与\Quote{厌恶}对立且相区别的。\par

% structure: p0010 source=h2 target=subsection reason=relative-heading-rank; base=h2

\subsection{五、田地与宝藏的诠释。}

这里我们必须分别探讨田地，以及藏在其中的宝藏，以及那个找到这隐藏宝藏的人如何满心欢喜地离开，卖掉他所有的一切，以便购买那块田地；并且我们还要探讨——他卖掉的是什么。根据这些，田地在我看来似乎就是圣经，它种植在历史、法律、先知以及其余思想的话中所明显表现出来的；因为整个圣经中话语的种植广大而多样；但藏在田里的宝藏是那些隐藏并潜伏在明显事物之下的思想，\Quote{奥秘中的智慧，}\BibleQuote{就是基督，在他内蕴藏着一切智慧和知识的宝藏。}\BibleRef{哥 2:3}但另一个人可能会说，田地就是那真正充满的、主所祝福的、天主的基督；而藏在其中的宝藏，是保禄所说\Quote{隐藏在基督内}的事物，保禄论及基督说：\BibleQuote{在他内蕴藏着一切智慧和知识的宝藏。}因此，天上的事物，甚至天国，正如经上以比喻的形式所写的——这些就是天国，或基督——他本身就是万世之王，就是那被比作藏在田地之宝藏的天国。\par

% structure: p0012 source=h2 target=subsection reason=relative-heading-rank; base=h2

\subsection{六、诠释继续。}

而且在此处你会询问，天国是否仅仅被比作藏在田里的宝藏，以致我们认为田与天国不同；还是被比作整个藏在田里的宝藏，以致天国按照比喻既包含田也包含藏在田里的宝藏。一个人来到田里，无论是来到圣经还是来到\Person{基督}——祂由有形和无形的事物构成——就发现了藏在\Person{基督}里或圣经中的智慧宝藏。因为，他四处巡视田里，查考圣经并寻求理解\Person{基督}，就在其中发现了宝藏；发现了之后，他就把它藏起来，认为把圣经的隐秘含义或\Person{基督}内的智慧和知识宝藏向所有人揭示并非没有危险。藏好之后，他就离开，努力谋划如何买下那块田，即圣经，好使它成为自己的财产，从天主子民那里接受天主的圣言，这些圣言当初是托付给犹太人的。\BibleRef{罗 3:2}当一个受\Person{基督}教导的人买下那块田时，按照另一个比喻作为葡萄园的天主之国，\BibleQuote{必从他们手中夺去，而交给一个能结果实的民族}\BibleRef{玛 21:43}——即交给那位凭信德买下田的人，作为他变卖所有一切的果实，不再保留任何从前属于自己的东西；因为那些东西曾是他作恶的来源。如果那含有隐藏宝藏的田是\Person{基督}，你也可以给出同样的应用，因为那些舍弃一切跟随祂的人，仿佛以另一种方式变卖了他们的财产，以便通过变卖和交出它们，并从天主——他们的助佑者——那里获得一个高尚的决心作为回报，他们就能以配得上那块田的巨大代价，购买那块含有隐藏宝藏的田。\par

% structure: p0014 source=h2 target=subsection reason=relative-heading-rank; base=h2

\subsection{七、宝贵珍珠的比喻。珍珠的形成与区别。}

\Quote{\Emph{天国又好像一个商人，寻找上好的珍珠。}}\BibleRef{玛 13:45} 有许多商人经营许多种商品，但天国不像其中任何一个，只像那寻找上好的珍珠，并找到一颗价值抵得许多颗的、非常贵重的珍珠，用许多颗换了这一颗的人。因此，我认为对珍珠的本性作一番探究是合理的。但请注意，基督并没有说：\Quote{他卖掉了所有的珍珠，}因为他不只卖掉了那寻找上好珍珠的人所买的，还卖掉了他所拥有的一切，来买那颗好珍珠。于是我们在那些撰写论石著作的人那里发现，关于珍珠的本性，有些珍珠是在陆地上找到的，有些是在海里。陆地珍珠只产于\Person{印度人}当中，适用于指环、项圈和项链；而海洋珍珠更为优越，也产于同一\Person{印度人}当中，最好的产在\Place{红海}。次好的珍珠是从\Place{不列颠}海中采集的；第三等的珍珠，不仅次于第一等，也次于第二等，是在\Place{斯基提亚}附近的\Place{博斯普鲁斯}找到的。关于印度珍珠，还有以下说法。它们产于蚌中，其性质如同非常大的螺旋蜗牛壳；这些蚌被描述为成群地以海为牧场，仿佛受某个领袖指引，该领袖在颜色和大小上显著，与服从它的那些不同，因此它的地位相当于所谓的蜂后。同样，关于捕捉最好的——即印度那些——珍珠的方法，有以下的叙述。当地人用网围住一大片海岸，然后潜下水，尽力抓住它们中领头的那个；因为据说，当这个被捉住时，捕捉属于它的那群蚌就不费吹灰之力，因为群中没有一只会静止不动，而是仿佛被皮带拴住一样跟着群的首领。据说，印度珍珠的形成需要一段时间，这种生物经历许多变化和改变直到完善。还进一步报告说，那蚌壳——我指的是产珍珠的动物的壳——张开并裂开，仿佛似的，张开后接收天上的露水；当它充满了纯净无障碍的露水时，就变得明亮，产出一颗又大又好的珍珠；但如果它某次接收了混浊、不均匀或冬季的露水，就会怀上一颗浑浊且带有瑕疵的珍珠。我们还发现，如果它正怀着珍珠即将完成时被闪电击中，就会闭合，并仿佛因恐惧而散开并涌出其子女，形成所谓的\Quote{\OriginalTerm{气泡}{physemata}。}有时，仿佛早产，它们生得小，虽形状好却有些混浊。与其他珍珠相比，印度珍珠具有这些特征：颜色洁白，透明度如银，闪耀着略带黄绿色的光辉，通常呈圆形；皮薄，比石头的本性更娇嫩；因此悦目，值得在较著名者中称颂，正如那位撰写论石著作的人常说的。这也是最佳珍珠的标志：外表圆润，颜色非常洁白，非常透明，尺寸非常大。关于印度珍珠就这些。但他们在\Place{不列颠}找到的珍珠，据说带金色，但有些混浊，光泽较暗。在\Place{博斯普鲁斯}海峡找到的比\Place{不列颠}的暗，发青，完全黯淡，柔软而小。而且\Place{博斯普鲁斯}海峡中产的珍珠并非产在作为珍珠贝的\Quote{\OriginalTerm{羽贝}{pinna}}中，而是在所谓的蚌中；它们的栖息地——我指的是\Place{博斯普鲁斯}的那些——是在沼泽中。据说还有第四类珍珠，在\Place{阿卡纳尼亚}的牡蛎\Quote{\OriginalTerm{羽贝壳}{pinnæ}}中。这些不太被追捧，但形状不规则，颜色非常暗且污秽；在同一个\Place{阿卡纳尼亚}还有另一些不同的珍珠，被抛弃在任何地方。\par

% structure: p0016 source=h2 target=subsection reason=relative-heading-rank; base=h2

\subsection{八、比喻的解释就是这些观点的光照。}

现在，从关于石头的论说中收集了这些事后，我说，救主深知珍珠的区别——有些在种类上是上好的，有些是无价值的——便说：\BibleQuote{天国好比一个商人寻找好珠子；} \BibleRef{玛 13:45} 因为，若有些珍珠不是无价值的，就不会说 \BibleQuote{一个商人寻找好珠子。} 现在，在一切自称宣告真理的言语中，以及在那些传述这些言语的人中，他寻找珍珠。并且让先知们，可以说，是那些孕育天露的蚌，因从天上来的真理之言而受孕，成了那好珠子，按此处的话说，商人所寻找的。而诸珍珠之首，找到它则其余也一并找到的，是那颗极贵重的珍珠，天主的基督，那超越法律和先知中宝贵文字和思想的圣言，找到它则其余一切都轻易获得。并且救主与众门徒交谈，视他们为不仅寻找好珠子、而且已经找到并拥有它们的商人，说：\BibleQuote{不要把你们的珍珠投在猪前。} \BibleRef{玛 7:6} 现在，这些话是对门徒说的，从祂话前所记的可以显明：\BibleQuote{祂看见群众，就上了山，坐下后，祂的门徒来到祂跟前；} \BibleRef{玛 5:1} 因为，在那段话中，祂说：\BibleQuote{不要把圣物给狗，也不要把你们的珍珠投在猪前。} \BibleRef{玛 7:6} 因此，或许他不是基督的门徒——那人不拥有珍珠或那颗极贵重的珍珠，我指那些好珠子，不是昏暗的，也不是黑暗的，如异端者的话，不是在日出时产生，而是在日落时或北方产生的，若有必要也将那些使我们发现不同地方出产的珍珠差异的事纳入比较的话。也许那些浑浊之词和与肉身行为相连的异端，就是黑暗的珍珠，那些出产于沼泽中的，不是好珠子。\par

% structure: p0018 source=h2 target=subsection reason=relative-heading-rank; base=h2

\subsection{基督——无价珍珠}

现在，你要将寻找好珠子的人与这句话相连：\BibleQuote{你们求，就会得到；} \BibleRef{玛 7:7} 以及这句——\BibleQuote{凡求的，就得到。} \BibleRef{玛 7:8} 你们求什么？或者说，凡求的，得到什么？我敢回答：珍珠，以及那人所拥有的那颗珍珠，他曾弃绝一切，并将它们视为损失；保禄说：\BibleQuote{为了祂，我将一切都看作损失，为赢得基督；} \BibleRef{斐 3:8} 以\BibleQuote{一切}指好珠子，\BibleQuote{为赢得基督，}那颗极其宝贵的珍珠。那么，灯在黑暗中为人是宝贵的，需要灯直到太阳升起；梅瑟脸上的荣耀，我想，先知脸上的荣耀，也是宝贵的，是美丽的景象，藉此我们被引入，得以看见基督的荣耀，圣父为这荣耀作证，说：\BibleQuote{这是我的爱子，我所喜悦的。} \BibleRef{玛 3:17} 但\BibleQuote{那曾得荣耀的，在这事上算不得荣耀，是因那超越的荣耀；} \BibleRef{格后 3:10} 我们先需要那可以被废去的荣耀，为了那超越的荣耀；正如我们需要那部分的认知，它将被废去，当那成全的来到时。\BibleRef{格前 13:9} 因此，每个灵魂达到童年，并走向成熟，直到时机的圆满临近，都需要监护人、管家和护卫，以致在这一切之后，他原先与奴仆并无区别，虽然他是万有的主人，可以摆脱了监护人、管家和护卫之后，领受与那颗极贵重的珍珠相对应的遗产，并与那成全的对应，那成全的来到时，就废去那部分的，当人能领受\BibleQuote{基督认知的卓越} \BibleRef{斐 3:8}，可以说，先前已在那些被基督认知所超越的认知形式中受过锻炼。但是群众，没有觉察到法律的许多珍珠之美，以及先知们那\Quote{部分}的认知（虽然如此），以为他们可以在没有对这些清晰阐述和理解的情况下，整体地找到那颗宝贵的珍珠，并看见\BibleQuote{基督认知的卓越}，与之相比，一切在这如此伟大认知之前的事，虽然在本性上并非渣滓，却显得是渣滓。这渣滓或许就是葡萄园看守扔在无花果树旁的\BibleQuote{粪}，使之结果实的原因。\BibleRef{路 13:8}\par

% structure: p0020 source=h2 target=subsection reason=relative-heading-rank; base=h2

\subsection{福音珍珠与旧约的关系}

\BibleQuote{所以，一切都有其季节，天下万物都有定时} \BibleRef{训 3:1}，有收集好珠子的时候，在收集之后寻找那颗宝贵珍珠的时候，那时人应该去变卖所有来买那颗珍珠。因为正如每个要在真理之言中成为智慧的人，必须先受启蒙，然后经过基础知识教导，并十分珍视它，但不拘泥于它，如同一个人在开始时尊崇它，但迈向成全，感谢那引导，因为它起初是有用的；同样，对法律和先知的完善理解，是对福音完善理解的基础训练，以及基督言行中一切含义的基础训练。\par

% structure: p0022 source=h2 target=subsection reason=relative-heading-rank; base=h2

\subsection{撒网的比喻}

\BibleQuote{\Emph{天国又好像撒在海里的网。}}\BibleRef{玛 13:47} 正如在图像和雕像的情形中，相似之处并非在所有方面都与它们所模拟的对象完全一致；例如，在木板的平面上用蜡绘制的图像，虽然表面与颜色相似，但不再保留凹陷和凸起，只保留其外观；而在雕像的铸造中，则努力在凹陷和凸起方面保持相似，但不保持颜色；如果铸件是用蜡制成的，它则努力同时保持两者，即颜色以及凹陷和凸起，但的确不是关于深度的图像；所以，你也跟我这样理解：在福音的比喻中，当天国被比作某物时，这种比较并不延伸到与天国相比的那物的一切特征，而只延伸到手头论证所需要的那些特征。因此，天国\Quote{像撒在海里的网}——并非如某些人认为的那样（他们以此指代那些进入网中的不同本性，即恶人和义人），仿佛因为\Quote{网聚拢各样水族}这句话，义人和恶人就有许多不同的本性；因为一切圣经都反对这种解释，圣经强调意志的自由，谴责那些犯罪的人，赞许那些行义的人；否则，对那本性如此的种类就不能恰当地加以责备，对那更好的种类也不能加以赞扬。因为鱼的好或坏的原因不在于鱼的灵魂，而基于圣言凭智慧所说的：\BibleQuote{水中要繁生蠕动的生物，具有生命的灵魂}，\BibleRef{创 1:20} 那时，\BibleQuote{天主还造了大鱼和水中各种蠕动的生物，各按其类。}\BibleRef{创 1:21} 因此，那里说\BibleQuote{水中繁生了各种蠕动的生物，各按其类}，其原因不在水中；但在这里，我们要为我们是好的种类、配得称为\Quote{器皿}，还是坏的种类、配得被丢在外面负责。因为在我们内造成邪恶的原因不是本性，而是那行恶的自由意志；同样，造成正义的原因也不是我们的本性，好像它不可能容纳不义，而是我们所接纳的原则使人成为义人；因为你也从未见过水中的种类从坏的种类变为好的，或从好的变为坏的；但你总能看见人中的义人或恶人，或者从恶行转向德行，或者从德行上的进步退回恶行的洪流。因此，在厄则克耳书中，关于那从不义转向遵守天主诫命的人，也是这样记载的：\BibleQuote{如果恶人离开他所行的一切恶事，}等等，直到\BibleQuote{使他离开恶道而生存}；\BibleRef{则 18:20} 但关于那从德行进步退回恶行洪流的人，则说：\BibleQuote{如果义人离开他的正义而行不义，}等等，直到\BibleQuote{他因所犯的罪过必须死亡。}\BibleRef{则 18:24} 那些从撒网的比喻引入不同本性学说的人，请告诉我们：关于那后来离开他所行的一切恶事、遵守天主的一切诫命、并做正义和仁慈之事的恶人，当他作恶时，他是什么本性？显然不是一种值得称赞的本性。如果确实是一种应受谴责的本性，那么当他离开他所犯的一切罪时，他又可以合理地被称为哪种本性？因为如果他因其先前行为而属于坏的那一类本性，他如何能变成更好的？或者，如果你因其后来的行为而说他是好的那一类，那么他既然本性是好的，又怎会变恶？关于那离开正义而行各样不义的义人，你也会遇到类似的难题。因为在他离开正义之前，他行正义之事，他并非属于坏的本性，因为坏的本性不可能在义中，因为坏树——即邪恶——不能结出好果子——即源于德行的果子。另一方面，如果他属于好的、不变的本性，他在被称为义人之后就不会离开善，以致在他所犯的一切罪中行不义。\par

% structure: p0024 source=h2 target=subsection reason=relative-heading-rank; base=h2

\subsection{十二、圣经比喻为网}

既然说了这些话，我们必须坚持\BibleQuote{天国好比网撒在海里，聚拢各样水族}\BibleRef{玛 13:47}，为要展示人类行为原则的多样性，这些原则彼此尽可能不同，所以\Quote{聚拢各样水族}这一说法包含了那些在美德或恶习的形式上值得称赞和值得谴责的倾向。而且天国被比作网的驳杂织法，这是指旧约和新约圣经，它由各种极其多样的思想编织而成。如同落入网中的鱼，有些在网的一部分，有些在另一部分，每条鱼都在它被捉住的地方；同样，对于那些进入圣经之网的人，你会发现有些被先知之网捉住，例如，按照这个说法，依撒意亚、耶肋米亚或达尼尔；有些在法律的网中，有些在福音的网中，有些在宗徒的网中；因为当一个人首次被圣言捕获或看似被捕获时，他是从整个网的某一部分被取出来的。若有些被捕的鱼被圣经中整个网的织法包围，四面受挤压而被捉住，以致无法逃脱，而且如同完全被奴役，不被允许逃出网外，这也毫不奇怪。这网已撒入海中——就是世界上各地人类波浪般动荡的生命，这生命在人生的苦难中游荡。在我们的救主耶稣基督之前，这网并未完全充满；因为法律和先知的网必须由祂来成全，祂说：\BibleQuote{你们不要以为我来是废除法律或先知，我来不是为废除，而是为成全}\BibleRef{玛 5:17}。而网的织法已在福音中，并在基督通过宗徒所说的话中得以完成。因此，为此缘故，\Quote{天国好比网撒在海里，聚拢各样水族}。而且，除了已经说的，\Quote{聚拢各样水族}这句话可以显示从各种族中对外邦人的召唤。而照料那撒在海中的网的是耶稣基督，网的主人，以及\BibleQuote{前来服事祂的天使}\BibleRef{玛 4:11}；他们不会把网从海中拉上来，也不会把它带到海那边的岸边——即超越此生的事物——除非网已装满，也就是说，除非\Quote{外邦人的全数}进入了网中。但当它来临时，他们就把网从下面的事物中拉上来，带到比喻上称为岸上的地方；在那里，拉网的人的工作是坐在岸边，在那里安定下来，以便将网中每一个好的按照此处所称的\Quote{器皿}放入自己的次序，而将那些品性相反、被称为坏的扔在外面，远远抛开。所谓\Quote{外面}，按救主的解释，是指火炉，祂说：\BibleQuote{今世的终结也将如此：天使要出来，把恶人从义人中分开，把他们扔进火炉里}\BibleRef{玛 13:49}。但必须注意，我们在莠子的比喻和所设的比喻中已受教导，天使被托付权力来辨别和分离恶人与义人；因为上面说过：\BibleQuote{人子要派遣祂的天使，从祂的国里收集一切使人跌倒的和作恶的事，把他们扔进火炉里：在那里要有哀号和切齿}\BibleRef{玛 13:42}。但这里说：\BibleQuote{天使要出来，把恶人从义人中分开，把他们扔进火炉里}。\par

% structure: p0026 source=h2 target=subsection reason=relative-heading-rank; base=h2

\subsection{十三、人与天使的关系。}

从这一点并不像有些人所认为的那样，在基督内得救的人甚至比圣天使还优越；因为那些被圣天使投入器皿的人，怎能与那些将他们投入器皿的天使相比呢？他们不是已置于天使的权下吗？我们说这话时，并非不知道在基督内得救的人会超越某些天使——即那些未被委以这项职务的天使——但并非超越所有的天使。因为我们读到：\BibleQuote{对于这些事，天使们渴望详细察看} \BibleRef{伯前 1:12}，这里并没有说\Quote{全部}天使。我们也知道这一点：\BibleQuote{我们要审判天使} \BibleRef{格前 6:3}，这里也没有说\Quote{全部}天使。既然关于渔网和网中人的事是这样记载的，那么，若有人希望在世代的终结之前，在天使来将恶人从义人中分别出来之前，网中不应有\Quote{各类}恶人，他似乎不明白圣经，并且是在追求不可能的事。因此，在天使为此目的被派遣来将恶人从义人中分别出来之前，即使我们看到我们的聚会中也充满了恶人，我们也不必惊讶。但愿那些将被投入火窑的人，数目不会比义人更多！但我们在开始时说过，比喻和比拟不能按它们所比喻或比较的一切方面来接受，而只能按某些方面来接受。因此，我们必须从即将要说的事中进一步确认：就鱼而言，就它们的生命来说，当它们被发现在网中时，对它们发生了一件坏事，因为它们被剥夺了本性的生命，无论它们被投入器皿还是被丢弃，它们所遭受的不过是作为鱼的生命丧失；但是，对于比喻所指的那些人，坏事却是留在海中而不进入网中，以便与好鱼一同被投入器皿。同样，坏鱼被丢在外面；但在我们面前的比喻中，恶人却被投入\BibleQuote{火窑}，以至于在\BibleBook{厄则克耳}{先知书}论及火窑时所说的也会降临到他们身上——\BibleQuote{主的话传给我说：人子，看，以色列家对我来说已完全混杂了铜和铁，}等等，直到这些话：\BibleQuote{你将知道是我，主，向你倾泻了我的怒火。} \BibleRef{则 18:17}\par

% structure: p0028 source=h2 target=subsection reason=relative-heading-rank; base=h2

\subsection{十四、门徒如同经师。}

\BibleQuote{\Emph{你们明白这一切事吗？他们说：是的。}}\BibleRef{玛 13:51} 基督耶稣知道人心里所存的事，\BibleRef{若 2:25} 正如若望也在 \WorkTitle{若望福音} 中论及祂所教导的，祂提出问题并非出于无知，而是因为祂一旦取了人性，便也使用人的一切特征，其中\Quote{询问}便是其一。救主这样做并不足为奇，因为宇宙的天主容忍人的行径，如同人容忍自己儿子的行径，祂查问说：\BibleQuote{亚当，你在哪里？}\BibleRef{创 3:9} 和\BibleQuote{你的兄弟亚伯尔在哪里？}\BibleRef{创 4:9}。但有人会强解此处，说\Quote{明白}一词不应视为疑问，而应视为肯定；并说门徒为祂的肯定作证，答说：\BibleQuote{是的。}唯独不论祂是发问还是肯定，都必然不是说\BibleQuote{这些事}——这是指示性的——也不仅仅说\BibleQuote{一切事}，而是说\BibleQuote{这一切事}。在此，祂似乎把门徒描绘为在天国之前曾是经师；\BibleRef{玛 13:52} 但 \WorkTitle{宗徒大事录} 中的话与此相反：\BibleQuote{他们看见伯多禄和若望的胆量，并发现他们是没有学问的平常人，就惊讶，并认明他们原是跟过耶稣的。}\BibleRef{宗 4:13}有人可能会就这些事发问：如果他们是经师，为何在 \WorkTitle{宗徒大事录} 中被说成是没有学问的平常人？或者如果他们是没学问的平常人，为何救主明明称他们为经师？对此可以回答说：事实上，\WorkTitle{宗徒大事录} 中只描述伯多禄和若望是没有学问的平常人，而非所有门徒；但还有其他门徒，因为他们明白一切事，故称\BibleQuote{每个经师}等等。或者可以说，凡按法律字面教训受教的人，都称为经师，因此那些没有学问、愚钝、并被法律字句俘虏的人，在特定的意义上被称为经师。这正是愚钝之人的特征，他们不擅长比喻解释，也不理解圣经奥秘的阐述，只相信纯字面，并为之辩护，而他们自称经师。因此，有人会解释\BibleQuote{祸哉，你们经师和法利塞假善人！}\BibleRef{玛 23:13} 这话是对那些只懂得字句的人说的。在此你要查考：福音的经师是否如同法律的经师？前者处理福音，如同后者处理法律，诵读、听闻并讲述\BibleQuote{那些含有寓意的事}\BibleRef{迦 4:24}，以便在保存事件历史真实的同时，理解应用于属灵事物的无误奥秘解释原则，使所学的不至于成为\BibleQuote{属邪恶的属灵事物}\BibleRef{弗 6:12}，而是完全相反，即属于良善的属灵事物。一个人在较简单的意义上是一个\BibleQuote{成为天国门徒的经师}，当他从犹太教而来，接受教会所界定的耶稣基督的教导；但在更深的意义上，当他通过圣经的字面接受了初步知识，进而上升到称为诸天之国的属灵事物时，他便是一个经师。并且根据每个思想被获得、被抽象掌握、并通过实例和绝对证明得以证实，人便能理解天国，以致那充满无误知识的人，就在此处所描绘的\Quote{诸天}的众多之国度中。同样，你也会寓意化\BibleQuote{你们悔改吧！因为天国临近了}\BibleRef{玛 3:2} 这句话，意思是经师——即满足于纯字面的人——可以悔改这种解释方法，并接受那通过耶稣基督、生活的圣言而被称为天国的属灵教导。因此，只要耶稣基督——\BibleQuote{在起初就与天主同在，是天主圣言}\BibleRef{若 1:1}——还没有在灵魂里安家，天国就不在其中；但当一个人接近接纳圣言时，天国就临近他了。但如果天国和天主的国在实际中是同一回事，即使概念上不同，那么显然对那些被说\BibleQuote{天主的国就在你们中间}\BibleRef{路 17:21} 的人，也可以对他们说\Quote{天国就在你们中间}；这尤其因为从字句到神恩的悔改，因为\BibleQuote{但当他们转向主时，帕子就除去了。主就是那神}\BibleRef{格后 3:16}。真正的一家之主既是自由又是富有的；富有是因为他由经师的职分成为天国的门徒，在旧约的每一句话中，及在关于基督耶稣新教导的全部知识中，并将这财富积存在自己的宝库中——在天上，他作为成为天国门徒的人在那里积蓄财宝——那里没有虫蛀，也没有贼挖洞透入。\BibleRef{玛 6:20}至于那如我们所说积蓄财宝在天上的人，我们可以真实地断定：没有一个情欲的蛀虫能触及他的属灵和天上的财产。我所说的\Quote{情欲的蛀虫}，是取自\WorkTitle{箴言}，其中写道：\BibleQuote{木中有蛀虫，忧伤同样伤害人心。}\BibleRef{箴 25:20}因为忧伤是虫和蛀，它伤害那心不在天上和属灵事物上的人；因为如果一个人的财宝在这些事物中——\BibleQuote{因为你的财宝在哪里，你的心也必在哪里}\BibleRef{玛 6:21}——他的心就在天上，为此他说：\Quote{虽有大军扎营攻击我，我的心也不害怕。}因此，救主论及的那些说\BibleQuote{凡在我以先来的，都是贼和强盗}\BibleRef{若 10:8} 的贼，也不能突破那些积蓄在天上的事物，也不能突破那在天上并因此说\BibleQuote{祂使我们与祂一同复活，在基督耶稣内使我们与祂一同坐在天上}\BibleRef{弗 2:6} 和\BibleQuote{我们的家乡原是在天上}\BibleRef{斐 3:20} 的心。\par

% structure: p0030 source=h2 target=subsection reason=relative-heading-rank; base=h2

\subsection{十五、家主及其宝库。}

既然\Quote{凡成为天国门徒的经师，就好像一个家主，从他的宝库里拿出新和旧的东西来}\BibleRef{玛 13:52}，那么就显然藉着所谓\OriginalTerm{命题的转换}{conversion of the proposition}，凡不从他的宝库里拿出新和旧的东西来的人，就不是成为天国门徒的经师。因此，我们必须竭尽全力，藉着\Quote{专务宣读、劝勉、教导}\BibleRef{弟前 4:13}，並\Quote{昼夜默思上主的法律}，在我们心里收藏的，不仅是福音、宗徒及其默示录的新神谕，还有法律中那\Quote{拥有未来美物的阴影}\BibleRef{希 10:1}的旧事，以及依照它们说预言的先知书。当我们阅读、知晓並记忆这些事，在适当的时候将属灵的事与属灵的事相比较时，这些事就会被收集起来——不比较那些无法相比的事，而是比较那些可以比较、且在措辞、思想和意见上具有某种相似性、表示相同意义的事，好藉着圣经两个、三个或更多见证人的口\BibleRef{玛 18:16}，我们得以确立並证实天主的每一句话。我们也必须藉着这些事驳斥那些，尽他们所能，将天主性一分为二，把新约从旧约割裂开来的人，以致他们远离那从宝库里拿出新和旧的东西来的家主。既然那被比作某人的人不同于他所比作的人，那么\Quote{成为天国门徒的经师}就是被比作的人，但不同于那\Quote{从他的宝库里拿出新和旧的东西来}的家主。但那位被比作他的人，如同效法他，愿意做那相似的事。或许，那作为家主的人就是耶稣自己，祂按照教导的时期，从祂的宝库里拿出新的、属灵的东西，这些东西也总是在义人的\OriginalTerm{内在的人}{inner man}中被祂更新——义人们自己天天在被更新\BibleRef{格后 4:16}——以及旧的东西，即\Quote{刻在石版上}\BibleRef{格后 3:7}和旧人石心上的东西，这样，藉着字句的比较和属神的彰显，祂便丰富了那成为天国门徒的经师，使他像祂自己一样，直到门徒如同师傅，首先效法那效法基督的人，然后效法基督本身，正如保禄所说的：\Quote{你们该效法我，如同我效法基督一样。}\BibleRef{格前 11:1}同样，在更简单的意义上，耶稣这位家主可以从祂的宝库里拿出新的东西——即福音的教导——和旧的东西——即取自法律和先知的言论的比较，我们在福音中可以找到这些例子。关于这些新和旧的东西，我们也必须注意\BibleBook{肋未}{记}中属灵的法律所说的：\Quote{你们要吃旧的，陈旧的，并从旧的中拿出新的来；我将在你们中间安置我的帐幕。}\BibleRef{肋 26:10}因为我们有福地吃旧的东西——先知的话——和旧中之旧的东西——法律的话；当新的、福音的话来到时，按福音生活，我们从新的面前拿出旧的文字，祂就在我们中间安置祂的帐幕，实现祂所应许的话：\Quote{我要在他们中间居住，在他们中行走。}\par

% structure: p0032 source=h2 target=subsection reason=relative-heading-rank; base=h2

\subsection{十六、比喻与相似。耶稣在自己的家乡。}

\Quote{ \Emph{耶稣讲完了这些比喻，就从那里离开了，来到自己的家乡。}} 见：\BibleRef{玛 13:53} 我们在上文曾探讨，对群众所说的是否是比喻，而对门徒所说的是否是\OriginalTerm{相似比喻}{similitudes}，并提出了在我看来并非微不足道的观察，你们必须知道，随后引用的句子\Quote{耶稣讲完了这些比喻，就从那里离开了}，似乎与所有这些论点相悖，因为它不仅适用于比喻，也适用于我们所阐释的相似比喻。因此，我们要追问，这一切是否都应被否决，或者我们必须谈到两种比喻：对群众所说的和向门徒宣告的；或者我们是否应认为比喻这个词是\OriginalTerm{一词多义}{equi-vocal}；或者那句\Quote{耶稣讲完了这些比喻}是否仅指上文在相似比喻之前的那些比喻。因为，由于这句话\Quote{天国的奥秘是给你们知道的，但对其他人，就用比喻}\BibleRef{玛 13:11}，不可能对门徒说——因为他们不属于外人——救主对他们说话用的是比喻。由此推论，那句\Quote{耶稣讲完了这些比喻，就从那里离开了}应指上文所说的比喻，或者比喻之名是一词多义，或者有两种比喻，或者我们称为相似比喻的根本不是比喻。请注意，他在自己的家乡之外讲说比喻，\Quote{他讲完了这些比喻，就从那里离开，来到自己的家乡，在会堂里教训他们。} 马尔谷说：\Quote{耶稣来到自己的家乡，门徒也跟随着他。}\BibleRef{谷 6:1}因此我们必须追问，\Quote{自己的家乡}是指纳匝肋还是白冷——纳匝肋，因为经上说\Quote{他要被称为纳匝肋人}\BibleRef{玛 2:23}，或者白冷，因为他诞生在那里。再者，我思量，福音作者们是否可能说\Quote{来到白冷}或\Quote{来到纳匝肋}，但他们没有这样做，而是称之为\Quote{自己的家乡}，因为在这段关于他家乡的经文中，以奥秘的方式宣告了某种事——即整个犹太地——在那里，他不受尊敬，正如经上所\Quote{先知除了在自己的家乡，没有不受尊敬的。}\BibleRef{玛 13:57}若有人想到耶稣基督是\Quote{为犹太人是绊脚石}\BibleRef{格前 1:23}，他在犹太人中被逼迫直到如今，却在外邦人中被宣讲并被相信——因为他的话传遍了整个世界——他将会看到，耶稣在自己的家乡没有尊荣，但在那些\Quote{远离盟约}\BibleRef{弗 2:12}的外邦人中，他却受到尊敬。但他在他们的会堂里所教训和所说的话，福音作者们没有记录，只记载了它们是如此伟大，如此性质，以致众人都惊讶。很可能所说的内容太过高深，无法记录下来。但须注意，他在他们的会堂里教训人，并未与之分离，也没有藐视它。\par

% structure: p0034 source=h2 target=subsection reason=relative-heading-rank; base=h2

\subsection{十七、耶稣的兄弟。}

并且这句话：\Quote{\Emph{这人从哪里得来了这样的智慧？}}\BibleRef{玛 13:54} 清楚表明，在\Person{耶稣}的言语中存在着一种伟大卓绝的智慧，配得起这句话：\Quote{看，这里有一位比\Person{撒罗满}更大的。}\BibleRef{玛 12:42} 祂常行比\Person{厄里亚}和\Person{厄里叟}所行的更大的奇迹，甚至比更早的\Person{梅瑟}和农的儿子\Person{若苏厄}所行的还大。他们惊奇地说（不知道祂是童贞之子，或者即使有人告诉他们也不相信，反倒以为祂是木匠\Person{若瑟}的儿子）：\Quote{这不就是木匠的儿子吗？}\BibleRef{玛 13:55} 并且贬低看来是祂最亲近的亲属全体，说：\Quote{祂的母亲不是叫\Person{玛利亚}吗？祂的兄弟不是\Person{雅各伯}、\Person{若瑟}、\Person{西满}和\Person{犹达}吗？祂的妹妹们不是都在我们这里吗？}\BibleRef{玛 13:55} 他们那时以为祂是\Person{若瑟}和\Person{玛利亚}的儿子。但有些人说，依据一个传统，即根据名为\WorkTitle{伯多禄福音}或\WorkTitle{雅各伯书}的传统，认为\Person{耶稣}的兄弟是\Person{若瑟}前妻所生的儿子，这前妻是在\Person{玛利亚}之前娶的。现在这样说的人，是希望保持\Person{玛利亚}直到末了的童贞尊荣，使她已被拣选来服侍那说\Quote{圣神要临于你，至高者的能力要庇荫你}\BibleRef{路 1:35}的圣言的身体，在圣神进入她、从上而来的能力庇荫她之后，不再认识与男人的交合。而且我认为，\Person{耶稣}在男人中、\Person{玛利亚}在女人中，成为由贞洁所构成的纯洁的初果，这是合乎理性的；因为将童贞的初果归于她以外的任何人，都不是虔敬的。\Person{雅各伯}就是\Person{保禄}在\WorkTitle{迦拉达书}中所说他所看见的：\Quote{至于别的宗徒，除了主的兄弟\Person{雅各伯}，我都没有看见。}\BibleRef{迦 1:19} 这位\Person{雅各伯}在人民中因义德而享有如此高的声誉，以致写了二十卷\WorkTitle{犹太人古史}的\Person{弗拉维·若瑟夫}，在想要说明人民为何遭受如此巨大的不幸、以致圣殿被夷为平地时，说这些事是依天主的义怒降在他们身上的，因为他们在被称为基督的\Person{耶稣}的兄弟\Person{雅各伯}身上胆大妄为。奇妙的是，尽管他不接受\Person{耶稣}为基督，却作证\Person{雅各伯}的义德如此伟大；并且他说，人民认为他们遭受这些事是因为\Person{雅各伯}。\Person{犹达}写了一封信，虽然只有寥寥数行，却充满了天上恩宠的健康言语，他在序言中说：\Quote{耶稣基督的仆人，\Person{雅各伯}的兄弟\Person{犹达}。}\BibleRef{犹 1:1} 至于\Person{若瑟}和\Person{西满}，我们没有什么可说的；但这句话：\Quote{祂的妹妹们不是都在我们这里吗？}\BibleRef{玛 13:56} 在我看来似乎意味着这样的事：他们关心我们的事，而不是\Person{耶稣}的事，并没有像\Person{耶稣}那样拥有不寻常的、超越的智慧份量。也许通过这些事，暗示了一种关于祂的新疑惑：\Person{耶稣}不是一个人，而是更神圣的什么，因为祂按他们所想，是\Person{若瑟}和\Person{玛利亚}的儿子，是四个兄弟以及那些妇女的兄弟，却与祂的亲属毫无相似之处，并且不是通过教育和教导达到如此高度的智慧和能力。因为他们在别处也说：\Quote{这人没有学过，怎么认识字呢？}\BibleRef{若 7:15} 这与这里所说的相似。只是，虽然他们说了这些话并如此困惑和惊奇，他们却不相信，反倒因祂而绊倒；好像他们心中的眼睛被那些在受难之际祂将要借着十字架凯旋率领的权势所制服了。\par

% structure: p0036 source=h2 target=subsection reason=relative-heading-rank; base=h2

\subsection{十八、先知在本乡。}

\BibleQuote{\Emph{但耶稣对他们说：先知除了在自己的家乡之外，没有不受尊敬的。}} \BibleRef{玛 13:57} 我们必须探究，这种说法当普遍应用于每一位先知时（似乎每一位先知都只在自己的家乡不受尊敬，但并非所有不受尊敬者都只在自己的家乡不受尊敬），抑或由于措辞是单数的，这些话是针对某一个说的。如果这些话是针对一个说的，那么，若我们将所写的应用于救主，以上所说已足够。但若是普遍性的，那么从历史看并不真实；因为\Person{厄里亚}在\Place{基肋阿得的提市贝}并未受羞辱，\Person{厄里叟}在\Place{阿贝耳默曷拉}也未受羞辱，\Person{撒慕尔}在\Place{辣玛塔因}也未受羞辱，\Person{耶肋米亚}在\Place{阿纳托特}也未受羞辱。然而，按喻意解释，这完全真实；因为我们必须将\Place{犹太}看作他们的家乡，将著名的\Place{以色列}看作他们的亲族，也许将身体看作房屋。因为所有先知在\Place{犹太}地，从那按血肉而生的\Place{以色列}手中遭受了羞辱，那时他们仍在身体内，正如在\WorkTitle{宗徒大事录}中所记载的，是对百姓的斥责说话：\BibleQuote{哪一位先知，你们的祖先没有迫害过？他们预先宣告了义人的来临。} \BibleRef{宗 7:52} \Person{保禄}在\WorkTitle{得撒洛尼前书}中也说了类似的话：\BibleQuote{弟兄们，你们曾效法天主在基督耶稣内，在犹太的各教会；因为你们也由同乡人受了苦难，如同他们由犹太人所受的一样，那些犹太人不但杀害了主耶稣和先知们，也驱逐了我们，他们不得天主的喜悦，并敌视众人。} \BibleRef{得前 2:14} 因此，先知在外邦人中并非没有尊荣；因为他们要么根本不认识他，要么在认识和接纳他为先知后尊敬他。教会的成员就是这样的人。先知受羞辱，首先，从历史事实看，是当他们受到百姓的迫害时；其次，是当他们的预言不被百姓相信时。因为如果他们相信了梅瑟和先知们，他们就会相信基督；基督表明，当人相信\Person{梅瑟}和先知时，对基督的信仰自然随之而来，而当人不信基督时，他们也不信\Person{梅瑟}。 \BibleRef{若 5:46} 此外，正如因违反法律而犯罪的人被称为羞辱天主，同样，因不相信预言，先知就被那不信预言的人所羞辱。就字面真理而言，回忆\Person{耶肋米亚}在百姓中所遭受的事是有益的，他曾说：\BibleQuote{我说：我不再提他，也不再以他的名发言。} \BibleRef{耶 20:9} 又在另一处说：\BibleQuote{我终日成为笑柄。} \BibleRef{耶 20:7} 他从当时的以色列王所受的苦难，在他的预言书中都有记载。也有记载说，百姓时常要砸死\Person{梅瑟}；因为他自己的家乡并非任何地方的石块，而是那跟随他的百姓，他在他们中间也受了羞辱。\Person{依撒意亚}据说是被百姓锯成两半的；若有人因这事出自\WorkTitle{伪经依撒意亚}而不接受，就让他相信\WorkTitle{希伯来书}所记载的：\BibleQuote{他们被石头砸死，被锯成两半，被试探；} \BibleRef{希 11:37} 因为\BibleQuote{被锯成两半}指\Person{依撒意亚}，就如\BibleQuote{被刀杀死}指\Person{匝加利亚}，他被杀于\BibleQuote{在圣所和祭台之间}，正如救主所教导的，我想这是为一部并非见于通行广泛流传的书籍的圣经作证，也许是在伪经中。他们也在自己的家乡，在犹太人中受了羞辱，他们\BibleQuote{披着绵羊皮、山羊皮，受着贫乏、磨难}等等； \BibleRef{希 11:37} \BibleQuote{凡是愿意在基督耶稣内虔敬生活的，都必遭受迫害。} \BibleRef{弟后 3:12} \Person{保禄}大概知道\BibleQuote{先知在自己的家乡没有尊荣}，所以他在许多地方宣讲圣言，却没有在\Place{塔尔索}宣讲。宗徒们也因此离开了\Place{以色列}，去执行救主所吩咐的：\BibleQuote{使万民成为门徒，} \BibleRef{玛 28:19} 以及\BibleQuote{你们要在\Place{耶路撒冷}、全\Place{犹太}和\Place{撒玛黎雅}，并直到地极，为我作证人。} \BibleRef{宗 1:8} 因为他们已在\Place{犹太}和\Place{耶路撒冷}做了所吩咐的事；但是，由于先知在自己的家乡没有尊荣，当犹太人不接纳圣言时，他们就去了外邦人那里。再想一想，由于\BibleQuote{我要将我的圣神倾注在一切有血肉的身上，他们要说预言} \BibleRef{岳 2:28} 这句话已在从外邦人中来的各教会中应验，你可以说，那些先前属于世界、因信不再属于世界的人，在自己的家乡（即世界）领受了圣神并说预言，却没有尊荣，反而受羞辱。为此，那些与先知遭受同样苦难的人是有福的，正如救主所说：\BibleQuote{因为他们的祖先也同样对待了先知。} \BibleRef{路 6:23} 如果细心留意这些事的人因生活严谨、责备罪人而受人憎恨和攻击，作为一个为义受迫害、受辱骂的人，他不仅不会悲伤，反而会欢喜踊跃，因为确信因这些事，他将在天上获得大赏报，那赏报来自那一位，因他遭受了同样的事而将他比作先知。因此，凡热切效法先知生活，并达到先知所具有的圣神的人，必定会在世界和罪人眼中受羞辱，因为义人的生活对他们是重担。\par

% structure: p0038 source=h2 target=subsection reason=relative-heading-rank; base=h2

\subsection{十九、信德与不信与耶稣超自然能力的关系}

继此之后，你可见到\BibleQuote{\Emph{祂因他们的不信，没有在那里行许多大能奇事。}}\BibleRef{玛 13:58} 我们由此学到，德能在信者身上显效，因为\BibleQuote{凡是有的，还要给他，使他富足}\BibleRef{玛 13:12}；但在不信者中，德能不但不起作用，而且如马尔谷所载，\BibleQuote{他们不能行。}\BibleRef{玛 17:19} 请留意经文：\Quote{祂不能在那里行任何大能奇事}；因为经文不是说\Quote{祂不愿意}，而是说\Quote{祂不能}；仿佛行使德能时，需要受惠者的信德与之合作，但此合作因不信而受阻。请看，对那些说\Quote{我们为什么不能赶出它？}的人，祂说：\BibleQuote{因为你们的小信德。}\BibleRef{玛 14:31} 对伯多禄，当他开始下沉时，祂说：\BibleQuote{小信德的人哪，你为什么怀疑？}\BibleRef{路 8:45} 此外，那患血漏的妇人，并未求医治，只心想若能触摸祂的衣边就必痊愈，当场就痊愈了。救主承认这医治方式，说：\BibleQuote{谁摸了我？我感觉有德能从我身上出去。}\BibleRef{玛 17:20} 或许，正如在物质事物中，某些事物对另一些事物有天然的吸引力，如磁石吸铁，所谓石脑油吸火，同样，信德对天主的德能也有吸引力，正如经上所记：\BibleQuote{如果你们有信德像芥子那样大，你们向这座山说：从这里移到那里去，它必移去。}\BibleRef{玛 13:58} 玛窦和马尔谷欲彰显天主权能的卓越，即在不信中也有权能，但不如在受益者信德中的权能那么大；我看他们说得准确：不是\Quote{任何}大能奇事，而是\Quote{许多}。\BibleRef{谷 6:5} 马尔谷也没有说祂不能在那里行任何大能奇事就止步，而是补充说：\BibleQuote{除给几个病人覆手，治愈了他们外}\BibleRef{谷 6:5}，内在于祂的权能如此克服了不信。在我看来，如同物质事物，耕作本身不足以保证收成，除非空气配合——不如说，祂以祂所愿的任何品质形成空气并使之成为所愿；空气也不能离开耕作，而是那位以祂的眷顾命定从地里生出的东西不能离开耕作而生——因为祂已在律法中一次而永远地命定了：\BibleQuote{地要生出青草，结种子的蔬菜，各从其类，并结果子的树，各从其类}\BibleRef{创 1:11}；同样，德能的运作，离开受治者的信德，不能显明完全的医治；信德不论多大，离开天主的德能也是如此。关于智慧所记载的，你也可应用于信德以及各项德行，以形成如下训诫：\BibleQuote{人子中若有人在智慧上成全，而缺了从祢来的德能，他就被算作虚无}\BibleRef{智 9:6}；或\Quote{人子中若有人在自制上成全，尽其所能，而缺了从祢来的自制，他就被算作虚无}；或\Quote{人子中若有人在义德及其余德行上成全，而缺了从祢来的义德及余德，他就被算作虚无。}因此，\BibleQuote{智慧人不要因他的智慧夸耀，勇士不要因他的勇力夸耀}\BibleRef{耶 9:23}，因为可夸耀的并非属于我们，而是天主的恩赐：智慧来自祂，勇力来自祂，其余亦然。\par

% structure: p0040 source=h2 target=subsection reason=relative-heading-rank; base=h2

\subsection{二十、关于洗者若翰的不同理解。}

\BibleQuote{\Emph{那时，分封侯黑落德听到耶稣的消息，就对臣仆说：这是洗者若翰。}}\BibleRef{玛 14:1} \BibleBook{马尔谷}{福音}\BibleRef{谷 6:14}和\BibleBook{路加}{福音}\BibleRef{路 9:7}也是如此记载。犹太人持有不同意见：有些是错误的，如撒杜塞人关于死人复活的主张——他们认为死人不会复活，关于天使，他们认为天使不存在，圣经中关于天使的记载只是比喻性的解释，实际上并不真实；有些则是正确的，如法利塞人所教导的——死人会复活。因此，我们在这里必须探究：黑落德和百姓中一些人错误地持有的关于灵魂的看法，是否类似这样——刚刚被他杀害的若翰，在斩首后复活了，只是换了名字，如今被称为耶稣，且具有从前在若翰内运行的能力。因为如果一个人被全体百姓如此熟知，名声传遍全犹大，人们称他为木匠和玛利亚的儿子，并有这样那样的兄弟姐妹，却被认为与若翰没有什么不同——若翰的父亲是匝加利亚，母亲是依撒伯尔，他们在百姓中也并非无名之辈——这种想法有何可信之处呢？但或许，百姓并非不知道他是匝加利亚的儿子，他们视若翰为真先知，人数众多，以致法利塞人为了避免说出令百姓不悦的话，不敢回答\Quote{他的洗礼是从天上来的，还是从人来的？}的问题。\BibleRef{玛 21:25} 或许，有些人也知道圣殿中所见的异象，即加俾额尔向匝加利亚显现的事。那么，黑落德或百姓中某些人的错误见解——认为若翰和耶稣不是两个人，而是同一个人若翰，在被斩首后复活，改名为耶稣——有何可信之处呢？然而，有人可能会说，黑落德和百姓中的一些人持有\OriginalTerm{灵魂转生}{transmigration of souls}进入身体的错误教义，因此他们认为先前的若翰通过新生再次显现，并从死亡中进入生命成为耶稣。但若翰的出生与耶稣的出生之间相隔不过六个月，这不允许这一错误见解被认为是可信的。或许，黑落德心中更可能是这样的想法：即在若翰内运行的能力转移到了耶稣身上，因此百姓认为耶稣就是洗者若翰。我们可以用以下论证：正如因着厄里亚的灵和能力（而非因他的灵魂），关于若翰有话说\Quote{他就是那将要来的厄里亚}\BibleRef{玛 11:14}，厄里亚内的灵和能力转移到了若翰身上——同样，黑落德认为，若翰内的能力在他身上施行了洗礼和教导的工作（因为若翰没有行过一件奇迹\BibleRef{若 10:41}），但在耶稣身上却显出了奇事。可以这样说，那些说厄里亚在耶稣身上显现，或说一位古先知复活了的人，\BibleRef{路 9:8}其想法与此类似。但那些说耶稣是\Quote{一位先知，如同众先知之一}\BibleRef{谷 6:15}的看法，与此无关。因此，关于耶稣的说法——无论是黑落德的看法，还是别人的说法——都是错误的。只是，\Quote{若翰在厄里亚的灵和能力中先行}\BibleRef{路 1:17}的说法，与当时他们对若翰和耶稣所怀的思想相符，在我看来更为可信。但既然我们首先知道，救主在受试探后听说若翰被交出去，就退到加里肋亚；其次，若翰在监狱中听到关于耶稣的事，就派两个门徒去问祂\Quote{你是那要来的，还是我们等候另一位？}\BibleRef{玛 11:2}；第三，大体上，黑落德关于耶稣说\Quote{这是洗者若翰，他从死人中复活了}\BibleRef{玛 14:2}；但我们先前未曾从任何途径得知洗者若翰被杀的具体情况，因此玛窦现在记下了此事，马尔谷也大致相同；而路加则对那记载中的大部分叙述保持沉默。\par

% structure: p0042 source=h2 target=subsection reason=relative-heading-rank; base=h2

\subsection{二十一、黑落德与洗者若翰。}

玛窦的叙述如下——\BibleQuote{\Emph{因为黑落德已抓住若翰，把他捆绑在监狱里。}}\BibleRef{玛 14:3}继此之事，在我看来，正如\BibleQuote{法律和先知到若翰为止}\BibleRef{路 16:16}，此后，从犹太人中间停止了预言的恩宠；同样，在人民中掌权者的权威，包括杀死他们认为该死之人的权力，也存在于若翰为止；当最后一位先知被黑落德非法杀害时，犹太人的王被剥夺了处死的权力；因为，如果黑落德没有被剥夺这权力，比拉多就不会判处耶稣死刑；但为此，黑落德连同为此聚集的司祭长和民间长老的公议会本是足够的。于是我想，雅各伯对犹大所说的这话应验了：\BibleQuote{权杖不离犹大，柄杖不离以色列，直到那所贮藏的来到，而他乃是外邦人所期待的。}\BibleRef{创 49:10}也许犹太人也被剥夺了这种权力，是天主上智安排，为使基督的教导在人民中传播，即使受到犹太人的阻挡，反对也不至于发展到杀害信徒的地步——这杀害在过去似乎是合法的。\BibleQuote{但是黑落德抓住若翰，把他捆绑起来，关在监狱里。}\BibleRef{玛 14:3}这一举动表明，就他的权力和人民的邪恶而言，他捆绑并囚禁了先知的话，阻止它像从前一样自由地宣告真理。但是黑落德这样做是为了他兄弟斐理伯的妻子黑落狄雅。因为若翰对他说：\BibleQuote{你占有她是不合法的。}\BibleRef{玛 14:3}这个斐理伯是依突勒雅和特辣曷尼地区的分封侯。有些人便推测，当斐理伯去世留下一个女儿黑落狄雅时，黑落德娶了他兄弟的妻子，虽然法律只允许在没有子女的情况下结婚。但是，由于我们找不到任何清楚的证据表明斐理伯已死，我们断定黑落德犯了更大的罪，即他在他的兄弟还活着的时候，诱使他的妻子背叛她的丈夫。\par

% structure: p0044 source=h2 target=subsection reason=relative-heading-rank; base=h2

\subsection{二十二、黑落狄雅的舞蹈。誓言的持守。}

因此，\Person{若翰}具有先知的胆量，不因\Person{黑落德}的王权而畏惧，也不因怕死而对此滔天大罪缄默不语，充满圣神向\Person{黑落德}说：\Quote{你娶她是不合法的；因为你占有你兄弟的妻子是不合法的。}因为\Person{黑落德}捉住\Person{若翰}，把他捆绑起来关在狱中，不敢公然杀他，从人民中除去先知之言；但\Place{特辣黑尼雅}王的妻子——这是一种邪恶的意见和恶毒的教导——生了一个同名的女儿，她的动作看似和谐，取悦了喜好生日庆典的\Person{黑落德}，成了人民中不再有先知之首的原因。至此，我认为犹太人人民的动作，看似合乎法律，无非是\Person{黑落狄雅}女儿的动作；但\Person{黑落狄雅}的舞蹈与那种神圣的舞蹈相对立，那些没有舞蹈的人将因听到\Quote{我们给你们吹笛，你们却不舞蹈}的话而受责备。在生日上，当不法之言统治他们时，他们舞蹈，以便他们的动作取悦那话语。我们之前有人注意到了\WorkTitle{创世记}中关于\Person{法郎}生日的记载，并指出那个喜爱生日之事的卑劣之人举行生日庆典；我们从这个提示中，在圣经中找不到任何义人过生日的记载。因为\Person{黑落德}比那位著名的\Person{法郎}更不义；后者在他的生日宴中杀了一个膳长；\BibleRef{创 40:20}但前者却杀了\Person{若翰}，\BibleQuote{女人所生的，没有比若翰更大的}，\BibleRef{玛 11:11}救主论及他说：\BibleQuote{你们出去究竟为什么？为看先知吗？我告诉你们：是的，而且他比先知还大。}\BibleRef{路 7:26}但感谢天主，即使预言的恩宠已从人民中夺去，一种比这一切更大的恩宠却由我们的救主\Person{耶稣基督}倾注在外邦人中，祂成了\Quote{死人中的自由者}；因为\BibleQuote{祂虽然由于软弱而被钉死，却由于天主的德能而活着。}\BibleRef{格后 13:4}也思考那涉及洁净与不洁食物的言论；但预言却被轻视，当它被放在碟子里代替食物时。但犹太人没有预言之首，因为他们否认一切预言之冠——\Person{基督耶稣}；先知被斩首，是因为一个誓言，在这种情况下，宁愿违背誓言也不愿遵守誓言；因为轻率发誓和因轻率而违背誓言的罪责，与先知之死相比，其罪过并不相同。不仅因此他被斩首，还因为\Quote{同席的人}，他们宁愿先知被杀而不愿他活着。他们同席吃饭，也与统治犹太人的邪恶之言一同宴乐，他们为他的出生而欢庆。有时你可以优雅地将这段经文应用于那些轻率发誓并希望坚持为不法行为所发誓言的人，说并非每一个守誓都是合宜的，就像\Person{黑落德}守誓一样。另外，请注意，\Person{黑落德}不是公然而是秘密在狱中杀了\Person{若翰}。因为犹太人现今的话语并不公开否认预言，而是实际上秘密地否认，并被证明不信预言。因为\BibleQuote{如果你们相信\Person{梅瑟}，你们也会相信\Person{耶稣}}，\BibleRef{若 5:46}同样，如果他们相信先知，他们也会接受那位曾是预言对象的\Person{基督}。但既然不信祂，他们也不信先知，并将预言之言割断、囚禁在狱中，视其为死而分裂的，一点也不健全，因为他们不懂。但我们却拥有完整的\Person{耶稣}，关于祂的预言得以应验：\BibleQuote{祂的骨头一根也不可折断。}\par

% structure: p0046 source=h2 target=subsection reason=relative-heading-rank; base=h2

\subsection{二十三、耶稣的退隐}

若翰的门徒前来收敛了他的遗体，就\BibleQuote{\Emph{便去告诉耶稣。}}\BibleRef{玛 14:12}耶稣便退到荒野地方——那就是外邦人中去——先知被杀后，群众从各城各处跟随他；看见这众多的人，他就怜悯他们，治好了他们的病人；随后用那被祝福而由几个饼倍增的饼，喂养了跟随他的人。\BibleQuote{\Emph{耶稣一听说，就从那里上船，退到荒野地方去。}}\BibleRef{玛 14:13}字面教导我们，要尽可能从迫害我们的人那里以及从预期通过言语而来的阴谋中退避；因为这是谨慎行事；当一个人能避开危险处境时，去迎头赶上就是鲁莽和任性。既然不仅耶稣在若翰所遇之事后撤退，而且还教导说：\BibleQuote{如果他们在这城中迫害你们，你们就逃到另一城}？\BibleRef{玛 10:23}那么谁还会犹豫不去避免这样的事呢？当试探来临而我们无法避免时，必须以极高的高尚和勇气忍受它；但如果我们可以避免却不这样做，就是鲁莽。然而，在字面之后，我们还需按神秘意义考察这个地方，我们必须说：当预言在犹太人中被阴谋反对并摧毁，因为他们重视生日之事，并接受那些虚妄的活动——尽管这些活动被邪恶的首领及其同席者视为规律和令他们喜悦，但若以真理为评判，却是不规律和不和谐的——，那时耶稣就从预言受攻击和定罪的地方退去；他退到外邦人中那曾经对天主荒芜的地方，为使天主圣言——当国度从犹太人那里被夺去，并\BibleQuote{交结果子的民族，}\BibleRef{玛 21:43}——能在外邦人中；并且，因此，\BibleQuote{荒凉者的子女，}他们未曾受法律和先知的教导，\BibleQuote{比有丈夫的生子更多，}那丈夫就是法律。当圣言从前在犹太人中间时，它并不像现在在外邦人中那样居留；因此，经上说，他\BibleQuote{在船上，}——即身体中——退到荒野地方，当他听到先知被杀的事。他来到荒野地方，就住在其中，因为圣言独自居住，他的教训与在外邦人中通行的习俗惯例相反。外邦人中的群众，听说耶稣来到他们的荒野居住，并且他是独处的，正如我们已经报告的，便从自己的城里跟随他，因为每个人都离开了父亲的迷信习俗，来到了基督的法律。他们从陆路跟随他，而不是在船上，因为他们不是以身体，而是以灵魂，以及他们被圣言说服的决心，跟随天主的肖像。耶稣出来到他们那里，因为他们不能到他那里去，目的是：他到了那些外面的人那里，可以把那些外面的人领到里面。那些外面的人群众很大，天主圣言出去到他们那里，倾注他\BibleQuote{眷顾，}的光，注视他们；看到他们更值得怜悯，因为他们处于这样的处境，那位不可动情者，作为爱人的，却动了怜悯之情，不仅怜悯，还治好了他们的病人，这些病人因他们的邪恶而有各种各类的疾病。\par

% structure: p0048 source=h2 target=subsection reason=relative-heading-rank; base=h2

\subsection{二十四、属灵疾病的多种形式。}

并且，若你愿意察看灵魂的疾病是何性质，请与我一同思量贪财者、贪恋虚荣者、好男色者，以及好女色者；因为祂看见这些人在群众中，就怜悯他们，治好了他们。因为并非每一种罪都应视为疾病，唯有那弥漫于整个灵魂的罪才是。因为如此，你看到贪财者完全专注于钱财及其保存与积聚，贪恋虚荣者完全专注于小小的荣耀，因为他们渴望来自群众和粗俗之人的赞美；以此类推，你也能理解我们所提及的其他情况，以及任何类似的情况。既然，当我们解释\BibleQuote{祂治好了他们的病人}\BibleRef{玛 14:14}时，我们说过并非每一种罪都是疾病，那么从圣经中探讨这些区别是合宜的。宗徒在写给哥林多人的书信中，论及那些有不同疾病的人，确实说：\BibleQuote{因此，在你们中间有许多软弱的、有病的，不少的人已沉睡了}\BibleRef{格前 11:30}。请在这些话中听他所言，他将不同的罪编织成一条带子并捻成辫子，根据有的软弱，有的比软弱更严重而有病，有的相较于两者则是沉睡。因为有些人由于灵魂的无能，倾向于滑入任何罪中，尽管他们可能并非完全被某种罪的形式所掌控（如那些有病的人），他们只是软弱；但另一些人，不按照\Quote{全心、全灵、全意}爱天主，却爱钱财、小小的荣耀、妻子或儿女，他们所遭受的比软弱更糟，而是有病。那些沉睡的人是那些本应小心谨慎、以灵魂警醒，却没有这样做，反而因极大的疏忽而在决意中打盹、在思虑中昏昏欲睡的人，例如\BibleQuote{他们在梦中玷污肉身，轻视主权者，亵渎众尊荣者}\BibleRef{犹 1:8}。这些人因为沉睡，生活在关于现实之空虚梦幻般的幻想氛围中，不接受真实的事，反被其虚妄想象所呈现的所欺骗；关于他们，依撒意亚说：\Quote{如一个口渴的人梦见自己在喝水，但醒来时仍然口渴，他的灵魂怀抱了虚妄的盼望；凡在耶路撒冷争战的所有民族的财富也是如此}。那么，若我们在叙述软弱者、有病者和沉睡者之间的区别时似乎离了题，那是因为宗徒在他写给哥林多人的书信中所说的话，我们已经解释过了；我们离题，是渴望表明我们从\BibleQuote{祂治好了他们的病人}\BibleRef{玛 14:14}这句话中所当了解的意思。\par

% structure: p0050 source=h2 target=subsection reason=relative-heading-rank; base=h2

\subsection{二十五、治愈先于参与耶稣的饼。}

此后，经上说：\BibleQuote{\Emph{当} \Emph{到了晚上，祂的门徒前来对祂说：\InnerQuote{这地方是荒野，时候已过；请遣散群众，叫他们往村庄里去，为自己买食物。}}}\BibleRef{玛 14:15} 首先应注意，当祂要将祝福的饼交给门徒，让他们放在群众面前时，祂先治好了病人，为使他们恢复健康后，得以参与祝福的饼；因为当他们还有病时，不能接受耶稣祝福的饼。但若有人，当他应听从\BibleQuote{但各人应省察自己，然后才吃这饼}等训导\BibleRef{格前 11:28}，却不服从这些话，而是随意参与主的饼和祂的杯，他便变得软弱或有病，甚至——容我这样说——因饼的能力而昏迷，进入沉睡。\par

\SourceRecord{Translated by John Patrick.；Ante-Nicene Fathers；Vol. 9.；Edited by Allan Menzies.；Buffalo, NY: Christian Literature Publishing Co.,；1896.；<http://www.newadvent.org/fathers/101610.htm>.}

% structure: p0001 source=h1 target=section reason=page-title

\section{\WorkTitle{玛窦福音释义}（第十一卷）}

% structure: p0002 source=h2 target=subsection reason=relative-heading-rank; base=h2

\subsection{一、饱食五千人的引言}

\Quote{\Emph{到了傍晚，祂的门徒来到祂跟前},}\BibleRef{玛 14:15} 也就是，在时期的终结，关于此我们适当地引用若望书信中的话：\Quote{现在是最后的时刻。}\BibleRef{若一 2:18} 他们尚不明白圣言将要做的事，就对祂说：\Quote{这地方是荒野，}\BibleRef{玛 14:15} 他们看到群众在有关天主、法律和圣言方面的荒芜状态；但他们又对祂说：\Quote{时候已经过了，}\BibleRef{玛 14:15} 仿佛法律和先知们的适当季节已经过去。或许他们说出这话，是引用耶稣的话语，即由于若翰被斩首，法律和先知直到若翰为止都停止了。\BibleRef{路 16:16} 因此他们说\Quote{时候已经过了，}并且没有食物在手，因为它的季节已不再存在，以致那些在荒野中跟随祢的人可以事奉法律和先知。此外，门徒们又说：\Quote{请打发他们走，}\BibleRef{玛 14:15} 好叫每人去买食物，如果不能在城里买，至少可以在村庄——更卑微的地方。门徒们说这些话，是因为在法律的字句已被废除、预言已停止之后，他们对于为群众找到意外的和新的食物感到绝望。但请看耶稣如何回答门徒，尽管祂没有大声喊叫，也没有明说：\Quote{你们以为，如果这大批群众没有食物离开我，他们会在村庄而不是在我这里找到食物，会在人群——不是市民而是村民——中找到，而不是留在我身边。但我告诉你们，关于你们以为他们缺乏的，他们并不缺乏，因为他们无需离开；而关于你们以为他们不缺——即缺少我——好像我不能喂饱他们，恰恰相反，他们需要。既然我已经训练了你们，并使你们适合给予那些需要的人理性的食物，你们就给我跟随的群众吃吧；因为你们有从我领受的能力，使群众吃饱；如果你们留意到这一点，你们就会明白我更有能力喂饱他们，就不会说：\InnerQuote{打发群众走，他们好去为自己买食物。}}\BibleRef{玛 14:15}\par

% structure: p0004 source=h2 target=subsection reason=relative-heading-rank; base=h2

\subsection{二、奇迹细节的阐释}

耶稣于是因祂赐给门徒的能力，即滋养他人的能力，说：\BibleQuote{你们给他们吃吧。}\newline{}\BibleRef{玛 14:16}但是（门徒并非否认他们能拿出饼，而是觉得饼太少，不足以喂养跟随耶稣的人，且没有考虑到，当耶稣拿起每一个饼——即\OriginalTerm{圣言}{Word}——祂便随意伸展它，使之足以供养所有祂愿意滋养的人），门徒说：\BibleQuote{我们这里只有五个饼和两条鱼。}\newline{}\BibleRef{玛 14:17}或许他们用五个饼暗指圣经中感性的话语，其数目因此与五种感官相应；而两条鱼则暗指表达之言与内在之言，它们可以说是圣经中所含感性事物的调味品；或者，也许指所传给他们的关于圣父与圣子的言语。因此，祂复活后也吃了烤鱼，\newline{}\BibleRef{路 24:42}从门徒那里取了一份，并接受了那关于圣父的神学，门徒能将其中一部分向祂宣告。这就是我们目前对五饼二鱼话语的解释所能提供的贡献；或许那些比我们更能将五饼二鱼集于自身的人，能给出更完备、更佳的意义诠释。然而必须注意，在玛窦、马尔谷和路加福音中，门徒说他们有五饼二鱼，未指明饼是麦面的还是大麦的；唯有\Person{若望}说，那些饼是大麦饼。\newline{}\BibleRef{若 6:9}因此，或许在\BibleBook{若望}{福音}中，门徒不承认饼在他们那里，却说：\BibleQuote{这里有一个孩童，带着五个大麦饼和两条鱼。}\newline{}\BibleRef{若 6:9}只要这五饼二鱼未被耶稣的门徒携带，它们就不会增加或增多，也不能滋养更多人；但当救主拿起它们，首先举目望天，借祂的目光如光束般从天吸引能力，与即将喂养五千人的饼和鱼混合；之后祝福五饼二鱼，以言语和祝福使之增多、加增；第三步，擘开分给门徒，让他们摆在群众面前，这时饼和鱼足够，使众人都吃了且饱足，而他们未能吃尽祝福之饼的部分。因为给群众剩下的如此多，这并不按群众的容量，而是按门徒的容量，他们能收集剩下的碎块，放入装满剩余物的篮子中，篮子的数目与以色列支派相同。关于\Person{若瑟}，圣咏上记载：\BibleQuote{他的手在篮子里服役。}但关于耶稣的门徒，他们收集了剩下的碎块，共十二篮，我以为是十二篮，不是半满，而是满的。我认为，直到如今，且直到时代的终结，与耶稣的门徒（他们优于群众）同在的，是那十二篮装满活面包碎块，群众不能吃。那些吃了十二篮剩余之前的五饼的人，其本性亲近数字五；因为吃的人已达到感觉之物的阶段，因为他们也被那位望天祝福并擘开的主滋养，他们不是孩童或妇女，而是男人。因为我认为，即使在感觉的食物中也有区别，以至于其中一些属于那\BibleQuote{已脱去幼稚之事，}\newline{}\BibleRef{格前 13:11}另一些则属于那些在基督内仍是婴孩和属血肉的人。\par

% structure: p0006 source=h2 target=subsection reason=relative-heading-rank; base=h2

\subsection{三、细节解释续。坐在草地上。分排坐下。}

我们这样说是因为下列话语：\BibleQuote{\Emph{吃了的人除了妇女孩子，约有五千男子}}，\BibleRef{玛 14:21}，这是一个模棱两可的说法；因为要么吃的人是五千男子，其中没有孩子或妇女；要么只有男子是五千，孩子和妇女不计在内。因此，如我们先前所言，有些人如此理解这段经文：当五饼二鱼增多和繁衍时，没有孩子或妇女在场。但有人可能会说，虽然许多人吃了并按照他们的功劳和容量分享了祝福的饼，但有些堪当被数算的人，相当于户籍纪中所数的二十岁以上的男子，\BibleRef{户 1:3}，是以色列男子；而其他不配如此计数和数算的人则是孩子和妇女。此外，请与我一同按寓意解释孩子，对应这段经文：\BibleQuote{我未能将你们当作属神的人，只得将你们当作属血肉的人，当作在基督内的婴孩}；\BibleRef{格前 3:1}；妇女则对应这句话：\BibleQuote{我愿意把你们众人如同贞洁的童女献给基督}；\BibleRef{格后 11:2}；男子对应这话：\BibleQuote{当我成了人时，便除去了孩子的事}。\BibleRef{格前 13:11}。我们不可不经阐释就略过以下话语：\BibleQuote{\Emph{祂吩咐群众坐在草地上，就拿起那五个饼和两条鱼，望着天祝福了，擘开，将饼交给门徒，门徒又递给群众。众人都吃了，并且吃饱了}}。\BibleRef{玛 14:19}。因为这句话是什么意思：\BibleQuote{祂吩咐所有群众坐在草地上}？我们在这段经文中当如何理解那堪当耶稣命令之事？如今，我认为祂吩咐群众坐在草地上，是因为依撒意亚所说：\BibleQuote{凡有血肉的都是草}；\BibleRef{依 40:6}；也就是说，祂吩咐他们制服肉体，使\BibleQuote{肉性的思念}\BibleRef{罗 8:6}居于服从的地位，好使任何人能分受耶稣祝福的饼。既然需要耶稣供应的食物的人有不同的等级，并非所有人都以相等的话语得到滋养，为此缘故，我认为马尔谷记载：\BibleQuote{祂就吩咐他们叫众人一帮一帮地坐在青草地上；他们就一排一排地坐下，有一百一排的，有五十一排的}；\BibleRef{谷 6:39}；而路加记载：\BibleQuote{祂对门徒说：叫他们一排一排地坐下，每排大约五十人}。\BibleRef{路 9:14}。因为那些要在耶稣的食物中得安息的人，要么必须属于一百的次序——这神圣的数字，因其中的一而归于天主；要么属于五十的次序——这数字包含罪的赦免，根据每五十年一次的\OriginalTerm{禧年}{Jubilee}的奥秘和\OriginalTerm{五旬节}{Pentecost}的庆节。我认为那十二筐为门徒所有，对他们曾说过：\BibleQuote{你们要坐在十二宝座上，审判以色列十二支派}。\BibleRef{玛 19:28}。正如审判勒乌本领支派的人的宝座可说是一个奥秘，审判西默盎支派的人的宝座是另一个，审判犹大支派的人的是另一个，其余类推；同样，可能有一个装着勒乌本食物的筐，另一个装着西默盎的，另一个装着肋未的。然而，现在偏离当前主题，去收集关于十二支派所说的话，以及关于每个支派分别所说的话，并说明以色列每一支派可能象征什么，这不符合我们当前的论述。\par

% structure: p0008 source=h2 target=subsection reason=relative-heading-rank; base=h2

\subsection{四、群众与门徒的对照。}

\BibleQuote{\Emph{祂即刻催迫门徒上船，先渡到对岸，等祂遣散群众。}}\BibleRef{玛 14:22} 应当留意，在同样的段落中，\Quote{群众}一词与\Quote{门徒}一词被提及的频次是何等频繁，以致通过观察并汇集与此事相关的段落，可以看出福音作者的目的，是要通过福音史事来表明那些来到耶稣面前之人的区别；其中有些是群众，并不被称为门徒，另有些是门徒，比群众优越。然而，对于目前而言，我们只需举出少数言说，以便凡受其感动者，皆可对全部福音书如法炮制。经上这样记载——仿佛群众在下面，而门徒却能在耶稣上山时到祂那里去，群众却不能在那里——如下：\BibleQuote{祂看见群众，就上了山，坐下以后，祂的门徒来到祂跟前；祂就开口教训他们说，神贫的人是有福的，}\BibleRef{玛 5:1} 又在另一处，因群众需要医治，经上说：\BibleQuote{许多群众跟随了祂，祂治好了他们。}\BibleRef{玛 12:15} 我们并未发现有关于门徒得医治的记载；因为如果有人已经是耶稣的门徒，他便是完整的，既已健康，便不需要耶稣作为医师，而是需要祂的其他权能。又在另一处，当祂向群众说话时，祂的母亲和兄弟站在外面，想同祂说话；有人将这事告诉了祂，祂便伸手指着——不是指向群众，而是指向门徒——说：\BibleQuote{看，我的母亲和我的兄弟！}\BibleRef{玛 14:46} 并因门徒履行了在天之父的旨意而为他们作证，祂又说：\BibleQuote{他就是我的兄弟、姊妹和母亲。}\BibleRef{玛 14:50} 又在另一处记载说：\BibleQuote{全体群众站在海滩上，祂就用比喻对他们讲了许多事。}\BibleRef{玛 13:2} 随后，在撒种的比喻之后，不再是群众，而是门徒来对祂说，不是\Quote{你为什么用比喻对我们讲话？}，而是\BibleQuote{你为什么用比喻对他们讲话？}\BibleRef{玛 13:10} 那时祂也回答说，不是对群众，而是对门徒说：\BibleQuote{天国的奥秘是赐给你们知道的，但对其余的人，却用比喻。}\BibleRef{玛 13:11} 因此，在那些来到耶稣名下的人中，有些认识天国奥秘的，可称为门徒；而那些未蒙此特恩的，则称为群众，被视作次于门徒。因为你要仔细留意，祂对门徒说：\BibleQuote{天国的奥秘是赐给你们知道的，}但关于群众，祂说：\Quote{对他们，却没有赐下}\BibleRef{玛 13:11} 又在另一处，祂确实遣散了群众，进了屋子，\BibleRef{玛 13:36} 但祂没有遣散门徒；到祂屋里来的，不是群众，而是祂的门徒，说：\BibleQuote{请将田间莠子的比喻解释给我们听。}\BibleRef{玛 13:36} 此外，又在另一处，当耶稣听到关于约翰的事，便乘船退到一处偏僻的荒野地方，群众却跟随了祂；祂出来看见一大群人，便对他们动了怜悯的心，治好了他们的病人——群众中的病人，而非门徒中的。\BibleRef{玛 14:13} \BibleQuote{到了傍晚，门徒来到祂跟前，}不是群众，而是门徒，以显出他们不同于群众，说：\BibleQuote{请遣散群众，好让他们往村庄里去，为自己买食物。}\BibleRef{玛 14:15} 并且，当耶稣拿起五个饼和两条鱼，举目望天，祝福了并掰开饼，祂没有给群众，而是给门徒，\BibleRef{玛 14:19} 好让门徒分给群众，群众不能直接从祂接受，而只能艰难地从门徒手中接受耶稣所祝福的饼，并且也没有吃完这些饼；因为群众吃饱了，剩下的碎块以满盛十二筐。\par

% structure: p0010 source=h2 target=subsection reason=relative-heading-rank; base=h2

\subsection{五、门徒在风浪中。耶稣步行于水上。}

我们之所以讨论这个主题，是因为所讨论的经文告诉我们，耶稣将门徒与群众分开，并强迫他们上船，先划到对岸去，直到祂自己打发群众散开；\BibleRef{玛 14:22} 因为群众不能到对岸去，他们在神秘的意义上不是希伯来人，希伯来文的意思是 \Quote{在彼岸居住}。但这是耶稣门徒的工作——我是指到对岸去，超越可见的、物质的事物，视之为暂时的，并趋向那不可见的、永恒的事物。被耶稣打发走是耶稣赐予群众的充分恩惠；只因他们是群众，就不能到对岸去；而这种打发，除了耶稣以外，没有人有能力做到，并且除非一个人先吃了耶稣所祝福的饼，否则他不可能被打发走。一个人也不可能吃耶稣祝福的饼，除非他照耶稣所吩咐的，坐在草地上，正如我们所说的。群众也不可能这样做，除非他们从自己的城跟随耶稣，那时祂退到一处僻静的地方。起初，当门徒请求耶稣打发群众时，祂没有打发他们，直到祂用祝福的饼餵饱了他们；但现在祂却打发了他们，首先强迫门徒上船；祂打发他们时，他们在某处低地——因为旷野是低地——但祂自己上山去祈祷。\BibleRef{玛 14:23} 你们必须注意这一点，即在五千人吃饱之后，耶稣立刻强迫门徒上船，先划到对岸去。只是，门徒不能先于耶稣到对岸去；但是，当他们到了海的中途，船被浪摇撼，\BibleQuote{因为风不顺}，\BibleRef{玛 14:24} 他们在夜间四更天，耶稣向他们走去时，就害怕了。如果耶稣没有上船，那么阻碍航行门徒的风也不会停止，航行的人也不能渡过到达对岸。也许，耶稣想通过经验教训他们，没有祂就不可能到对岸去，于是强迫他们上船，先到对岸去；但是，当他们不能前进超过海的中途时，祂向他们显现，做了所记载的事，\BibleRef{玛 14:25} 并表明到达对岸的人之所以到达，是因为耶稣与他同航。但耶稣强迫门徒上的那条船是什么？也许就是试探与困难的争战，任何人被圣言强迫进入其中，不情愿地，仿佛救主要藉着操练来训练门徒在这艘被波浪和逆风摇撼的船中。但由于马尔谷稍微改变了读法，他写的是\BibleQuote{祂立刻强迫门徒上船，先划到对岸，往贝特赛达去}，\BibleRef{谷 6:45} 而不是\BibleQuote{祂立刻强迫门徒上船，先划到对岸}，我们必须注意\Quote{祂强迫}这个词，首先我们看到了马尔谷的细微变化，他通过增加代词指明了更确定的意思；因为同一件事不是用\Quote{强迫门徒}来表达。马尔谷写的不仅是\Quote{门徒}，而是\Quote{祂的门徒}。因此，或许要注意措辞，那些难以离开耶稣、不能因任何普通原因与祂分离的门徒，愿意与祂同在；但祂判断他们应当尝试波浪和逆风（如果与耶稣同在，这些就不会是逆风），便将他们与祂分离并上船的必要性强加给他们。救主于是强迫门徒进入试探之船，先划到对岸，藉着胜过试探而超越关键的困难；但是，当他们进入了海的中途，在试探的波浪中，并被阻止他们到对岸去的逆风阻挡时，他们无法在没有耶稣的情况下挣扎着胜过波浪和逆风而到达对岸。因此，圣言怜悯那些竭尽全力要到对岸去的人，便步行在海面上来到他们那里，对祂来说，没有波浪或风能够阻挡，如果祂愿意的话；因为经上不是写着\BibleQuote{祂步行在波浪上来到他们那里}，而是写着\BibleQuote{在水上}；\BibleRef{玛 14:25} 正如伯多禄，起初当耶稣对他说\BibleQuote{来吧}时，他下了船，不是步行在\BibleQuote{波浪上}，而是\BibleQuote{水上}，\BibleRef{玛 14:29} 来到耶稣那里；但当他疑惑时，他看到风很大，那风对于放下小信德和疑惑的人并不大。但是，当耶稣与伯多禄一起上了船时，风就停了，因为当耶稣上了船时，风就没有力量再与船对抗了。\par

% structure: p0012 source=h2 target=subsection reason=relative-heading-rank; base=h2

\subsection{六、叙述细节的诠释。其对所有门徒的应用。}

然后门徒们\Quote{\Emph{渡到对岸，来到革尼撒勒地方}，}\BibleRef{玛 14:34} 若我们知道这词的释义，也许对解释本段会有所帮助。请看，既然天主是信实的，且不容许群众受试探超过他们所能承担的，那么天主子是如何强迫门徒们上船——因为他们更强壮，能够到达海的中部并忍受波浪的考验，直到他们配得上神圣的援助，并看见耶稣，听见祂上岸后的话，然后渡到对岸，来到革尼撒勒地方；但至于群众，因他们较为软弱，没有经历船、波浪和逆风，祂就遣散了他们，独自上山祈祷。\BibleRef{玛 14:22} 祂为谁祈祷？也许是为群众祈祷，使他们被遣散后，在祝福的饼之后，不致做出违背耶稣遣散他们的事？也为门徒祈祷，使他们被祂强迫上船、先行到对岸去时，不致在海中或逆风中受害？我敢说，因耶稣为门徒向圣父的祈祷，当海、波浪和逆风与他们对抗时，他们毫发无伤。因此，较为质朴的门徒可以满足于单纯的叙述；但让我们记住，若我们陷入痛苦的试探，是耶稣强迫我们上他们的船，希望我们比祂先到对岸；因为我们不可能到达对岸，除非我们经历了波浪和逆风的考验。然后，当我们看到许多困难包围我们，并通过适度的挣扎勉强游过一些时，让我们认为我们的船在海中，被波浪困苦，这些波浪企图使我们在信德或某种德行上遭遇船难；但当我们看到恶者的灵与我们对抗时，让我们意识到那时风是逆着我们的。于是，当我们在此痛苦中经历了三更——即试探中的黑暗——用全力英勇奋斗并警醒自己，以免在信德或某种德行上遭遇船难——第一更对抗黑暗与邪恶之父，第二更对抗他的儿子\Quote{他敌对并高举自己在各种称为神或受人崇拜者以上，}\BibleRef{得后 2:4} 第三更对抗那敌对圣神之神，那么我们就相信，当第四更临近时，当\Quote{黑夜已深，白昼临近，}\BibleRef{罗 13:12} 天主子会来到我们这里，为我们预备海，并在其上行走。当我们看见圣言向我们显现时，我们确实会惊慌，直到我们清楚明白是救主来到我们这里，以为我们仍在看见幻象，因恐惧而呼喊；但祂自己立刻会对我们说：\Quote{放心！是我，不要怕！}\BibleRef{玛 14:27} 并且，若有人因祂的\Quote{放心！}而热切感动，在我们中间有某个伯多禄，他正走向成全但尚未成为成全的，从船上下到水面上，仿佛从他困苦的试探中出来，起初确会行走，想要踏水来见耶稣；但因他小信德且尚存疑虑，看见风大就害怕，将要下沉；但他不会沉没，因为他要大声向耶稣呼求，对祂说：\Quote{主，救我！}\BibleRef{玛 14:30} 随即，当这样一个伯多禄还在说话，说\Quote{主，救我！}时，圣言会伸出祂的手，向此人提供援助，在他开始下沉时拉住他，并责备他的小信德和疑虑。\BibleRef{玛 14:31} 只是，注意祂并没有说：\Quote{你这无信德的人！}而是说：\Quote{你这小信德的人！}并且问：\Quote{你为什么疑虑？}因为他仍有几分信德，但也有趋向与信德相反之事的倾向。\par

% structure: p0014 source=h2 target=subsection reason=relative-heading-rank; base=h2

\subsection{七、对岸病人的医治。医治的方法。}

但此后，\Person{耶稣}和\Person{伯多禄}将上船，风将平息；船上的人觉察到从大危险中得救，就朝拜祂，说，不只是像两个附魔的人所说\Quote{祢是天主子，}而是\Quote{祢真是天主子。}\BibleRef{玛 14:33}这是船上的门徒说的，因为我不认为除门徒外还有别人这样说。当我们经历了这一切，渡过对岸后，我们就会来到耶稣吩咐我们在他之前先去的陆地。也许，关于某些被耶稣拯救的人，还有一些隐秘的奥迹由这些话指出：\Quote{那地方的人一认出祂来，}——显然是指对岸那地方——\Quote{就打发人向周围整个地区，}——对岸的周围地区，不是对岸本身，而是其周围——\Quote{把一切有病的人都带到祂跟前。}\BibleRef{玛 14:35}这里要注意，他们不仅把许多病人带来，而且带来了周围整个地区的所有病人；被带来的人恳求祂，只求能摸一摸祂的衣边，\BibleRef{玛 14:36}向祂恳求这份恩宠，因为他们不像\Quote{那个患血漏十二年的妇人，她来到祂背后，摸了祂的衣边，心里说：我只要一摸祂的衣裳，就会痊愈。}\BibleRef{玛 9:20}请看关于祂衣边所说的话，因着这句话，她的血漏立刻止住了。但来自\Place{革乃撒勒}周围地区的人——\Person{耶稣}和祂的门徒渡海来到这地方——并没有自己来到耶稣跟前，而是被那些报信的人带来，因为他们极其软弱，无法自己前来。他们也不仅只是摸一摸衣裳，像那个患血漏的妇人那样，而是恳求之后才摸。只是这些人中，\Quote{凡摸到的都痊愈了。}\BibleRef{玛 14:36}至于他们身上所说的\Quote{痊愈了}与\Quote{得救}之间是否有区别——因为曾对那患血漏的妇人说：\Quote{你的信德救了你}\BibleRef{玛 9:22}——你可以自己思考。\par

% structure: p0016 source=h2 target=subsection reason=relative-heading-rank; base=h2

\subsection{八、论前来询问的法利塞人和经师：为什么祢的门徒违犯先人的传授？}

\BibleQuote{\Emph{那时，有从耶路撒冷来到祂那里的} \Emph{法利塞人和经师说：为什么祢的门徒违犯先人的传统？因为他们吃饭时并不洗手。}} \BibleRef{玛 15:1} 谁若留意法利塞人和经师从耶路撒冷来到耶稣那里，说：\BibleQuote{为什么祢的门徒违犯先人的传统，}等等，就会察觉玛窦必然不是简单地写说法利塞人和经师从耶路撒冷来到救主那里询问我们眼前的事，而是这样写：\BibleQuote{那时，他们从耶路撒冷来到祂那里。} 那么，我们应如何理解\BibleQuote{那时}所指的时间呢？就是在耶稣和门徒渡海，乘船来到革乃撒勒地方的时候，自从耶稣上了船，风就停了，并且当\BibleQuote{那地方的人认出祂来，就派人走遍周围整个地区，把一切患病的人都带到祂那里，恳求他们使他们至少能摸一摸祂的衣边；凡摸着的，都痊愈了。} \BibleRef{玛 14:35} 那时，有法利塞人和经师从耶路撒冷来到祂那里，他们并不惊叹耶稣内蕴的能力——这能力治愈了那些即使只摸到祂衣边的人——而是以吹毛求疵的态度，在师傅面前控告门徒，所指责的不是触犯天主的诫命，而只是犹太先人的一条传统。而且很可能，这些吹毛求疵的人所提出的指控恰恰证明了耶稣门徒的虔诚，因为法利塞人和经师未能找到门徒违反天主诫命的把柄；如果他们能责备他们所控告的人，并指出他们违犯了天主的诫命，他们就不会控告门徒违犯先人的传统了。但不要以为这些事证实了必须按字面遵守梅瑟法律，因为耶稣的门徒直到那时还遵守它；因为祂是在受苦之前就\BibleQuote{赎我们脱离了法律的诅咒，} \BibleRef{迦 3:13} 祂为了人而受苦，\BibleQuote{为我们成了诅咒。} 然而，正如保禄合宜地成了犹太人中的犹太人，为要赢得犹太人\BibleRef{格前 9:20}，宗徒们的生活方式是在犹太人中间度过，即使他们明白法律中属灵的事物，他们采取迁就的态度，又有什么奇怪呢？保禄也曾这样做：给弟茂德行割礼\BibleRef{迦 2:3}，并按照某种法律誓愿献祭，正如\BibleBook{}{宗徒大事录}所记载的。只是，他们似乎又喜爱控告，因为他们无法以天主的诫命来指控耶稣的门徒，只能以一条先人的传统来控告。这种爱控告的性情尤其显明在这一点上：他们是在那些被治愈的病人面前提出控告；表面上是针对门徒，实际上却意在诽谤他们的老师，因为洗手是先人的传统，被认为是虔敬所必需的。因为他们认为，那些在吃饭前不洗手的手是不洁不净的，而那些用水洗过的手则变得纯洁神圣——这不是比喻意义上的，而是按字面与梅瑟法律相应。但是，让我们不按犹太先人的传统，而是按正确的理性，努力净化我们自己的行为，从而洗净我们灵魂的双手，当我们即将吃那三个饼时——这饼是我们向耶稣祈求的，祂愿意作我们的朋友；因为手若不洁、未洗、不净，我们就不该分享这饼。\par

% structure: p0018 source=h2 target=subsection reason=relative-heading-rank; base=h2

\subsection{九、\OriginalTerm{科尔班}{Corban}的解释}

耶稣却并未根据犹太长老的传统控告他们，而是针对天主的两个最重要的诫命。其中一条是十诫中的第五诫，内容如下：\BibleQuote{应孝敬你的父亲和你的母亲，好使你能得福，并在上主你的天主赐给你的地上延年益寿；} \BibleRef{出 20:12}另一条在\BibleBook{肋未}{纪}中如此记载：\BibleQuote{若有人咒骂自己的父亲或母亲，应处死；他咒骂了父亲或母亲，应自负血债。} \BibleRef{肋 20:9}但是，当我们想要查考\BibleBook{玛窦}{福音}所记的字句时，\BibleQuote{咒骂父亲或母亲的，应处死。} \BibleRef{玛 15:4}就要考虑这是否取自经上记载的地方：\BibleQuote{凡殴打父亲或母亲的，应处死；咒骂父亲或母亲的，应处死。}因为这就是法律中关于这两条诫命的准确字句；但玛窦只引用了部分且形式简略，并非原文字句。但我们必须思考，救主指控从耶路撒冷来的法利塞人和经师说他们因自己的传统而违犯了天主的诫命，这指控的性质为何。天主曾说：\BibleQuote{应孝敬你的父亲和你的母亲，} \BibleRef{出 20:12}教导子女应给父母当得的孝敬。这种对父母的孝敬，一部分是与父母分享生活必需品，如食物、衣服，以及其他任何可能向父母表示恩惠的事物。但法利塞人和经师却颁布了一条与法律相对立的传统，这传统在福音中相当隐晦地出现，我们自己也想不到，除非有一个希伯来人告诉我们以下与这段经文相关的事实。他说，有时放债人遇到顽固的债务人——他们有偿还能力却不愿偿还——便将所欠的款项奉献给穷人账户，因为每个愿意按自己的能力将部分财物施舍给穷人的人，都会将钱投入库中。因此，他们有时用本地方言对债务人说：\Quote{你所欠我的，是\OriginalTerm{科尔班}{Corban}，}——即奉献——\Quote{因为我已将它奉献给穷人，归于对天主的虔敬账上。}于是债务人，既然不再欠人，而是欠天主和天主的虔敬，就仿佛被束缚，即使不情愿，也不得不偿还债务，不再还给放债人，而是以放债人的名义还给天主，归入穷人账上。放债人对债务人所做的，有时儿子们也对自己的父母这样做，对他们说：\Quote{父亲或母亲，你本可从我得利的那份，要知道你将从\OriginalTerm{科尔班}{Corban}获得，} \BibleRef{玛 15:4}即从奉献给天主的穷人账上获得。于是，父母听说本应给他们的是\OriginalTerm{科尔班}{Corban}——奉献给天主——便不再愿意从儿子那里收取，即使他们极度缺乏生活必需品。于是，长老们向百姓宣布了这样一条传统：\Quote{无论谁对父亲或母亲说：本应给他们的任何东西都是\OriginalTerm{科尔班}{Corban}，是奉献，这人便不再有义务给父母提供生活必需品。}救主谴责这条传统，认为它不正，且违背天主的诫命。因为如果天主说：\BibleQuote{应孝敬你的父亲和你的母亲，}而传统却说：凡将本应给父母的东西奉献给天主作为\OriginalTerm{科尔班}{Corban}的人，便没有义务以礼物孝敬父母，那么显然，天主关于孝敬父母的诫命被法利塞人和经师的这条传统无效化了，他们说：凡是一次性地将父母可能收到的奉献给天主的人，便不再有义务孝敬父母。法利塞人既贪爱钱财，为了假借穷人的名义甚至收取本应给任何人的父母的东西，便这样教导。福音也证明他们贪爱钱财，说：\BibleQuote{法利塞人既贪爱钱财，听了这一切，便嗤笑祂。} \BibleRef{路 16:14}因此，如果我们中间被称为长老的人，或任何以某种方式做人民首领的人，声称以公益之名施舍穷人，而不愿施舍给偶然需要生活必需品的亲属，且施舍者不能两者兼顾，那么此人完全有理由被称为那些法利塞人的同道，他们因自己的传统使天主的话失效，并被救主斥为伪善者。作为一个强有力的劝阻，防止任何人急于从穷人账上取用，并认为\BibleQuote{别人的虔敬是获利之道} \BibleRef{弟前 6:5}，我们不仅有这些，还有关于叛徒犹达斯的记载，他表面上维护穷人的利益，愤慨地说：\BibleQuote{这香液原可卖三百银钱，施舍给穷人！}但实际上\BibleQuote{他是个贼，且掌管钱囊，常取其中所存。} \BibleRef{若 12:6}因此，如果在我们时代，有人掌管教会的钱囊，像犹达斯一样为穷人说话，却取走其中所存，愿他得与行这些事的犹达斯同分一份；因为这些事如毒疮侵蚀他的灵魂，魔鬼便将出卖救主的念头投入他心中；当他已领受\BibleQuote{火箭} \BibleRef{弗 6:16}之后，魔鬼随后亲自进入他的灵魂，完全占据了他。也许，当宗徒说：\BibleQuote{贪爱钱财是万恶的根源} \BibleRef{弟前 6:10}时，他这样说是因为犹达斯贪爱钱财，这是对耶稣所行一切恶事的根源。\par

% structure: p0020 source=h2 target=subsection reason=relative-heading-rank; base=h2

\subsection{十、长老传统与天主法律的冲突。}

但让我们回到眼前的主题，其中救主简化和阐述了法律中的两条诫命，一条出自出谷纪的十诫，另一条出自肋未纪，或另一条出自梅瑟五书中的某卷。既然我们已经说明了他们如何藉着说\Quote{你们不可孝敬你的父亲或母亲}来废了天主所说\Quote{孝敬你的父亲和母亲}的话，——凡是对父亲或母亲说：\Quote{你由我所得必成供物}，——那么有人可能会问，\Quote{咒骂父亲或母亲的，必须死} \BibleRef{玛 15:4}这句话是否无关呢？因为，的确，若有人将本应用于孝敬父母的东西奉献给所谓的\OriginalTerm{供物}{Corban}，他就不孝敬父母；那么，法利塞人的传统是如何废了\Quote{咒骂父亲或母亲的，必须死}这句话的呢？但是，也许，当任何人对父亲或母亲说：\Quote{你由我所得必成供物} \BibleRef{玛 15:5}时，他仿佛是在辱骂父亲或母亲，如同称呼父母为亵渎者，因为他们从已奉献给供物的人那里拿取了已奉献给供物的东西。犹太人于是按法律惩罚他们的儿子，视他们为咒骂父亲或母亲，当他们向父亲或母亲说：\Quote{你由我所得必成供物}时；但你们藉着你们的一个传统，废了天主的两个诫命。而你们却不以为羞地控告我的门徒，他们不违犯任何诫命，因为他们\Quote{在祂的一切诫命和法令上无瑕可指}，但违犯了长老的传统，为的是不违犯天主的诫命。如果你们事先持守了这个目标，你们必会遵守关于孝敬父母的诫命，以及那句\Quote{咒骂父亲或母亲的，必须死}的话；但你们必不会遵守与这些诫命相悖的长老传统。\par

% structure: p0022 source=h2 target=subsection reason=relative-heading-rank; base=h2

\subsection{十一、耶稣所引依撒意亚预言的阐释。}

此後，為要從先知的話語徹底駁斥猶太人中的這些長老傳統，祂便向他們引用依撒意亞的一句話，其準確措辭如下：\Quote{上主說：這民族口頭上親近我，}等等；\BibleRef{依 29:13}而如同我們先前所說，瑪竇並未逐字寫出先知的話。若因福音中使用此段而需按其力所能及解釋，我們將補充前面的一段；依我判斷，這段經文對我們解釋福音中取自先知的那段經文頗有裨益。依撒意亞從開頭的經文是這樣的：\Quote{你們要昏迷，且要瘋狂；你們要沉醉，但不是因烈酒或葡萄酒；因為上主將沉睡之靈傾注給你們，祂將封閉他們的眼睛，就是他們的先知和那些觀看隱密事的首領的眼睛。這一切話語對你們將如同被封住的書卷上的話；若將它交給一個識字的人說：讀這書，他會回答：我不能讀，因為它被封住了。這書將交給一個不識字的人，有人對他說：讀這書，他會說：我不識字。上主說：這民族親近我，等等，直到這些話：\InnerQuote{禍哉！那些在暗中策劃陰謀，他們的行為必在黑暗中。}}\BibleRef{依 29:9}於是，我們拿起福音中在我們面前的這段經文，我列出了它前面的一些經節和後面的一些經節，以表明聖言如何威脅要封閉那些驚奇、沉醉並被灌下沉睡之靈的百姓的眼睛。它也威脅要封閉他們的先知和那些自稱能看見隱密事的首領的眼睛——我認為這些事發生在救主來到這民族之後；因為全部聖經的話，連依撒意亞的話，對他們都如同被封住的書卷上的話。現在，\Quote{封住}一詞用於一本因其晦澀而封閉、不因其清晰而敞開的書；它對於完全因不識字而不能讀的人，與那些自稱識字卻不明白所寫之事的意義的人，同樣晦澀。那麼，祂是否補充說：當這民族因自己的罪而暈厥，並處於瘋狂狀態，藉著他們將要沉醉於其中的那些罪惡，以沉睡之靈向祂發怒——這沉睡之靈是上主在封閉他們的眼睛時傾注給他們喝的，因為他們不配看見，以及他們的先知和那些自稱能看見聖經中奧秘隱藏之事的首領的眼睛；當他們的眼睛被封閉時，先知的話對他們就成了封住和隱藏的，正如那些不信耶穌是基督的人所遭遇的。當先知的話語成了封住的書卷上的話，不僅對不識字的人，也對自稱識字的人，那時上主說，猶太民族只用口親近天主，他們用嘴唇尊敬祂，但他們的心因不信耶穌而遠離上主。如今，特別是在他們否認我們的救主之後，天主可以論及他們說：\Quote{但他們徒然敬拜我；}\BibleRef{瑪 15:9}因為他們不再教導天主的誡命，而是教導人的誡命，教義是屬人的，不再是智慧之神的教義。因此，當這些事臨到他們時，天主便撤走了猶太民族，並使他們中間智慧人的智慧消亡；因為他們中間不再有智慧，如同不再有預言；但天主徹底摧毀了精明人的精明，將它隱藏起來，\BibleRef{依 29:14}它不再輝煌奪目。因此，即使他們似乎以深奧的方式策劃某種計謀，但因不是藉著上主而行，他們被稱為可憐的；即使他們自稱透露一些天主計畫的奧秘，他們也是說謊，因為他們的行為不是光明之舉，而是黑暗和黑夜之舉。\BibleRef{依 29:15}我認為適宜簡要陳述這預言，並在某種程度上闡明其意義，因為瑪竇提到了它。馬爾谷也提到了它，我們可以從他那裡有益地記下以下有關長老違背的話——長老們認為猶太人吃餅時必須洗手：\Quote{因為法利塞人和所有的猶太人，若不仔細洗手，就不吃，持守長老的傳統；從市場回來，若不洗浴，就不吃。還有其他他們所領受持守的事：杯、壺、銅器及床榻的洗滌。}\BibleRef{谷 7:3}\par

% structure: p0024 source=h2 target=subsection reason=relative-heading-rank; base=h2

\subsection{十二、按法律和福音的潔與不潔。}

\Quote{\Emph{祂便叫群众过来，对他们说：你们听，且要明白}}等等。\BibleRef{玛 15:10} 在这些话里，救主清楚地教导我们，当我们读\WorkTitle{肋未纪}和\WorkTitle{申命纪}中关于洁净与不洁净肉食的戒律时，我们因违犯这些戒律而被世俗的犹太人和与他们相差无几的厄彼雍派所指责，但我们不要认为圣经的要旨在于对它们的肤浅理解。因为如果\BibleQuote{不是入于口的使人污秽，而是出于口的才使人污秽}，\BibleRef{玛 15:11} 而且特别是按照马尔谷的记载，救主说这些事\BibleQuote{洁净了一切食物}，\BibleRef{谷 7:19} 那么显然，当我们吃那些愿意受制于法律字句的犹太人称为不洁净的东西时，我们并不受玷污；但当我们——尽管我们的嘴唇应当用理性束缚，并且我们应当\BibleQuote{为它们制造所谓的天平和砝码}，\BibleRef{德 28:25}——信口开河、讨论不当之事时，我们才受玷污，罪恶的源头就从那里临到我们。诚然，天主的法律应当禁止出于邪恶的事，命令趋于德行的事；但对于那些本身无所谓善恶的事，则让它们保持原状，因为它们根据我们的选择和内在的理性，若我们在其上犯罪，便行得恶；若我们正确引导，便行得善。凡细心思考这些事的人都会看到，即使在被认为是善的事上，人若以邪恶的方式、在情欲冲动下接受它们，也可能犯罪；而这些被称为不洁的东西，如果我们按理性使用，就可以被视为洁净。因此，正如犹太人犯罪时，他的割损被视为未割损；而外邦人正直行事时，他的未割损被视为割损，\BibleRef{罗 2:25} 同样，那些被认为是洁净的东西，若使用不当——不在适当的时候、不按适当的方式、不为适当的理由——将被视为不洁。至于那些被称为不洁的东西，\BibleQuote{对洁净的人，一切都是洁净的}；因为，\BibleQuote{对败坏和不信的人，什么也不洁净，连他们的理智和良心也败坏了}。\BibleRef{铎 1:15} 当这些东西败坏了，它们所触及的一切都变得败坏；反之，清洁的理智和良心使一切变得洁净，即使它们看似不洁；因为义人食用肉食或饮料，并非出于无节制、贪图享乐或使人犹豫的疑虑，而是谨记诫命：\BibleQuote{你们或吃或喝，或无论做什么，一切都要为光荣天主而做}。\BibleRef{格前 10:31} 如果有必要描述福音中不洁的食物，我们将说：它们是由贪婪供给的，是卑劣营利的结果，出于享乐之爱，以及神化备受尊崇的肚腹——当肚腹及其欲望，而非理性，统治我们的灵魂时。至于我们，我们知道有些东西被魔鬼使用，或者如果我们不知道，但怀疑并犹豫，如果我们使用这些东西，那我们就不是为光荣天主，也不是因基督的名而使用；因为不仅怀疑祭过偶像的食物会定那吃的人为有罪，就连这一点上的怀疑也是如此；因为，按宗徒的说法，\BibleQuote{那怀疑的，若吃了，就被定罪，因为不是出于信德；凡不出于信德的，就是罪}。\BibleRef{罗 14:23} 他是在信德中吃，因为他相信所吃的没有在偶像的庙中祭过，不是勒死的，也不是血；但他不是在信德中吃，若他对这些事中任何一件有怀疑。而那人知道它们已祭过魔鬼，却仍然使用，就成了魔鬼的同伙，同时他的想象也被魔鬼参与祭献所玷污。然而，宗徒知道肉食的本性既不会对使用它的人造成损害，也不会对戒用的人带来益处，而是见解和其中的理性，便说：\BibleQuote{其实食物不能使我们更接近天主，因为我们不吃也无损，吃也无益}。\BibleRef{格前 8:8} 而且，他知道那些对什么是洁净、什么是不洁有更高理解的人，放弃了关于使用洁净与不洁之物的区别（我想，还放弃了事物有别这迷信），从而对肉食的使用变得漠不关心，并因此被犹太人谴责为违犯法律，所以他曾在某处说：\BibleQuote{所以，不要让任何人在饮食上判断你们}等等，\BibleRef{哥 2:16} 教导我们：按文字的事是影子，但法律中蕴藏的真义才是未来美好的事物，在其中我们可以找到什么是灵魂的洁净属灵肉食，什么是在虚假和矛盾的话语中伤害受滋养之人的不洁食物，\BibleQuote{因为法律带有未来美好事物的影子}。\BibleRef{希 10:1}\par

% structure: p0026 source=h2 target=subsection reason=relative-heading-rank; base=h2

\subsection{十三、法利塞人的冒犯。}

正如在许多情况下，我们必须思考犹太人对救主的话感到惊奇，因为祂讲话带有权威，同样也适用于此处的话。于是祂叫了群众来，对他们说：\BibleQuote{你们听，且要明白！}\BibleRef{玛 15:10}等等。祂说这话，是因为法利塞人对这话起了反感，因为他们怀着恶念和对法律的曲解，他们不是祂在天上的圣父所栽种的，因此正被连根拔除；\BibleRef{玛 15:13}因为他们没有接受真葡萄树而被连根拔除——就是那由圣父所栽培的耶稣基督。\BibleRef{若 15:1}因为怎能算是祂圣父的栽种呢？那些人对耶稣的话起了反感，而耶稣的话使人远离那\BibleQuote{不可拿、不可尝、不可摸——这一切都是注定要随着使用而消亡的——随从人的诫命和教理}\BibleRef{哥 2:21}的规条，却促使有智慧的听者寻求那些上面的事，而不是像犹太人那样寻求地上的事？\BibleRef{哥 3:2}并且，因为法利塞人由于恶念而不是祂在天上的圣父所栽种的，因此，如同对待那些不可救药的人，祂对门徒说：\BibleQuote{由他们去吧！}\BibleRef{玛 15:14}\BibleQuote{由他们去吧！}祂这样说，是因为既然他们是瞎子，就应当意识到自己的瞎眼而寻求向导；但他们对自己的瞎眼毫无察觉，却自称为瞎子的向导，没有料到自己会掉进坑里，关于此，在\BibleBook{}{圣咏}中写着：\BibleQuote{他挖了坑，又掘了穴，却会掉进自己所掘的坑里。}此外，别处写着：\BibleQuote{祂一见群众，就上了山，坐下后，祂的门徒来到祂跟前；}\BibleRef{玛 5:1}但在这里，祂却伸手向群众，把他们叫到跟前，并将他们的思想从对法律问题的字面解释转向，祂首先对那些尚未明白所听见的事的群众说：\BibleQuote{你们听，且要明白！}然后像讲比喻那样对他们说：\BibleQuote{不是入于口的使人污秽，而是出于口的才使人污秽。}\BibleRef{玛 15:10}\par

% structure: p0028 source=h2 target=subsection reason=relative-heading-rank; base=h2

\subsection{十四、为什么法利塞人不是天主所栽种的。奥利振关于\Quote{主的饼}的教导。}

\Quote{此后，值得考察一下这个说法，它被那些声称法律的天主与耶稣基督福音的天主并非同一位的人以诡辩的方式攻击；因为他们说，耶稣基督的圣父并非那些自以为按照梅瑟法律崇拜天主的人的栽培者。耶稣亲自说过，那些崇拜创造世界和法律的天主的法利塞人，并非祂的圣父所栽种的植物，因此它将被连根拔起。\BibleRef{玛 15:13} 但你也许还可以这样说：即使那\Quote{带领并栽种了百姓，}当他们出埃及后，\Quote{到祂自己产业的山上，到祂为自己预备居住的地方}的，是耶稣的圣父，\BibleRef{出 15:17} 然而耶稣关于法利塞人会这样说：\BibleQuote{凡不是我圣父所栽种的植物，必要被连根拔起。} 对此，我们将说：那些因对法律之事的曲解而并非祂天上的圣父所栽种的人，他们的心思被蒙蔽了，因为他们不信真理，反而喜爱不义，\BibleRef{得后 2:12} 被那被此世之子神化、因此被保禄称为\Quote{此世之神}的人所蒙蔽。\BibleRef{格后 4:4} 不要以为保禄说他真是天主；因为正如肚腹，虽然它不是那些过于看重享乐之人的神——他们是爱享乐胜于爱天主的人——却被保禄称为他们的神，\BibleRef{斐 3:19} 同样，此世的元首，救主论及他说：\BibleQuote{如今这此世的元首已受了审判，}\BibleRef{若 16:11} 尽管他不是神，却被说成是那些不愿领受义子的心神的人的神，好使他们能成为此世之子及从死者中复活之子，并因此停留在此世之子身份中。因着那句\BibleQuote{他们是盲人给盲人领路，}\BibleRef{玛 15:14} 我虽以插话的方式说了这些，也认为是必要的。谁是这样的人？法利塞人，此世之神蒙蔽了他们不信的心思，因为他们不信耶稣基督；他蒙蔽了他们，以致\BibleQuote{天主荣耀的福音之光在基督的面容上不得照耀他们。}\BibleRef{格后 4:4} 但我们不但要避免被那些盲人引导——他们自觉需要向导，因为他们尚未领受自身的视觉能力——即使对所有声称以纯正道理指导我们的人，我们也必须谨慎听从，并对所说的话运用正确的判断，免得我们被那些盲人（他们看不见纯正道理之事）的无知所引导，自己也被视为盲人，因为我们看不见圣经的意义，以致引导者和被引导者都掉进我们先前说过的坑里。紧接着，经上记载伯多禄如何回答并问救主，仿佛他未明白\BibleQuote{不是入口的使人污秽，而是出口的}这句话：\BibleQuote{请给我们解释这个比喻。}\BibleRef{玛 15:11} 救主说：\BibleQuote{连你们也还是不明白吗？}\BibleRef{玛 15:16} 仿佛祂说：\BibleQuote{你们与我相处这么久，还不明白所说的话的意思吗？你们不明白因此入口的不能使人污秽，因为它进入肚腹，又排泄出去落入茅厕？}\BibleRef{玛 15:17} 法利塞人不是耶稣的圣父的植物，不是就他们表面上相信的法律而言，而是就他们对法律及其中所写之事的曲解而言。因为关于法律，有两件事需要理解：刻在石板上、与圣神无关的\OriginalTerm{死的职务}{the ministration of death}，和在属神法律中所理解的\OriginalTerm{生的职务}{the ministration of life}。那些能以诚实之心说\BibleQuote{我们知道法律是属神的，}\BibleRef{罗 7:14} 并因此\BibleQuote{法律是圣的，诫命也是圣的、公义的、善的}\BibleRef{罗 7:12} 的人，乃是圣父所栽种的植物；但那些并非如此，反而谨慎持守仅能杀死的字句的人，不是天主的植物，而是那使他们的心刚硬、并给他们盖上帕子者的植物，这帕子在他们未转向主时一直辖制他们；\BibleQuote{因为谁若转向主，帕子就除去了，主就是那神。}\BibleRef{格后 3:16} 现在，有人处理这段经文时可能会说，正如\BibleQuote{不是入口的使人污秽，}\BibleRef{玛 15:11} 即使犹太人认为不洁的事，同样，也不是入口的使人成圣，即使被称为主的饼的物，在较简单的门徒看来可能使人成圣。我认为这话不可轻视，因此需要清楚的解释，在我看来是这样的：正如不是食物，而是那存疑而吃者的良心污秽了他，因为\BibleQuote{那存疑的若吃了，就被定罪，因为他不是出于信德而吃，}\BibleRef{罗 14:23} 并且对于污秽不信的人，万物本身并非不洁，而是因他的污秽和不信而不洁。同样，那藉天主的话和祈祷而祝圣的物，就其本性而言并不使使用者成圣，因为若是如此，它也会使不配地吃主的饼的人成圣，并且没有人会因这食物而软弱、生病或死亡，因为保禄说了某种类似的事：\BibleQuote{因此，你们中间有许多人软弱、患病，也不少死了的。}\BibleRef{格前 11:30} 因此，关于主的饼，当人以无玷的心思和纯洁的良心领受时，对使用者有益。所以，既不因不吃——我指的是不吃那藉天主的话和祈祷祝圣的饼——我们就丧失任何好处，也不因吃就增进任何好处；因为我们缺乏的原因是邪恶和罪恶，我们富足的原因是公义和善行；所以保禄所说的话\BibleQuote{其实食物不能使我们悦纳天主，因为我们不吃也无损，吃也无益}\BibleRef{格前 8:8} 就是这个意思。现在，如果\BibleQuote{凡入口的，都进入肚腹，又排泄出去落入茅厕}\BibleRef{玛 15:17}，即使是那藉天主的话和祈祷祝圣的食物，就其物质本性而言，也进入肚腹，又排泄出去落入茅厕；但就临到它上面的祈祷而言，按照信德的比例，它成为益处，并为那注视有益之事的头脑提供清晰的视觉；使之有益的并非饼的物质，而是那祝谢的话，给那不配地吃主饼的人带来益处。以上所说的是关于\OriginalTerm{预像和象征的身体}{the typical and symbolic body}。但关于那成了血肉的圣言本身，\BibleRef{若 1:14} 以及真正的食粮，凡吃它的必永远活着，没有不配的人能吃它；因为如果一个持续不配的人能吃那成了血肉的圣言——祂是圣言和生命之粮——就不会写\BibleQuote{凡吃这食粮的，必永远活着}\BibleRef{若 6:51}。}\par

% structure: p0030 source=h2 target=subsection reason=relative-heading-rank; base=h2

\subsection{十五、以未洗的心进食污秽人。}

接下来，让我们看看那些从人里面出来并污秽人的事物，如何不是因为它们从口里出来而污秽人，而是因为它们在心中有污秽的根源，当恶念从心里出来（先于那些从口里出来的事物）时，这些恶念的种类有：谋杀、奸淫、行淫、偷窃、假见证、毁谤。\BibleRef{玛 15:18} 因为当这些从心里出来，并离开心而从口中发出时，正是这些污秽了人。这样，如果它们不从心里出来，而是被留在心中的某个地方，不允许从口中说出来，它们就会很快消失，而人就不再被污秽。因此，一切罪的泉源和源头都是恶念；因为除非这些恶念得胜，否则就没有谋杀、奸淫或任何其他此类事情。因此，各人必须全神贯注地守护自己的心；\BibleRef{箴 4:23} 因为当主在审判之日来临时，\BibleQuote{祂将揭露暗中的隐密并显明心中的计谋，} \BibleRef{格前 4:5} \BibleQuote{人的思想此时互相控告或辩解，} \BibleRef{罗 2:15} \BibleQuote{当他们的计谋环绕他们时。} \BibleRef{欧 7:2} 但恶念的性质是这样的：有时它们使那些看起来良善的事——就大众的判断而言值得称赞的事——也变得该受责备。因此，我们若在众人面前施舍，心里存着要显得爱人并因爱心受尊荣的意图，我们便从人那里得到了赏报；\BibleRef{玛 6:1} 并且，普遍来说，凡行事者意识到自己会被人荣耀的事，就没有从那在暗中察看并暗中报答纯洁之人的那一位那里得到赏报。所以，表面的纯洁也是如此，如果它受到虚荣或贪利的考虑的影响；那被认为是教会教导的教义，如果它通过谄媚的言语成为奴性的，无论是作为贪婪的借口，还是有人因自己的教导而寻求人的荣耀，就不被认为是那\BibleQuote{被天主安置在教会中的：首先是宗徒，其次是先知，第三是教师。} \BibleRef{格前 12:28} 的教导。对于那为了人的荣耀、人的谄媚、或因那些因虔诚之名而皈依道的人所给予的收益而寻求主教职位的人，你也会这样说；因为这样的主教根本不\BibleQuote{渴望善工，} \BibleRef{弟前 3:1} 也不能无可指摘，不能节制，不能清醒，因为他被荣耀所陶醉并毫无节制地沉溺其中。对于司铎和执事，你也会同样说。如果我们似乎在某些人看来在这些谈论中离题了，请思考一下是否必须说这些话，因为恶念是一切罪的泉源，并且能玷污那些如果没有恶念就会使行义之人成义的行为。我们已按自己的能力探究了那些污秽人的事物；但用未洗的手吃饭并不污秽人；但若我们必须大胆地说，用未洗的心吃任何我们理性的自然食物，就污秽人。\par

% structure: p0032 source=h2 target=subsection reason=relative-heading-rank; base=h2

\subsection{十六、关于客纳罕妇人。\Quote{提洛和漆冬的边境}的意义。}

\BibleQuote{\Emph{耶稣从那里出去，退到提洛和漆冬一带。看，有一个客纳罕妇人。}} \BibleRef{玛 15:21} \Quote{那里}指哪里？是从\Place{革乃撒勒}地方吗？关于这事，先前曾记载说：\BibleQuote{他们渡过了海，来到革乃撒勒地方？} \BibleRef{玛 14:34} 但祂退却，也许是因为\Person{法利塞人}听见了\BibleQuote{不是入口的，而是出口的，才污秽人} \BibleRef{玛 15:11} 而觉得反感；而且，因为他们有谋害祂的嫌疑，经上记载说：\BibleQuote{祂退却了} 这从这段经文可以明显看出：\BibleQuote{当祂听说\Person{若翰}被交付了，祂就退到\Place{加里肋亚}去了。} \BibleRef{玛 4:12} 也许也因为这个缘故，当描述此地之事时，\Person{马尔谷}说：\BibleQuote{祂起身往提洛边界去了，进了屋子，不愿意任何人知道。} \BibleRef{谷 7:24} 很可能祂是要避开那些因祂的教导而反感的法利塞人，等待祂受苦的时机，那是更合宜且正确指定的。但有人可能说，\Place{提洛}和\Place{漆冬}用来指代外邦人；因此当祂从\Place{以色列}退却时，祂来到外邦人的地方。在希伯来人中，提洛被称为Sor，解释为\Quote{痛苦}。漆冬，也是希伯来名称，解释为\Quote{猎人}。在外邦人中，那猎人同样也是恶势力，在他们中间有大苦难，即存在于邪恶与情欲中的苦难。所以，当\Person{耶稣}从革乃撒勒出去时，祂确实从以色列退却，来到提洛和漆冬的\Quote{地方}，而不是提洛和漆冬本身，结果那些外邦人现在部分相信了；因此，如果祂访问了整个提洛和漆冬，其中就不会留下任何不信者了。现在，根据马尔谷的记载：\BibleQuote{耶稣起身，往提洛边界去了。} \BibleRef{谷 7:24}——即外邦人的苦难——为使他们也得以从这些边界中出来，相信而得救；因为请注意：\BibleQuote{看，有一个客纳罕妇人从那些边界出来，喊说：主啊，\Person{达味}之子，可怜我吧！我的女儿被魔鬼纠缠得甚苦。} \BibleRef{玛 15:22} 我认为，如果她没有从那些边界出来，她就不能以那被见证的大信德向耶稣呼求；按照信德的比例，人从外邦人的边界出来，这些边界是\BibleQuote{当\Person{至高者}划分万民时，按以色列子孙的数目设立} \BibleRef{申 32:8}，并阻止了他们进一步前进的。这里，某些边界被称为提洛和漆冬的边界，但在\BibleBook{出谷}{纪}中，有\Person{法郎}的边界，\BibleRef{出 8:2}，他们说，那里形成了对\Place{埃及}人的灾祸。我们必须认为，我们每个人犯罪时，都是在提洛或漆冬或法郎和埃及的边界，或者是在天主所分配产业之外的某个边界；但当他从邪恶转变为美德时，他就从邪恶的边界出来，来到天主分地的边界，其中也有区别，这将对那些能理解有关以色列分地和产业之事的人显明，与属灵法律一致。并且也注意，可以这么说，耶稣与客纳罕妇人的相遇；因为祂来到提洛和漆冬的地方，而她从那些地方出来，喊说：\BibleQuote{主啊，达味之子，可怜我吧。} \BibleRef{玛 15:22} 那妇人是客纳罕人，这解释为预备受辱。义人确实预备进入天国和在天主国内被高举；但罪人则预备承受他们内在邪恶的羞辱，以及由此产生且为此预备他们的行为，还有在他们必朽身体中作王的罪的羞辱。只是，客纳罕妇人从那些边界出来，离开了预备受辱的状态，喊说：\BibleQuote{主啊，达味之子，可怜我吧。}\par

% structure: p0034 source=h2 target=subsection reason=relative-heading-rank; base=h2

\subsection{十七、叙述细节的阐明。}

现在请将福音中称呼祂为达味之子的人聚集起来，如她（即客纳罕妇人）和耶里哥的瞎子们\BibleRef{玛 20:30}；以及称呼祂为天主子的人，且不加\Quote{确实}一词，如附魔者所说：\Quote{我们与祢有什么相干？天主子。}\BibleRef{玛 8:29}；以及加上\Quote{确实}一词称呼祂的人，如船上朝拜祂的人说：\Quote{祢真是天主子。}\BibleRef{玛 14:33}因为将这些经文汇集一处，我认为将有助于你看出那些前来（就耶稣者）的差异：有些人前来，视祂为\Quote{按血统是达味后裔所生的}；\BibleRef{罗 1:3}另有些人前来就祂，视祂为\Quote{按圣洁的神性，以大能宣告为天主子}；\BibleRef{罗 1:4}而这些人中，有的加上\Quote{确实}，有的不加。此外，请注意，客纳罕妇人为一个女儿恳求祂，她似乎根本没有生过儿子，那女儿被魔鬼苦苦折磨；但另一位母亲却将已抬出去死的儿子活着接了回来。\BibleRef{路 7:12}又会堂长为一位十二岁的女儿祈求，以为她已死了，\BibleRef{玛 9:18}但那位贵族为儿子祈求，以为他还病着，且濒临死亡。\BibleRef{若 4:46}因此，被魔鬼折磨的女儿和已死的儿子来自两位母亲；而死了的女儿和病危的儿子来自两位父亲，其中一位是会堂长，另一位是贵族。我深信这些事含有关于耶稣所赋予生命和医治的不同种类灵魂的道理。祂在人群中所行的所有治愈，尤其是福音作者所记载的那些，在当时发生，为叫那些若不看见神迹奇事就不相信的人可以相信；\BibleRef{若 4:48}因为先前的事是那些藉耶稣的能力而始终成就之事的预象；因没有一刻，每一件所记载的事不是凭耶稣的能力按各人的功过而行。因此，客纳罕妇人因她的种族，连从耶稣那里得到回答也不配，耶稣承认祂被圣父派遣，只为了以色列家迷失的羊\BibleRef{玛 15:24}——一群具有清晰视野的迷失灵魂；但因她的决心，且因她朝拜耶稣为天主子，她得到了回答，这回答责备她出身卑微，并显明她配得的程度，即她配得碎屑如小狗，却不配得饼。但她以更坚定的决心，接受耶稣的话，提出要求，甚至如小狗般取得碎屑，并承认主人们是更高贵的种族，于是她得到第二次回答，这回答证明她的信德是大的，并应许照她所愿的给她成就。\BibleRef{玛 15:28}我认为，与\Quote{那在上自由的耶路撒冷，它是保禄和与他相似之人的母亲}相应\BibleRef{迦 4:26}，我们必须将客纳罕妇人、那位被魔鬼苦苦折磨之女的母亲设想为这种灵魂之母亲的预象。并考虑，是否合乎健全理智，也有许多父亲和母亲，与圣祖亚巴郎所往归的先祖们\BibleRef{创 15:15}以及保禄论及自己和与自己相似之人时所说的\Quote{母亲}耶路撒冷相对应。很可能，客纳罕妇人作为其预象的那位灵魂，从提洛和漆冬的边界（地上的地方是这些的预象）出来，来到救主面前，恳求祂，甚至如今仍恳求祂说：\Quote{主，达味之子，可怜我吧！我的女儿被魔鬼苦苦折磨。}\BibleRef{玛 15:22}然后，祂也对外人和门徒们，在必要时回答说：\Quote{我未被差遣，}\BibleRef{玛 15:24}教导我们，有一些迷失的灵魂，格外富有理智且视野清晰，比喻称为以色列家的羊；这些事，我想，那些认为这些话系指按血统的以色列而言的头脑简单之人，将不得不承认，即我们的救主由圣父所派遣，只为了那些迷失的犹太人，而非其他。但我们，能诚实地夸耀说：\Quote{即使我们曾按血统认识基督，如今却不再这样认识祂了}，\BibleRef{格后 5:16}确信拯救更具智慧者乃是圣言首要的工作，因为这些人比那些较迟钝者与祂更相似。但由于以色列家迷失的羊，除了\Quote{按照恩宠所选的遗民}外\BibleRef{罗 11:5}，都不信从圣言，因此\Quote{天主拣选了世上愚拙的}\BibleRef{格前 1:27}，即那不是以色列、也无清晰视野的，为叫以色列有智慧的羞愧；祂又召选了\Quote{那不存在的}\BibleRef{格前 1:28}，将一个能接纳\Quote{宣讲的愚妄}的智慧民族交给他们\BibleRef{格前 1:21}，并凭祂的美意拯救了那些相信这愚妄的人，为要驳倒\Quote{那存在的}，当它们与真理为敌时，祂从\Quote{婴孩和吃奶者口中}为自己完成了赞美。如今，客纳罕妇人前来，朝拜耶稣为天主，说：\Quote{主，帮助我。}祂却回答说：\Quote{拿孩子的饼扔给小狗是不可以的。}\BibleRef{玛 15:25}但有人或许也会探询这句话的含义，因为——既然饼的一个量度如此，使家中的孩子和狗都不能吃饼，除非狗吃的是另一种饼，而非那些精制的饼——那么按照正确的理智，将孩子精制的饼作为食物给小狗是不可能的。但在耶稣的能力方面，却不见这样的事，因为孩子们和那些称为小狗的，都能分享这能力。那么，考虑是否或许关于那句话：\Quote{拿孩子的饼是不可能的}，我们应当说：\Quote{祂空虚自己，取了奴仆的形状}，\BibleRef{斐 2:7}带来了一个世界所能接受的能力的量度，祂也知道这能力中有一定量从祂身上出去，正如从这话所表明的：\Quote{有人摸了我，因为我感觉有能力从我身上出去了。}\BibleRef{路 8:46}因此，祂从这个能力的量度中施与，将较大的一份给予那些卓越的、称为儿子的人，将较小的一份给予那些不如此的人，如同给小狗。但尽管事已如此，然而哪里有大信德，祂就赐给那因在客纳罕地出身卑微而为小狗的妇人，如同赐给一个孩子那样，孩子的饼。或许，在耶稣的话中，也有一些饼可以给那些较有理性的人，如同只给孩子们；而另一些话，则正如从贵人和主人们的大家和桌子上掉下的碎屑，可为一些灵魂所用，如同狗一样。按照梅瑟法律，论及某些事物写着：\Quote{你们要丢给狗吃。}\BibleRef{出 22:31}并且圣神关心，对某些食物给予指示，应留给狗。那么，让那些不熟悉教会道理的人假定，灵魂按照邪恶的不同程度，从人的身体转到狗的身体；但我们，在圣经中根本找不到这一点，认为较为理性的状态转化为较为非理性的状态，是因极大的懒惰和疏忽而经历这种变化。但同样，一个较为非理性的意志，因忽略理性，有时转变为理性，以致那曾是一条喜欢从主人桌子上吃碎屑的狗，进入了儿子的状态。因为美德大大有助于使人成为天主的儿子，但邪恶、在放荡言论和无耻中的疯狂怒气，则有助于使人被称为狗，正如圣经的话所说。\BibleRef{撒下 16:9}在用于无理性动物的其他名称上，你也会明白类似的事。唯有那被斥为狗，却不因被称为不配得孩子的饼而愤慨，并以全部宽容重复那客纳罕妇人之言的人说：\Quote{主，可是小狗也吃主人桌子上掉下来的碎屑。}\BibleRef{玛 15:27}他必将得到耶稣极其温和的回答，对他说：\Quote{你的信德真大！}——当他已接受了如此大的信德——并说：\Quote{照你所要的，给你成就。}\BibleRef{玛 15:28}以致他自己可得痊愈，并且若他产生了任何需要医治的果实，这果实也得医治。\par

% structure: p0036 source=h2 target=subsection reason=relative-heading-rank; base=h2

\subsection{十八、论被治愈的群众。\Person{耶稣}所坐的山与教会的比较。}

\BibleQuote{\Emph{耶稣从那里离开}}——显然，从上文可知，是从\Place{提洛}和\Place{漆冬}一带——\BibleQuote{来到\Place{加里肋亚海}边}，\BibleRef{玛 15:29}即通常称为\Place{革乃撒勒湖}的地方，又上山，在那里坐下。我们可以说，\Person{耶稣}所坐的这座山，不仅是健康的人上去，连同那些患有各种疾病的人也一同上去。也许，\Person{耶稣}上去并坐下的这座山，就是通常被称为教会的山，它借着\OriginalTerm{天主圣言}{the word of God}建立在世界及世上众人之上；上山的不仅是门徒（如同在真福八端时那样，他们留下群众），还有大批群众——这些人自己并非聋子或患有任何病痛，但他们带着这样的人同来。因为你可以看到，与来到这山上——天主圣子坐于其上——的群众一起，有些人对所应许的事成了聋子，有些人灵魂瞎眼，不看真光，有些人瘸腿，不能按理性行走，还有些人残废，不能按理性工作。那些灵魂遭受这些痛苦的人，虽然与群众一同上山到\Person{耶稣}所在之处，但若他们处于\Person{耶稣}的脚之外，便不被祂治愈；然而，当他们因患有这类疾病而被群众抛到\Person{耶稣}脚前\BibleRef{玛 15:30}——即基督身体的末端——时，他们自己本不配获得这些，却被祂治愈了。当你在通常称为教会的集会中，看见慕道者被抛在那些处于教会最末端的人后面，如同在\Person{耶稣}身体——教会——的脚前，带着自己的聋、瞎、瘸、跛来到教会，并按时按\OriginalTerm{圣言}{the Word}得到医治时，你若不犯错，就会说这样的人是与教会的群众一同上山到\Person{耶稣}所在之处，被抛在祂脚前而得到医治；以致教会群众惊讶于目睹从如此大恶转变为更好的变化，从而可以说：那些先前哑巴的后来讲论天主圣言，瘸子行走，\Person{依撒意亚}的预言得以应验，不仅是在身体上，也是在属灵上——这预言说：\BibleQuote{那时瘸子必跳跃如鹿，哑巴的舌头必歌唱。}\BibleRef{依 35:6}除非\BibleQuote{瘸子必跳跃如鹿}这一表述是偶然的，否则我们将说，那些先前瘸腿、现在借着\Person{耶稣}的能力跳跃如鹿的人，并非无目的地被比作鹿，因为鹿是洁净的动物，与蛇为敌，丝毫不受蛇毒伤害。而且，就哑巴被看见说话一事，预言也应验了，那预言说：\BibleQuote{哑巴的舌头必歌唱}，或者更确切地说，那预言说：\BibleQuote{聋子啊，你们要听；瞎子啊，你们要睁眼看，以便看见。}\BibleRef{依 42:18}如今，瞎子看见，是指他们看见世界，并从受造物的极大之美中，体悟到造物主在伟大和美善上与它们相称；并且他们清晰地看见\BibleQuote{天主那看不见的事，从创世以来，借着所造之物就可以明了}\BibleRef{罗 1:20}；也就是说，他们仔细而清晰地看见并理解。群众看见这些事，就光荣了以色列的天主\BibleRef{玛 15:31}，并且坚信那医治前述之人的父与以色列的天主是同一位天主。因为祂不但是犹太人的天主，也是外邦人的天主。\BibleRef{罗 3:29}那么，让我们带着那些愿意随同上山到\Person{耶稣}所坐之处——祂的教会——的人，将聋子、瞎子、瘸子、残废者以及许多其他人一同带上，并将他们抛到\Person{耶稣}脚前，让祂治愈他们，以致群众因他们的治愈而惊讶；因为经上并未记载门徒对这些事感到惊奇，尽管那时他们与\Person{耶稣}同在，这从以下的话可以显明：\BibleQuote{\Person{耶稣}叫门徒来，说：\InnerQuote{我怜悯这批群众}}\BibleRef{玛 15:32}等等；也许，你若仔细留意\BibleQuote{有许多群众来到祂跟前}\BibleRef{玛 15:30}这句话，会发现那时门徒并没有来到祂跟前，而是早已开始跟随祂，并跟随祂上了山。但来到祂跟前的，是那些不及门徒的人，他们那时是初次接近祂，没有经历过与那些一同上山的人相同的事。此外，你要留意福音中哪些人被记载为跟随\Person{耶稣}，哪些人来到祂跟前，哪些人被带到祂跟前，以及那些在前和那些在后的区分；还有那些来到祂跟前的人，哪些是在屋里来到祂跟前，哪些是在别处。因为借着观察和以属灵的事比较属灵的事，你会发现许多与福音中精准智慧相称的事。\par

% structure: p0038 source=h2 target=subsection reason=relative-heading-rank; base=h2

\subsection{十九、论七个饼。四千人吃饱的叙述与五千人的比较。}

\BibleQuote{\Emph{耶稣便叫了自己的门徒来，说。}} \BibleRef{玛 15:32} 上面有关饼的类似记载中，在提及饼之前，就有：\BibleQuote{耶稣出来，看见一大群人，就怜悯了他们，治好了他们的病人。到了晚上，门徒来到他跟前说：\InnerQuote{这地方是荒野，时候已过去了，请打发他们走吧。}} \BibleRef{玛 14:15} 等等。但现在，在治好聋子等人之后，他对那已经跟随他三天且没有东西吃的群众动了怜悯之心。在那里，门徒为五千人求情；\BibleRef{玛 14:15} 但这里，他主动提起四千人。\BibleRef{玛 15:32} 那些人也是在傍晚与他共度一天后得喂养；但这些人，被证实已跟随他三天，吃饼免得在路上昏倒。在那里，门徒在耶稣未询问时对他说，他们只有五个饼和两条鱼；但这里，他们在耶稣询问时回答了关于七个饼和几条小鱼的情况。在那里，他命令群众坐下或躺在草地上；因为路加也写道：\BibleQuote{叫他们坐下。}\BibleRef{路 9:14} 马尔谷说：\BibleQuote{他吩咐众人坐下。}\BibleRef{谷 6:39} 但这里，他不是命令，而是宣告让群众坐下。此外，在那里，三位福音作者用同样的话说：\BibleQuote{他拿起五个饼和两条鱼，望天祝福了；} 但这里，如玛窦和马尔谷所记：\BibleQuote{耶稣感谢了，就掰开；} 在那里，他们斜靠在草地上，但这里，他们坐在地上。此外，你可以考察各处记载中若望所记的差异，他关于那件事写道：耶稣说：\BibleQuote{叫众人坐下。}\BibleRef{若 6:10} 并且他感谢后，把饼分给了坐下的人，但他完全没有提及这个奇迹。那么，注意到各处有关饼的记载的不同，我认为这些属于与那些不同的类别；因此这些是在山上得喂养，而那些是在荒野之地；这些是在跟随耶稣三天之后，而那些只是一天，在当天傍晚得喂养。再者，除非耶稣主动行事和听了门徒的话后行事是同一回事，那么请思考：耶稣即时喂养他们以示仁慈，这些蒙受仁慈的人是否并不更优越？而且，如果按若望所记，\BibleRef{若 6:13} 他们吃的是大麦饼，剩下十二筐，而关于这些饼并无此类说法，那么这些饼岂不优于前者？并且他治好了那些病人，\BibleRef{玛 14:14} 但这里，他和群众一同治好了这些不是病人而是瞎眼、瘸腿、聋子、残废的人；因此关于这些人，四千人惊奇，\BibleRef{玛 15:31} 但关于病人却没有这样说。我认为，这些吃那七饼并为此感恩的人，优于那些吃那五饼并被祝福的人；这些吃那几条小鱼的，优于那些吃那两条鱼的；或许这些坐在地上的，也优于那些坐在草地上的。那些从较少饼中剩下十二筐，但这些从较多饼中剩下七筐，因为他们能领受更多。也许这些践踏一切属地之物并坐在其上，而那些坐在草地上——只坐在肉体上——因为\BibleQuote{凡有血肉的，尽都如草。}\BibleRef{依 40:6} 此后也要思考，耶稣不愿他们饥饿离去，以免他们在路上昏倒，因为没有耶稣的饼，而他们还在路上——走向自己事务的路上——可能会受伤害。也要留意耶稣曾打发人离开的实例，好让你看到那些被他喂养后打发走的人与那些以其他方式被打发走的人之间的区别；至于以其他方式被打发走的例子，请取\Quote{妇人，你的病已脱离了。} 但进一步，经常与耶稣在一起的门徒并不被他打发走；但群众在吃过后则被打发走。同样，门徒们对客纳罕妇人没有大想法，便说：\BibleQuote{打发她走吧，因为她在我们后面喊叫。}\BibleRef{玛 15:23} 但救主似乎完全没有打发她走；因他对她说：\BibleQuote{妇人，你的信德真大！照你所愿的，给你成就吧！}\BibleRef{玛 15:28} 从那一刻起，她的女儿就好了：然而并未记载他打发她走。到目前为止，我们已能考察并探究眼前这段经文。\par

\SourceRecord{Translated by John Patrick.；Ante-Nicene Fathers；Vol. 9.；Edited by Allan Menzies.；Buffalo, NY: Christian Literature Publishing Co.,；1896.；<http://www.newadvent.org/fathers/101611.htm>.}

% structure: p0001 source=h1 target=section reason=page-title

\section{《玛窦福音注释（第十二卷）》}

% structure: p0002 source=h2 target=subsection reason=relative-heading-rank; base=h2

\subsection{一、论那些求祂从天上显一个征兆的人}

\Quote{\Emph{法利塞人和撒杜塞人前来，试探祂，求祂从天上显一个征兆给他们。}} \BibleRef{玛 16:1} 法利塞人和撒杜塞人在最重要的真理上彼此意见不合——因为法利塞人拥护死人复活的教义，盼望有来世；而撒杜塞人则认为人死后，无论他是在德行上有所进步，还是根本没有努力摆脱邪恶的群山，都一无所存——这些人，我说，竟因试探耶稣而达成一致。现在，正如路加所记述的，\BibleRef{路 23:12} 在黑落德和比拉多之间也发生了类似的事，他们彼此成为朋友，以便杀害耶稣；因为，也许他们彼此之间的敌意原本会阻止黑落德为了取悦那些喊着\Quote{钉祂十字架，钉祂十字架}\BibleRef{路 23:21} 的民众而要求处死祂，并会对比拉多产生影响——他先前有点倾向于不判祂有罪，而他对黑落德的敌意又使他先前想要释放耶稣的意向得到了新的刺激。但他们表面上的友谊使黑落德在比拉多面前对耶稣的要求更加强硬，而比拉多或许也因为新建立的友谊，想要做些什么来取悦黑落德和全体犹太民族。甚至在今天，你也常常在日常中看到那些在希腊哲学或其他思想体系中持有最相左意见的人，为了嘲笑并攻击在祂门徒身上的耶稣基督而显得同心一意。从这些事，我想你可以通过理性推论来思考：当彼此不和的势力联合起来反对时——例如法郎与拿步高、\BibleRef{列下 24:7} 雇士王提尔哈卡与散乃黑黎布\BibleRef{列下 19:9}——是否就会形成反对耶稣和祂子民的联盟。所以，也许也正如经上所记：\Quote{地上的君王起来，官长聚集商议}，\BibleRef{咏 2:2} 尽管他们以前并不和睦，却共同商议敌对主和祂的基督，以便杀害荣耀之主。\par

% structure: p0004 source=h2 target=subsection reason=relative-heading-rank; base=h2

\subsection{二、论法利塞人为何求从天上来的征兆}

如今，我们在论述中已到此点，因为法利塞人和撒杜塞人一同来到耶稣面前，他们在有关复活的事上意见分歧，但却仿佛达成一致，为试探我们的救主，并请求祂向他们显示一个来自天上的记号。原来，他们不满足于在百姓中显示的奇妙记号——即医治各种疾病和灾殃，以及我们的救主在许多人知晓下所行的其余奇迹，他们希望祂也向他们显示一个来自天上的记号。我推测，他们怀疑地上的记号可能不是出于天主；因为他们确实毫不迟疑地说：\Quote{耶稣靠魔王\OriginalTerm{贝尔则步}{Beelzebub}驱魔。}在他们看来，来自天上的记号不可能源自贝尔则步或任何邪恶势力。但他们在两方面都错了，无论是地上的记号还是天上的记号，因为他们不是\Quote{合格的兑换银钱者}，也不知道如何分辨那些运作的神灵，哪些是出于天主，哪些是背离了祂。他们本应知道，甚至在梅瑟时代为攻击埃及所行的许多异兆，虽不是来自天上，却明显出于天主；而自天降下烧毁约伯羊群的火，并非来自天主；\BibleRef{约 1:16}因为那火属于同一位，即那掳走并三次派兵攻击约伯牲畜者也属于他。此外，我认为在依撒意亚书中——似乎记号能同时从地和天上被显示，真记号出于天主，而\Quote{各种能力、奇迹和虚假的奇事}\BibleRef{得后 2:9}出于恶者的记号——就曾对阿哈次说：\Quote{你向主你的天主求一个记号，或在下深处，或在上高处。}\BibleRef{依 7:11}因为，除非在深处或高处有些记号不是出于主天主，否则就不会说：\Quote{你向主你的天主求一个记号，或在下深处，或在上高处。}但我深知，对经文\Quote{你向主你的天主求一个记号}的如此解释，对某些人而言将显得牵强；但请注意宗徒论及\OriginalTerm{罪恶之人}{man of sin}、\OriginalTerm{丧亡之子}{son of perdition}时所说的话：\Quote{带着各种能力、奇迹和虚假的奇事，以及各种不义的欺骗，}\BibleRef{得后 2:9}他要在那些丧亡的人中间显现，效法各种奇事，即真理的奇事。正如埃及的术士和巫师，虽次于那罪恶之人和丧亡之子，却模仿某些能力，即真理的记号与奇事，行虚假的奇事，以致真者不被相信；同样，我认为那罪恶之人也将模仿记号与能力。或许，法利塞人也因关于祂的预言而怀疑这些事；但我要问，是否撒杜塞人在试探祂时，也请求耶稣向他们显示来自天上的记号。因为，除非我们说他们怀疑这一点，我们怎能描述他们与耶稣所行异兆的关系？他们一直心硬，不为所行的奇事感到羞愧。但若有人以为，我们给了法利塞人和撒杜塞人辩护的借口，既在他们说耶稣靠贝尔则步驱魔时，也在他们试探祂而问及天上记号时，愿他知道，我们合情理地说，他们被引向那样的结局，为使他们不信耶稣的奇迹；但并非因此配得宽恕；因为他们没有留意先知的话——那些话正在耶稣的作为中应验，邪恶势力根本无法模仿。然而，使一个已离去的灵魂返回，以致那已经发臭、过了四天的尸体从坟墓中出来，\BibleRef{若 11:39}这工作唯有祂能行，祂曾听见圣父的话：\Quote{让我们按我们的肖像和模样造人。}\BibleRef{创 1:26}而且，命令风，并以一言平息海上的狂暴，这工作也唯有祂能行，万物——海洋本身和风——都是藉着祂而受造的。此外，至于那激发人对造物主之爱、与法律和先知相谐、抑制情欲并依虔诚塑造道德的教导，它向那些能看见的人指明了什么？岂非那行如此大能之事者确是天主之子吗？关于这些事，祂也向若翰的门徒说：\Quote{你们去，把你们所见所闻的大事告诉若翰：瞎子看见，}\BibleRef{玛 11:4}等等。\par

% structure: p0006 source=h2 target=subsection reason=relative-heading-rank; base=h2

\subsection{三、耶稣对他们请求的答复。}

接下来让我们注意，当法利塞人和撒杜塞人问祂，求祂从天上显一个征兆时，祂如何回答并说：\Quote{\Emph{一个邪恶淫乱的世代寻求征兆，除了约纳先知的征兆外，再没有征兆给他们。}} 当时祂也\Quote{离开他们走了。}\BibleRef{玛 16:4} 但事实上，按着他们的要求，约纳的征兆不仅是一个征兆，而且是从天上来的征兆；因此，即使对那些试探祂、求祂从天上显征兆的人，祂仍然出于自己的大善，给了他们这个征兆。因为正如约纳在鲸鱼腹中过了三天三夜，人子也在地心里过了三天三夜，并从那里复活——我们怎能不说，基督复活的征兆是从天上来的呢？尤其是在受难时，祂成了那获得祂恩宠得以进入天主乐园的强盗的征兆；此后，我认为，祂下到阴府，到了死者那里，\Quote{在死者中自由。} 在我看来，救主似乎把那要来自祂自己的征兆与关于约纳的征兆的理由联结起来，因为祂说，祂所赐的不仅是一个类似的征兆，而是那个征兆本身；因为请听这话：\Quote{再没有征兆给它，除了约纳先知的征兆。}\BibleRef{玛 16:4} 因此，那个征兆就是这征兆，因为那个成了这个的预表，以致那表面晦涩的征兆的阐明，可以在救主受苦并在地心里度过三天三夜这一事实中找到。同时我们也学到一项普遍原则：如果征兆象征什么，那么所记载的每一个征兆，无论是实际历史中的，还是以诫命形式的，都预示后来应验的某事；例如，约纳三天后从鲸鱼腹中出来的征兆，预示了我们的救主死后三天三夜复活；而所谓的割损，乃是保禄用这话所指明的记号：\Quote{我们才是受割损的。}\BibleRef{斐 3:3} 你们也当在旧约圣经中寻找每一个记号，作为新约圣经某段经文的预表，并且在新约中称为记号的，也预表着将要来临的世代，甚或在此记号发生后的后世中的某事。\par

% structure: p0008 source=h2 target=subsection reason=relative-heading-rank; base=h2

\subsection{四、为什么耶稣称他们为淫乱的世代。法律如同丈夫。}

祂果然称他们为\BibleQuote{邪恶的世代}，因为他们内在所产生的邪恶品质，因为邪恶是自愿的恶行；又称他们为\BibleQuote{奸淫的世代}，因为当法利塞人和撒杜塞人离弃那被比喻地称为丈夫的——那就是真理之言或法律，他们就被虚伪和罪恶的法律所败坏。因为如果有两种法律，即在我们的肢体中与思想的法律作战的法律，和思想的法律，\BibleRef{罗 7:23}我们必须说，思想的法律——即属灵的法律——是丈夫，天主将灵魂作为妻子赐给了它，即赐给了作为法律的丈夫，正如经上所记载的：\BibleQuote{妻子是由天主许配给丈夫的。}\BibleRef{箴 19:14}而另一种法律则是灵魂的奸夫，灵魂受制于它，也因此被称为奸妇。现在，保禄在\WorkTitle{罗马书}中清楚地表明，法律是灵魂的丈夫，他说：\BibleQuote{法律管辖人，只在人活着的时候；因为已婚的妇人，在法律上她被丈夫束缚，只要丈夫活着，她就被那作为丈夫的法律所束缚。}等等。请思考这些事：法律管辖人，只在那法律活着的时候——如同丈夫管辖妻子。因为\BibleQuote{那已婚的妇人}，即处于法律之下的灵魂，\BibleQuote{只要丈夫活着，就被束缚}于作为丈夫的法律；但如果丈夫——即法律——死了，她就从法律——她的丈夫——中获得释放。现在，法律对于那些上升至有福境界的人已经死了，他不再活在法律之下，而是如同基督一样行事，祂虽为法律之下的人成了法律之下的人，为要赢得那法律之下的人，\BibleRef{格前 9:10}却没有继续处于法律之下，也没有使那些被祂释放的人仍受法律束缚；因为祂带领他们与祂一同进入高于法律的属天公民身份，这身份对于不完善且仍是罪人的人，包含赎罪的祭献。因此，那无罪且不再需要法律祭献的人，也许当他成为完善时，甚至已经超越了属灵的法律，而到达超越它的圣言——祂为活在肉体内的人成了肉身，但对那些不再按肉体作战的人，祂被理解为圣言，正如祂在起初就是与天主同在的天主，并启示了圣父。因此，与这处境相关，有三件事要考虑：那有丈夫的妇人，即在丈夫——法律——之下的妇人；那成为奸妇的妇人，即灵魂，当她的丈夫法律活着时，她却与另一个丈夫，即肉体的法律结合；以及那嫁给已死丈夫的兄弟的妇人，即嫁给那活着且不死、\BibleQuote{从死者中复活，不再死，死亡再不能管辖祂}的圣言\BibleRef{罗 6:9}。到目前为止，因为这句话：\BibleQuote{但如果丈夫死了，她就从法律——丈夫——中被释放}，以及这句话：\BibleQuote{所以，如果丈夫活着，她若与别的男人结合，便被称为奸妇}，以及这句话：\BibleQuote{但如果丈夫死了，她就从法律中获得自由，因而她虽与别的男人结合，也不是奸妇。}\BibleRef{罗 7:2}我们引用了这句话：\BibleQuote{所以，只要丈夫活着，她就被称为奸妇}，想清楚地表明，为什么耶稣在回答试探祂、求祂从天上显个征兆的法利塞人和撒杜塞人时，不仅说\BibleQuote{一个邪恶的世代}，还说\BibleQuote{一个奸淫的世代}\BibleRef{玛 16:4}。一般而言，那在我们的肢体中与思想的法律作战的法律\BibleRef{罗 7:23}，作为一个犯奸淫的男人，是灵魂的奸夫。但现在，每一个敌对的势力，就是掌控人的灵魂并与它交合的，也与那曾被天主赐予新郎——即圣言——的灵魂犯奸淫。此后再经上记载：\BibleQuote{祂离开了他们，就去了。}因为新郎——圣言——怎能不离开那奸淫的世代而离去呢？但你可以说，天主的圣言离开犹太人的会堂这个奸淫者，离弃了它，并娶了一个娼妓为妻\BibleRef{欧 1:2}，即那些来自外邦的人；因为那些曾是以\BibleQuote{信任的城邑}\BibleRef{依 1:21}的人成了娼妓；但这些外邦人却成了像娼妓辣哈布一样的人，她接待了若苏厄的探子，并与她全家一同得救\BibleRef{苏 6:25}；此后她不再行淫，反而来到耶稣脚前，用忏悔的泪水沾湿了祂的脚，又用圣洁言谈的香膏涂抹了祂的脚；为了她的缘故，祂责备了患癞病的西满——即先前的人民——说了一些经上记载的话。\par

% structure: p0010 source=h2 target=subsection reason=relative-heading-rank; base=h2

\subsection{五、关于法利塞人的酵母。}

\BibleQuote{\Emph{门徒渡到对岸，忘了带饼。}}\BibleRef{玛 16:5} 门徒在渡到对岸之前所拥有的饼，在他们到了对岸之后就不再有用，因为他们渡海前需要一种饼，渡海后需要另一种饼——因此，他们去对岸时，因不在意带饼，就忘了带饼。于是，耶稣的门徒渡到了对岸，他们从物质的事过渡到了属神的事，从可感觉的事过渡到理智的事。或许，为了阻止那些藉着渡到对岸而\Quote{开始于神}的人退回属血肉的事，耶稣在对岸对他们说：\BibleQuote{你们应当谨慎，提防。}\BibleRef{玛 16:6} 因为有一种教训的团和真正古老的酵母——即按着字面的教训，因此未能摆脱出于邪恶的事——这是法利塞人和撒杜塞人所提供的，耶稣不愿自己的门徒再吃这酵母，祂为他们造了一个新的属神的团，将自己赐给那些放弃了法利塞人和撒杜塞人的酵母、来到祂跟前的人——\BibleQuote{从天上降下来、赐生命给世界的活粮。}\BibleRef{若 6:33} 但是，对于不再用法利塞人和撒杜塞人的酵母和团与教训的人来说，首先要\Quote{看}，然后要\Quote{提防}，这样就没有人因不看和不谨慎而沾染他们那被禁止的酵母——因此，祂先对门徒说\Quote{看}，然后说\Quote{提防}。明察而谨慎的人的特征是，将法利塞人和撒杜塞人的酵母以及一切不是\BibleQuote{真诚和真理的无酵饼}\BibleRef{格前 5:8}的食物与活粮——即从天上降下来的——分开，这样，凡吃的人就不致采纳法利塞人和撒杜塞人的东西，而是藉着吃活而真正的饼来坚固自己的灵魂。我们也很合时宜地将这句话应用于那些在基督徒生活方式之外，还愿意像犹太人一样物质地生活的人；因为这些人既不看也不提防法利塞人和撒杜塞人的酵母，反而违背耶稣禁止的旨意，吃法利塞人的饼。是的，还有所有不愿意理解法律是属神的、且是未来美好事物的影子\BibleRef{希 10:1}、是将来事物的影子\BibleRef{哥 2:17}的人，他们既不探究每条法律是何种将来好事的影子，也不看也不提防法利塞人的酵母；而那些拒绝死人复活教义的人，也没有防范撒杜塞人的酵母。在异端者中有许多人，由于他们不信死人复活，而沾染了撒杜塞人的酵母。耶稣说这些话时，门徒心中议论，并非出声，而是在心中说：\BibleQuote{我们没有带饼。}\BibleRef{玛 16:7} 他们大概这样说：如果我们有饼，就不必用法利塞人和撒杜塞人的酵母了；但因为我们没有饼，有沾染他们酵母的危险，救主不愿我们退回他们的教训，因此祂对我们说：\BibleQuote{你们应当谨慎，提防法利塞人和撒杜塞人的酵母。}\BibleRef{玛 16:6} 于是他们就这样议论；耶稣洞察他们心中所想，听见他们心里的理由，作为人心真正的监察者，祂责备他们，因为他们没有看、没有记住从祂那里领受的饼；因此，即使他们看似缺少饼，也不需要法利塞人和撒杜塞人的酵母。\par

% structure: p0012 source=h2 target=subsection reason=relative-heading-rank; base=h2

\subsection{六、酵母的寓意。耶稣对心的认识。}

然后，祂清楚地解释并向那些因\Quote{饼}和\Quote{酵}的双重含义而困惑的人毫无掩饰地说明，祂并非在谈论感官的饼，而是关于教导中的酵，祂接着说：\Quote{\Emph{你们怎么不明白我不是对你们讲关于饼的事呢？但你们要防备法利塞人和撒杜塞人的酵。}}\BibleRef{玛 16:11} 虽然祂没有揭开解释，但仍然继续使用比喻性语言，门徒们会理解救主的言论是关于教导——比喻性地称为酵——即法利塞人和撒杜塞人所教导的。那么，只要我们拥有耶稣与我们同在，履行那应许：\Quote{看，我同你们天天在一起，直到今世的终结。}\BibleRef{玛 28:20} 我们就不能禁食，也不会缺少食物，以至于因缺乏而想要取食并吞吃那被禁止的酵，即使是从法利塞人和撒杜塞人而来的。有时候，祂与我们同在时，我们可能没有食物，正如上面经文中提到的：\Quote{他们同我在一起已经三天，也没有什么可吃的。}\BibleRef{玛 15:32} 但是，即使如此，祂不愿遣散我们空腹，免得我们在路上昏倒，便祝谢了门徒们所有的饼，使我们从七个饼中有了七个筐子剩余，正如我们记载的。此外，针对那些认为救主的神性从玛窦福音中丝毫不能证明的人，也要注意这一点：当门徒们彼此议论说：\Quote{我们没有饼了}，耶稣知道他们的议论，便说：\Quote{小信德的人哪！你们为什么彼此议论没有饼呢？}\BibleRef{玛 16:8} 这超出人的能力；因为只有主，正如撒罗满在\WorkTitle{列王纪上}中所说，知道人的心。\BibleRef{列上 8:39} 既然门徒们明白了，当耶稣说：\Quote{你们要防备酵}\BibleRef{玛 16:6} 时，祂并非告诉他们要防备饼，而是要防备法利塞人和撒杜塞人的教导，你就会明白，每当提到\Quote{酵}，它都是比喻性地指代教导，无论是在法律中，还是在法律之后的圣经中；所以，酵也许不被献在祭台上；因为祈祷不应采取教导的形式，而只应是对天主所赐美善之事的恳求。但有人可能会因关于来到对岸的门徒们所说的话而追问：是否任何到达对岸的人仍会被责备为小信德、且尚未理解或记起耶稣所行的事。然而，我认为针对这一点不难回答说：相对于那完美的，当那完美的来到时，\Quote{那局部的就必消失}\BibleRef{格前 13:10}，我们在此的全部信德都是小信德，而相对于那完美的，我们这些局部认识的人既不知道也不记起；因为我们无法获得足够且能触及思辨本质之大的记忆。\par

% structure: p0014 source=h2 target=subsection reason=relative-heading-rank; base=h2

\subsection{七、内心之罪与实际行动之罪的相对大小。}

但我们也可以从这里学到：仅凭我们内心所起的议论，我们有时就被定罪和责备为小信德。我认为，正如一个人只在心里犯奸淫，虽然完全没有走向公开行动，同样，他也以心里犯下其余被禁止的事。那么，正如在心里犯奸淫的人将按此类奸淫受相应惩罚，同样，在心里犯下任何一件被禁止之事之人——例如，只在心里偷窃，或只在心里作假见证——将不会像实际上偷窃或完成假见证行为的人那样受罚，而只会像在心里做这些事的人那样受罚。还有一种情况：一个人没有达到邪恶行动，但尽管他愿意，却未能做到。因为如果除了愿意之外，他还企图去做，但未能实施，他将不会像仅在心里犯罪的人那样受罚，而是像在行为上犯罪的人那样受罚。对于这类问题，有人可能会问：是否有人在心里犯奸淫，即使他实际上没有犯奸淫的行为，而只是在心里缺乏自制？关于其他值得称赞的事，你也会这样说。但这段经文可能包含一个看似合理却错误的推理，我想必须这样澄清：在心里发生的奸淫，其罪小于若加上行为的情况。但不可能存在一种贞洁在心里阻止贞洁行为——除非有人举出法律中在野地里被强奸的童贞女\BibleRef{申 22:25}作为实例；因为可以承认，一个人的心可能非常纯洁，但淫乱中的暴力导致了她贞洁身体的败坏。事实上，在我看来，她在隐秘的心里是完全贞洁的，但身体不再像受暴之前那样纯洁；然而，虽然在外表上她不再纯洁，她因此也不再是贞洁的吗？我说这些是因为那句话：\Quote{他们彼此议论说：我们没有带饼}接着是\Quote{耶稣知道了，就说：小信德的人哪！你们为什么彼此议论……}\BibleRef{玛 16:7} 等等；因为有必要就隐秘之事的责备和相应地隐秘之事的赞美进行探究。\par

% structure: p0016 source=h2 target=subsection reason=relative-heading-rank; base=h2

\subsection{八、酵是比喻性的，如同耶稣向撒玛黎雅妇人所讲的水一样。}

但我感到惊奇，门徒们在耶稣向他们解释这话之前，是否认为他们的导师和主禁止他们提防法利塞人或撒杜塞人的有形酵母，视为不洁，并因此禁止，以免他们因没有带饼而使用那酵母。在其他事上我们也可以作类似的探究；但举例来说，关于撒玛黎雅妇人的叙述就足够了：\BibleQuote{凡喝这水的，还要再渴；但谁若喝了我所赐给的水，将永远不渴。}\BibleRef{若 4:13} 因为在那里，仅就表达形式而言，撒玛黎雅妇人似乎认为救主在应许关于物质的水，当祂说：\BibleQuote{谁若喝了我所赐给的水，将永远不渴。} 那些事必须比喻性地解释，我们必须考察和比较雅各伯泉的水（撒玛黎雅妇人从中取水）与耶稣的水；这里也必须同样做；因为也许饼没有被烘烤，而只是一种生酵母，即法利塞人和撒杜塞人的教导。\par

% structure: p0018 source=h2 target=subsection reason=relative-heading-rank; base=h2

\subsection{九、关于耶稣在凯撒勒雅的提问：\Quote{人们说我是谁？}不同观念中的耶稣。}

\BibleQuote{\Emph{现在耶稣来到斐理伯的凯撒勒雅境内，就问祂的门徒。}}\BibleRef{玛 16:13} 耶稣问门徒：\Quote{人们说我是谁？} 这是为了使我们从宗徒的回答中得知当时犹太人中关于我们救主的各种观念；也许也是为了让耶稣的门徒学会关心人们关于他们所说的话；因为这对做这事的人有益，如果有什么恶言，就断绝一切作恶的机会；如果有什么好话，就增加行善的动机。不过，请注意，由于犹太人中关于耶稣的舆论动荡，有些人受错误理论的影响，说祂是洗者若翰，如同分封侯黑落德对他的仆人所言：\BibleQuote{这是洗者若翰，他从死者中复活了，因此奇能在他身上运行；}\BibleRef{玛 14:2} 另有人说，那如今称为耶稣的，是厄里亚，或是再次出生，或是自那时起一直活在肉身中，如今显现。但那些说耶稣是耶肋米亚的人，而不是说耶肋米亚是基督的预像，也许是受耶肋米亚书开头关于基督的记载影响，那段话当时并未应验在那位先知身上，而是在耶稣身上开始应验，\BibleQuote{天主立祂于万民和万国之上，为拔除、为拆毁、为毁灭、为建立、为栽植，}\BibleRef{耶 1:10} 使祂成为向外邦人传话的先知，向他们宣讲圣言。此外，那些说\BibleQuote{祂是某一位先知}\BibleRef{玛 16:14} 的人，之所以有此想法，是因为先知们所说的那些话是针对他们的，但在他们身上并未应验。犹太人，由于心中蒙着帕子，也持有关于耶稣的错误观点；而伯多禄，并非\BibleQuote{血肉}\BibleRef{玛 16:17} 的门徒，而是堪受圣父在天启示的人，承认祂是基督。伯多禄对救主说：\Quote{祢是基督，}当犹太人不认识祂是基督时，这确实是大事；但更伟大的是，他不仅知道祂是基督，还知道祂是\BibleQuote{永生天主之子，}\BibleRef{玛 16:16} 祂也曾通过先知说：\BibleQuote{我生活着，}\BibleRef{耶 22:24} 以及\BibleQuote{他们离弃了我这活水的泉源；}\BibleRef{耶 2:13} 而且祂也是生命，正如从圣父而来，祂是生命的泉源，祂说：\BibleQuote{我是生命；}\BibleRef{若 14:6} 并且要仔细考虑：正如河的源头与河不是同一回事，生命的源头是否与生命不同？我们添加这些，是因为在\BibleQuote{祢是基督，天主子}这句话之后，附有\BibleQuote{永生}\BibleRef{玛 16:16} 一词；因为有必要阐明关于天主和万物之父被称作永生的一些值得注意的事，无论是关乎祂绝对的生命，还是关乎那些分享生命的事物。既然我们说，那些宣称耶稣是洗者若翰或上述任何一位的人受错误理论影响，那么让我们证明：如果他们曾遇见耶稣去若翰那里受洗，或遇见若翰为耶稣施洗，或从别人那里听说，他们就不会说耶稣是若翰。同样，如果他们理解了耶稣说\BibleQuote{如果你们愿意接受，这厄里亚就是要来的那位}\BibleRef{玛 11:14} 的语境，并像有耳朵的人一样去听，有些人就不会说祂是厄里亚。如果那些说祂是耶肋米亚的人意识到大多数先知都承担了象征祂的某些特征，他们就不会说祂是耶肋米亚；同样，其他人也不会说祂是某一位先知。\par

% structure: p0020 source=h2 target=subsection reason=relative-heading-rank; base=h2

\subsection{十、伯多禄的回答。}

也许，西满伯多禄所回答并说的：\BibleQuote{\Emph{你是默西亚，永生天主之子}}\BibleRef{玛 16:16}，如果我们像伯多禄那样说这话，不是因血肉启示了我们，而是因来自天上圣父的光照在我们心中，我们也成为伯多禄，被宣称为有福的，如同他一样；因为使他被称为有福的理由也适用于我们：血肉没有启示我们关于耶稣说祂是默西亚、永生天主之子，而是天上的圣父，从诸天启示我们，使我们的家乡在天上，\BibleRef{斐 3:20}启示那提升到天上的人的启示，他们除去心中每一层帕子，并领受\BibleQuote{智慧和启示的神}\BibleRef{弗 1:17}。如果我们像伯多禄那样说：\BibleQuote{你是默西亚，永生天主之子}，不是因血肉启示了我们，而是因来自天上圣父的光照在我们心中，我们就成为伯多禄，并且圣言会对我们说：\BibleQuote{你是伯多禄}等等\BibleRef{玛 16:18}。因为，凡是基督的门徒，那些喝了那跟随他们的神性盘石所涌出的水的人，\BibleRef{格前 10:4}就是盘石；在这样的盘石上，教会的每一句话以及与之相符的生活方式都被建立；因为在每位成全的人身上——他们拥有言语、行为与思想相结合，充满真福——教会是由天主建立的。\par

% structure: p0022 source=h2 target=subsection reason=relative-heading-rank; base=h2

\subsection{十一、给予伯多禄的应许不限于他，也适用于所有像他一样的门徒。}

但如果你以为整个教会是由天主仅仅建立在那一个伯多禄身上，那么关于雷霆之子若望或其他每位宗徒，你会说什么？我们岂敢说，阴府的门唯独不能胜过伯多禄，却能胜过其他宗徒和成全的人？先前所说的话：\BibleQuote{阴府的门不能胜过她}\BibleRef{玛 16:18}，难道不是对所有人、对其中每一位都有效吗？还有那句话：\BibleQuote{在这盘石上，我要建立我的教会}？\BibleRef{玛 16:18}难道天国的钥匙只由主赐给伯多禄，而其他有福之人不会领受它们吗？如果这个应许：\BibleQuote{我要将天国的钥匙交给你}\BibleRef{玛 16:19}是与其他门徒共享的，那么先前所说的一切，以及接着对伯多禄所说的话，岂不也是与众人共享吗？因为在这里，这些话似乎只是对伯多禄说的：\BibleQuote{凡你在地上所束缚的，在天上也要被束缚}，等等；\BibleRef{玛 16:19}但在若望福音中，救主向门徒们嘘气，赐予他们圣神，说：\BibleQuote{你们领受圣神吧}，等等。\BibleRef{若 20:22}许多人会对救主说：\BibleQuote{你是默西亚，永生天主之子}；但并非所有说这话的人都是在没有血肉的启示、而是藉着天上的圣父亲自除去他们心中的帕子的情况下向祂说的，以便此后\BibleQuote{他们以揭开的脸面反映主的荣耀}，\BibleRef{格后 3:18}藉着天主之神说话，论及祂说：\Quote{主耶稣}，并向祂说：\BibleQuote{你是默西亚，永生天主之子。}\BibleRef{玛 16:16}如果有人这样向祂说，不是因血肉启示他，而是藉着天上的圣父，他将会得到按照福音字面所讲给那伯多禄的话，但正如福音的精神所教导的，是给每一个成为像那伯多禄一样的人。因为所有效法基督、即那跟随得救者的神性盘石的人，\BibleRef{格前 10:4}都得了\Quote{盘石}的称号，以便从它饮用神性的饮品。他们像基督一样受了盘石的名号。而且，作为基督的肢体，他们从祂得名，被称为基督徒；从盘石，他们得名为伯多禄。由此你可以说，义人承受基督之名，基督就是正义；智慧人承受基督之名，基督就是智慧。\BibleRef{格前 1:30}同样，关于祂的其他一切名号，你都可以作为名号应用于圣徒；对所有这些人都可说救主的话：\BibleQuote{你是伯多禄}等等，直到\BibleQuote{不能胜过她}。但\Quote{她}指什么？是基督建立教会的盘石，还是教会？因为这话是歧义的。或者，盘石和教会是同一个？我认为这是真的；因为阴府的门既不能胜过基督建立教会的盘石，也不能胜过教会；正如在箴言中所写，蛇在盘石上无路可寻。如果阴府的门胜过某个人，这样的人就不能成为基督建立教会的盘石，也不能成为耶稣建立在盘石上的教会；因为盘石不能被蛇触及，它比反对它的阴府之门更坚固，因此阴府的门不能胜过它；而教会，作为基督的建筑物——祂聪明地将自己的房屋建立在盘石上——，\BibleRef{玛 7:24}不能容纳阴府的门，因为阴府的门能胜过每一个在盘石和教会之外的人，却不能胜过它。\par

% structure: p0024 source=h2 target=subsection reason=relative-heading-rank; base=h2

\subsection{十二、每一种罪——每一种虚假的教义都是一道\Quote{阴府的门}。}

但当我们明白了每一件引至阴府的罪都是阴府的一扇门时，我们就应领悟到，那有\BibleQuote{斑点、皱纹或任何这类东西}\BibleRef{弗 5:27}，且因邪恶而不圣洁、无可指摘的灵魂，既非基督在其上建造的磐石，也非教会，亦非基督在磐石上建造的教会的一部分。若有人因那些被认为相信教会的绝大多数人，要以此事羞辱我们，必须对他说：不仅\BibleQuote{许多蒙召，但少有被选}\BibleRef{玛 22:14}，还有救主对来到祂面前的人所说的话，如\BibleBook{路加}{福音}记载的：\BibleQuote{你们要从窄门尽力进入，因为许多，我告诉你们，想要从窄门进去，却不能}\BibleRef{路 13:24}；以及\BibleBook{玛窦}{福音}所写的：\BibleQuote{因为窄门和狭路通往生命，找到它的人很少}\BibleRef{玛 7:14}。你若留意\BibleQuote{许多，我告诉你们，想要进去却不能}\BibleRef{路 13:24}这句话，便会明白这是指那些自夸属于教会，却活得软弱且违背圣言的人。因此，在那些试图进入的人中，那些不能进入的人将无法进入，因为阴府之门胜过他们；但对于那些阴府之门不能胜过的人，试图进入的人将强而有力，在基督耶稣内能成就一切，祂加强他们的力量\BibleRef{斐 4:13}。同样，每一个邪恶意见的创始者，都成了某一阴府之门的建造者；而那些与这类教导合作的人，则是仆人和管家，是那建造不敬的邪恶教义的奴仆。虽然阴府之门众多且几乎无数，但没有一扇阴府之门能胜过那磐石，或胜过基督在其上建造的教会。然而，这些门具有某种力量，藉此它们制胜一些不抵抗、不奋战的人；但另一些人却能胜过它们，因为他们未曾偏离那说\BibleQuote{我就是门}\BibleRef{若 10:9}的那一位，并从自己的灵魂中铲除了所有阴府之门。我们也应当知道，如同城市之门各有其名，同样，阴府之门也可按罪的种类命名；因此，一扇阴府之门被称为\Quote{淫乱}，通过它，行淫者进入阴府；另一扇被称为\Quote{否认}，通过它，否认天主者下到阴府。同样，每一位异端者，以及任何生出\BibleQuote{伪称为知识的}\BibleRef{弟前 6:20}的人，都已建造了一扇阴府之门——\Person{马西昂}造一扇，\Person{巴西里得}造另一扇，\Person{瓦伦提努斯}又造一扇。\par

% structure: p0026 source=h2 target=subsection reason=relative-heading-rank; base=h2

\subsection{十三、\Quote{阴府之门}与\Quote{熙雍之门}的对比}

在此处，所说的是阴府的门；但在圣咏中，先知感谢说：\Quote{他从死亡的门中提拔我，使我能 在熙雍女儿的门内宣扬你的一切赞美。}由此我们得知，除非一个人从死亡的门中被提拔，来到熙雍的门，否则他永远不可能适合宣扬天主的赞美。熙雍的门可以被视为与死亡的门相对，因此有一个死亡的门是放荡，而熙雍的一个门是节制；同样，有一个死亡的门是不义，而熙雍的一个门是正义，先知如此显示并说：\Quote{这是上主的门，义人将要进入其中。}再者，胆怯是死亡的门，而勇敢是熙雍的门；不智是死亡的门，而与之相反的明智是熙雍的门。但对于所有\Quote{假冒的知识} \BibleRef{弟前 6:20}的门，有一个门与之相对，即毫无虚假的知识之门。但试想，因着\Quote{我们战斗不是对抗血肉之躯} \BibleRef{弗 6:12}等话，你能说每一个权势和这黑暗的霸主，以及每个\Quote{天界里邪恶的属神军队} \BibleRef{弗 6:12}都是阴府的门和死亡的门吗？那么，让我们称之为阴府的门，而\Quote{服役的灵} \BibleRef{希 1:14}是正义的门。但正如在更好的事物中，首先提到许多门，在门之后，有一个门，在经文中：\Quote{给我敞开正义的门，我要进去，向上主完整宣认}，以及\Quote{这是上主的门，义人由它进入}；同样，在那些相对的门中，有许多阴府和死亡的门，每一个都是一个权势；但凌驾所有这些之上的，是邪恶者本身。我们要留意每一件罪，仿佛若犯罪就下到某个死亡的门；但当我们从死亡的门被提拔时，让我们在熙雍女儿的门内宣扬上主的一切赞美。例如，在熙雍女儿的一个门——即称为节制的门——我们以我们的节制宣扬天主的赞美；在另一个称为正义的门中，我们以正义宣扬天主的赞美；总之，在我们所从事的一切值得赞美的事中，我们就在熙雍女儿的某个门，在每个门宣扬天主的一些赞美。但我们必须研究，是否在十二卷小先知书之一中说：\Quote{他们在城门口恨那指摘的人，厌恶那说圣言的。} \BibleRef{亚 5:10}那么，也许在城门口指摘的人属于熙雍女儿的门，指摘那些处在与此门相反的罪中的人，即阴府或死亡的门。但如果你不这样理解\Quote{他们在城门口恨那指摘的人}这句话，那么\Quote{在城门口}这个说法要么被视为多余，要么你要探究这话如何与先知精神相称。\par

% structure: p0028 source=h2 target=subsection reason=relative-heading-rank; base=h2

\subsection{十四、在何种意义上\Quote{钥匙}赐给了\Person{伯多禄}，以及每一个伯多禄。此权力的限制。}

此后，让我们看看在何种意义上对伯多禄，以及对每一位伯多禄说：\Quote{\Emph{我要把天国的钥匙交给你。}}\BibleRef{玛 16:19}首先，我认为说\Quote{我要把天国的钥匙交给你}这句话，是与\Quote{阴府的门不能胜过她}这句话一致的。\BibleRef{玛 16:18}因为凡是在阴府的门前坚固自卫，使其不能胜过他的人，就配从同一圣言领受天国的钥匙，仿佛作为奖赏而接受天国的钥匙，因为阴府的门对他没有权势，使他能为自己打开那些对曾被阴府之门征服的人关闭的门。他作为一个节制的人，藉着打开节制的钥匙，进入一扇敞开的门——节制之门；作为一个义人，藉着另一扇门——正义之门，它由正义之钥匙开启；其余德行也是如此。因为我认为，对于每一种知识的德行，与德行种类相应的某些智慧奥秘，会向那按德行生活的人开启；救主赐给那些不被阴府之门所制胜的人与德行数目相等的钥匙，这些钥匙开启数量相等的门，这些门按照奥秘的启示与每一种德行相应。也许，每一种德行都是一个天国，而所有德行合起来是一个诸天之国；因此，按照这种理解，那按德行生活的人已经是在诸天之国里，所以\Quote{你们要悔改，因为天国近了}这句话，不应指向时间，而应指向行为和心性；因为基督，祂就是一切德行，已经来临，并发言，因此天主的国就在祂的门徒们心中，而不在这里或那里。\BibleRef{路 17:21}但是请想一想，基督在其上建立教会的磐石具有何等大的能力，以及每一个说\Quote{你是基督，永生天主之子}的人具有何等大的能力，以致此人的判决坚定可靠，仿佛天主在他内审判，以至于在审判之际，阴府的门不能胜过他。但是，当一个人不义地审判，不按照天主的话在地上捆绑，也不按照祂的旨意在地上释放时，阴府的门就胜过他；然而，对于任何阴府的门不能胜过的他，这人便正义地审判。因此他拥有天国的钥匙，向那些在地上已被释放的人打开，使他们也在天上被释放并得自由；并向那些因他的正义审判在地上被捆绑的人关闭，使他们在天上也被捆绑并受定罪。但是，当那些担任主教职的人，如同伯多禄一样使用这句话时，他们从救主那里接受了天国的钥匙，教导说由他们捆绑（即定罪）之物在天上也捆绑，由他们获得赦免之物在天上也释放时，我们必须说，如果他们拥有那使得对那位伯多禄说\Quote{你是伯多禄}的生活方式，\BibleRef{玛 16:18}并且他们是基督在其上建立教会的这种人，并且此语能合理地指向他们，那么他们说话是正确的；并且当他想要捆绑和释放时，阴府的门不应胜过他。但如果他被自己的罪绳紧紧捆绑，\BibleRef{箴 5:22}那么他捆绑和释放就毫无意义。也许你可以说，在智者的诸天（即诸德行）中，恶人被捆绑；同样，在这些德行中，善人被释放，并为其在德行之前所犯的罪获得了赦免。但是，正如那个没有罪绳和不义之物（这些罪和不义被比作\Quote{长绳或牛轭的带子}\BibleRef{依 5:18}）的人，连天主也不能捆绑他；同样，任何一位伯多禄，无论他是谁，也不能；如果有人并不是伯多禄，也不拥有此处所说的事物，却自以为像伯多禄一样，以为他在地上如此捆绑，以致捆绑的在天上也绑定，在地上如此释放，以致释放的在天上也释放，那么他是自高自大，不明白圣经的含义，并且因自高自大而陷入了魔鬼的毁灭。\BibleRef{弟前 3:10}\par

% structure: p0030 source=h2 target=subsection reason=relative-heading-rank; base=h2

\subsection{耶稣先前给门徒的使命与现在沉默命令的关系。信德与知识的对比。}

\BibleQuote{那时祂嘱咐门徒，不要对任何人说祂是基督。}\BibleRef{玛 16:20} 上面记载，耶稣派遣这十二人，吩咐他们说：\BibleQuote{外邦人的路，你们不要走。}\BibleRef{玛 10:5} 以及记录的其他话语，是在祂派遣他们去执行宗徒使命时对他们说的。那么，当门徒们已经执行宗徒职务时，祂是否希望他们宣告祂是基督呢？因为，如果祂希望如此，就应当追问，为何现在祂命令门徒不要说祂是基督？如果祂不希望，那么关于宗徒职务的事怎能稳妥地维持呢？在此处也可以追问这些问题：当祂差遣十二人时，难道祂不是差遣他们带着\Quote{祂是基督}的认识吗？但如果十二人有这样的认识，显然伯多禄也有；那么，他现在为何被称为有福呢？因为这里的表述清楚地表明，伯多禄现在首次承认基督是永生天主之子。玛窦福音在一些手抄本中写道：\BibleQuote{那时祂嘱咐门徒，不要对任何人说祂是基督。}\BibleRef{玛 16:20} 但马尔谷说：\BibleQuote{祂嘱咐他们，不要对任何人谈及祂。}\BibleRef{谷 8:30} 路加说：\BibleQuote{祂嘱咐并命令他们，不要把这事告诉任何人。}\BibleRef{路 9:21} 但这\Quote{这事}是什么？是否根据路加，伯多禄在回答\Quote{你们说我是谁？}的问题时，说\BibleQuote{是天主的基督？还是永生天主之子？}\BibleRef{玛 16:15} 然而，你们要知道，有些玛窦福音的手抄本写的是\BibleQuote{祂嘱咐了。}\BibleRef{玛 16:20} 这样引起的困难在我看来确实很大；但让我们寻求一个无可辩驳的解答，并且如果它比我们将要提出的较为温和的观点更可信，就让找到它的人在众人面前提出来。所以，请考虑，你们是否可以说：相信耶稣是基督，低于对所信之事的认识。或许在认识耶稣是基督方面也有差异，因为认识祂的人都并非同样地认识祂。从若望的话中：\BibleQuote{你们若存留我的话，就必认识真理，真理必使你们自由。}\BibleRef{若 8:31} 很明显，没有知识的信仰低于知道；但在认识耶稣是基督方面存在差异，因为认识祂的人都并非同样地认识祂，这是事实，对稍加思考的人都是显而易见的。例如，谁不会承认，弟茂德虽然知道耶稣是基督，但在对祂的认识上并没有像宗徒那样被光照？谁不会也承认这一点：虽然许多人说实话，论及天主说：\Quote{祂赐给我对存有之物的真知。}但他们不会以同样的洞察和领悟说出所认识的事，也不会知道同样多的事。然而，不仅在认识程度上有差异，而且根据知识的来源也有不同；因此，根据这一点，那藉着圣父的启示认识圣子的人，如伯多禄被证实认识的那样，享有最高的真福。现在，如果我们的这些观点是正确的，你们将考虑，十二人先前相信但并不知道；但相信之后，他们也获得了知识的开端，知道了关于祂的一些事；之后他们继续在知识上进步，以致能够接受来自那启示子的圣父的知识；伯多禄当时就处于这种状态，因此他被宣布为有福；因为他被宣布为有福，不仅仅是因为他说了\BibleQuote{你是基督}，而且加上\BibleQuote{永生天主之子}。因此，马尔谷和路加记载伯多禄回答说\BibleQuote{你是基督}，但未记载玛窦所加的内容，也没有记载他因所说的话而被宣布为有福，以及随之而来的祝福：\BibleQuote{你是伯多禄}\BibleRef{玛 16:18} 等等。\par

% structure: p0032 source=h2 target=subsection reason=relative-heading-rank; base=h2

\subsection{十六、门徒知识的逐步增长}

但现在我们必须首先调查一个事实：他们曾宣告关于祂的其他伟大而奇妙的事，但尚未宣布祂是基督，免得救主似乎剥夺了他们原本可宣布祂是基督的权威。或许有人会支持这样一种论点，说当犹太人被引入基督的学校时，门徒们向他们传授关于耶稣的光荣之事，以便在适当的时候，将这些作为基础，建立关于耶稣是基督的事；也许对他们所说的许多事，实际上是对所有实质上的信徒说的；因为不是只有宗徒才适用这句话：\BibleQuote{为了我的缘故，你们也要被带到总督和君王面前，对他们和外邦人作证。}\BibleRef{玛 10:18} 而且也许也不是绝对地对宗徒，而是对所有将要相信这句话的人：\BibleQuote{兄弟要把兄弟置于死地。}\BibleRef{玛 10:21} 等等；而\BibleQuote{凡在人前承认我的。}\BibleRef{玛 10:32} 等等，不是特别地对宗徒说的，也是对所有信徒说的。因此，根据这一点，通过对宗徒所说的话，预先给出了一个教学大纲，这教学大纲日后将对他们和每位教师都有益。\par

% structure: p0034 source=h2 target=subsection reason=relative-heading-rank; base=h2

\subsection{十七、这种逐步知识的原因}

同样，凡认为宗徒们早已宣报基督这一事实的，当他们听到\BibleQuote{我在暗中告诉你们的，你们要在光明中说出来；你们在耳边听到的，要在屋顶上宣扬出来}\BibleRef{玛 10:27}这句话时，就会说：祂愿意首先给那些将要听到基督之名的宗徒们某种教理讲授，然后容许这，可以这么说，在听众的脑海里消化，以便在沉默一段时间后，关于祂的这类宣布，在更合宜的时候，可以在先前的基础上建立\Quote{基督耶稣被钉十字架并从死者中复活}，这在起初连宗徒们也不知道；因为在本段经文中写道：\BibleQuote{从那时起，耶稣开始向门徒指示，祂必须往耶路撒冷去}\BibleRef{玛 16:21}，并受这些苦。但如果宗徒们现在才第一次从耶稣那里得知将要发生在祂身上的事，即长老们要谋害祂，祂将被杀死，此后第三天祂要从死者中复活——那么有什么必要假定那些由宗徒们教导有关耶稣的人先前就知道这些，或者虽然基督被宣告给他们，但只是以一种没有清楚阐明关于祂的事的方式介绍的呢？因为我们的救主吩咐门徒不可告诉任何人祂是基督时，祂愿意将关于祂的更完善的教导保留到更合适的时候，那时那些看见祂被钉十字架的人，以及看见祂被钉十字架并复活的宗徒们，能够为祂的复活作证。因为如果宗徒们常与祂同在，看见了祂所行的一切奇事，并为祂的话作证说它们是永生的话\BibleRef{若 6:68}，却在祂被出卖的那夜跌倒——那么你想那些早先就得知祂是基督的人会作何感想？我认为，祂赐下这道命令是为了体恤他们。\par

% structure: p0036 source=h2 target=subsection reason=relative-heading-rank; base=h2

\subsection{十八、耶稣起初仅被十二位宗徒宣报为行异能者和师傅。}

但若有人认为对十二位宗徒所说的话指的是此后的事，并且宗徒们尚未向听众宣布祂是基督，那么他会说：祂愿意将与耶稣之名相连的基督概念保留给那更完善且带来救恩的宣讲，正如保禄所知道的，他对格林多人说：\BibleQuote{我决定在你们中只知道耶稣基督，并被钉十字架的那一位。}\BibleRef{格前 2:2}因此，先前他们宣报耶稣是某些行动的施行者和某些教义的导师；但现在，当伯多禄承认祂是基督，永生天主之子时，由于祂不愿意让人已经宣布祂是基督，以便在更合适的时候，并作为被钉十字架的那一位被宣报，祂命令门徒们不可告诉任何人祂是基督。当祂禁止宣布祂是基督时，这一含义在一定程度上由以下话语得到证实：\BibleQuote{从那时起，耶稣开始向门徒指示，祂必须往耶路撒冷去，并从长老们那里受许多苦}等等；\BibleRef{玛 16:21}因为那时，在合适的时候，祂仿佛向那些知道耶稣是基督、永生天主之子（父已启示给他们）的门徒宣布：他们不要相信已被钉十字架的耶稣基督，而要相信即将被钉十字架的耶稣基督。同样，不要相信基督耶稣并从死者中复活，祂教导他们相信基督耶稣并即将从死者中复活。既然\BibleQuote{祂脱去了率领者和掌权者，公开羞辱了它们，借着十字架取得了胜利}\BibleRef{哥 2:15}，若有人以基督的十字架为耻，他就以这因之而胜过那些权势的救恩计划为耻；因此，凡既相信又知道这些事的人，应当以我们的主耶稣基督的十字架为夸耀，\BibleRef{迦 6:14}借着这十字架，当基督被钉时，率领者——其中我想也有这世界的元首——被示众并在相信的世界面前被战胜。因此，当祂的受难临近时，祂说：\BibleQuote{如今这世界的元首已受了审判}\BibleRef{若 16:11}，并且\BibleQuote{如今这世界的元首要被赶出去}，\BibleQuote{我若从地上被举起来，就要吸引众人归向我}\BibleRef{若 12:31}，因为他不再有足够的能力阻止那些被祂吸引的人去到耶稣那里。\par

% structure: p0038 source=h2 target=subsection reason=relative-heading-rank; base=h2

\subsection{十九、宣报耶稣为被钉十字架者的重要性。}

因此，为了宣报耶稣为基督，必须宣报祂被钉十字架；在我看来，宣报耶稣是基督，若沉默不谈祂的其他任何奇迹，都不如沉默不谈祂的十字架更不完善。因此，为了保留宗徒们关于祂的更完善的宣讲，祂命令门徒不可告诉任何人祂是基督；并且祂预备他们说出祂是钉十字架并复活的基督，\BibleQuote{当祂开始}不仅说，甚至不止于教导，而是\BibleQuote{指示}\BibleRef{玛 16:21}给门徒祂必须往耶路撒冷等等；因为请注意\Quote{指示}这个词；正如有形之物被说成被显示一样，耶稣对门徒所说的话也被说成由耶稣指示。而且我不认为，当那些看见祂在身体上从民间长老们受许多苦的人所看见的每件有形之物，其清晰程度能比得上关于祂的门徒所领受的理性证明。\par

% structure: p0040 source=h2 target=subsection reason=relative-heading-rank; base=h2

\subsection{二十、为何耶稣必须前往耶路撒冷。}

\Emph{\BibleQuote{然后祂开始显示}}；\BibleRef{玛 16:21} 之后，当他们能够接受时，祂可能更清楚地显示，不再像向那些学习入门者那样开始显示，而是在显示上已有长进；若合理认为耶稣完全成就了祂所开始的，那么，在某时刻，祂完全成就了祂开始向门徒显示的关于祂必须受所记载之事的必要性。因为，当任何人从圣言领悟到这些事的圆满知识时，就必须说，出于理性的展示（心灵看见所显示之事），展示对于那有志有能力且实际默观这些事的人成为圆满。但由于\BibleQuote{先知不在耶路撒冷丧命是不能的}\BibleRef{路 13:33}——这丧命对应\BibleQuote{为我的缘故丧失生命的，必要得着生命}\BibleRef{玛 10:39}——为此祂必须上耶路撒冷去，在那耶路撒冷受许多苦，好使祂在上的耶路撒冷成为从死者中复活的\BibleQuote{初果}\BibleRef{格前 15:20}，废除并瓦解地上的城及其中的一切敬礼。因为只要基督\BibleQuote{没有复活，成为睡了之人初果}\BibleRef{格前 15:20}，那些成为符合祂死亡与复活的人尚未与祂一同复活，天主的城、圣殿、洁净礼等一切便仍在地上寻求；但当此事发生后，不再寻求地上的事，而是天上的事；为建立这些，祂必须往下到耶路撒冷，在那里受长老、司祭长和民间的经师许多苦，为要受天上的长老——那些能接受祂恩惠者——以及更神圣的、在唯一大司祭之下被立的大司祭的荣耀，并为要受那从事\BibleQuote{不是用墨写的}\BibleRef{格后 3:3}、而是由永生天主的圣神所明示的文字的民间的经师的荣耀，并在下的耶路撒冷被杀，从死者中复活，在熙雍山、永生天主的城——天上的耶路撒冷\BibleRef{希 12:22}——为王。但第三日祂从死者中复活，为要拯救他们脱离恶者及其子（其中是谎言、不义、战争及一切与基督相反之事），也脱离那伪装为圣神的亵渎之神，使被拯救者获得在灵、魂、体上受洗归入圣父、圣子、圣神之名——这代表那三日期，永恒地同时呈现在那些借此成为光明之子的人面前。\par

% structure: p0042 source=h2 target=subsection reason=relative-heading-rank; base=h2

\subsection{二十一、伯多禄的责斥与耶稣的答复。}

\BibleQuote{ \Emph{伯多禄便拉过祂来，开始责斥祂说：主，愿天主怜悯祢，这事绝不会临到祢身上。}} \BibleRef{玛 16:22} 耶稣却对他说：\BibleQuote{ \Emph{撒殚，退到我后面去！你是我的绊脚石，因为你所体会的不是天主的事，而是人的事。}} \BibleRef{玛 16:23} 既然耶稣开始向门徒显示祂必须上耶路撒冷去，受许多苦，伯多禄至此学到了所显示之事的开端。但他认为受难与基督、永生天主之子不相称，也低于曾向他启示如此伟大之事的天父的尊严——因为关于祂将要受难的事尚未启示给他——为此他拉过祂来，如同忘记了对基督应有的尊敬，也忘记了永生天主之子既不做什么也不说什么值得责斥的事，便开始责斥祂；并如同向一个需要赎罪祭的人——因为他尚不知道\BibleQuote{天主设立了耶稣为赎罪祭，是凭着祂的血，藉着信德}\BibleRef{罗 3:25}——他说：\BibleQuote{主，愿天主怜悯祢。}\BibleRef{玛 16:22} 耶稣赞许他的善意，却斥责他的无知，因他的善意是正确的，便对他说：\BibleQuote{退到我后面去}\BibleRef{玛 16:23}，如同一个因无知而说错话的人，已放弃了跟随耶稣；但因他的无知，如同一个与天主的事敌对的人，祂说：\BibleQuote{撒殚}，这在希伯来语中意为\Quote{仇敌}。但若伯多禄不是出于无知而说话，也不是责斥永生天主之子，对他说：\BibleQuote{主，愿天主怜悯祢，这事绝不会临到祢身上}，基督就不会对他说：\BibleQuote{退到我后面去}，如同一个已放弃在祂后面跟随祂的人；也不会如同一个说了与祂所说相反的话的人说：\BibleQuote{撒殚}。但现在撒殚胜过了那跟随耶稣并走在祂后面的人，使他偏离跟随，不再在天主之子后面，并因他无知所说的话，使他有资格被称为\BibleQuote{撒殚}、天主之子的绊脚石，并\Quote{不体贴天主的事，只体贴人的事}。但伯多禄先前是在天主之子后面的，在他犯这罪之前，这从以下话明显可见：\BibleQuote{来跟随我，我要使你们成为渔人的渔夫。}\BibleRef{玛 4:19}\par

% structure: p0044 source=h2 target=subsection reason=relative-heading-rank; base=h2

\subsection{二十二、\Quote{后面}与\Quote{转过身来}这些表达的重要性。}

但你要将祂对伯多禄所说的\BibleQuote{撒殚，退到我后面去！}\BibleRef{玛 16:23}与祂对魔鬼（他对祂说：\BibleQuote{你若俯伏朝拜我，我就把这一切都给你}\BibleRef{玛 4:9}）所说的\BibleQuote{退去！}\BibleRef{玛 4:10}而没有加上\BibleQuote{到我后面}比较；因为到耶稣后面是一件好事。因此经上记载：\BibleQuote{来，到我后面来，我要使你们成为渔人的渔夫。}\BibleRef{玛 4:19}同样的话也说：\BibleQuote{凡不背起自己的十字架追随在我后面的，不配是我的。}\BibleRef{玛 10:38}作为一个一般原则，请注意\Quote{后面}这个词；因为当一个人走在主天主的后面、在基督的后面是一件好事；但当一个人把天主的话抛在脑后，或违反诫命说：\BibleQuote{不要跟在你的私欲后面}\BibleRef{德 18:30}，那就不然。厄里亚在列王纪上对民众说：\BibleQuote{你们心持两意要到几时呢？如果天主是主，就跟随祂；如果巴耳是主，就跟随他。}\BibleRef{列上 18:21}耶稣对伯多禄说这话时，祂\Quote{转过身来}，这是施恩的方式。所以，如果你收集更多关于\Quote{转过身来}的例子，特别是那些归给耶稣的，并将它们互相比较，你会发现这个词并非多余。但目前从若望福音中举出这一点就够了：\BibleQuote{耶稣转过身来，看见他们跟着——明显是伯多禄和安德肋——就问他们说：你们找什么？}\BibleRef{若 1:38}请注意，当祂\Quote{转过身来}时，是为了那些祂转向的人的利益。\par

% structure: p0046 source=h2 target=subsection reason=relative-heading-rank; base=h2

\subsection{二十三、\Person{伯多禄}作为耶稣的绊脚石。}

接下来我们必须探究祂如何对伯多禄说：\BibleQuote{你是我的绊脚石}\BibleRef{玛 16:23}，尤其当达味说：\BibleQuote{爱慕你法律的人，必得平安，对他们没有绊脚石。}因为有人会说，如果这在先知书中这样说了，是因为那些有爱的人的坚定，他们不可能跌倒，因为\BibleQuote{爱包容一切，相信一切，希望一切，忍耐一切，爱永不衰败}\BibleRef{格前 13:7}，那么主亲自\BibleQuote{扶持所有跌倒的，扶起所有被压弯的}，怎能对伯多禄说\BibleQuote{你是我的绊脚石}？但必须说：不仅救主，就连那在爱中成全的人，也不会跌倒。但就他自己而言，说或做这样事的人，即使对那不会跌倒的人，也是绊脚石；除非耶稣称那犯罪的门徒为绊脚石，甚至是祂自己的绊脚石，如同耶稣远比保禄更会凭着爱而说：\BibleQuote{谁软弱，我不软弱？谁跌倒，我不焦急？}\BibleRef{格后 11:29}与此一致，我们也可以说：\BibleQuote{谁跌倒，我不跌倒？}但如果伯多禄当时因为说\BibleQuote{主，愿天主怜悯你，这事绝不会临到你身上}\BibleRef{玛 16:22}而被耶稣称为绊脚石，因为他在所说的话中不思念天主的事，只思念人的事，那么对那些自称是耶稣门徒却不思念天主的事、不关注不可见和永恒的事，只思念人的事、关注可见和暂时的事的人，更当如何说？\BibleRef{格后 4:18}但这样的人更要被耶稣指责为祂的绊脚石，并且因为作了祂的绊脚石，也就作了祂弟兄们的绊脚石？正如祂关于他们说的：\BibleQuote{我渴了，你们没有给我水喝}\BibleRef{玛 25:42}等等，同样祂也可以说：\BibleQuote{当我奔跑时，你们使我跌倒。}因此我们不要认为思念人的事是小罪，因为我们应当在一切事上思念天主的事。这对每一个偏离天主教训、教会圣言和正确思想的人也是恰当的；例如，对于那些把\Person{巴西理德}、\Person{瓦伦廷}或\Person{马西翁}的教导当作真理，以及任何将人的事当作天主的事来教导的人。\par

% structure: p0048 source=h2 target=subsection reason=relative-heading-rank; base=h2

\subsection{二十四、克己与背十字架。}

\Quote{\Emph{于是耶稣对祂的门徒说：\InnerQuote{谁若愿意跟随我}},}等等。\BibleRef{玛 16:24}祂用这些话显示，愿意跟随耶稣并跟从祂，不是出于平常的男子气概，而且没有舍弃自己的人不能跟随耶稣。一个人若通过一次惊人的革命抹去自己先前在邪恶中度过的生活，他就是舍弃了自己；例如，一个曾经放荡的人否认了自己放荡的自我，成为持续自制的人。但可能有人会提出异议，他是否在舍弃自己时也承认自己？当他舍弃自己——那不正义的——时，并承认自己——那正义的？但如果基督是正义，那接受了正义的人承认的不是自己，而是基督；同样，那找到了智慧的人，因拥有智慧本身而承认基督。确实，正如\BibleQuote{心里相信以致成义，口里承认以致得救}，\BibleRef{罗 10:10}并为基督的作为作证，藉这一切在众人面前承认基督，他也要被基督在祂在天的父面前承认。\BibleRef{玛 10:32}同样，那没有舍弃自己却否认基督的人，将要经历这句话：\BibleQuote{我也要否认他。}\BibleRef{玛 10:33}因此，让每一个思想、每一个意念、每一句话、每一个行动都成为对自己的舍弃，但成为关于基督并在基督内的见证；因为我深信，成全之人的每一个行动都是对基督耶稣的见证，而远离一切罪恶就是对自己的舍弃，引导他跟随基督。这样的人与基督同钉十字架，背起自己的十字架跟随那为我们背负自己十字架的主，正如\BibleBook{若望}{福音}所记载的：\BibleQuote{他们于是把耶稣带走，把十字架放在他身上，}等等，直到\BibleQuote{他们把他钉在十字架上。}\BibleRef{若 19:17}但可以说，根据\BibleBook{若望}{福音}所记载的耶稣，祂自己背着十字架，背着走了出去；而根据\BibleBook{玛窦}{福音}、\BibleBook{马尔谷}{福音}和\BibleBook{路加}{福音}所记载的耶稣，祂没有背自己的十字架，而是基勒乃人西满背着。也许这个人指的是我们，我们因耶稣而背起耶稣的十字架，但耶稣自己背负了它；因为可以说有两个十字架的概念：一个是基勒乃人西满所背的，另一个是耶稣自己所背的。\par

% structure: p0050 source=h2 target=subsection reason=relative-heading-rank; base=h2

\subsection{二十五、论及保禄关于与基督同钉十字架的话。}

此外，关于这句话：\BibleQuote{让他舍弃自己，}\BibleRef{玛 16:24}保禄的话似乎合适，他舍弃了自己，说：\BibleQuote{我活着，却不再是我，而是基督在我内活着；}\BibleRef{迦 2:20}因为\BibleQuote{我活着，却不再是我}这句话，是舍弃自己之人的声音，如同一个人放弃了自己的生命，接受了基督，为叫祂作为正义、智慧、圣化、和平以及天主的大能住在他内，祂在他内成就一切。此外，还要注意这点：虽然死亡有多种形式，天主子被钉在十字架上，悬挂在木头上，为叫所有死于罪恶的人，除了藉十字架这条路，别无他法死于罪恶。因此他们会说：\BibleQuote{我已与基督同钉在十字架上，}以及\BibleQuote{我决不夸耀，只夸耀我们主耶稣基督的十字架，因基督的十字架，世界于我已被钉死，我于世界也被钉死。}因为也许每个与基督同钉十字架的人，都\BibleQuote{脱去自身的率领者和掌权者，公然羞辱他们，在十字架上凯旋}\BibleRef{哥 2:15}；或者更确切地说，基督在他们内成就了这些事。\par

% structure: p0052 source=h2 target=subsection reason=relative-heading-rank; base=h2

\subsection{二十六、生命的丧失；与它的得救。}

\BibleQuote{\Emph{因为谁若愿意救自己的性命，必要丧失性命。}}\BibleRef{玛 16:25}第一个表达方式有两种可能；可以这样理解：如果有人因贪生怕死，认为现世生命是好的，为了活在肉身中而小心翼翼地照料自己的生命，害怕死亡，好像死亡会失去它，此人因这样愿意救自己的生命，反而会丧失它，将其置于福乐之外。但如果有人因为我的话语而轻视现世生命，我的话促使他为永生而奋斗，甚至为真理而死，他丧失了自己的生命，为了虔敬而放弃那通常被称为死亡的东西，此人因我为缘故丧失了生命，反而会救得生命，并保持拥有它。根据第二种方式，我们可以这样解释这句话：如果有人真正理解了救恩是什么，希望获得自己生命的救恩，让他告别此生，舍弃自己，背起自己的十字架，跟随我，向世界丧失自己的生命；因为为我并为我的全部教导而丧失了生命，他将获得这种丧失的结局——救恩。\par

% structure: p0054 source=h2 target=subsection reason=relative-heading-rank; base=h2

\subsection{二十七、向世界丧失的生命得以得救。}

但同时也请注意，起初说的是：\BibleQuote{凡愿意的}，但后来是说：\BibleQuote{凡丧失的}。\BibleRef{玛 16:25} 如果我们希望它得救，便应在世俗中丧失它，如同那些与\Person{基督}同钉十字架的人，并以我们主\Person{耶稣基督}的十字架为夸耀，因为藉着这十字架，世界于我们已被钉在十字架上，我们于世界也被钉在十字架上。\BibleRef{迦 6:14} 好使我们达到目的，即我们生命的救恩；这救恩从我们为圣言的缘故丧失生命时开始。但如果我们认为生命的救恩是蒙福的事，且是指向那在天主内的救恩及与祂同在的福乐，那么任何生命的丧失都应当是一件好事，并且为了\Person{基督}的缘故，必须证明是这蒙福救恩的前奏。因此，在我看来，根据前面所说的，循着自我弃绝的类比，每个人都应丧失自己的生命。所以，让每个人丧失自己有罪的生命，好使在失去那有罪的之后，能领受那因正义行为而得救的生命；但人若赚得全世界，却毫无益处。我认为，那赚得世界的人，就是世界没有被他钉在十字架上的人；而世界没有被钉在十字架上的人，他将丧失自己的生命。但当两件事摆在我们面前：或赢得生命而丧失世界，或赢得世界而丧失生命，更可取的选择是：我们应丧失世界，并因\Person{基督}的缘故丧失生命而赢得生命。\par

% structure: p0056 source=h2 target=subsection reason=relative-heading-rank; base=h2

\subsection{二十八、生命的交换}

但是，那句\BibleQuote{\Emph{人还能拿什么作为自己生命的交换？}},\BibleRef{玛 16:26} 如果以反问的方式说出，似乎能表明：一个人在他犯罪后，若放弃了全部财产，以他的财物喂养穷人，彷彿要借此获得救恩，那么他就是为自己生命做了交换；但若以肯定语态说出，我认为这表明：人没有任何东西可以给出，作为自己被死亡战胜的生命之交换，来将生命从死亡手中赎回。所以，人不能为自己的生命给出任何交换，但天主为我们众人的生命给出了交换，就是\BibleQuote{\Person{基督耶稣}的宝血}，\BibleRef{伯前 1:19} 正如\BibleQuote{你们是用高价买来的}，\BibleRef{格前 6:20} \BibleQuote{不是用能朽坏的金银等物，而是用宝血，即无玷无瑕羔羊的血}，即\Person{基督}的宝血。\BibleRef{伯前 1:18} 在依撒意亚书中也向以色列说：\BibleQuote{为了你，我给了雇士作为交换，埃及和色巴也给了你；从何时起你在我眼中成为宝贵，你就被荣耀了。} \BibleRef{依 43:3} 例如，以色列长子的交换是埃及人的长子，以色列的交换是那些在埃及最后灾祸中死去的埃及人，以及灾祸后的淹没。但是，从这些事情上，让能够的人探究：天主赐给真以色列的交换，即\Quote{那救赎以色列脱离一切过犯的}，是否是真正的雇士，以及所谓的属灵的埃及和埃及的色巴；并且更大胆地探究：也许色巴是耶路撒冷的交换，埃及是犹太的交换，雇士是那些敬畏者的交换，他们不同于以色列、\Person{肋未}家族和\Person{亚郎}家族。\par

% structure: p0058 source=h2 target=subsection reason=relative-heading-rank; base=h2

\subsection{二十九、人子带着光荣来临}

\BibleQuote{\Emph{因为人子要在他父的光荣中，同他的天使一起降来。}} \BibleRef{玛 16:27} 然而，人子确实还没有在他的光荣中来临；\BibleQuote{因为我们看见了他，他没有俊美，也没有华丽；但他的容貌被损伤，比起世人来有所缺陷；他是一个受苦辛劳的人，熟悉疾病的忍受，因为他掩面不顾，被侮辱而不被尊重。} \BibleRef{依 53:2} 而且他必须以此形态来临，好担当我们的罪过\BibleRef{依 53:4}并为我们受苦；因为他在光荣中担当我们的罪过并为我们受苦是不合适的。但他也将在光荣中来临，先通过他那没有俊美、没有华丽的显现预备了门徒；并且，他成了像他们一样，好使他们能成为像他一样，\BibleQuote{与他光荣的肖像同化}，\BibleRef{罗 8:29} 既然他先前曾取用了\BibleQuote{我们卑微的身体}，\BibleRef{斐 3:21} 当他\BibleQuote{空虚了自己，取了奴仆的形体}时，\BibleRef{斐 2:7} 他便恢复了天主的肖像，并也使他们与这肖像同化。\par

% structure: p0060 source=h2 target=subsection reason=relative-heading-rank; base=h2

\subsection{三十、圣言以不同形态显现；他光荣来临的时刻}

但若你愿意理解圣言的差异——圣言藉着\BibleQuote{宣讲的愚昧}\BibleRef{格前 1:21}向相信的人宣讲，并以智慧向完全人讲述——你就会看到圣言如何以奴仆的形体对待学习初步道理的人，以致他们说：\BibleQuote{我们看见了他，他没有形貌，也没有美丽。}\BibleRef{依 53:2}但圣言在光荣中降临到完全人面前，\BibleQuote{在他父的光荣中}\BibleRef{玛 16:27}，他们可以说：\BibleQuote{我们看见了他的光荣，正如父独生者的光荣，充满恩宠和真理。}\BibleRef{若 1:14}因为完全人确实看见了圣言的光荣，以及天主父的独生子，恩宠和真理的圆满，而那种人无法领悟，他为了相信，需要\BibleQuote{宣讲的愚昧}。但\BibleQuote{人子要在自己父的光荣中降临}，并非独自一人，而是\BibleQuote{同他的天使一起}。如果你能设想所有在圣言的光荣中，以及在启示那即基督的智慧中共同协助的人，与祂一同来临，你就会看到人子如何在自己父的光荣中同自己的天使一同降临。并请思考，你是否能在此联系说：那些以前因自己的言语而受苦、具有\Quote{没有形貌也没有美丽}的先知，与那具有\Quote{没有形貌也没有美丽}的圣言处于相似的地位？并且，正如人子在自己父的光荣中降临，同样，那些在先知中的言语之天使，也与祂同在，保持着自己光荣的分寸。但当圣言以这种形式同自己的天使降临时，祂将根据各人的行为，把祂自己光荣的一部分以及祂天使的光辉的一部分分赐给各人。但我们说这些，并不否定对天主圣子第二次来临的更简单理解。但这些事何时发生呢？是否当那宗徒的神谕实现时，它说：\BibleQuote{因为我们众人都必须出现在基督的审判台前，为使每人按照自身所行的，或善或恶，领取相应的事。}\BibleRef{格后 5:10}但如果祂按各人的行为报应各人，不仅是善行，也不只是恶行单独，那么显然祂是按每一个恶行和每一个善行报应各人。但我以为——在此也跟随宗徒，同时也比较厄则克耳的话，其中完全悔改者的罪过被涂抹，而完全堕落者先前的义德不被记念——即对于完全人且已彻底抛弃邪恶者，罪过被涂抹；但对于完全背离虔敬者，即使他先前做过任何善事，也不被记念。\BibleRef{则 18:21}但对我们这些处于完全人和背教者中间地位的人，当我们站在基督的审判台前时\BibleRef{格后 5:10}，我们所做的无论善恶都将被报应；因为我们还没有纯洁到恶行不被归算于我们，也没有堕落至此而使我们的善行被遗忘。\par

% structure: p0062 source=h2 target=subsection reason=relative-heading-rank; base=h2

\subsection{三十一、关于不尝死亡的应许的简单解释}

\BibleQuote{\Emph{我实在告诉你们：站在这里的有些人，在还没有尝到死亡以前，必要看见人子带着自己的国和自己的光荣来临。}}\BibleRef{玛 16:28}有些人将这些话联系到六天后——或如路加所说，八天后\BibleRef{路 9:28}——耶稣带领三个门徒单独上了一座高山。采用这种解释的人说：伯多禄和另外两人在看见人子带着自己的国和自己的光荣来临之前，没有尝到死亡。因为当他们看见耶稣在他们面前显圣容，以致\BibleQuote{他的面貌发光}等等，以及\BibleQuote{他们看见了天主的国带着威能来临}。\BibleRef{谷 9:1}因为正如一些持戟卫士围立在君王周围，同样，梅瑟和厄里亚出现在那些上山的人面前，与耶稣谈话。但值得考虑的是，在祂的国中坐在救主左右是否指他们，以致\BibleQuote{但给谁预备了}这句话\BibleRef{玛 20:23}是因他们而说的。现在，这个关于三位宗徒直到看见耶稣显圣容而未尝死亡的诠释，适合于那些被伯多禄称为\BibleQuote{新生的婴儿贪求合理的、无诡诈的乳汁}\BibleRef{伯前 2:2}的人，保禄也对他们说：\BibleQuote{我给你们喝的是奶，不是饭}\BibleRef{格前 3:2}等等。同样，任何经文诠释，若能建立那些不能领受更大真理的人，都可合理地被称为乳，是从圣经的圣土流出，那里流着奶和蜜。但那些像依撒格一样\BibleRef{创 21:8}已经断奶的人，配得上亚巴郎在他儿子断奶时所设的筵席和款待，他们在这里和在每部圣经中寻求的食粮，我想，是不同于那确实是肉但并非固体食物的，也不同于被比喻称为蔬菜的，那是给已经断奶但仍不强壮、软弱之人的食物，正如经上所说：\BibleQuote{软弱的吃蔬菜}\BibleRef{罗 14:2}。同样，那些像撒慕尔一样已经断奶，并被其母亲奉献给天主的人——他的母亲是亚纳，意为恩宠——他也就是恩宠之子，如同在圣殿中长大的人，寻求天主的肉，那是同时身为完全人和司铎的神圣食物。\par

% structure: p0064 source=h2 target=subsection reason=relative-heading-rank; base=h2

\subsection{三十二、站在救主旁}

对于眼前这段经文，我们此刻所想到的反思如下：有些人站在耶稣所在之处，他们的灵魂之足坚定地与耶稣同在，他们的立足与梅瑟在经文中所说的立足相似：\BibleQuote{我在山上住了四十天四十夜}\BibleRef{申 10:10}，他被天主视为配得听天主对他说的话：\BibleQuote{你同我站在这里}\BibleRef{申 5:31}那些真正站在耶稣——即天主圣言——旁边的人，并不都站得平等；因为站在耶稣旁边的人彼此之间有差别。因此，并非所有站在救主旁边的人，而只有其中那些站得更好的人，不会尝到死味，直到他们看见那曾与世人同住、因此被称为人子的圣言，在自己的王国中来临；因为耶稣来临并不总是带着自己的国度，对于初信者而言，祂是如此，以至于他们看着圣言本身不光荣、不伟大，反而比他们中的许多人更为卑微，会说：\BibleQuote{我们看见了祂，祂没有威仪，也没有美貌；祂的容貌被侮辱，比世人更憔悴}\BibleRef{依 53:2}这些话将由那些在早先时候见过祂荣耀的人说出，那时圣言在会堂中被理解，对他们而言既无威仪也无美貌。因此，对于那已最明显地承受了超越一切言语之大能的圣言来说，有一种王者的尊严，为某些站在耶稣旁边的人所看见，当他们能够跟随祂前行，登上祂显现的高山时。这种殊荣，只有那些站在耶稣旁边的人中配得的人才能得到，他们或是伯多禄——阴间的门不能胜过他，或是雷霆之子\BibleRef{谷 3:17}，由天主那大能的声音所生，这声音向那些有耳且智慧的人从天上传出伟大的事。这样的人至少不会尝到死味。\par

% structure: p0066 source=h2 target=subsection reason=relative-heading-rank; base=h2

\subsection{三十三、对\Quote{尝死味}的解释}

但我们必须寻求理解\Quote{尝死味}的含义。祂就是生命，祂说：\BibleQuote{我是生命}\BibleRef{若 14:6}，这生命确实与基督一同藏在天主内；并且\BibleQuote{当基督，我们的生命，显现时，那时你们也要与祂一同}\BibleRef{哥 3:3}显现在荣耀中。但生命的仇敌，也就是祂所有仇敌中最后被消灭的，就是死亡\BibleRef{格前 15:26}，犯罪的灵魂所经历的死亡，其倾向与那正直生活的灵魂中所发生的相反，并因正直生活而生活。当法律中说：\BibleQuote{我将生命摆在你面前}\BibleRef{申 30:15}时，圣经指的是那说\BibleQuote{我是生命}的那一位，以及祂的仇敌死亡；我们每个人总是通过自己的行为选择其中之一。当我们面对生命而犯罪时，那针对我们的诅咒就应验了，它说：\BibleQuote{你的性命必悬在你面前}等等，直到\BibleQuote{和你眼中所见的情景}\BibleRef{申 28:66}。因此，正如生命也是从天上降下来并赐生命给世界的生命之粮\BibleRef{若 6:33}，同样地，祂的仇敌死亡是死粮。每个理性的灵魂要么以生命之粮为食，要么以死粮为食，取决于它所接受的好或坏的见解。正如在较普通的食物中，有时只是尝一尝，有时则更多地吃；同样，对于这些饼，一个人只是不足地尝一尝，而另一个人则吃饱——善人或走向善的人以从天上降下来的生命之粮为食，恶人以死粮即死亡为食；有些人或许只是少量地尝一尝死亡，稍微犯罪；但那些达到美德的人甚至不尝它，总是以生命之粮为食。因此，对于伯多禄——阴间的门不能胜过他——来说，他自然不会尝到死味，因为任何人尝死味并吃死亡，都是在阴间的门胜过他的时候；一个人根据阴间的门胜过他的程度——程度更大或更小、数量更多或更少——而吃或尝死味。但对于雷霆之子——他们由属天的雷霆所生——来说，尝死味是不可能的，因为死味与他们的母亲雷霆相距极远。然而，圣言预言了这些事，是对那些将要被成全、并因与圣言一同站立而进步到如此程度的人所说：他们没有尝过死味，直到他们看见天主圣言的显现、荣耀、国度和卓越——凭借这卓越，祂超越每一言语，这言语借着真理的表象引诱并拖住那些不能摆脱分心之束缚、不能登上真理之言卓越之高度的人。\par

% structure: p0068 source=h2 target=subsection reason=relative-heading-rank; base=h2

\subsection{三十四、\Quote{直到}的意义。应许不受限制。}

但或许有人认为，救主的应许为其不尝到死亡规定了时间的界限，即他们不会尝到死亡\Quote{直到}\BibleRef{玛 16:28}他们看见人子在自己的国里来临，但在这之后便会尝到死亡。我们要表明，根据圣经的用法，\Quote{直到}这个词表示所指之事的时间紧迫，但并非被如此限定，以至于在\Quote{直到}之后，与所指之事相反的事就一定会发生。现在，救主在复活后对十一位门徒说了许多话，其中包括：\Quote{看，我天天与你们在一起，直到时代的终结。}\BibleRef{玛 28:20}祂说这话时，是否许诺要与他们在一起直到时代的终结，但在时代的终结之后，当另一个时代临近（这时代被称为\Quote{将来的时代}），祂就不再与他们同在？——那么据此，门徒在时代终结之前的状态反而比时代终结之后更好？但我想，没有人敢说，在时代终结之后，天主子不再与门徒同在，因为这话宣告祂要与他们同在如此之久，直到时代的终结临近；因为显然，所质问的问题是：天主子是否要在将来的时代和所盼望的天主应许（作为赏报）之前，立即与祂的门徒同在？但可能存在一个问题——虽然肯定祂会与他们同在——祂是否有时临在，有时不临在。因此，为了消除可能因疑惑而生的怀疑，祂宣布：从现在起，甚至天天，祂都要与门徒同在，并且不会离开那些已成为祂门徒的人，直到时代的终结；（因为祂说\Quote{天天}，祂并没有否认，夜晚日落时，祂也会临在他们中间。）但如果\Quote{直到时代的终结}这话具有这样的力量，那么显然，我们就不必被迫承认，那些看见人子在自己的国里来临的人，在配得以如此形态瞻仰祂之后，还会尝到死亡。但是，正如我们所引用的经文那样，急需教导我们的是：祂\Quote{直到时代的终结}不会离开我们，而是天天与我们同在；同样，在此情况下，我认为对于懂得观察事物逻辑一致性的人来说，清楚的是：那一次性地看见了\Quote{人子在自己的国里来临}，并看见祂\Quote{在自己的荣耀里}，并看见\Quote{天主的国带着权能来临}的人，在沉思如此美好伟大之事后，不可能尝到死亡。然而，除了耶稣应许的话之外，我们并非没有理由推测，只要我们尚未被配得看见\Quote{天主的国带着权能来临}和\Quote{人子在自己的荣耀和自己的国里来临}，我们就会尝到死亡。\par

% structure: p0070 source=h2 target=subsection reason=relative-heading-rank; base=h2

\subsection{三十五、圣经中关于死亡的引用。}

但既然在此处三部福音中记载着\Quote{他们必不死尝死味,} 而在其他作者笔下关于死亡有不同的记载，那么将这些段落连同\Quote{尝}一并提出并审视，或许并非不合时宜。在\WorkTitle{圣咏集}中，这样说道：\Quote{有谁能存活而看不见死亡?} 又在另一处说：\Quote{愿死亡降临他们，愿他们活着下到阴府;} 但在一先知书中说：\Quote{死亡成了强大的吞灭者，吞没了他们,} \BibleRef{依 25:8} 而在\WorkTitle{默示录}中：\Quote{死亡和阴府跟随着某些人.} \BibleRef{默 6:10} 在我看来，在这些段落中，尝死味是一回事，看见死亡是另一回事，死亡降临某人又是另一回事，而第四种不同于前述的，由\Quote{死亡成了强大的吞灭者，吞没了他们,} 所指明，第五种不同于这些的，由死亡和阴府跟随他们所指明。如果你收集起来，或许还会发现其他我们未曾提到的差异，通过相互比较和正确的探究，你会发现每处所指的事物。但在此我探究的是：看见死亡是较小的恶，而尝死味是比看见更大的恶，但比这更糟的是死亡跟随某人，不仅跟随，而且现在降临并抓住它先前所跟随的人；而被吞没似乎比所言的一切都更痛苦。但是，留意所说的话以及所犯罪过的差异，我想你不会迟迟不承认这类事物是由那使这些事写在天主圣言中的圣神所意旨的。但是，如果必须对看见人子在自己的国度中或在自己的荣耀中来临，以及看见天主的国带着权能来临所指的意义，比前述更清楚地阐释，这些事——无论是那些在我们心中闪耀的，还是那些被寻求者寻见的，或是那些逐渐进入我们思想的——各人可随己意判断——我们将陈明。那注视并领悟圣言之卓越的人，当他破除并驳倒一切真实谎言却自称真理的事物之看似合理的形态时，就看见了人子（按\Person{若望}的话，\Quote{天主的圣言,}）在自己的国度中来临；但如果这样一个人不仅看见圣言破除似是而非的反对，而且以完全清晰的方式呈现自己的真理，那么他除了看见祂的国度之外，还会看见祂的荣耀。这样一个人确实会在祂身上看见天主的国带着权能来临；他将以此身份看见，即他不再处于\Quote{那在犯罪者必死身体中作王的罪,} \BibleRef{罗 6:12} 的统治之下，而是永远服从那位是万有之天主的王的命令，祂的国度确实潜在\Quote{在我们内,} \BibleRef{路 17:21} 但实际上，并如\Person{马尔谷}所称\Quote{带着权能,} 且毫无软弱地只存在于完美者内。因此，耶稣向当时站立的门徒应许了这些事，预言不是关于他们全体，而是关于其中一些人。\par

% structure: p0072 source=h2 target=subsection reason=relative-heading-rank; base=h2

\subsection{三十六、论救主的显容。}

\Quote{ \Emph{过了六天},} 按\Person{玛窦}和\Person{马尔谷}记载，\Quote{ \Emph{耶稣带着伯多禄、雅各伯和若望，单独领他们上了一座高山，在他们面前变了容貌.} } 现在，让我们也承认，在我们想到的关于这些事的解释之前，这件事是过去发生的，是按字面。但在我看来，那些被耶稣领上高山，并被恩准单独看见祂显容的人，并非无故在先前讲论的六天之后被领上去。因为既然在六天——这完美的数字——之内，整个世界——这完美的艺术作品——被造成，因此我认为，那超越世间一切，不再看那些可见的，因为它们是暂时的，而已经看那些不可见的，并且只看那些不可见的，因为它们是永恒的，这样的人由\Quote{六天后耶稣带了几个人与祂一同}这句话所代表。因此，如果我们中任何人愿意被耶稣带走，并被祂领上高山，并被恩准单独看见祂的显容，就让他超越这六天，因为他不再看可见之物，也不再爱世界和世界上的事物 \BibleRef{若一 2:15}，也不贪图任何世俗的贪欲，即肉身的贪欲、财物的贪欲、肉身的荣耀，以及任何本性是分散并拖曳灵魂远离更好更神圣事物，并将它拉下来紧紧束缚于这世代之欺骗的事物，在财富、荣耀以及其他与真理为敌的贪欲中。因为当他度过了这六天，如我们所说，他将守一个新的安息日，在高山上欢欣，因为他看见耶稣在他面前变了容貌；因为圣言有不同的形态，祂按各人看见的便利而显现，且不对超过观看者能力者显现。\par

% structure: p0074 source=h2 target=subsection reason=relative-heading-rank; base=h2

\subsection{三十七、\Quote{在他们面前} 这些话的力量。}

但你会问：当祂在那些被祂带上高山的人面前显圣容时，是否向他们显示为天主原有的形象，以致祂对山下的人保有仆人的形象，但对那些六天后跟随祂上高山的人，祂没有那形象，而是天主的形象？不过，你若能同时属神地留意，就请听这些事：经文并非简单地说\Quote{祂变了容貌}，而是带有某种必要的补充，玛窦和马尔谷都记载了；因为按两位所说，\Quote{祂在他们面前变了容貌}。据此，你确实会说，耶稣可能在一些人面前以这种显圣容出现，但在另一些人面前同时未显圣容。但若你想看见耶稣在那些与祂一同单独上高山的人面前的显圣容，请与我一同注视福音中的耶稣：祂被那些没有藉着提升人的言行登上智慧高山的人，更简单地理解为，可以说，\Quote{按肉体}认识；但祂不再按肉体被认识，而是藉着所有福音在祂的天主性中被认识，并按照他们的知识以天主的形象显现；因为耶稣在他们面前显圣容，而不在下面任何人面前。但当祂显圣容时，祂的面貌也发光如太阳，为要显给那些脱去黑暗作为、穿上光明武器的光明之子\BibleRef{罗 13:12}，他们不再是黑暗或黑夜之子，而是成了白昼之子，并像在白昼一样端庄行事；当祂显现时，祂要光照他们，不单像太阳，而是显现为正义的太阳。\par

% structure: p0076 source=h2 target=subsection reason=relative-heading-rank; base=h2

\subsection{三十八、衣裳洁白如光。}

祂不仅在这样的门徒面前显圣容，也不仅使祂的面貌如太阳发光；更甚，对那些被祂带上高山私下的人，祂的衣裳也洁白如光。\BibleRef{玛 17:2}但耶稣的衣裳是祂所穿戴的福音的表达和文字。然而，我认为宗徒著作中指示关于祂真理的话语也是耶稣的衣裳，它们对那些与耶稣一同上高山的人变得洁白。但由于白色也有差异，祂的衣裳变得像所有白色中最明亮纯粹的那样洁白，那就是光。因此，当你看见某人不仅透彻理解关于耶稣的神学，而且阐明了福音的每一句话，你不妨说，耶稣的衣裳在他面前变得洁白如光。但当天主圣子在祂的显圣容中被如此理解和看见时，祂的面貌成为太阳，衣裳洁白如光，那么，看见耶稣如此形象的人，会立刻看见梅瑟——法律——和厄里亚——以提喻的方式，不仅是一位先知，而是所有先知——与耶稣交谈；因为\Quote{与祂谈话}这话的力量如此；\BibleRef{玛 17:3}但按路加所说，\Quote{梅瑟和厄里亚在光荣中显现}，直到\Quote{在耶路撒冷}的话。\BibleRef{路 9:30}但若有人看见梅瑟的光荣，理解了属神的法律是与耶稣和谐的道理，以及先知们隐藏在奥秘中的智慧\BibleRef{格前 2:7}，那么当他看见他们与耶稣在一起时，就看见了梅瑟和厄里亚在光荣中。\par

% structure: p0078 source=h2 target=subsection reason=relative-heading-rank; base=h2

\subsection{三十九、耶稣显圣容——\Quote{祂正在祈祷时}}

然后，由于必须讲解马尔谷所记的段落：\Quote{\Emph{祂正在祈祷时，祂在他们面前变了容貌}}，我们必须说，也许我们特别能看见圣言在我们面前显圣容，如果我们做了前述之事，上了高山，并看见了绝对的圣言与圣父交谈，为那些真正的大司祭可能向唯一真天主祈祷的事而向祂祈祷。但为能与天主如此共融并向圣父祈祷，祂上了高山；然后，按马尔谷所记，\Quote{祂的衣裳变得洁白发光，如光一样，世上没有漂布者能漂得那样白。}\BibleRef{谷 9:3}或许地上的漂布者是此世的智者，他们注重辞藻，认为那是明亮纯洁的，以致连他们卑劣的思想和虚假的教条，也因所谓漂洗而显得美丽；但那位向登上高山的人显示自己衣裳闪烁发光、比它们所能漂的更明亮者，就是圣言；祂在许多人轻视的圣经表达中显露出思想的闪烁，那时耶稣的衣裳，按路加所说，变得洁白灿烂。\BibleRef{路 9:29}\par

% structure: p0080 source=h2 target=subsection reason=relative-heading-rank; base=h2

\subsection{四十、论伯多禄的话}

但让我们接下来看看\Person{伯多禄}在回答\Person{耶稣}时的心思是什么，他说：\BibleQuote{主啊，我们在这里真好；\Emph{让我们搭三座帐棚}}等等。正因如此，这些话需要非常特别的检查，因为\Person{马尔谷}以他自己的注释补充说：\BibleQuote{因为他不知道回答什么，}\BibleRef{谷 9:6}但\Person{路加}说：\BibleQuote{不知道，}他说：\BibleQuote{他所说的是什么。}\BibleRef{路 9:33}因此，你要考虑，他是否是在恍惚中说了这些话，充满了那感动他说话的那个灵，而这不能是圣神；因为\Person{若望}在\BibleBook{若望}{福音}中教导说，在救主复活之前，没有人领受过圣神，他说：\BibleQuote{因为圣神还没有赐下，因为耶稣还没有受到光荣。}\BibleRef{若 7:39}但如果圣神还没有赐下，而他，不知道自己在说什么，是在某个灵的影响下说话，那导致这些话被说出的灵，是那些尚未在十字架上被胜过、也未被公开羞辱的灵之一，关于这些灵有记载说：\BibleQuote{祂解除了\OriginalTerm{率领者与掌权者}{principalities and powers}的武装，公开羞辱了它们，仗着十字架战胜了它们。}\BibleRef{哥 2:15}但这灵或许就是那被\Person{耶稣}称为\OriginalTerm{绊脚石}{stumbling-block}的，也是在下列经文中被称为\Person{撒殚}的那一位：\BibleQuote{撒殚，退到我后面去！你是我的绊脚石。}\BibleRef{玛 16:23}但我深知，这样的话会得罪许多遇到它们的人，因为他们认为，这违背健全的理性：一个不久前还被\Person{耶稣}宣布为有福的人，因为天上的父向他启示了关于救主的事，即祂确实是\Person{耶稣}、是基督、是永生天主之子，竟会受到这样贬损的评述。但让这样的人更准确地留意有关\Person{伯多禄}和其余宗徒的陈述，看看他们如何提出请求，仿佛他们仍与那位要从仇敌手中救赎他们、并用祂自己的宝血购买他们的那一位相疏远；或者，也让那些坚持认为在\Person{耶稣}受难之前宗徒们就已成全的人，告诉我们，为什么\BibleQuote{伯多禄和他同在一起的人都因困倦而沉睡。}\BibleRef{路 9:32}但为了预先考虑随后的一些事并将之应用于当前的主题，我要反过来提出这些问题——是否有人能在\Person{耶稣}身上找到绊倒的缘由，而不涉及那使他绊倒的魔鬼的工作；以及是否有人能否认\Person{耶稣}，且是在一个小使女、一个看门人和最卑贱的人面前否认，除非有一个与他否认同在的灵，敌视那赐予的圣神和智慧（这智慧是赐给那些因自己的某种功绩而蒙天主帮助作宣认的人）。但那学会了将罪的根源归咎于罪之父魔鬼的人，不会说，没有魔鬼，宗徒们就会绊倒，或者\Person{伯多禄}会在那著名的鸡叫之前三次否认基督。但如果这样，请考虑，或许为了按自己的能力使\Person{耶稣}绊倒，并使他偏离那以带来救恩的苦难为特征的\OriginalTerm{救恩计划}{dispensation}，而这救恩计划是祂心甘情愿接受的，寻求达成那些看似有助于此目标的事，他本人在这里也仿佛想用诡计诱惑\Person{耶稣}，好像召唤祂不再屈尊俯就人、来到他们中间、为他们受死，而是与\Person{梅瑟}和\Person{厄里亚}一起留在高山上。他还许诺要搭三座帐棚，一座为\Person{耶稣}，一座为\Person{梅瑟}，一座为\Person{厄里亚}，好像一个帐棚不足以容纳三者，如果他们有需要住在帐棚里并在高山上。也许在此他也是存心不良，当他煽动那\Quote{不知道自己在说什么的人}时，他并不希望\Person{耶稣}、\Person{梅瑟}和\Person{厄里亚}在一起，而是想借口搭三座帐棚使他们彼此分开。同样，这也是一个谎言：\BibleQuote{我们在这里真好}\BibleRef{玛 17:4}，因为如果真是好事，他们也会留在那里。但如果这是谎言，你会想知道是谁导致说了这谎言；尤其是因为按照\Person{若望}所说：\BibleQuote{他说谎时，是出于自己的本性，因为他是说谎者，而且就是谎言之父}\BibleRef{若 8:44}；并且，正如没有真理能离开那说\BibleQuote{我是真理}\BibleRef{若 14:6}的那一位的工作，同样，也没有谎言能离开那真理之敌。因此，这些相反的特质，真理和谎言，仍存在于\Person{伯多禄}内；出于真理，他说：\BibleQuote{你是基督，永生天主之子}\BibleRef{玛 16:16}，但出于谎言，他说：\BibleQuote{主，愿天主可怜祢，这事绝不会临到祢身上}\BibleRef{玛 16:20}，以及\BibleQuote{我们在这里真好}\BibleRef{玛 17:4}。但如果有人不承认\Person{伯多禄}说这些话是出于任何邪恶的灵感，而是出于他自己的自由选择，并且问他将如何解释\BibleQuote{不知道他所说的是什么}\BibleRef{路 9:33}和\BibleQuote{因为他不知道回答什么}\BibleRef{谷 9:6}，他会说，在前一种情况下，\Person{伯多禄}认为承认永生天主之子、基督，已由圣父启示给他的那一位，应该被杀害，是可耻的、不配\Person{耶稣}的事；而在目前的情况下，因为他看见了\Person{耶稣}的两种形态，以及在显圣容时那更卓越得多的形态，对此非常满意，就说留在这山上真好，为的是使他本人和与他在一起的人，能因目睹\Person{耶稣}显圣容、祂的面容发光如太阳、祂的衣服洁白如光而欣喜，除此之外，还能永远在荣耀中看见那些他们曾一度在荣耀中看见的人，即\Person{梅瑟}和\Person{厄里亚}；并且他们能因听见他们彼此交谈和来往时所说的话而欣喜，即\Person{梅瑟}和\Person{厄里亚}与\Person{耶稣}，以及\Person{耶稣}与他们之间的交谈。\par

% structure: p0082 source=h2 target=subsection reason=relative-heading-rank; base=h2

\subsection{四十一、同一段落的喻意解释。}

但是，既然我们尚未将精力用在喻意解释这一段，而只是按字面探究了这些事，那么我们就应当据此考虑：前述的伯多禄和雷霆之子被带上真理教义之山，看见了耶稣以及与他一同显现在荣耀中的梅瑟和厄里亚的显容，他们是否愿意在自己内为即将居住在他们内的天主圣言、为其被看见的荣耀中的法律、以及为预言耶稣将要完成的逝世，搭设帐幕。\BibleRef{路 9:31}伯多禄作为热爱沉思生活的人，并认为沉思生活中的乐趣优于人群中的喧嚣生活，出于有益于那些渴望它的人的意愿，说：\BibleQuote{我们在这里真好！}\BibleRef{玛 17:4}但是，既然\BibleQuote{爱不求自己的益处}\BibleRef{格前 13:5}，耶稣没有做伯多禄以为好的事；因此他从山上下来，到那些不能上山看见他显容的人那里，好让他们能以他们能看见他的那种形态看见他。因此，一个拥有\BibleQuote{不求自己益处的爱}\BibleRef{格前 13:5}的义人，应当摆脱一切束缚，却使自己成为所有下面的人的奴仆，好能赢得更多的人。\BibleRef{格前 9:19}但是，有人可能会针对我们就伯多禄的恍惚和恶神的作用所提出的关于\Quote{他不知道所说的是什么}\BibleRef{路 9:33}这话的解释，不接受我们的理解，而说保禄提到有些人\BibleQuote{想要做法律的教师}\BibleRef{弟前 1:7}，他们不知道自己所谈论的，虽然他们不清楚地阐明所说之事的性质，也不理解其意义，却对自己不知道的事妄下断言。伯多禄的情感也是如此，因为他没有领悟到就耶稣和山上显现的梅瑟与厄里亚的安排而言什么是好的，就说\BibleQuote{我们在这里真好}等等，\Quote{他不知道所说的是什么}，\Quote{因为他不知道该说什么}\BibleRef{路 9:33}；因为如果\BibleQuote{智慧人必从自己的口中明白事理，并在嘴唇上携带谨慎}\BibleRef{箴 16:23}，那么不是这样的人就不从自己的口中明白事理，也不理解他所说之事的性质。\par

% structure: p0084 source=h2 target=subsection reason=relative-heading-rank; base=h2

\subsection{四十二、\Quote{光明的云彩}的意义。}

接着这些话是：\BibleQuote{\Emph{他还在说话的时候，看啊，也有一片光明的云彩遮蔽了他们}}\BibleRef{玛 17:5}等。现在，我认为天主想要劝阻伯多禄不要搭三座帐幕（按他的选择，他打算住在里面），便展示了一个更好的、可以说更卓越的帐幕，即云彩。因为帐幕的功能是遮蔽其中的人，保护他；光明的云彩遮蔽了他们，天主便造了一个更神圣的帐幕，它是光明的，好成为他们未来复活的预像；因为光明的云彩遮蔽义人，他们立即被它保护、照亮并照耀。但是，这遮蔽义人的光明云彩是什么呢？或许是父的德能，从中发出圣父的声音，为圣子作证，称他是爱子和喜悦的，并劝勉那些在它荫下的人听他，而不听其他的人？正如他自古以来说话，也总是借着他所愿意的说话。或许，圣神也是光明的云彩，遮蔽义人，并预言在他内做工的天主之事，说：\BibleQuote{这是我的爱子，我所喜悦的；}但我也敢说，我们的救主是一朵光明的云彩。因此，当伯多禄说\BibleQuote{我们在这里搭三座帐幕}时……一座来自圣父本身，一座来自圣子，以及一座来自圣神。因为来自圣父、圣子和圣神的光明云彩遮蔽耶稣的真门徒；或者，云彩遮蔽福音、法律和先知，对于能看见它在福音、法律和先知中的光的人，它是光明的。但或许从云彩中发出的声音对梅瑟和厄里亚说：\BibleQuote{这是我的爱子，我所喜悦的，你们要听他，}因为他们渴望看见人子，听他，并看见他如同在荣耀中一样。或许它也教导门徒：按字面意义是天主子、他的爱子、他所喜悦的那一位——门徒尤其应当听从他——就是当时被看见、被显容、他的面容发光如太阳、衣服洁白如光的那一位。\par

% structure: p0086 source=h2 target=subsection reason=relative-heading-rank; base=h2

\subsection{四十三、梅瑟和厄里亚与耶稣的关系。沉默的诫命。}

但这些事以后，经上记载说：当三位宗徒听见云中有声音为圣子作证时，他们承担不住那声音的荣耀和其上所寄托的能力，就\Quote{俯伏在地} \BibleRef{玛 17:6}，恳求天主；因为他们因这超自然的景象和从景象中所说的话而极其害怕。但请思考，你是否也可以就这段经文的细节说：门徒们明白了天主子曾与梅瑟交谈，并且正是祂曾说过\Quote{没有人看见我的面还能存活} \BibleRef{出 30:20}；此外，他们领受了天主关于祂的见证，因不能承受圣言的光辉，就在天主大能的手下自卑；\BibleRef{伯前 5:6}但在圣言触摸之后，他们举目，只见耶稣，别无他人。\BibleRef{玛 17:8}梅瑟——法律，和厄里亚——先知，只与耶稣的福音成为一体；他们并不像从前三者那样各自存留，而是三者成为一体。但请你与我一起就奥秘之事思考这些；因为就字面的单纯意义而言，梅瑟和厄里亚在光荣中显现并与耶稣交谈后，便回到了他们来的地方，也许是为了将耶稣与他们所说的话传达给那些将要因祂受益的人，几乎是立即地，即在受难之时，当许多已长眠的圣人的身体，坟墓被打开，要前往那座真正神圣的城——不是耶稣为之哭泣的耶路撒冷——并在那里显现给许多人。\BibleRef{玛 27:52}但在山上的事工之后，当门徒们从山上下来，好能来到群众中为天主子服务关乎百姓的救恩时，耶稣命令门徒说：\Quote{在人子从死者中复活以前，不要将所见告诉任何人。} \BibleRef{玛 17:9}但这句\Quote{不要将所见告诉任何人}的说话，与上文所查考的相似，即\Quote{祂嘱咐门徒不要告诉任何人祂是基督。} \BibleRef{玛 16:20}因此，在那段经文中所说的话，对我们当前的段落也有用；因为耶稣也愿意，依照这些，祂荣耀的事在受难之后的荣耀以前不可被传讲；因为那些听见的人，尤其是群众，若看见曾被如此光荣化的祂被钉十字架，会受到伤害。因此，既然祂在复活中的光荣化与祂的显圣容相似，也与祂的面貌如同太阳的异象相似，为此祂愿意这些事在祂从死者中复活后，才由宗徒们传讲。\par

\SourceRecord{Translated by John Patrick.；Ante-Nicene Fathers；Vol. 9.；Edited by Allan Menzies.；Buffalo, NY: Christian Literature Publishing Co.,；1896.；<http://www.newadvent.org/fathers/101612.htm>.}

% structure: p0001 source=h1 target=section reason=page-title

\section{玛窦福音释义（第十三卷）}

% structure: p0002 source=h2 target=subsection reason=relative-heading-rank; base=h2

\subsection{一、洗者与厄里亚的关系。论轮回说}

\Quote{\Emph{门徒问祂说：\InnerQuote{那么，为什么经师说厄里亚必须先来呢？}}}\BibleRef{玛 17:10}那时与耶稣一同上山的门徒记起了经师关于厄里亚的传说，即在基督来临之前，厄里亚会先来，为那些将要接受祂的人预备心灵。但在山上显现的景象中，厄里亚出现了，这与所说的话似乎不一致，因为他们看来厄里亚是在耶稣之后而不是之前来的；所以他们问这些话，认为经师说谎了。救主回答说，祂并没有否定关于厄里亚的传统，而是说在基督来临之前厄里亚另有第一次来临，是经师所不知道的；并且关于这一点，由于他们不认识他，他们\BibleQuote{随意待他}\BibleRef{玛 17:12}，好像他们参与了将他下在监狱并由黑落德杀害的事；然后祂说，正如他们对待厄里亚一样，祂也要在他们手中受害。\BibleRef{玛 17:12}门徒问的是关于厄里亚的事，救主回答，但他们听了以后，就明白\BibleQuote{厄里亚已经来了}这句话以及救主随后说的话是指着洗者若翰说的。\BibleRef{玛 17:13}以上是关于这段文字的解释。现在，按我们的能力，让我们也探究其中深藏的奥义。在我看来，此处并非指厄里亚的灵魂，以免我陷入\OriginalTerm{轮回}{transmigration}的教义中，这教义与天主的教会是格格不入的，并非宗徒所传授，也非圣经任何地方所教导的；因为它也与说\BibleQuote{所见的是暂时的}\BibleRef{格后 4:18}以及\BibleQuote{这个时代将有终结}的经文相悖，也与\BibleQuote{天地要过去}\BibleRef{玛 24:35}和\BibleQuote{这世界的景象正在逝去}\BibleRef{格前 7:31}以及\Quote{诸天将要毁灭}等经文相违背。因为假设从开始到世界末了的事物秩序中，同一个灵魂可以两次进入身体，那么它进入身体的缘由是什么？如果是因为罪而两次进入身体，为什么不会是三次，甚至多次呢？因为惩罚，就今生和今生所犯的罪而言，只能通过轮回的方式予以执行。但如果这个推论成立，那么也许永远不会有灵魂不再轮回的时候：因为它总会因先前的罪而居住在身体里；这样，世界便不会有败坏的时候，那时\BibleQuote{天地要过去}。\BibleRef{玛 24:35}并且，如果按照这个假设，完全无罪的人不会通过出生进入身体，那么你认为经过多长时间才能找到一个完全纯洁、无需轮回的灵魂呢？但尽管如此，如果某个灵魂总是从固定的灵魂数目中被移除，不再回到身体中，那么经过无限的时间后，出生将会停止；世界将只剩下一个、两个或少数几个灵魂，在他们完美之后世界将毁灭，因为进入身体的灵魂供应已经断绝。但这不符合圣经；因为圣经知道在世界毁灭时会有许多罪人。这从\BibleQuote{然而人子来的时候，遇得见世上有信德吗？}\BibleRef{路 18:8}这句话中显明。我们在玛窦福音中也读到：\BibleQuote{诺厄的日子怎样，人子降临也要怎样；因为洪水以前的日子，他们怎样}等等\BibleRef{玛 24:37}。对于那时存在的人，他们将因自己的罪受到惩罚，但并非通过轮回；因为如果他们还在犯罪时被抓住，他们要么在死后受到另一种形式的惩罚——这样就有两种普遍的惩罚形式，一种是通过轮回，另一种是在这种身体之外，并且他们必须说明这些惩罚的原因和区别——要么他们不会受到惩罚，好像世界末日时留下的人立刻丢弃了他们的罪；或者，更可能的是，对于在身体内犯罪的人只有一种惩罚形式，即在他们之外受苦，也就是在这个生命的构造之外，按照他们罪行的应得。但洞察事物本质的人都知道，这些说法都足以推翻轮回的教义。但如果那些引进轮回教义的希腊人，出于必然，遵循与此符合的学说，不承认世界将归于败坏，那么当他们直面明确宣称世界将毁灭的圣经时，要么他们不信圣经，要么他们会发明一系列关于末世解释的论点；即使他们愿意，他们也无法做到。此外，我们还要对那些敢于断言世界不会毁灭的人说：如果世界不毁灭，而要存在无限的时间，那么就不会有天主知道一切尚未发生的事。但如果祂或许只部分知道，要么祂在每件事发生前就知道了，要么祂知道某些事，然后知道其他事；因为无限的事物不可能被那种以限定已知事物为本性的知识所把握。由此得出结论，不可能有关于所有事物的预言，因为所有事物都是无限的。\par

% structure: p0004 source=h2 target=subsection reason=relative-heading-rank; base=h2

\subsection{二、\Quote{厄里亚的精神和能力}——而非灵魂——在洗者若翰身上。}

我认为有必要花一些时间考察灵魂转世的教义，因为有些人怀疑所讨论的灵魂在厄里亚和若翰中是同一个，在前者被称为厄里亚，在后者被称为若翰；并且，他被称为若翰并非出于天主，正如向匝加利亚显现的天使所说的话所表明的：\Quote{匝加利亚，不要害怕！因为你的祈祷已蒙应允，你的妻子依撒伯尔要给你生一个儿子，你要给他起名叫若翰。}\BibleRef{路 1:13}以及匝加利亚在写字板上写下那出生者应叫若翰后，他的口立时开了。\BibleRef{路 1:63}但如果那是厄里亚的灵魂，那么当他第二次降生时，就应该被称为厄里亚；或者，名字的改变应当有某种理由，如同亚巴郎和亚伯拉罕、撒辣和撒辣、雅各伯和以色列、西满和伯多禄的情形。然而，即使这样，他们的论点也站不住脚；因为在上述情形中，名字的改变发生在同一生命中。但有人可能会问，如果厄里亚的灵魂并非先在提斯比人厄里亚身上，其次在若翰身上，那么救主在两个中所称为厄里亚的究竟是什么？我说，加俾厄尔在向匝加利亚说话时，已经暗示了厄里亚和若翰中相同的东西是什么；因为他说：\Quote{他要使许多以色列子民转向上主、他们的天主；他要以厄里亚的精神和能力，在他前面先行。}\BibleRef{路 1:16}请注意，他没有说以厄里亚的\Quote{灵魂}，那样灵魂转世之说或许有些根据，而是说\Quote{以厄里亚的精神和能力}。因为圣经清楚区分精神和灵魂，如\Quote{愿赐平安的天主亲自完全圣化你们，将你们整个人的灵魂和肉身，在我们的主耶稣基督来临时，保持得无瑕可指；}\BibleRef{得前 5:23}以及\BibleBook{达尼尔}{书}中根据\WorkTitle{七十士译本}的经文\Quote{义人的诸精神和灵魂，请赞美上主}，这显示了精神和灵魂的区别。因此，厄里亚被称为若翰并非因为灵魂，而是因为精神和能力，这丝毫不与教会的教导相冲突，尽管它们先前在厄里亚身上，后来在若翰身上；并且\Quote{先知的神魂是顺从先知的，}\BibleRef{格前 14:32}但先知的灵魂并不顺从先知，而\Quote{厄里亚的精神停在厄里叟身上。}\BibleRef{列下 2:15}但我们应当查考厄里亚的精神是否与厄里亚身上的天主圣神相同，或者它们彼此不同，并且厄里亚身上的精神是否是一种超自然的东西，不同于每个人内在的精神；因为宗徒清楚表明，即使天主圣神在我们内，它也与每个人内在的精神不同，他在某处说：\Quote{圣神亲自和我们的心神一同作证：我们是天主的子女；}\BibleRef{罗 8:16} 又在别处说：\Quote{除了人内里的心神外，有谁能知道那人的事呢？同样，除了天主圣神外，也没有人知道天主的事。}\BibleRef{格前 2:11} 但对于厄里亚所说的事，不要惊奇，如果正如他所经历的、与所有记载的圣人们不同的事那样，他乘旋风被提升天，\BibleRef{列下 2:11} 他的精神也具有某种特选的卓越，以致不仅停在厄里叟身上，而且在他出生时与若翰一同降临；并且若翰一方面\Quote{自母胎中就充满了圣神}，另一方面\Quote{以厄里亚的精神和能力在基督前面先行}。因为同一个人的身上可能同时有多个精神，不只是较坏的，也有较好的。达味因此恳求被自由的精神所扶持，并求在他内里更新正直的精神。但如果为了使救主将\Quote{智慧与聪敏的精神，超见与刚毅的精神，知识与敬畏的精神}\BibleRef{依 11:2} 赐给我们，他也充满了敬畏上主的精神；那么也可以认为这些不同的善灵同在一个人身上。我们提出这一点，也是因为若翰\Quote{以厄里亚的精神和能力}\BibleRef{路 1:17} 在基督前面先行，以便\Quote{厄里亚已经来了}\BibleRef{玛 17:12} 这句话可以指若翰内里的厄里亚的精神；正如与祂一同上山的三个门徒所理解的，祂对他们所说的是关于洗者若翰。\BibleRef{玛 17:13} 厄里叟身上只停下了厄里亚的精神，但若翰不仅以精神，而且以厄里亚的能力先行。因此，厄里叟不能被称为厄里亚，但若翰就是厄里亚本人。但如果需要引述经师们说厄里亚必须先来的经文，请听玛拉基亚说：\Quote{看，我要派遣先知厄里亚到你们这里来}等等，直到\Quote{免得我来击打这地，完全毁灭}\BibleRef{拉 4:5} 这些话语似乎表明，厄里亚要以某些圣言和灵魂中的意向来预备基督光荣的来临，使那些为此最合适的人得以预备，因为若没有厄里亚预先预备，地上的人因那荣耀的卓越性无法承受。同样，在此处，我理解的厄里亚不是那先知的灵魂，而是他的精神和能力；因为藉着这些，一切将要恢复，\BibleRef{玛 17:11} 以致当一切恢复后，并因那恢复而能接受基督的荣耀时，那位将要在荣耀中显现的天主子才能与他们同住。但如果厄里亚在某种意义上也是一低于\Quote{那在起初就与天主同在的圣言，天主圣言}\BibleRef{若 1:1} 的言语，那么这言语也可能作为预备性的训练临到那由它预备的人民，使他们得以受训接纳圆满的圣言。但有人可能会提出疑问：厄里亚的精神和能力在若翰身上是否遭受了所遭受的事，正如所说：\Quote{他们任意待了他。}\BibleRef{玛 17:12} 对此，一方面，更简单地回答是，那施助者因爱而与被助者一同受苦，这并不奇怪；耶稣确实说：\Quote{因了软弱的人，我成了软弱的；因了饥饿的人，我成了饥饿的；因了口渴的人，我成了口渴的。} 另一方面，从更深的意义上说，经文不是\Quote{他们任意待了他内里的他}，因为那遭受的事依赖于厄里亚的精神和能力，若翰的灵魂绝非厄里亚；并且可能身体也依赖于它们。因为灵魂在身体中以一种方式，精神和能力以另一种方式；而义人的身体在这些更好部分中以一种方式，如同依赖并附着于它们；但\Quote{属于血肉的人不能中悦天主；你们却不属于血肉，而是属于神，如果天主的圣神住在你们内；}\BibleRef{罗 8:8} 罪人的灵魂在血肉中，但义人的灵魂在神中。同样地，还可以进一步查考：\Quote{但他们任意待了他}\BibleRef{玛 17:12} 这句话是指谁？是指门徒们所询问并说的经师吗？\Quote{那么，为什么经师们说厄里亚必须先来？}\BibleRef{玛 17:10} 但若翰是否从经师们手中遭受了什么，完全不清楚，除非他们不信他；或者，如我们之前所说，他们是黑落德敢于加害于他的共犯。但另一个人可能会说：\Quote{但他们任意待了他} 这句话不是指经师，而是指黑落狄雅和她的女儿，以及黑落德，他们任意待了他。接着的\Quote{人子也要同样受他们的害}\BibleRef{玛 17:12} 如果前者指经师，也可指经师；但如果前者指黑落德、黑落狄雅和她的女儿，那么第二个段落也指他们；因为黑落德似乎也投票赞成耶稣该死，也许他的妻子也参与谋害祂。\par

% structure: p0006 source=h2 target=subsection reason=relative-heading-rank; base=h2

\subsection{三、论癫癇者。}

\BibleQuote{\Emph{当他们来到群众那里时，有一个人前来跪在祂跟前说：主，可怜我的儿子罢。}} \BibleRef{玛 17:14} 那些受苦的人，或受苦者的亲属，与群众在一起；因此，当祂分施了超越群众的事物之后，祂下降到他们中间，好让那些因压制他们灵魂的疾病而无法上升的人，当圣言从更高之处降到他们那里时，能获得益处。但我们应当查考，在哪些疾病中，受苦者自己相信并为自己的治愈祈祷，在哪些疾病中，别人为他们这样做，例如，百夫长为他的仆人，王臣为他的儿子，会堂长为女儿，客纳罕妇人为她被魔鬼折磨的小女儿，以及现在这个人跪在祂面前为他的癫癇儿子。与此一并，你当探究何时救主自己主动医治，不经任何人请求，例如瘫子；因为这些治愈，为此目的彼此比较，并一同考察，将向能听见\BibleQuote{在奥秘中隐藏的天主智慧} \BibleRef{格前 2:7}的人显示许多关于不同灵魂疾病的教义，以及治愈它们的方法。\par

% structure: p0008 source=h2 target=subsection reason=relative-heading-rank; base=h2

\subsection{四、属灵的癫癇者。}

但由于我们目前的目的不是查考每一个病例，而是查考眼前的经文，让我们采用寓意解释，思考谁可能是那癫癇者，谁是为他祈祷的父亲，以及患者不是常常而是屡次跌倒，有时跌入火中，有时跌入水中的含义，以及他不能被门徒治愈而只能被耶稣亲自治愈的事实。因为如果我们的救主当时在百姓中治愈的每一种疾病和每一种虚弱，都对应灵魂中的不同失调，那么合理的是：瘫子象征灵魂瘫痪的人，他们使灵魂瘫痪在身体中；盲人象征那些在唯有灵魂可见的事物上盲目的人，他们才是真正的盲人；聋子象征那些在接纳救恩之言上耳聋的人。按同样原则，也须查考关于癫癇者的事。这种病每隔相当长的时间发作，患者在不发作时似乎与健康人无异。你可以在某些灵魂中发现类似的失调，它们常被认为在节制和其他德行上是健康的；然后，有时，仿佛被一种源于情欲的癫癇所攫住，它们从似乎站立的位置跌倒，被此世的欺骗和其他私欲所拖走。因此，如果你说这些人可以说是属灵的癫癇者，被\BibleQuote{天空属灵的邪恶势力}\BibleRef{弗 6:12}打倒，并常常生病，当情欲攻击他们的灵魂时；他们时而跌入燃烧的火中，如欧瑟亚所说，他们成了奸淫者，\BibleQuote{像从燃烧的火焰加热的烤盘}\BibleRef{欧 7:4}；时而跌入水中，当水中一切巨龙之王将他们从他们似乎自由呼吸的领域抛下，使他们落入人生命海洋的波涛深处。我们在智慧书上关于义人的均衡气质所说的话将支持我们对癫癇者的这种解释：\BibleQuote{虔敬人的言论常是智慧}，但关于我们所说的，\BibleQuote{愚人如月亮一样变化}\BibleRef{德 27:11}。有时你甚至可以在这样的人身上看到冲动，这些冲动可能使那些不注意他们缺乏压舱物的人称赞他们，以致他们会说在他们身上是满月或近乎满月。你又会看到他们里面的光似乎减弱——因为它不是白昼的光，而是夜晚的光——衰退到如此程度，以致在他们身上似乎看到的光不再存在。但至于那些最初为事物命名的人是否因此将癫癇称为月癇，你自己判断。\par

% structure: p0010 source=h2 target=subsection reason=relative-heading-rank; base=h2

\subsection{五、聋哑的魔鬼。}

现在，那癫癇者的父亲——或许是被分配给他的天使，如果我们说每个灵魂都隶属于某位天使的话——向灵魂的医生为他的儿子祈祷，求祂医治那些不能通过门徒中低级言语治愈的失调。但被圣言赶出的聋哑魔鬼，必须寓意地理解为非理性的冲动，甚至向那看似良善之物的冲动，以致一个人曾凭非理性冲动所做而看似善的事，他不再非理性地做，而是按照耶稣教导的理性。受此启发，保禄也说：\BibleQuote{我若有全备的信德，甚至能移山} \BibleRef{格前 13:2}；因为凡有像芥子那样全备信德的人\BibleRef{玛 17:20}，不仅移一座山，也移几座类似的山；因为虽然信德被人轻视，看起来非常微小和可鄙，但当它遇到好地，即那能合适接受这种子的灵魂时，它就成为一棵大树，以致没有翅膀的东西，而只有属灵有翅膀的天空飞鸟，才能栖在这如此大的信德枝上。\par

% structure: p0012 source=h2 target=subsection reason=relative-heading-rank; base=h2

\subsection{六、月亮星辰对人的影响。}

现在，让我们仔细留意这段经文的字面意思，首先探究：那被投入黑暗、被不洁且聋哑之灵压制的人，为何被称为\Quote{月癫}？\Quote{月癫}这一称呼，又从何得名于天上仅次于太阳的巨光，就是天主命定\BibleQuote{掌管黑夜}的？\BibleRef{创 1:16} 让医生们从生理角度讨论此事吧，因为他们认为这并非不洁之灵作祟，而是身体失调；他们探究事物的本性，会说：头部湿润的体液因与月光（月光本性湿润）有一种感应而运动。而我们，既相信福音所载：此病由不洁、聋哑之灵在患者身上所致，又看到那些惯于像埃及术士般承诺治愈此类病症的人，有时似乎真能奏效，我们将说：或许，为诽谤天主的创造，为使\BibleQuote{不义得以高谈阔论，使他们口出狂言攻击高天}，这不洁之灵便观察月亮的某些位相，从而使人根据患者在某一月相发作的观察，显得如此大恶的根源并非聋哑恶魔，而是天上那被命定\BibleQuote{掌管黑夜}的巨光——而这光根本无力在人中引起这般紊乱。然而，凡说地上一切紊乱——无论是总体还是每一细节——皆源于星辰运行的人，都是\BibleQuote{高声谈论不义}；他们实在\BibleQuote{口出狂言攻击高天}，因为他们声称有些星辰有恶意影响，有些有益影响。但宇宙的天主并未创造任何星辰来行恶，正如耶肋米亚在\WorkTitle{哀歌}中所记：\BibleQuote{从主的口中会出高贵与良善的事。}很可能，这所谓产生月癫的不洁之灵，观察月相，以便作用于那因某些原因被交付于它、且不使自己配受天使守护的人；同样，也有其他邪灵和魔鬼，在其余星辰的某些位相发作，如此一来，不仅月亮，连其余星辰也被那些\BibleQuote{高声谈论不义}的人诽谤了。因此，值得听听那些占星家，他们把各种疯狂和附魔的起源都归因于月相。显然，那些患所谓月癫的人有时落水，有时也落火（虽较少见，却确会发生），而且此病极难治愈，以至于能驱魔的人有时对此无能为力，有时却藉禁食、祈祷和更多苦工而成功。但你要探究：在灵体中和在人中是否有这样的失调？以致有的灵说话，有的缄默无声；有的听见，有的耳聋。正如他们不洁的原因在于他们自身，同样，因自由意志，他们也被判定为无声和耳聋。因为某些人若留心先知在圣神感动下所说的祈祷（其中论及某些罪人说：\BibleQuote{愿说谎的嘴唇沉默不语}），也会遭此定罪。或许，那些滥用听觉、接纳虚妄之言的人，会被那曾说\BibleQuote{谁使人口哑耳聋}的\BibleRef{出 4:11} 弄成耳聋，好使他们不再听虚妄之事。\par

% structure: p0014 source=h2 target=subsection reason=relative-heading-rank; base=h2

\subsection{七、信德的能力。}

但当救主说：\BibleQuote{无信且败坏的世代哪}\BibleRef{玛 17:17}，祂表明那违反本性的邪恶，从败坏中悄然侵入，使我们变得败坏。然而，我想，救主针对地上全人类（因他们的邪恶和祂的滞留而受压迫）说：\Quote{我同你们在一起要到几时呢？}我们已部分讨论了\BibleQuote{如果你们有信德像芥子那样大，你们向这座山说}等等的话\BibleRef{玛 17:20}；但在此处，我们仍要说说那些看来能增进明晰的事。依我之见，此处所说的\Quote{山}，是指那些存在于极大邪恶洪流中的敌对势力，它们如同盘踞在某些人的灵魂中。因此，什么时候有人具有全备的信德，以致不再怀疑圣经所载的任何事，并拥有像亚巴郎那样的信德（他如此信天主，以至于他的信德被算为正义），他就有了像芥子那样大的全备信德；那时，他将向这座山——我指的是那被称为月癫者身上的聋哑之灵——说：\BibleQuote{移开此地}，显然是指从受苦者身上移开，也许移到深渊，它就会移开。宗徒（我想，是从此处出发）权威地说：\BibleQuote{如果我有全备的信德，以致能移山}\BibleRef{格前 13:2}；因为不只一座山，而是好几座类似的山，那拥有像芥子那样大之全备信德的人都能移开；对拥有如此大信德的人，没有不可能的事。\BibleRef{玛 17:20} 我们也要注意：\BibleQuote{这类魔鬼，若不藉祈祷和禁食，是赶不出去的。}\BibleRef{玛 17:21} 这样，若有时我们必须从事治愈此类患者，我们就不必驱魔、提问或对不洁之灵说话，好像它听得见一样，而要专心祈祷和禁食，藉着为患者祈祷而成功，并藉着自己的禁食，将不洁之灵从他身上赶出去。\par

% structure: p0016 source=h2 target=subsection reason=relative-heading-rank; base=h2

\subsection{八、\Person{耶稣}预言祂被\Quote{交付}于人手。}

\BibleQuote{ \Emph{他们住在加里肋亚时，耶稣对他们说：\InnerQuote{人子将被交于人手。}} } \BibleRef{玛 17:22} 但这些事似乎与那些事有相同的效果，即\BibleQuote{耶稣开始指示祂的门徒说：祂必须往耶路撒冷去，由长老、司祭长和经师受许多苦。} \BibleRef{玛 16:21} 但并非如此；因为\BibleQuote{指示门徒说祂必须往耶路撒冷去，由长老、司祭长和经师受许多苦}，以及受苦之后\BibleQuote{被杀}，被杀之后\BibleQuote{第三日复活}，与他们在加里肋亚时对他们所说的——我们之前并未得知——即人子\BibleQuote{将被交付}，并不是同一件事；因为\InnerQuote{被交付}的事前面没有提及，但现在又说祂\BibleQuote{将被交于人手}。\BibleRef{玛 17:22} 至于这些事，让我们探究祂将被交于何人之手；因为那里我们得知祂将受谁之苦，以及将在何处受苦；但在这里，我们还得悉，虽然祂的诸多苦难是由前述之人造成的，但他们并非祂受苦的主要缘由，而是那将祂交于人手中的一位或几位。因为有人会说，宗徒解释此事时，论及天主说：\BibleQuote{祂没有怜惜自己的儿子，反而为我们众人把祂交出来；} \BibleRef{罗 8:32} 但圣子也为我们舍身至死，所以祂不仅被圣父交付，也被祂自己所交付。但另有人不仅如此说，且汇集各处经文，将会说圣子首先被天主交付——然后即将受试探，然后交战，然后为人受苦，甚至为全世界受苦，以除掉世人的罪 \BibleRef{若 1:29} ——交付给这世代的王侯及其他王侯，然后由他们交在要将祂杀害的人手中。以约伯的事为例。\BibleQuote{看，凡他所有的，我都交在你手中，只是不可加害他本人；} \BibleRef{约 1:12} 此后，他好像被魔鬼交付给他的王侯们，即那些掳掠战俘的、骑兵、从天降下的火、从旷野刮来摧毁他房屋的大风。\BibleRef{约 1:15} 但你要考虑，正如他将约伯的财产交付给那些掳掠他们的人和骑兵，他也将其交付给某种力量，这力量隶属于\BibleQuote{那在空中掌权的首领，即现今在悖逆之子身上运行的邪神}，\BibleRef{弗 2:2} 为要使那从那里降在约伯羊群上的火，对那向约伯报告\BibleQuote{ \InnerQuote{火从天降，烧灭了羊群和牧人} }\BibleRef{约 1:16}的人而言，似乎是自天而降的。同样，你也要查考那从旷野而来、袭击房屋四角的大风，是否也是属于那些在魔鬼手下、魔鬼将约伯子女的宴席交付给他们的势力之一，为要使房屋倒塌在义人子女身上，使他们死亡。那么，我们就承认，如同约伯的例子，圣父首先将圣子交付给敌对势力，然后他们将祂交付于人手，这些人中也包括犹达斯，在吃了那饼以后 \BibleRef{若 13:27} 撒殚进入了犹达斯，他以比犹达斯更具权威的方式将祂交付。但你要小心，免得将圣父将圣子交付给敌对势力，与敌对势力将救主交付于人手这两件事相提并论，就以为所谓的\InnerQuote{交付}在两方是相同的。因为你要明白，圣父因祂对人的爱，为了我们众人将祂交付出来；而敌对势力将救主交付于人手时，并非意图为拯救某些人，而是尽他们所能——因为他们谁也不认识\BibleQuote{那隐藏在奥秘中的天主智慧}，\BibleRef{格前 2:7}——他们将祂交出，使之被处死，好让祂的仇敌死亡将祂收归权下，如同那些在亚当内死去的人一样；\BibleRef{格前 15:22} 而那些杀害祂的人之所以这样做，也是因为他们受那些确愿耶稣被死亡辖制者的意愿所塑造。我认为也有必要探究这些事，因为当耶稣被交于人手中时，祂并不是由人交于人手，而是由那些圣父为了我们众人将圣子交付给它们的势力所交；并且在祂被交付、落在那接受祂之势力的权下时，祂毁灭了那掌有死亡权柄的；因为\BibleQuote{祂借着死亡，毁灭了那握有死亡权势的，即魔鬼，并释放了那些因死亡的恐惧，一生处于束缚中的人们。} \BibleRef{希 2:14}\par

% structure: p0018 source=h2 target=subsection reason=relative-heading-rank; base=h2

\subsection{九、\Person{撒殚}与\Person{耶稣}的\Quote{交付}。}

现在我们应当认为，魔鬼拥有死亡的权柄——并非那普遍而无差别的死亡，按照这死亡，由灵魂与肉体所组成的人，当其灵魂与肉体分离时就死去；而是那与说\BibleQuote{我是生命}\BibleRef{若 14:6} 的一位相反且敌对的死亡，按照这死亡，\BibleQuote{犯罪的人，它必要死亡}\BibleRef{则 18:4}。但将祂交在人手中的并非天主，救主清楚宣告说：\BibleQuote{我的国度不属于这世界；若是我的国度属于这世界，我的仆役必要争战，使我不致被交给犹太人}\BibleRef{若 18:36}。因为当祂被交给犹太人时，祂并非被自己的仆役所交，而是被这世代的君王所交，那君王对属于无形领域的权势、即敌对人类的国度宣称：\BibleQuote{这一切我都要赐给你，如果你俯伏朝拜我}\BibleRef{玛 4:9}。因此我们也应当认为，关于它们曾论及说：\Quote{地上的君王齐集，首领聚在一起，反对主和祂的受傅者}。诚然，那些君王和那些首领齐集并聚在一起反对主和祂的受傅者；但我们，因祂被他们交在人手中并被杀害而得益，就说：\Quote{让我们挣断他们的捆绑，从我们身上抛开他们的轭}。因为当我们与基督的死亡同化时，我们便不再受地上君王的捆绑，如我们所说，也不受聚在一起反对主的这世代首领的轭。为此，\BibleQuote{圣父没有怜惜自己的圣子，反而为我们众人把祂交出来}\BibleRef{罗 8:32}，为使那些拿住祂并交在人手中的，被住在高天之上的嘲笑，被主讥讽，因他们出乎意料地，从圣父那里接受了那第三日复活、消灭了祂的仇敌死亡、并使我们不仅相似祂死亡的肖像也相似祂复活的肖像的圣子；这圣子却反倒毁灭了他们自己的国度和权能；藉着祂，我们走在新生命中\BibleRef{罗 6:4}，不再坐在\BibleQuote{黑暗和死影之地}\BibleRef{玛 4:16}，因为我们有上升的光照在我们身上。但当救主说：\BibleQuote{人子将被交在人手中，他们要杀害祂，第三日祂要复活}\BibleRef{玛 17:22}时，他们\BibleQuote{非常忧伤}\BibleRef{玛 17:22}，只注意祂将交在人手中并被杀害，视为阴郁而应哀伤之事，却不注意祂第三日复活，因为祂无需更长时间\BibleQuote{藉着死亡来废除那掌有死亡权柄的}\BibleRef{希 2:14}。\par

% structure: p0020 source=h2 target=subsection reason=relative-heading-rank; base=h2

\subsection{十、关于那些索要半银元的人。}

\BibleQuote{他们来到葛法翁时，那些收半银元的人来到伯多禄跟前}\BibleRef{玛 17:24}。有些地上君王，他们的儿子不纳丁税或贡赋；另有别的人，与他们的儿子不同，是地上君王的陌生人，地上君王从他们那里收取丁税或贡赋。在地上君王中，他们的儿子如同在父亲面前一般自由；但那些对他们是陌生人的人，虽在尘世之外的事物上是自由的，在那些管辖并奴役他们的人面前却如同奴隶；如同埃及人管辖以色列子民，严重苦害他们的生命并强暴地奴役他们\BibleRef{出 1:13}。为了那些处于与希伯来人奴役相应奴役中的人，天主子只取了奴隶的形状\BibleRef{斐 2:7}，未作任何污秽或奴役的工作。因此，既然祂具有那奴隶的形状，祂所缴纳的丁税和贡赋与其门徒所缴纳的并无不同；因为同一样钱币就足够了，正是为耶稣和祂的门徒所付的那一枚。但这钱币并不在耶稣的家中，而是在海中，在海里一条鱼的口中；按我的判断，这鱼当它上来并被捕入伯多禄的网中就得了益处，伯多禄成了捕人的渔夫；在那网中有这比喻称为鱼的东西，为要从它取得带有凯撒像的钱币，它便列于那些被学会成为捕人的渔夫者所捕之中。那么，谁有属于凯撒的，就应还给凯撒，为以后他才能将属于天主的还给天主。但耶稣是\BibleQuote{不可见的天主的肖像}\BibleRef{哥 1:15}，祂并没有凯撒的肖像，因为\BibleQuote{这世代的君王在祂内一无所有}\BibleRef{若 14:31}，为此祂从自己的地方，即海中，取了凯撒的肖像，为将自己和门徒的交给地上君王，使那些收半银元的人不致以为耶稣欠他们和地上君王的债；因为祂付了债，并非借入，也非拥有，也非获得，也从未使之成为自己的财产，以免凯撒的肖像与不可见的天主的肖像并存。\par

% structure: p0022 source=h2 target=subsection reason=relative-heading-rank; base=h2

\subsection{十一、儿子的自由。}

这也可以用另一种方式来表达。有一些人在地上是君王的儿子，但他们并非那些君王的儿子，而是儿子，且是绝对的子；另有一些人，因他们不亲近地上君王的儿子，且不是地上任何人的儿子，但正因如此，他们是儿子，无论是天主的儿子，还是祂的圣子的儿子，或是某位属于天主者的儿子。既然如此，救主询问\Person{伯多禄}说：\BibleQuote{地上的君王向谁征收关税或丁税？向自己的儿子，还是向外人？}\BibleRef{玛 17:25}伯多禄回答说，不是向自己的儿子，而是\BibleQuote{向外人}，于是耶稣论到那些不亲近地上君王的人，并因自由而成为儿子的人说：\BibleQuote{那么儿子是自由的。}\BibleRef{玛 17:26}因为地上君王的儿子并不自由，因为\BibleQuote{凡犯罪的，就是罪恶的奴仆}\BibleRef{若 8:34}，但那些常在\OriginalTerm{天主的话}{the word of God}的真理中的人乃是自由的，正因如此，他们认识真理，好使他们也能从罪中得自由。所以，若有任何人只是一个儿子，且在此事上全然不是地上君王的儿子，他就是自由的。然而，尽管他是自由的，他仍注意不得罪地上君王和他们的儿子，以及那些收取\OriginalTerm{半舍客勒}{half-shekel}的人；因此他说：\BibleQuote{我们不要使他们绊倒，但去投下你的网，捉住那首先上来的鱼}\BibleRef{玛 17:27}等等。但我愿向那些喜欢编造不同本性神话的人询问：这些人的本性如何，无论是地上的君王，还是他们的儿子，或是那些收取半舍客勒的人——救主不愿得罪他们；\OriginalTerm{根据假设}{ex hypothesi}，他们显然不是值得赞扬的本性，然而祂仍小心避免使他们绊倒，并阻止任何绊脚石放在他们路上，以免他们犯更重的罪，并且旨在使他们得救——如果他们愿意的话——甚至通过接受那位已使他们免于绊倒的祂。正如在一个确实安慰的地方（按解释，即\Place{葛法翁}），祂安慰那位既自由又是儿子的门徒，赐给他先捉住鱼的能力，好使当鱼上来时，\Person{伯多禄}因鱼的上来和被捉住，以及从鱼口中取出\OriginalTerm{斯塔特}{stater}来付给那些应得斯塔特的人而得到安慰，那些人要求这块银钱作为自己的。\par

% structure: p0024 source=h2 target=subsection reason=relative-heading-rank; base=h2

\subsection{十二、\OriginalTerm{斯塔特}{stater}寓意解}

但有时你可以优雅地将这段话应用于贪财者，他口中除了关于银钱的事别无他物，当你看见他被某位\Person{伯多禄}治愈时，这位伯多禄不仅从他的口中和话语中，而且从他整个品格中取走了\OriginalTerm{斯塔特}{stater}——这是他一切贪婪的象征。因为你会说，这样一个人曾在大海中，在人生的苦海中，在贪婪的忧虑和焦虑的波涛中，当他尚未相信且贪婪时，口中含着斯塔特；但他从大海中上来，被理性的网捕获，并因某位\Person{伯多禄}教导他真理而受益，不再口中含着斯塔特，取而代之的是那些包含祂肖像的事物，即\OriginalTerm{天主的话}{the oracles of God}。\par

% structure: p0026 source=h2 target=subsection reason=relative-heading-rank; base=h2

\subsection{十三、\OriginalTerm{圣半舍客勒}{Sacred Half-Shekel}}

此外，针对那句话：\BibleQuote{那些收取半舍客勒的人来到\Person{伯多禄}面前}\BibleRef{玛 17:24}，你将从\BibleBook{}{户籍纪}中引证：为那些按照天主的法律成为圣洁的人，所付的不是简单的半舍客勒，而是\OriginalTerm{圣半舍客勒}{sacred half-shekel}。因为经上记载：\BibleQuote{你要按圣半舍客勒，每人取五个舍客勒}\BibleRef{户 3:47}。而且也为全体以色列子民按人头献上圣半舍客勒。既然天主的圣者不可能同时拥有圣半舍客勒和所谓的俗舍客勒，因此，对于那些不接受圣半舍客勒的人，他们问\Person{伯多禄}说：\BibleQuote{你们的老师不纳半舍客勒吗？}救主命令付给一块\OriginalTerm{斯塔特}{stater}，其中包含那首先上来的鱼口中发现的半舍客勒，以便为老师和门徒付出。\par

% structure: p0028 source=h2 target=subsection reason=relative-heading-rank; base=h2

\subsection{十四、论那些说\Quote{谁最大？}的人，以及论\Person{耶稣}所召来的孩子}

\Quote{\Emph{那一天，门徒来到耶稣跟前说：那么，谁在天国里是最大的呢？}} \BibleRef{玛 18:1} 为了使我们可以得知门徒来到耶稣跟前向祂求问的是什么，以及祂如何回答他们的疑问，玛窦——虽然本可以只记述这件事本身——根据某些手稿，加上了\Quote{那时，门徒来到耶稣跟前}，但根据另一些手稿，却是\Quote{那一天}。我们必须好好研究圣史这句话的意思。因此，关注\Quote{那一天}或\Quote{那一刻}之前的话语，让我们看看是否可能从中找到一种方法，来理解加上\Quote{那一天}或\Quote{那一刻}的必要性。那时耶稣和祂的门徒来到\Place{葛法翁}，在那里，那些收\OriginalTerm{半舍客勒}{half-shekel}的人来到\Person{伯多禄}跟前，问道：\Quote{你的老师不纳半舍客勒吗？}然后，当伯多禄回答他们说\Quote{纳}时，耶稣进一步为纳半舍客勒的事辩护，打发伯多禄去用网拖上一条鱼，说在鱼嘴里会找到一枚\OriginalTerm{斯塔特}{stater}，这钱要用来为祂和伯多禄缴纳。因此，在我看来，门徒认为耶稣对伯多禄施予了极大的尊荣，祂断定伯多禄比祂其他的朋友更伟大，所以他们希望通过询问耶稣并亲耳听见祂的回答，来准确了解他们的猜疑是否属实，即祂是否如他们所想，断定伯多禄比他们更伟大；同时，他们也希望了解伯多禄被优待胜过其他门徒的理由。我想，玛窦为了使这一点清楚明了，就把\Quote{拿这个}——即斯塔特——\Quote{交给他们，作为你和我的殿税}这句话之后，紧接着写了\Quote{那一天，门徒来到耶稣跟前，说：那么，谁在天国里是最大的呢？}或许，他们因为在那显圣容时被给予那三人的优待而心存疑惑，他们疑惑的是——主认为这三人的哪一个是最大的。因为\Person{若望}出于爱曾靠在祂的胸膛上，而且我们可以推断，在最后晚餐之前，他们已经看到耶稣给予若望许多特殊尊荣的迹象；但伯多禄因他的宣认，在他们的耳中被称是有福的，因为他说：\Quote{你是基督，永生天主之子。}\BibleRef{玛 16:16}但又因为祂说：\Quote{退到我后面去，撒殚！你是我的绊脚石，因为你所体会的，不是天主的事，而是人的事。}\BibleRef{玛 16:23}他们心中就困惑起来，不知道究竟是不是他，或者是不是\Person{载伯德}的儿子中的一个，才是最大的。关于\Quote{那一天}或\Quote{那一刻}的话就说到这里，正是在那时发生了与斯塔特有关的事。\par

% structure: p0030 source=h2 target=subsection reason=relative-heading-rank; base=h2

\subsection{十五、伟大之程度有别。}

但接下来我们必须寻求理解这点：门徒们来到祂跟前，如同门徒向一位老师提出难题，询问说，\Quote{那么在天国里谁最大？}\BibleRef{玛 18:1} 并且，在这方面，我们必须效法耶稣的门徒；因为如果任何时候，在我们中间有任何探讨的题目未能查清，让我们就所争议的问题同心合意地到耶稣那里去，祂就在因祂的名聚集两三个人的地方，\BibleRef{玛 18:20} 并且祂以祂的能力临在，随时准备光照那些真正渴望成为祂门徒的心，以使他们领悟所探讨的事。同样，如果我们到任何一位被天主立为教会导师的人那里，提出一个与此类似的问题：\Quote{那么，在天国里谁最大？} 这也并不奇怪。那么，门徒们对于与此问题相关的事已经知道了什么？他们所探究的要点是什么？他们已经领悟到，那些被认为配得上天国的人之间并非平等；既然不平等，就有某人是最大的，然后依次一直到最小的；但他们进一步想知道，那最大的性质如何，那最小的生活之道如何，以及那居于中间的是谁？除非，更准确地说，他们从以下的话知道了谁是最小的：\Quote{谁若废除这些诫命中最小的一条，也这样教训人，在天国里他将被称为最小的}；但对于谁是最大的，他们却不知道，即使他们已领会了以下话的意思：\Quote{谁若实行并教导这些诫命，在天国里他将被称为大的}；\BibleRef{玛 5:19} 因为由于有许多伟大的人，按照人的标准，他们不清楚谁是伟大者中最伟大的。并且，有许多人是伟大的，但伟大者并不相等，这一点可以从\Quote{伟大}这一称号被用于依撒格，以及\Quote{他渐渐强大，成为极其强大}\BibleRef{创 26:13} 这句话，以及关于梅瑟、洗者若翰和救主所说的话中明显看出。人人都将承认，即使所有这些人都按圣经来说是伟大的，但救主比他们更伟大。但是，若翰（在妇人所生的人中没有比他更大的）\BibleRef{玛 11:11} 是否比依撒格和梅瑟更大，或者他并不更大，而是与两者或其中之一相等，这是难以断定的。并且从这句话：\Quote{依撒格渐渐强大，变得更大}\BibleRef{创 26:13}，直到他不仅成为强大，而且加上两次重复的\Quote{极其}，我们可以知道，在伟大者中间有区别：一个是伟大的，另一个是极其伟大的，另一个是极其极其伟大的。因此，门徒们来到耶稣跟前，想要知道在天国里谁最大；也许他们想听到祂有时这样说：\Quote{某人在天国里是最大的}；但祂使论述具有普遍性，表明那在天国里最大的人具有怎样的特质。让我们尽己所能，从所记载的来理解这是谁。\Quote{因为耶稣叫了一个小孩子来，}\BibleRef{玛 18:2} 等等。\par

% structure: p0032 source=h2 target=subsection reason=relative-heading-rank; base=h2

\subsection{十六、为何伟大者被比作小孩子。}

但首先我们可以用简单的方式解释它。有人若按简单方法解释救主在这里的话，可能会说：如果一个人，作为成人，克制成人的私欲，藉着灵治死身体的恶行，并且\BibleQuote{时常带着耶稣的死在身上}\BibleRef{格后 4:10}，以致达到那样一种状态，即具有没有尝过感官快乐、也没有成人冲动观念的小孩子的状态，那么这样的人就悔改了，并变成了像小孩子一样。而且，他在这方面越接近于小孩子的状态，与那些仍在修行、尚未达到如此崇高自制境界的人相比，他在天国中就越大。但是，关于小孩子在情欲快乐方面所说的话，同样也可以论及灵魂的其他情感、软弱和病患，这些是尚未完全拥有理性的小孩子本性上不会陷入的；例如，若有人悔改，且虽是成人，却变得像小孩子一样不轻易发怒；又如，小孩子对于悲伤的态度是，有时在父亲、母亲或兄弟去世时，他也会笑和玩耍；悔改的人就会变得像小孩子一样；并从圣言领受了不能悲伤的气质，以致在悲伤方面变得像小孩子。至于所谓的快乐，恶人无理地沉醉于其中，而小孩子却不受其扰，悔改并变成像小孩子一样的人也不受其扰，你也可以这样说。因此，正如别人也精确地证明过的，没有情欲会发生在尚未完全拥有理性的小孩子身上；如果没有情欲，显然也没有恐惧；但如果有与情欲相应的东西，这些东西在小孩子身上是微弱的，非常快速被压制并痊愈，所以那个悔改如同小孩子的人，达到了仿佛使自己的情欲像小孩子那样服从的程度，是值得爱的。因此，关于恐惧，也可以想到相似的事情：小孩子不会经历恶人的恐惧，而是另一种东西，那些精确了解情欲及其名称问题的人称之为恐惧；例如，在小孩子身上，他们在哭泣的时候就会忘记自己的痛苦，因为他们在瞬间转变，与那些被认为曾给他们带来悲伤和惊吓、但实际上并未在他们内引起这样情感的人一起欢笑玩耍。同样地，一个人也将谦卑自己如同耶稣所叫来的小孩子一样；因为无论是傲慢，还是对高贵出身、财富或任何那些被认为是好而实际不是的事物的自负，都不会临到小孩子。因此，你可以看到那些不是完全婴孩、直到三四岁的孩子，就像出身卑微的人一样，尽管他们可能看起来出身高贵，并且根本不比贫穷的孩子更爱富有的孩子。所以，如果按照他们的年龄，孩子们对于那些使愚昧人高升的情欲的态度是，因着理性，耶稣的门徒谦卑自己如同耶稣所指示的小孩子一样，不因虚荣而高抬自己，不因财富或衣服而自负，也不因高贵出身而骄傲，那么特别要因耶稣的名接待和效仿那些按照圣言所指示的、如同耶稣所带到身边的小孩子一样悔改的人；因为基督特别住在这样的人内，因此祂说：\BibleQuote{谁因我的名，接待一个这样的小孩子，就是接待我。}\BibleRef{玛 18:5}\par

% structure: p0034 source=h2 target=subsection reason=relative-heading-rank; base=h2

\subsection{十七、小孩子和他们的绊脚石。}

但要将下文与上文和谐一致地解释清楚，是一项艰巨的工作；因为有人可能会说：那皈依并变得像小孩的人，在相信耶稣的人中是小的，并且能够被绊倒，这怎么可能呢？同样，让我们尝试连贯地解释这一点。每个依照关于祂的真实历史而坚持耶稣是天主子，并按照福音行事的人，正走在按照德行生活、皈依并走向变得像小孩的路上；他不可能不进入天国。的确，这样的人有很多；但并非所有为了变得像小孩而皈依的人都达到了变得像小孩的地步；每个人在多大程度上缺少小孩般的相似性，就在多大程度上缺乏我们所说的小孩对情欲的态度。因此，在整个信徒群体中，也有那些好像刚刚为变得像小孩而皈依的人，在他们皈依的那一刻（为了变得像小孩），他们被称为小孩；那些皈依以变得像小孩，但远未真正变得像小孩的人，是能够被绊倒的；每个人在多大程度上缺乏对他们的相似性，就在多大程度上缺乏我们所说的小孩对情欲的态度，我们不应该给他们绊倒的机会；但是，倘若不然，那绊倒他的人将需要（作为有助于治疗他的方法）把\BibleQuote{驴的磨石挂在他的颈上，并沉入海的深处。}\BibleRef{玛 18:6}因为，这样，当他在海中（那里有\Quote{天主所造来在其中嬉戏的龙}）付清应受的惩罚，并为此目的而受罚受苦之后，他将参与那些属于海深的痛苦，这些痛苦是他被驴的磨石拖下去时所忍受的。因为磨石也有区别，所以其中的一个可以（姑且如此称呼）是人的磨石，另一个是驴的磨石；人的磨石，关于它经上记载说：\BibleQuote{两个女人推磨；一个被拿去，一个被留下；}\BibleRef{玛 24:41}但驴的磨石是要套在那给出绊脚石机会的人身上的。但有人可能会说——我不知道他说得正确还是错误——驴的磨石是恶人的沉重身体，它向下沉，他在复活时将要领受它，以便沉入被称为海深的深渊，那里\Quote{有天主所造来在其中嬉戏的龙}。但另一个人将绊倒其中一个小子归因于人类不可见的能力；因为从这些能力中，对耶稣所指出的那些小子产生了许多绊脚石。但当他们绊倒耶稣所指出的那些小子中的一个（他们是相信祂的人），那人将承受驴的磨石，即可朽坏的身体，它重重地压迫着灵魂，它本身挂在颈上，被拖向今生的事务，好藉着这些事除去他们的骄傲，并付清惩罚后，他们将通过驴的磨石达到对他们有益的境况。\par

% structure: p0036 source=h2 target=subsection reason=relative-heading-rank; base=h2

\subsection{十八、\Person{耶稣}所叫的那个小孩是谁？}

现在，或许可以提出另一种解释，不同于所谓较简单的解释；暂且不论这是作为教义，还是为了操练，让我们也探究一下：那被\Person{耶稣}叫来、放在门徒中间的小孩子，究竟是谁？现在思想一下：你是否可以说，那被\Person{耶稣}叫来的小孩子，就是圣神；当祂被救主呼召时，祂谦卑了自己，并被放置在\Person{耶稣}门徒的理智中间；的确，祂愿意我们远离一切其余事物，转向圣神所提示的榜样，好使我们自己也变得像那些小孩子——那些小孩子本身也被转化，并与圣神相似。这些小孩子是\Person{天主}赐给救主的，正如\BibleBook{}{依撒意亚}书上所说：\BibleQuote{看，我和\Person{天主}所赐给我的小孩子。}任何人若不远离世俗事务，并变得像那些拥有圣神的小孩子，就不可能进入天国；这圣神被\Person{耶稣}呼召，从祂自身的完美降下，成为好像小孩子一样，被\Person{耶稣}放在门徒中间。因此，那远离世俗欲望的人，必须谦卑自己，不仅要像一个小孩子，更要按照经上所记：\BibleQuote{如同这个小孩子。}\BibleRef{玛 18:4} 但是像那个小孩子一样谦卑自己，就是效法圣神，祂为了人的救恩而谦卑了自己。现在，关于救主和圣神是由\Person{圣父}为人的救恩所派遣的，已在\BibleBook{}{依撒意亚}书上由救主亲口宣告：\BibleQuote{现今，\Person{主}派遣了我和祂的圣神。}\BibleRef{依 48:16}然而，你必须知道，这句话有歧义；因为要么是\Person{天主}派遣了救主，而圣神也派遣了；要么，如我们所理解的，是\Person{圣父}派遣了二者——救主和圣神。因此，那谦卑自己超过所有效法那个小孩子而谦卑的人，就是天国中最大的。因为有许多人愿意像那个小孩子一样谦卑自己；但那个人，若在各方面都变得像那谦卑自己的小孩子，并且是奉\Person{耶稣}的名——尤其是在\Person{耶稣}自己内——实际上，将被发现是那被称为在天国中比众人都大的。但是，正如那奉\Person{耶稣}的名接待一个这样的小孩子的人，就是接待\Person{耶稣}；同样，那不愿奉\Person{耶稣}的名接待一个这样的小孩子的人，就是拒绝\Person{耶稣}，并把他赶出去。但是，如果那些被认为配得圣神的人也有差别，因为信徒领受圣神或多或少，那么在信\Person{天主}的人中，会有一些小子可以被绊倒。为了报复那些使人绊倒的人，圣言论到那些使他们绊倒的人说：\BibleQuote{为他更好，若将驴拉的磨石拴在他的颈上，沉在海底。}\BibleRef{玛 18:6}关于我们面前的这段\BibleBook{玛窦}{福音}经文，就讲到这里。\par

% structure: p0038 source=h2 target=subsection reason=relative-heading-rank; base=h2

\subsection{十九、\BibleBook{玛尔谷}{福音}与\BibleBook{路加}{福音}中的平行经文}

但让我们也思量其他福音作者相似的记载。马尔谷，\BibleRef{谷 9:33} 于是说，那十二位在路上争论谁最大。因此祂坐下，叫他们来，教导谁最大，说：那藉着谦和与温良成为众人中最后的，将作为最大的得到首位，这样他取得的不是受服事者的地位，而是服事者的地位，且不是对某些人而不对另一些人，而是对所有人绝对地服事；因为请听这话：\Quote{若有人愿意作首先的，他必作众人中最后的，并作众人的仆役。} \BibleRef{谷 9:35} 接着祂说，\Quote{祂}——即耶稣——\Quote{拿起一个小孩，把他放在祂自己的门徒中间，并且抱起祂，对他们说：凡因我的名接待一个这样的小孩，就是接待我。} \BibleRef{谷 9:36} 但按照经文的深意，耶稣所拿起并抱在怀中的小孩是什么呢？是圣神吗？的确，有些人被比作这个小孩，关于他们祂说过：\Quote{凡因我的名接待一个这样的小孩，就是接待我。} 按路加记载，这思念并非门徒自发产生，而是由这个问题提示给他们：\Quote{他们中谁最大。} \BibleRef{路 9:46} 耶稣看透了他们心中的思念，正如祂有着能看透人心思念的眼目——看透他们心中的思念——不待他们发问，按路加记载，\Quote{拿起那小孩，叫他站在自己身边，} 不像玛窦和马尔谷所说的只放在中间，如今也\Quote{站在祂身边，} 对门徒说，不仅\Quote{凡因我的名接待这样一个小孩，} 或\Quote{凡因我的名接待这样一个小孩子，就是接待我，} 如今更进一步：\Quote{凡因我的名接待这个小孩，就是接待我。} \BibleRef{路 9:47} 因此，按路加，必须在耶稣的名内接待耶稣所拿起并放在身边的那一个小孩。我不知道是否有人能寓意解释这句话：\Quote{凡因我的名接待这个小孩。} 因为我们每个人都必须在耶稣的名内接待耶稣那时所拿起并放在身边的那小孩；因为他永远活着，是不朽的，我们必须从耶稣自己手中在耶稣的名内接待他；并且耶稣与那接待小孩的人不可分离，因此说：\Quote{凡因我的名接待这个小孩，就是接待我。} 既然圣父与圣子不可分离，祂便与那接待圣子的人同在。因此说：\Quote{凡接待我的，就是接待那派遣我的。} \BibleRef{路 9:48} 但那接待了小孩、救主和派遣祂的那一位的人，是耶稣众门徒中最小的，使自己卑微。然而，他越使自己卑微，就越成为伟大；因为正是那促使他更加卑微的事，有助于他增进伟大；请听这话：\Quote{你们众人中最小的，这人才是大的；} 但在其他手抄本中我们读到：\Quote{这人将要成为大的。} 如今，按路加：\Quote{凡不像小孩子一样接受天主的国，决不能进去。} \BibleRef{路 18:17} 这句话含混；或者指接受天主之国的人可能变成如同小孩子，或者指他可能接受那已为他变成如同小孩子一样的天主之国。或许在此世，那些接受天主之国的人是在它如同小孩子时接受它，但在来世将不再像小孩子；并且他们接受在属灵成人的完美上的伟大，可以说，这种完美在现今时代向所有接受它的人彰显，当它在此世如同小孩子时。\par

% structure: p0040 source=h2 target=subsection reason=relative-heading-rank; base=h2

\subsection{二十、世界与绊脚石。世界的多种含义。}

\BibleQuote{\Emph{祸哉，世界因了绊倒人的事！}}\BibleRef{玛 18:7} \Quote{\OriginalTerm{世界}{cosmos}}这个表述，在\BibleQuote{祂在世界，世界却不认识祂}这段经文中，是就其本身且绝对地使用的，\BibleRef{若 1:10} 但它也被相对地、就其与所关联的事物的关系来使用，如\BibleQuote{恐怕你举目望天，看见太阳、月亮和天上的一切万象，便偏离正道去叩拜事奉它们}。\BibleRef{申 4:19} 在\BibleBook{艾斯德尔}{传}中，关于她也有类似用法，写着\Quote{脱下她所有的装饰}。因为\Quote{世界}一词，单纯而言，并不等同于天上的\Quote{万象}或艾斯德尔的\Quote{装饰}；我们现在所探讨的又是另一个意义。我认为，按照圣经，世界并非这个天地组合的整体，而仅仅是围绕大地的那片区域，并且这区域不是就整个大地而言，而只是就我们所居住的那部分而言；因为真光\BibleQuote{在世界上}，即在那片围绕的区域，就我们的大地部分而言；\BibleQuote{世界却不认识祂}，\BibleRef{若 1:10} 即那区域的人，或许还包括与这地方有亲缘的权势。因为若把这里的\Quote{世界}理解为天地及其中的一切所组成的整体，那就荒谬了；这样便可说太阳、月亮、星辰的众天使和这整个世界中的天使都不认识那真光，并虽不认识它，却仍遵守天主为他们所定的秩序。但是当救主在向圣父的祷文中说：\BibleQuote{父啊，现在求祢以祢自己，就是我在未有世界以前，同祢所有的光荣，来光荣我}，\BibleRef{若 17:5} 你们必须按\Quote{世界}理解为我们所居住的大地；因为圣父将人赐给圣子，正是从这世界中，救主仅为这些人求父，而不是为全人类的世界。再者，当救主说：\BibleQuote{我往祢那里去，我不再在世界上}，\BibleRef{若 17:11} 祂指的是地上的世界；因为不能认为祂说\BibleQuote{我往祢那里去，我不再在世界上}和\BibleQuote{我在世界上}，是自相矛盾的。同样，在\BibleQuote{我还在世界上说这些话}这一句，\BibleRef{若 17:13} 我们也要理解为围绕大地的区域。这由\BibleQuote{世界恨他们，因为他们不属于世界}\BibleRef{若 17:14} 这话也清楚表明。因为自从我们因耶稣的教导，不再\BibleQuote{顾念所见的，而是顾念所不见的}以来，\BibleRef{格后 4:18} 世界就恨我们；这世界不是天地及其中的一切所组成的整体，而是与我们同在世上的人。而\BibleQuote{他们不属于世界}\BibleRef{若 17:21} 这句话，等同于说他们不属于围绕大地的区域。同样，耶稣的门徒也不属于这世界，正如祂不属于世界。此外，在\BibleBook{若望}{福音}中两次出现的\BibleQuote{叫世界相信祢派遣了我}，也不是指超乎人类的事物，而是指需要相信圣父派遣圣子到这世界上来的人。是的，宗徒也说：\BibleQuote{你们的信德传遍了全世界}。\BibleRef{罗 1:8}\par

% structure: p0042 source=h2 target=subsection reason=relative-heading-rank; base=h2

\subsection{二十一、\Quote{祸哉}并不适用于耶稣的门徒。}

但如果地上各处的人因绊倒人的事而遭祸哉——那些事会抓住人；但门徒不属于世界，因为他们不顾念所见的，正如主不属于这世界；那么对任何一位耶稣的门徒，都不适用\Quote{因绊倒人的事而有的祸哉}，因为\BibleQuote{爱天主法律的人，大有平安，他们不会失足}。但如果有人似乎被称为门徒，却仍属于世界，因为他爱世界和世界上的事——我是指围绕大地的区域中的生命、财产、占有物或任何形式的财富——以致\BibleQuote{他们不属于世界}\BibleRef{若 17:16} 这话不适合他；那么，作为真正属于世界的人，那临到世界的事，即\Quote{因绊倒人的事而有的祸哉}，就要降到他身上。但愿意逃避这祸哉的人，不要贪爱生命，而要同\Person{保禄}一起说：\BibleQuote{世界于我已被钉在十字架上，我于世界也被钉在十字架上}。\BibleRef{迦 6:14} 因为圣人们\BibleQuote{在这帐棚里，叹息呻吟，深感重负}，\BibleRef{格后 5:4} 带着\Quote{卑贱的身体}，他们做一切事，为能配得在复活的奥秘中被找到，那时天主将不再为所有人，而是为那些真正成了基督门徒的人，更新那卑贱的身体，使之相似祂光荣的身体。\BibleRef{斐 3:21} 因为正如没有一个\Quote{祸哉}临到基督的任何门徒，同样这个\Quote{因绊倒人的事而有的祸哉}也不会临到他们；因为，即使有成千的绊脚石出现，它们也不会触及那些不再属于世界的人。但如果有人，因信德缺乏根基，对天主圣言的服从不够坚定，而可能跌倒，就让他知道，他并未被耶稣称为祂的门徒。现在我们应当设想，有如此多的绊脚石出现，结果祸哉不再只临到地上某些部分，而是临到其中整个的\Quote{世界}。\par

% structure: p0044 source=h2 target=subsection reason=relative-heading-rank; base=h2

\subsection{二十二、\Quote{绊倒人的事}是什么？}

\Quote{ \Emph{而且必须要有绊倒人的事来到}，}\BibleRef{玛 18:7} 我认为这与那些带来绊倒的人不同。那么，来到的绊倒事是魔鬼的军队、他的天使和一队不洁的邪灵，它们寻找工具来工作，常常找到完全与虔敬疏远的人，有时甚至找到那些被认为相信天主圣言的人，对于这些人，有比那被绊倒的人更坏的祸患，正如在审判之日，\BibleRef{玛 11:22} 对提洛和漆冬要比对耶稣行奇迹奇事却仍不被相信的地方更可容忍。但是，正如有人可以从圣经中收集那些被称为有福的人以及他们被称为有福的事，同样他也可以收集所写的祸患以及那些受祸患的人。然而，那导致绊倒的人的祸患比那被绊倒的人更坏，你可以通过经文证明：\Quote{无论谁使这些信我的小子中的一个绊倒，对他更好，}\BibleRef{玛 18:6} 等等；因为那被绊倒的小子从使他绊倒的人那里得到报应，而那导致绊倒的人遭受所写的严厉且不可容忍的惩罚是合适的。但如果我们更仔细地考虑这些事，我们就应当提防得罪弟兄们，并伤害他们软弱的良心，免得我们得罪基督；\BibleRef{格前 8:11} 正如我们周围的弟兄们，\Quote{基督为他而死}，常常不仅因我们的知识而丧亡，也因与我们有关的其他原因；在这种情况下，我们得罪基督，将受惩罚，那因我们而丧亡的灵魂的代价将从我们手中索要。\par

% structure: p0046 source=h2 target=subsection reason=relative-heading-rank; base=h2

\subsection{二十三、在何种意义上\Quote{必要}}

接下来我们必须准确检验\Quote{必要}一词在经文\Quote{ \Emph{因为必须要有绊倒的事来到}，}\BibleRef{玛 18:7} 中的意义，类似地在路加中，\Quote{是\InnerQuote{不可能不发生}，但绊倒的事应该来到}，\BibleRef{路 18:1} 代替\Quote{不可能}。正如那必死的一定会死，且不可能不死，正如那在身体里的人必须进食，因为不进食的人不可能活，同样绊倒的事必然且不可能不发生，因为在人身上，恶必然先于德性而存在，而绊倒事源自恶；因为不可能找到一个人完全无罪，且没有罪而达到德性。因为邪恶势力中的恶，是人类恶的源头，完全渴望通过某些工具对抗世上的人。而且，或许邪恶势力在它们被耶稣的言语驱逐时更加激怒，它们的敬拜减少，通常的祭献不再献给它们；因此这些冒犯事必然发生；但并非必然通过任何特定的人发生；因此\Quote{祸}落在带来绊倒的人身上，因为他给那意图制造绊倒的邪恶势力提供了位置。但不要以为某些绊倒事是出于本性和构造，寻找人作为工具；因为正如天主没有创造死亡，祂也没有创造绊倒事；而是自由意志在那些不愿为德性忍受劳苦的人中产生了绊倒事。\par

% structure: p0048 source=h2 target=subsection reason=relative-heading-rank; base=h2

\subsection{二十四、那绊倒人的手、脚或眼睛}

因此，如果眼睛和手值得称赞，那么眼睛就不能合理地对手说：\Quote{我不需要你。}\BibleRef{格前 12:21} 但如果教会集会整个身体中的某个人，因其实际恩赐而被称为手，却变成一只绊倒人的手，就让眼睛对这样的手说：\Quote{我不需要你，}并且说着，就把它砍下来丢掉。\BibleRef{玛 18:8} 同样，如果有一个头是有福的，而脚配得上有福的头，使头遵守适合本身的事，就不能对脚说：\Quote{我不需要你。}然而，如果发现任何一只脚成为整个身体的绊脚石，就让头对这样的脚说：\Quote{我不需要你，}并把它丢掉，从自己身上扔掉；因为即使没有那绊倒的脚或手，身体其余部分进入永生，也比当绊脚石蔓延到全身时，带着两只脚或两只手被丢进地狱的火好得多。同样，那能成为整个身体眼睛的人，应当配得上基督和整个身体；但如果这样的眼睛曾经改变，成为整个身体的绊脚石，最好把它取出来丢到整个身体之外，并且身体其余部分没有那只眼睛而得救，比连同它一起，当全身被败坏时，整个身体被丢进地狱的火好。因为灵魂的实践能力，如果倾向于犯罪，和灵魂的行走能力，可以这么说，如果倾向于犯罪，以及灵魂的清晰视力，如果倾向于犯罪，可能就是那绊倒人的手、绊倒人的脚和绊倒人的眼睛；这些东西最好丢掉，在摒弃它们之后，没有它们而进入生命，如同一个跛子、残缺或独眼的人，总比连同它们失去整个灵魂要好。同样，在灵魂的情况下，善用其能力达到最高贵的目的是一件好而蒙福的事；但如果我们因任何原因将要失去其中一种能力，那么失去其用途更好，以便连同其他能力我们得以得救。\par

% structure: p0050 source=h2 target=subsection reason=relative-heading-rank; base=h2

\subsection{二十五、对眼睛或手的寓意解释}

这些话也可以应用在我们最近的亲属身上，他们可以说是我们的肢体；因为亲密的关系，他们被视为我们的肢体；无论是由于血缘，还是由于习惯性的友谊，可以这么说；如果他们伤害我们的灵魂，我们就不该宽容他们。因为如果我们有一个父亲或母亲希望我们做出违背虔敬的事，有一个儿子或女儿想要我们（就他们所能而言）背弃基督的教会和对祂的爱，我们就要像砍掉一只手或一只脚或一只眼睛一样，将他们从我们身上砍去。但是，即使是我们怀中的妻子，或一个灵魂相通的友人，成了我们的绊脚石，我们也不该宽容他们，而要将其从我们身上割除，从我们的灵魂中驱逐出去，因为他们不是我们真正的亲属，而是我们救恩的仇敌；因为\Quote{谁不恨自己的父亲、母亲，}\BibleRef{路 14:26}以及随后提到的那些人，当应该像对待仇敌和攻击者一样憎恨他们，以便能够赢得基督时，这人就不配作为天主子。关于这些人，我们可以说，从一个关键的位置上，任何一个（可以说是）瘸子都得救了，当他失去了一只脚——比如说一个兄弟——而独自获得了天主的国的产业；一个残废的人得救了，当他的父亲没有得救，而他们丧亡了，他却与他们分离，以便他独自获得祝福。同样，任何一只眼睛得救的人，就是砍掉了自己家里的眼睛——他的妻子，如果她犯了奸淫，免得他有两只眼睛而被投入地狱之火。\par

% structure: p0052 source=h2 target=subsection reason=relative-heading-rank; base=h2

\subsection{二十六、小子和他们的天使。}

\Quote{\Emph{你们小心，不要轻视这些小子中的一个。}}\BibleRef{玛 18:10} 在我看来，如同人的身体在大小方面有差别——有些人矮小，有些人高大，有些人中等身材，而且在矮小者中也有差异，有的更小或稍小，高大者和中等身材也是如此——同样，在人的灵魂中，有些事物给他们打上渺小的印记，有些打上伟大的印记（姑且这样说），一般地，按照身体事物的类比，其他事物打上平庸的印记。但在身体方面，一个人矮小或高大或中等，并非出于人的行为，而是由于\OriginalTerm{种子原理}{spermatic principles}；但在灵魂方面，是我们的\OriginalTerm{自由意志}{free-will}、如此的行为和如此的习惯提供了理由，让人成为伟大、渺小或平庸；而或是在身量上长进以增大体积，或是不长进而矮小，这都取决于我们的自由意志。因此，我这样理解关于耶稣取了人性灵魂的话：\BibleQuote{耶稣前进}\BibleRef{路 2:52} 因为正如由于自由意志，祂的灵魂在智慧和恩宠上有所增长，同样也在身量上。宗徒也说：\BibleQuote{直到我们众人达到一个成年的人，达到基督圆满年龄的程度。}\BibleRef{弗 4:13} 因为我们必须认为，一个人达到成年，按照内在的人，他经历了孩童的事，达到了成人的阶段，并丢弃了孩童的事，一般地，成全了成人的事。因此，我们必须设想，存在着某种灵性身量的尺度，最完美的灵魂能够借着赞美上主达到这个尺度，并成为伟大。这样，那些被称为伟大的人，如\Person{依撒格}、\Person{梅瑟}、\Person{若翰}，以及超越一切的救主本身，就成了伟大；因为关于祂，\Person{加俾额尔}说：\BibleQuote{他将为大。}\BibleRef{路 1:32} 但小子们是\BibleQuote{新生的婴孩，渴慕纯洁的灵性之乳，}\BibleRef{伯前 2:2} 他们需要\OriginalTerm{乳养之父}{nursing-fathers}和\OriginalTerm{乳养之母}{nursing-mothers}，\WorkTitle{依撒意亚}论到外邦人的召叫时说：\BibleQuote{他们必将你的儿子抱在怀中，将你的女儿放在肩上；君王要作你的乳养之父，王后要作你的乳养之母。}\BibleRef{依 49:22} 因此，你将会留意这句话：\Quote{不要轻视这些小子中的一个，}\BibleRef{玛 18:10} 并考虑是否他们的天使把他们抱在怀中，因为他们已经成了儿子，并且将所谓女儿放在肩上，是否从那称为君王的乳养之父和称为王后的乳养之母而来。既然我们的救主所指明的小子们在乳养之父和乳养之母的\OriginalTerm{管家职分}{stewardship}之下，基于此，我认为\Person{梅瑟}，他相信自己已经被置于伟人的行列中，关于这应许说：\Quote{我的天使要在你前面引路，}\BibleRef{出 32:34} \Quote{如果你自己不与我同行，就不要从这里带我上去。}\BibleRef{出 33:15} 因为即使小子也是继承人，但作为孩童，他与奴仆毫无差别，当他还是孩童时，\BibleRef{迦 4:1} 并且在他渺小到\BibleQuote{得着了奴役之灵，仍有所恐惧；}\BibleRef{罗 8:15} 的程度上，但那些不再如此的人，不再有奴役之灵，而是已经有了儿子名分之灵，那时\BibleQuote{完美的爱驱逐恐惧；}\BibleRef{若一 4:18} 这样，你会明白如何根据这些事，上主的天使被说成\Quote{在敬畏祂的人四周安营，并拯救他们。}但你也将根据这些事思考：是否这些确实是小子们的天使，\Quote{他们被奴役之灵引导而恐惧，}\Quote{当初上主的天使在敬畏祂的人四周安营并拯救他们；}但就伟大者而言，是否上主（祂比天使更大）关于他们每个人可以说：\Quote{我在困苦中与他同在；}并且，只要我们还不完美，需要一位帮助我们脱离邪恶者，我们就需要一位天使，就是\Person{雅各伯}所说的：\Quote{从所有邪恶中拯救我的天使}\BibleRef{创 48:16} 但是，当我们成了成全的，并经历了从属于乳养之父、乳养之母、监护人和管家的阶段，\BibleRef{迦 4:4} 我们就适合由主亲自治理了。\par

% structure: p0054 source=h2 target=subsection reason=relative-heading-rank; base=h2

\subsection{二十七、当小子们被托付给天使时。}

然后还可以追问，那些被称为他们天使的，是从什么时候开始承担起基督所指明的小子们的监护职分；他们是否从\OriginalTerm{重生之洗}{laver of regeneration}\BibleRef{铎 3:5}开始接受了这个使命，就是他们作为\BibleQuote{新生的婴孩，渴慕纯洁的灵性之乳，}\BibleRef{伯前 2:2} 不再受任何邪恶权势的管辖；或者，是否从出生开始，他们就按照上主的\OriginalTerm{预知和预定}{foreknowledge and predestination}被指派给那些上主也预知并预定要肖似基督荣耀的人。\BibleRef{罗 8:29} 关于天使从出生就有的观点，有人可以引用：\BibleQuote{他从母胎中就分别了我，}\BibleRef{迦 1:15} 以及\Quote{从母胎中你就作了我的保护者，}以及\Quote{他从母胎中就帮助了我，}以及\Quote{我从母胎中就被投靠于你，}以及在犹达书信中：\BibleQuote{蒙天主父所爱，为耶稣基督而保存，且蒙召的人，}\BibleRef{犹 1:1} ——完全被守护他们的天使所保存。\par

% structure: p0056 source=h2 target=subsection reason=relative-heading-rank; base=h2

\subsection{二十八、天使与他们的\Quote{小子们}的密切关系。}

论及\Quote{\Emph{当我透过洗濯池在基督内成为婴孩时}}这句话，可以说，没有圣天使与仍在邪恶中的人同在，但在不信的时期，他们处于撒殚的天使之下；然而，在重生之后，那以自己的血赎回了我们的那位，将我们托付给一位圣天使，这位天使因他的纯洁而常目睹天主的圣容。关于这段经文，还可以有第三种解释，类似如下：正如一个人能从不信转变为信，从放纵转变为节制，从邪恶转变为德行，同样，一个自出生就被托付给某个灵魂的天使，起初可能是邪恶的，但后来，当人相信时，他也按比例相信，并且进步到成为那常在天上瞻仰圣父面容的天使之一\BibleRef{玛 18:10}，从那时起，他与那被预知和预定在那时相信的人联合，天主那不可言说、不可探究、如同深渊的审判，恰到好处地促成这整个和谐的关系——天使与人。也有可能，正如一对夫妻同为不信者，有时丈夫先相信，随后拯救妻子；有时妻子先开始，然后说服丈夫，天使和人之间也是如此。然而，若其他天使也是如此，或不是，你可自行探究。但请考虑，是否可以对每位按照救主的话如此受尊荣、被称为常在天上圣父面容的天使，说出类似的话。由于我们在上面所说，小子有天使，但大人已经超越这种地位，有人会引用《宗徒大事录》来反对我们，其中记载：有一个名叫洛弟的使女，当伯多禄敲门时，她前来应门，认出伯多禄的声音，就跑进去报告说伯多禄站在门前；但在屋内聚集的人感到惊讶，认为伯多禄确实不可能站在门前，便说：\BibleQuote{那是他的天使}\BibleRef{宗 12:13}。反对者会说，既然他们一次性地学到每个信徒都有某个特定的天使，他们就知道伯多禄也有一个。但坚持我们先前所说的人会说，洛弟的话不一定是教条，而且那些没有准确知道的人的话也可能如此，因为一个人何时作为小子而敬畏天主，受天使管理，何时又由主自己亲自管理。此后，为了确立我们提出的关于小子的概念，可以说，我们不需要关于\BibleQuote{不轻视}大人的命令，但需要关于小子的命令；因此，并非只说\BibleQuote{不要轻视这些中的一个}，指着所有门徒，而是\BibleQuote{这些小子中的一个}\BibleRef{玛 18:10}，由那位看见灵魂的小与大者所指明。\par

% structure: p0058 source=h2 target=subsection reason=relative-heading-rank; base=h2

\subsection{二十九、小子与成全者。}

但另有人可能会说，成全的人在此被称为小子，引用\BibleQuote{因为你们中间最小的，这人才是大的}\BibleRef{路 9:48}，并断言那谦卑自下、在众信徒中成为小孩的人，即使他是宗徒或主教，且成为\BibleQuote{如同母亲乳养自己的孩子}\BibleRef{得前 2:7}，就是耶稣所指的小子，他的天使配得目睹天主的圣容。因为根据\BibleQuote{你们中间最小的，这人才是大的}\BibleRef{路 9:48}这段经文，以及保禄所说\BibleQuote{我原是最小的宗徒，不配称为宗徒}\BibleRef{弗 3:8}，称小子在这里指成全的人，似乎与\BibleQuote{无论谁使这些小子中的一个跌倒}\BibleRef{玛 18:6}和\BibleQuote{你们在天之父的旨意，是使这些小子中一个也不丧亡}\BibleRef{玛 18:14}和谐一致。因为，如前所述，那如今是小子的人，既不能跌倒也不能丧亡，因为\BibleQuote{爱慕天主法律的人，享有极大的平安，他们没有绊脚石}；而那在基督众门徒中是最小的人，并因此成为大的人，不会丧亡；既然他不会丧亡，他就能说\BibleQuote{谁能使我们与基督的爱隔绝？}\BibleRef{罗 8:35}等等。但坚持这最后一种解释的人会说，就连义人的灵魂也是可改变的，正如厄则克尔也作证说，义人可能离弃天主的诫命，以致他先前的义德不被记念\BibleRef{则 33:12}；因此经上记着\BibleQuote{谁使这些小子中的一个跌倒}\BibleRef{玛 18:6}和\BibleQuote{你们在天之父的旨意，是使这些小子中一个也不丧亡}\BibleRef{玛 18:14}。\par

[关于\BibleQuote{一百只羊}\BibleRef{玛 18:12}相关的解释，你可参阅\BibleBook{路加}{福音}的讲道。]\par

% structure: p0061 source=h2 target=subsection reason=relative-heading-rank; base=h2

\subsection{三十、犯罪的弟兄。}

\BibleQuote{\Emph{如果你的弟兄得罪了你，去，只在你们两人之间指证他的过错。}}\BibleRef{玛 18:15} 那么，仔细考察这句话，以证明耶稣超凡仁爱的人会说，因这话并未暗示罪的区别，他们便会以一种独特的方式行事，并且违背耶稣的良善——那些提供这种想法的人认为，这话应理解为仅适用于较轻的罪。但另一个人，也仔细考察这句话，不愿引入这些外来的想法，也不承认这句话是针对一切罪说的，便会说：犯那些大罪的人不是弟兄，即使他被称作弟兄，正如宗徒所说：\BibleQuote{如果有人称为弟兄，却是行淫乱的、或贪婪的、或拜偶像的等，这样的人，不要与他共食。}\BibleRef{格前 5:11} 因为没有一个拜偶像的、或行淫乱的、或贪婪的人是弟兄；因为如果一个人，他看似带有基督的名号，虽被称为弟兄，却有这些人的某些特征，他便不能正确地被称为弟兄。因此，正如那说这样的话是针对一切罪的人——无论那罪是谋杀、下毒、娈童或诸如此类——会给基督的至善造成损害的机会；相反，那区分弟兄和被称为弟兄的人，可能会教导说，在人的最小罪的情形中，那在被告知过犯后仍未悔改的人，应被视为外邦人和税吏，因为这些罪是\BibleQuote{不至於死的}\BibleRef{若一 5:16}，或者，正如法律在户籍纪中所描述的，不是\BibleQuote{致死的}。\BibleRef{户 18:22} 这似乎非常严苛；因为我不认为有谁能轻易找到一个人，未曾因同一种形式的罪而受三次谴责，比如说，辱骂（辱骂者以此辱骂邻人，或被激情冲昏头脑的人）、或酗酒、或撒谎与闲言，或在群众中存在的任何这些事。因此，你要探究，是否这段经文的一些观察点，被那些受其关于圣言之良善的概念影响的人忽略了，他们宽恕犯最大罪的人；同时也被那些教导说，在最小的罪的情形中，一个人在犯了三次极小的过犯后，就要被算为外邦人和税吏，使他成为教会的外人的人忽略。但我认为两者都忽略了以下这一点，即\BibleQuote{你已赢得了你的弟兄}这句话\BibleRef{玛 18:15}。圣言将它只归于那聆听的人，不再应用于那跌倒数次并被谴责的人；但关于那被谴责两三次的人应当说的话，与\BibleQuote{你已赢得了你的弟兄}相对应，祂却悬而未言，可谓如此。因此，他并没有完全被赢得，也不会完全丧亡，或者他将受鞭打。且要仔细注意第一段：\BibleQuote{如果他听从你，你便赢得了你的弟兄}；以及第二段，其原文是：\BibleQuote{如果不听从你，就另带一两个人去，以便凭两三个证人的口，可确立一切话。}\BibleRef{玛 18:15} 那么，当一切话已由两三个证人确立后，那第二次被谴责的人将遭受什么，祂留给我们去思考。再者，\BibleQuote{如果他不听他们}——显然是指所带的那些证人——\BibleQuote{就告诉教会}；祂说\BibleRef{玛 18:17}；祂没有说若他不听从教会会受什么惩罚，但祂教导说，如果他拒绝听从教会，那已经三次劝诫而未被听从的人，今后要将他看作外邦人和税吏\BibleRef{玛 18:17}。因此，他并非完全被赢得，也非完全丧亡。但总的说来，那起初未听从、却要求证人、甚至拒绝听从证人、而被带到教会的人，将遭受什么，天主知道；因为我们不加以断言，按诫命\BibleQuote{你们不要判断，免得你们受判断}\BibleRef{玛 7:1}，\BibleQuote{直等到主来，祂要照亮黑暗中的隐密事，并显明心中的意图}\BibleRef{格前 4:5}。但是，关于在犯较小罪的人身上看似严苛的情况，有人可能会说，那连续两次未听从的人，不可能第三次听从，以致因此不再被视为外邦人或税吏，或不再需要全体教会面前的谴责。因为我们必须记住：\BibleQuote{你们的在天之父不愿这些小子中的一个丧亡。}\BibleRef{玛 18:14} 因为\BibleQuote{我们众人必须出现在基督的审判台前，为使每人按照他所行的，或善或恶，领取他藉身体所做的事的报应}\BibleRef{格后 5:10}，愿各人尽己所能，为使自己不因身体所行的更多恶事而受惩罚，即使他将要为自己所行的一切不义得到报应；但我们的雄心应是获取更多善行的赏报，因为\BibleQuote{你们用什么升斗量，也要用什么升斗量给你们}\BibleRef{玛 7:2}，并且\BibleQuote{按我们手所作的，必要临到我们}\BibleRef{依 3:11}，不是无限地，而是罪人从主的手中为自己的罪接受双倍或七倍的报应；因为祂不按人手所作的回报任何人，而是超过他所行的，因为依撒意亚教导说：\BibleQuote{耶路撒冷从主的手中为自己的一切罪接受了双倍}\BibleRef{依 40:2}；但以色列的邻邦，不论是谁，要接受七倍，按照圣咏中如下的话：\BibleQuote{主啊，求您将邻邦加给您的辱骂，七倍报应到他们怀中。}还可以找到其他形式的回报；如果我们明白了这些，就会知道，在任何罪之后悔改，无论罪有多大，都是有益的，为使我们不仅不因更多过犯而受罚，而且对于以后某个时候所行的善事，还留有一些希望，即使在此之前我们中任何人犯下了成千的过失。因为如果恶事被记在某人的账上，而恶之后所做的善事却毫无益处，那才是奇怪的。这一点也可从厄则克耳学到，只要留意关于这类情况所说的话。\EditorialNote{厄则克耳第三十三章}\par

% structure: p0063 source=h2 target=subsection reason=relative-heading-rank; base=h2

\subsection{三十一、在地上与天上捆绑的权柄。}

但在我看来，那在三次告诫后被视为像外邦人和税吏一样的人的案件，适宜地接上了这句：\Quote{\Emph{我实在告诉你们}}——即对那些将任何人视为\OriginalTerm{外邦人}{Gentile}和\OriginalTerm{税吏}{publican}的人——\Quote{\Emph{并且你们在地上所捆绑的}}\BibleRef{玛 18:18}等等；因为那三次告诫而未被听从的人，正义地捆绑了被视为外邦人和税吏的人；因此，当这样的人被一个如此性格的人捆绑和定罪时，他仍被捆绑，因为天上没有人推翻那捆绑者所作的判决。同样，那一次被告诫并行配得被赢回之事的人，因那赢回他的人的告诫而得释放，不再被他自己罪恶的绳索所捆绑，\BibleRef{箴 5:22} 为此他受到了告诫，应被判定为由天上的人释放了。不过，似乎暗示：那些先前只赐予伯多禄的事，现在赐予所有向所有犯罪者给予三次告诫的人；这样，如果未被听从，他们将在世上捆绑被视为外邦人和税吏的人，正如这样的人在天上已被捆绑。但由于有必要，即使在对伯多禄和那些三次告诫弟兄的人所说的话中有一些共同之处，伯多禄应该拥有某些优于那些三次告诫者的元素，在伯多禄的情况中，这句话\Quote{我要将天国的钥匙给你}\BibleRef{玛 16:19} 被特意放在\Quote{凡你在地上所捆绑的}等话之前。而且，如果我们仔细留意福音书，我们也会在这里发现，就那些看似伯多禄和三次告诫弟兄者所共有的事而言，所说与伯多禄有关的话与第二类人相比，存在着巨大的差异和卓越之处。因为伯多禄接受的\OriginalTerm{钥匙}{keys}不是一个天，而是多个天，并且他在地上所捆绑的在\OriginalTerm{诸天}{heavens}中也被捆绑，这与许多在地上捆绑和释放的人相比，是一个不小的差异；这些人在天上不是像伯多禄那样被捆绑和释放，而只是在一个天上；因为他们没有达到那样高的境界，拥有像伯多禄那样在所有诸天中捆绑和释放的能力。\BibleRef{玛 16:19} 因此，捆绑者越好，被释放者就越有福，以致在诸天的每一部分他的释放都得以完成。\par

\SourceRecord{Translated by John Patrick.；Ante-Nicene Fathers；Vol. 9.；Edited by Allan Menzies.；Buffalo, NY: Christian Literature Publishing Co.,；1896.；<http://www.newadvent.org/fathers/101613.htm>.}

% structure: p0001 source=h1 target=section reason=page-title

\section{\WorkTitle{玛窦福音释义（第十四卷）}}

% structure: p0002 source=h2 target=subsection reason=relative-heading-rank; base=h2

\subsection{一、和谐的力量与祈祷}

\BibleQuote{\Emph{我再给你们说：如果你们中两人在地上同心合意}，\Emph{无论为什么事祈求，他们在天之父必要成全他们。}} \BibleRef{玛 18:19} \OriginalTerm{交响}{symphony}一词严格用于音乐中声音的和谐。在音乐的声音中，确实有些是协调的，有些是不协调的。但福音圣经却熟用此词指音乐方面，如在经文：\BibleQuote{他听见交响乐和跳舞的声音。} \BibleRef{路 15:25} 这是因为：当那失而复得的儿子藉着悔改与父亲和睦时，在这家庭的欢乐喜宴上听到交响乐是合适的。但恶人\Person{拉班}却不识\Quote{交响}一词，他对\Person{雅各}说：\BibleQuote{你若告诉我，我必欢欢喜喜地打发你走，有音乐、有鼓、有竖琴。} \BibleRef{创 31:27} 与此类交响相似的是在\BibleBook{撒慕尔}{纪下}所记载的：\BibleQuote{阿米纳达布的弟兄们走在约柜前，大卫和他的儿子在天主面前弹奏精心制作的乐器，大有能力，并配有歌曲；} \BibleRef{撒下 6:4} 因为这些精心制作、大有能力并配有歌曲的乐器，本身具有音乐的交响，其力量如此之大，以至于当只有两人带着与神圣属灵音乐相关的交响，向天上的父祈求任何事时，父就会应允那些带着地上的交响祈求的人——这极其奇妙——他们可能已经求得了那些筹成前述交响之人所祈求的事物。同样，我也理解宗徒的话：\BibleQuote{你们不要彼此亏负，除非两相情愿，暂时分房，为专务祈祷。} \BibleRef{格前 7:5} 因为\Quote{和谐}一词应用于那些按照天主所结的婚姻，正如\BibleBook{}{箴言}中的话：\BibleQuote{父亲将房屋和财产分给儿子，但妻子是由天主赐予丈夫的。} 那么，和谐来自天主，理当使名与实共享为祈祷而达成的合一，正如经上所说：\BibleQuote{除非两相情愿。} \BibleRef{格前 7:5} 然后，圣言重复说，在地上两人同心合意，就是与基督同心合意，并补充道：\BibleQuote{因为哪里有两个或三个人因我的名字聚集在一起。} \BibleRef{玛 18:20} 因此，因基督的名字聚集在一起的两个人或三个人，就是在地上同心合意的人，不单是两人，有时也是三人。但有权能的人将会考虑：这种同心合意以及基督在其间的这种聚集，是否能在更多的人中找到？因为\BibleQuote{通往生命的门是窄的，路是狭的，找到的人很少。} \BibleRef{玛 7:14} 但或许不单是少数，而是两三个人构成一个交响，如\Person{伯多禄}、\Person{雅各伯}和\Person{若望}，天主的圣言向这几位构成交响的人显示了自己的荣耀。但有两人构成交响：\Person{保禄}和\Person{索斯特乃}，在写\BibleBook{格林多}{前书}时\BibleRef{格前 1:1}；之后，\Person{保禄}和\Person{弟茂德}在致同一教会的第二封信中\BibleRef{格后 1:1}。甚至三人构成交响：\Person{保禄}、\Person{息耳瓦诺}和\Person{弟茂德}在致得撒洛尼人的信中训诲时\BibleRef{得前 1:1}。但若也有必要从古经中举出在地上构成交响的三人，使得圣言在他们中间使他们合而为一，请留意 \BibleBook{}{圣咏集} 的标题，例如第四十一篇，如下：\BibleQuote{交与伶长，训诲，科辣黑子孙的诗。} 因为\Person{科辣黑}有三个儿子，我们在\BibleBook{出谷}{纪}中找到他们的名字\BibleRef{出 6:24}：\Person{阿色尔}，意为\Quote{训诲}；次子\Person{厄耳卡纳}，意为\Quote{天主的产业}；三子\Person{阿彼雅撒夫}，希腊语可译为\Quote{父亲之集会}。然而预言并未分开，而是由同一个灵、同一种声音、同一个灵魂说出并写下，这灵魂以真正的和谐运作，三人如同一人说道：\BibleQuote{如同心渴慕水泉，我的灵魂渴慕你，天主。} 但在第四十四篇中，他们也用复数说：\BibleQuote{天主，我们亲耳听见了。} 但若你愿进一步看到那些在地上构成交响的人，请看那些听从劝勉的人：\BibleQuote{好使你们同心合意，有同样的心思，同样的判断。} \BibleRef{格前 1:10} 以及追求目标的人：\BibleQuote{全体信徒一心一意。} \BibleRef{宗 4:32} 他们已成为这样——如果可能的话，在多于两三人中找到这种状态——他们之间没有不和谐，正如十弦瑟的弦之间没有不和谐一样。但那些说\BibleQuote{我是属保禄的，我是属阿颇罗的，我是属刻法的，我是属基督的}\BibleRef{格前 1:12}的人，在地上并没有交响，他们中间有分裂。等到分裂瓦解后，他们与\Person{保禄}的灵、主耶稣基督的能力一同聚集\BibleRef{格前 5:4}，以便不再\BibleQuote{彼此咬食，以致互相消灭。} \BibleRef{迦 5:15} 因为不和谐消灭，和谐则聚集，并接纳那位来到已和睦者中间的天主子。严格来说，和睦确实发生在两类事物中：一是通过共同的成全（如宗徒所称），即因心智把握相同的意见而有同样的心思；二是通过共同的成全，即因相似的生活方式而有同样的判断。但如果任何时候我们中的两人在地上同心合意，无论为什么事祈求，他们在天之父（即\Person{耶稣}之父）必要成全他们\BibleRef{玛 18:19}，那么显然，当他们在天之父没有成全他们祈求的任何事时，那两人在地上并未同心合意；这就是我们祈祷不蒙垂听的原因：我们在地上不彼此同心合意，无论在见解还是生活上。此外，如果我们就是基督的身体，天主将每一个肢体安置在身体内，使肢体彼此同样关怀、互相和睦，并且如果一个肢体受苦，所有肢体一同受苦；如果一个肢体得荣耀，所有肢体一同喜乐，那么我们就应当操练那源于神圣音乐的交响，以便当我们因基督的名字聚集时，祂就在我们中间——天主的圣言、天主的智慧和祂的能力\BibleRef{格前 1:24}。\par

% structure: p0004 source=h2 target=subsection reason=relative-heading-rank; base=h2

\subsection{二、夫妻的和谐。}

以上是关于圣言所劝勉要同心合意的\Quote{两个人或三个人}的较普遍理解。现在让我们也触及我们某位先辈所提出的另一种解释，他劝勉已婚者追求圣洁与纯洁；因为，他说，圣言愿意在地上同心的两个人，应理解为丈夫与妻子，他们藉着同意彼此免除肉身的结合，以便专心祈祷；\BibleRef{格前 7:5}这样，如果他们为任何所求之事祈祷，就必得着，因着这样的同心合意，耶稣基督的天上之父必赐予他们所求的。在我看来，这种解释并不导致婚姻的解除，反而是对同心的激励；如此，若一方愿意守贞，而另一方不愿意，那么那愿意且能成就更好部分的一方，便俯就那无力或不愿的一方，这样，他们二人必能从耶稣基督的天上之父那里，得着他们任何所求之事。\par

% structure: p0006 source=h2 target=subsection reason=relative-heading-rank; base=h2

\subsection{三、肉身、灵魂与心神的和谐。}

在恶人身上，罪掌管灵魂，安居在这可朽的肉身中如同在自己的宝座上，以致灵魂顺从肉性的私欲；\BibleRef{罗 6:12}但在那些已经激起了从前在肉身中掌权的罪、并与之争战的人身上，\Quote{肉性相反心神，心神相反肉性；}\BibleRef{迦 5:17}但在那些已成全的人身上，心神已得胜，治死了肉身的作为，并将自己的生命赐予肉身，以致应验了这话：\Quote{祂必使你们必死的身体活过来，因为祂的圣神住在你们内；}\BibleRef{罗 8:11}于是肉身与心神二者在地上产生了和谐；在此和谐成功达成之后，那\Quote{心里相信成义，口里宣认得救}\BibleRef{罗 10:10}的人也会发出和谐的祈祷，以致心不再远离天主，同时义人也以自己的嘴唇和口亲近天主。但更蒙福的是，若三者都因耶稣的名聚集，以致应验这话：\Quote{愿天主完全圣化你们，并愿你们的心神、灵魂和肉身，在我们主耶稣基督来临时，得以完全无可指摘地保存。}\BibleRef{得前 5:23}但有人可能会问：关于心神与肉身的和谐，若没有第三者——即灵魂——的和谐，是否可能？以及，是否当二者在地上和谐，并在基督的名内聚集后，三者就已经在祂的名内聚集，在它们中间天主子来到，因为三者都已奉献给祂，没有人反对祂，不仅心神没有敌对，连灵魂也没有，肉身也更没有。\par

% structure: p0008 source=h2 target=subsection reason=relative-heading-rank; base=h2

\subsection{四、旧盟约与新盟约的和谐。}

同样，努力理解并展示两个盟约——即救主肉身来临之前的盟约与新盟约——的和谐事实，也是一件愉快的事；因为凡在两个盟约和谐一致、彼此无冲突之处，必有祈祷，以致他们无论求什么，都必从天上的父那里得着成就。若你还想要那联合二者的第三者，不妨说那就是圣神，因为\Quote{智者的话，}无论它们是来临之前的、来临之时的、还是来临之后的，\Quote{就如刺棒，又像钉牢的钉子，都是由一位牧者所赐，出于同一的协定。}\BibleRef{训 12:11}此外，也不可忽略这一点：祂并没有说，哪里有两三个人因我的名聚集，那里\BibleQuote{我将在}他们中间，而是\BibleQuote{我就在}他们中间，\BibleRef{玛 18:20}不是将要，不是延迟，而是在同心合意的那一刻，祂亲自临在，在他们中间。\par

% structure: p0010 source=h2 target=subsection reason=relative-heading-rank; base=h2

\subsection{五、宽恕的限度。}

\BibleQuote{\Emph{那时\Person{伯多禄}前来对祂说：主啊，若我的弟兄得罪了我，我该原谅他多少次？}}\BibleRef{玛 18:21}认为这些话是按字面简单由\Person{伯多禄}所说，仿佛他在询问当他的弟兄得罪他时，他应原谅他七次，但若他第八次犯罪，则不再原谅；而救主则认为人应静坐计算邻人得罪他的次数，以便原谅他七十个七次，但从第七十八次起，他就不应原谅那亏待他的人。——这种理解在我看来完全是愚蠢的，既与\Person{伯多禄}在与\Person{耶稣}相处中所取得的进步不符，也与\Person{耶稣}神圣的宽宏大量不相称。因此，这些话可能也接近类似于\BibleQuote{你们\Person{拉默客}的妻子，请听我的声音}\BibleRef{创 4:23}等话语的晦涩。若有人已成了\Person{耶稣}的朋友，以至受祂的圣神教导，这圣神光照那按功绩已进展到如此地步之人的理智，那么他就能知道关于这些事的真义，即\Person{耶稣}自己会清楚解释的那样；但我们这些不及\Person{耶稣}友谊之伟大的人，若能对这经文略发呓语，也就知足了。因此，数字六似乎代表辛劳工作，而数字七则包含安息的概念。请你想想，能否说那爱世界并做世界之事、做属物质之事的人，犯了六次罪，而数字七在他身上是罪的终结，以致\Person{伯多禄}带着类似想法，愿意原谅他的弟兄对他犯下的七种罪。但既然个位数的十和百在比例上有某种共同尺度，而\Person{耶稣}知道这个数字可能被超越，因此，我想，祂在数字七上又加了七十，\BibleRef{玛 18:22}并说在这里对弟兄们，以及对那些在关于这里的事情上犯罪的人，应给予宽恕。但若有人超越关于世界和今世的事情而犯罪，即使那是微不足道的，他也不能再合理地获得罪的宽恕；因为宽恕延伸到这里的事，并关系到在这里所犯的罪，无论宽恕来得迟或早；但对于那犯了超过七十七次罪的弟兄，甚至没有宽恕。但你可以说，那如此犯罪的人，无论是得罪了他的弟兄\Person{伯多禄}，还是得罪了那阴府的门不能战胜的\Person{伯多禄}，按此类的罪，他属于较小的数目，但按更坏的罪，他属于没有罪之宽恕的数目。\par

% structure: p0012 source=h2 target=subsection reason=relative-heading-rank; base=h2

\subsection{六、关于那与自己的仆人算账的王，有人带了一个欠一万塔冷通的人来见他。}

\BibleQuote{\Emph{因此我对你们说：天国好比一个国王，要同他的仆人算账。}} \BibleRef{玛 18:23} 这个比喻的一般用意是教导我们，应当宽恕那些得罪我们的人所犯的过错，尤其当犯错者在行恶之后向受害者恳求宽恕时。比喻还希望教导我们：即使天主已经宽恕了我们所犯的罪债，但若我们不宽恕得罪我们的人的罪，那么即使得赦免之后，仍会要求我们偿还债款，直到我们心中不再存有丝毫对恶事的记忆，而是全心全意，藉着宽恕之德——这并非寻常的美德——原谅那对我们行恶、甚至背信弃义之人。但紧接比喻的一般用意，也应按字面更简单地审视全文，使慎重研究前述各细节的人能从所言的审查中获益。因为可能存在一种超验而难以捉摸的解释，带有某种神秘性，按照圣史所解释的比喻的类比，人们会探究其中的每一个细节；例如，国王是谁，仆人是谁，他开始算账的缘由是什么，那欠了许多塔冷通的债务人是谁，他的妻子和儿女是谁，除了国王命令变卖以偿还债务之物外所谓\Quote{一切}是什么，那被宽恕了许多塔冷通的人出去意味着什么，被发现且欠债的不是对家主而是对被宽恕者的人是谁，那一百个德纳的数目是什么意思，\Quote{他扼住他的喉咙说：还你所欠的}这句话是什么意思，那被宽恕了一切塔冷通的人出去将同伴投入的监狱是什么，那些忧闷并告知主人所发生之事的同伴仆人是谁，那被打入监狱的人被交给的刑吏是谁，被交给刑吏的人如何还清了一切债务不再亏欠。但一个更有能力的探究者或许还能添加其他事情，它们的阐述和解释我认为超乎人的能力，需要讲述它们的基督的神，使基督按他所讲的而被理解；因为\BibleQuote{除了人内里的心神外，谁能知道人的事呢？}并且\BibleQuote{除了天主的神外，谁也不能知道天主的事} \BibleRef{格前 2:11}，同样，除了基督的神外，没有人能知道基督在箴言和比喻中所讲的天主的事；谁在基督内分享不只是作为神的他，而是作为智慧、作为圣言的基督，就必能看到这段经文中启示给他的事。至于最高层次的解释，我们不敢自称；另一方面，藉着作为天主智慧的基督的帮助，我们也不失望于理解比喻中所意指的事；但关于这段圣经是否会有这样的内容赐予我们，愿天主在基督内提示做他所喜悦的事，只愿我们也藉着圣神由天主赐下智慧之言，并按照圣神赐下知识之言 \BibleRef{格前 12:8}。\par

% structure: p0014 source=h2 target=subsection reason=relative-heading-rank; base=h2

\subsection{七、讲解继续：国王与仆人。}

祂说：\Quote{\Emph{天国}好比}，\BibleRef{玛 18:23}等等。但若它好比这样一位君王，一位行了这样的事的君王，我们怎能不说它就是天主子呢？因为祂是诸天的君王，并且既然祂是绝对的智慧、绝对的正义和绝对的真理，难道祂不也是绝对的王国吗？然而，这王国不是任何地上事物的王国，也不是天上部分事物的王国，而是所有天上事物——这些事物被称为诸天——的王国。但若你探究这句话的含义：\Quote{天国是他们的}，\BibleRef{玛 5:3}你可以说基督是他们的，就祂是绝对的王国而言，祂在每一个思想中为王，这思想属于一个不再受罪统治的人——罪在那些已顺服于罪的人那必死的身体中为王。\BibleRef{罗 6:12}若我说祂在每一个思想中为王，我的意思是这样的：祂作为正义、智慧、真理以及其他美德，在已变成天的人心中为王——因为这人心怀天上的肖像——也在每一种能力中为王，无论是天使的能力，还是在今世和来世被称为圣者的其余能力，这些能力都配得这样的王国。因此，这天国（当它被造成\Quote{罪恶肉身的形状}\BibleRef{罗 8:3}，为要因罪而定罪罪，当天主使那\Quote{不知罪的，替我们成了罪}\BibleRef{格后 5:21}——我们这带着罪恶肉身的人），好比一位特定的君王，这君王被理解为与耶稣相连，与祂合而为一，若我们敢这样说，祂比那与主结合、成为与祂一灵的人\BibleRef{格前 6:17}，更有能力与那\Quote{一切受造物的首生者}\BibleRef{哥 1:15}结合并完全成为一体。如今，这天国——好比一位特定的君王，按对耶稣的理解与祂相连并合而为一——被提前说成祂愿意与祂的仆人算账。但祂要与他们算账，为要显明各人是如何运用了家主的合格货币和理性的钱币。比喻中的图像确实取自主人与自己的仆人算账；但若我们定睛于那些管理主人财物、并被要求就此交账的仆人所行的事，我们就能更准确地理解这部分比喻的含义。因为每个仆人从主人的财物中领受了不同分量，有的正确使用，使主人的财物增多；有的则浪费在不当之处，毫无判断与审慎地挥霍了交托给他们的东西。然而，有些人智慧地管理了这些财物——而且是这样大的财物——却失去了其他，每当主人与他们算账、他们交账时，就汇集出每人遭受了多少损失，计算出每人带来了多少收益，并根据他管理财物的方式是否相称，或得尊荣，或受惩罚；有时债务被豁免，有时塔冷通被夺走。那么，从以上所说，让我们先看看理性的钱币和家主的合格货币：有人收获得多，有人收获得少——因为按各人的能力，有人领了五个塔冷通，因为他有能力管理那么多；另一个人领了两个，因为他不能承受前者的量；还有一人领了一个，因他比第二个人更差。\BibleRef{玛 25:15}这些是仅有的区别吗？还是我们要在福音接着谈论的某些人身上认出这些区别，同时还有其他区别？在其他比喻中也能找到某些人，比如两个债务人，一个欠五百德纳，另一个欠五十；\BibleRef{路 7:41}但这些人是否曾受托管理财物并管理不善（像那领了一个塔冷通的人能力较差），或者他们曾领受什么，我们不得而知；但比喻似乎教导我们，他们欠了那么多。还有另外十个仆人，每人分别受托一米纳。\BibleRef{路 19:13}若有人理解了人类灵魂的多样性以及彼此间在天赋上的巨大差异——有的人对更多或更少的美德有或无天赋，对某些美德或另一些美德有天赋——他或许就能领会每个灵魂是如何带着家主的某些钱币而来的，这些钱币在理性完全成熟时，以及在理性完全成熟后随之而来的关注和正确事物的操练中，或是在其他事物（无论是作为追求有用，或部分有用部分无用，如那些不完全真也不完全假的意见）的勤奋与操练中显露出来。\par

% structure: p0016 source=h2 target=subsection reason=relative-heading-rank; base=h2

\subsection{八、清算的原则。}

但在此你定要探究，是否所有的人皆可称为国王的仆人，抑或有些人是蒙他预知并预定的仆人，另有些人则与仆人进行交易，而被称为银行家。\BibleRef{玛 25:27} 同样地，你也要探究，是否有那些不属奴仆之列的人，户主声称要连本带利向他们收回自己的东西，这些人不仅是远离虔敬的人，也包括一些信徒。如今唯独仆人是圣言的管家，但国王与仆人结账时，要向那些向仆人借过东西的人追索，无论是一百桶小麦还是一百桶油，\BibleRef{路 16:6}或实际上那些不属国王家室的人所领取的任何东西；因为那欠了一百桶小麦或一百桶油的人，依照比喻并非不义管家的同仆，这一点从问题——你欠我主人多少？\BibleQuote{你欠我主人多少？}\BibleRef{路 16:5}即可看出。但请与我一同注意，每一件善的或合宜的行为，都好比一种收益和增加；而恶行好比一种亏损。正如钱多时有一种收益，钱少时有另一种收益，且有多寡之别，同样，按照善行，也可说有一种或多或少的收益估值。要断定何种工作是大收益、何种是较小收益、何种是最小收益，此乃唯独他享有的特权，因为唯有他知道如何审视这些事，从意向、言语和行为的角度观察，并考虑那些非我们能力所及之事与能力所及之事相互配合；同样，在相反的事上，由他来说，当与仆人结账时，何种罪过被认定为大亏损、何种为较小亏损，以及——容我们如此说——何种是最后一个小钱，\BibleQuote{最后一个小钱}\BibleRef{路 12:59}或最后一块铜板，\BibleQuote{最后一块铜板}\BibleRef{玛 5:26}。因此，那被称为天国、可比作国王的，要追索整个人生的总账，那时\BibleQuote{我们众人都应出现在基督的审判座前，为使每人按他肉身所行的，或善或恶，领取相应报应；}\BibleRef{格后 5:10}然后当结账之时，凡人们所说的每一句废话也都将被纳入结算，\BibleQuote{每一句废话}\BibleRef{玛 12:36}以及任何人只因为门徒的名分，给一杯凉水喝的事，\BibleQuote{给一杯凉水喝的事}\BibleRef{玛 10:42}。\par

% structure: p0018 source=h2 target=subsection reason=relative-heading-rank; base=h2

\subsection{九、结账所占的时间。}

这些事将要发生，每当那写在达尼尔书上的事迹成就时：\BibleQuote{书卷展开了，审判已设立；}\BibleRef{达 7:10} 因为，可以说，凡所说、所做、所思想的一切都作了记录，且借着天主的能力，我们一切隐藏的事都将被显明，一切遮盖的事都将被揭露，为要使那被发现未曾\BibleQuote{尽力从对头手中脱身}的人，依次经过官长、审判官、差役之手，被投入监牢，直到他还清最后一分钱；\BibleRef{路 12:58} 但若有人尽力从中脱身，且不欠任何人分文，又已经使那一弥纳变成十弥纳或五弥纳，或使五塔冷通加倍，或使两塔冷通变成四塔冷通，他便可得到应得的赏报，进入他主人的福乐中，或被委派管理祂全部财产，\BibleRef{玛 24:47} 或听到\BibleQuote{你掌管十座城} \BibleRef{路 19:17} 或\BibleQuote{你掌管五座城}的言语。然而我们以为，这些话说得仿若需要很长时间，为使我们能勾销尘世生命的全部时光，以致我们可能设想，当国王与他的众多仆人中每一位结账时，这件事将需要如此漫长的时间，直到那些从世界伊始直至世代终结——不是一代，而是许多世代——所存在的事物结束。但事实并非如此；因为当天主意欲一下子在所有人的记忆中重新燃起每个人在整个时间中所做的一切，为使人人都能意识到自己的行为是善是恶时，祂会以其不可言喻的能力做到。因为天主不像我们；我们若想记起某些事，需要充裕的时间来详细叙述我们所说的话，并唤起我们想要记住的事；但若祂愿意唤起我们记忆中所行过的事，为使我们在意识到自己所行之后能明白为何受罚或受赏，祂是能够做到的。但若有人不信天主在这些事上的能力之迅速，他尚未对那位创造宇宙的天主有正确的认识，那位天主并不需要时间就创造了浩大的天地及其中的万物；因为，尽管祂似乎是在六天内造了这些，仍需理解\Quote{在六天内}一语是在何种意义上说的，并因此而要知道\BibleQuote{这是天地万物的来历}，\BibleRef{迦 2:4}等等。因此可以大胆断言，那所期待的审判之时并不需要时间，而是正如复活据称是\BibleQuote{在一刹那，一眨眼之间}发生的，\BibleRef{格前 15:52}照样，我认为审判也会如此。\par

% structure: p0020 source=h2 target=subsection reason=relative-heading-rank; base=h2

\subsection{十、欠了许多塔冷通的人。}

接下来，我们必须就此讲论：\Quote{\Emph{及至他开始算账时，有人带了一个欠他一万塔冷通的人来。}}\BibleRef{玛 18:24}我认为其寓意如下：审判的季节始于天主的家，祂对那些被任命执行惩罚的说，正如在\BibleBook{厄则克耳}{书}所记载：\BibleQuote{从我的圣者开始。}\BibleRef{则 9:6}这就像\Quote{眨眼的一瞬；}而算账的时间也包含这同一\Quote{一瞬，}在观念上被领悟，因为我们并没有忘记先前对那些欠更多债者的论述。因此，经上并未记载当他算账时，而是说\Quote{当他开始算账时}，便在开始算账之初，有人带了一个欠他一万塔冷通的人来；他丧失了数以万计的塔冷通，原本受托付大事，有许多事物交给他管理，但他没有为主人带来利益，反而亏损了万计，以致欠下许多塔冷通；或许正因如此，他欠了许多塔冷通，因为他时常追随那坐在铅塔冷通上的妇人，她的名字叫邪恶。\BibleRef{匝 5:7}但请注意，在这里，每一项大罪都是对家主塔冷通的亏损，犯下这类罪的是奸淫者、通奸者、男色者、柔弱的人、拜偶像者、杀人者。那么，那个被带到王前欠许多塔冷通的，所犯的并非小罪，而是所有重大且可憎的罪；如果你在人类中寻找他，或许会发现他是\BibleQuote{那罪恶的人，丧亡之子，他敌对并高举自己超过一切被称为神或受崇拜的对象。}\BibleRef{得后 2:3}但若你在人类之外寻找他，除了那使许多接待他的人败亡、在他们内行罪的魔鬼，还能是谁呢？因为\BibleQuote{人是一件大事，一个可怜的人是宝贵的。}\BibleRef{箴 20:6}宝贵到足以堪当一塔冷通，无论是金塔冷通——像那等于一金塔冷通的灯台\BibleRef{出 25:39}——或银塔冷通，或任何按理智理解的其他质料，其象征记载于\BibleBook{}{编年纪}中，\BibleRef{编上 22:14}当达味获得许多塔冷通时，那里提到了数目：金塔冷通若干，银塔冷通若干，以及其余有所命名的质料，用以建造天主的圣殿。\par

% structure: p0022 source=h2 target=subsection reason=relative-heading-rank; base=h2

\subsection{十一、欠一百银钱的仆人。}

只是，虽然他无法偿还那些塔冷通，因为他已经失去了它们，但他有妻子、儿女和其他东西，经上记载说：\BibleQuote{凡他所有的。} \BibleRef{玛 18:25} 而且，当他和他的家人一起被卖时，如果有人买了他，他可能会兴旺起来，并凭借他的价值和属于他的东西，全额还清债务；并且他可能不再是王的仆人，而成为了买主的仆人。他恳求不要与他的家人一起被卖，而是继续住在王的家中；因此他俯伏朝拜他，知道王就是天主，并说：\BibleQuote{求你给我宽限，我必全还给你；} \BibleRef{玛 18:26} 因为他大概是个勤奋的人，知道可以通过第二次行动来填补先前损失许多塔冷通的全部亏空。这位真正良善的王动了慈心，免除了那欠他许多塔冷通的人的债，赐予了他比所求更大的恩惠；因为欠债人向忍耐的主人承诺偿还所有债务，但主动了慈心，不仅因为忍耐而赦免了他以收回自己的钱财，甚至完全释放了他，免除了他全部债务。但这个恶仆，曾恳求主人为他许多塔冷通忍耐，行事却毫无怜悯，因为他遇见一个同僚仆人欠他一百个德纳，就抓住他，掐住他的喉咙，说：\BibleQuote{你所欠的，还我。} \BibleRef{玛 18:28} 他难道不是表现出极端的邪恶吗？他为了一百个德纳抓住同僚仆人，掐住他的喉咙，剥夺了他呼吸的自由，而他自己却因许多塔冷通既未被抓住，也未被掐住喉咙，起初只是被吩咐与妻子、儿女和所有他的一切一同被卖；但后来，当他朝拜主人时，主人动了慈心，他被释放，并被免除了全部债务。然而，根据耶稣的理解，要说明谁是那个被发现欠一百个德纳的同僚仆人（不是欠他自己的主人，而是欠那欠许多塔冷通的人），以及谁是那些看见一人掐住喉咙、另一人被掐而极其忧愁，并向他们自己的主人清楚报告了一切所发生之事的同僚仆人，这确实是艰巨的任务。但我断言，除非耶稣——祂曾私下向自己的门徒解释一切——住在这人的理性中，并打开比喻中所有隐秘、隐藏、不见的宝藏，用清楚的证明坚定祂所愿光照的人，使他能立即明白谁是被带到王面前欠许多塔冷通的人，以及谁是那个欠他一百个德纳的另一个人，等等。他是否是先前所提到的那罪恶的人 \BibleRef{得后 2:3}，或是魔鬼，或二者都不是，而是别的某人，或是一个受到魔鬼支配的人；因为借由圣神所记载的，关于那些本质上具有某种本性或具有如此这般特质的人（无论是有形界的能力还是某些人）的预言之事，正是天主的智慧之工。但由于我们尚未领受那能与基督的心智结合、能领悟如此伟大之事、并能与圣神一同 \BibleQuote{洞察一切，甚至天主深奥的事} \BibleRef{格前 2:10} 的足够心智，我们对于这段经文仍只有模糊的印象，认为比喻中所指的那恶仆（此处代表许多塔冷通的债务）指的是某一个特定的人。\par

% structure: p0024 source=h2 target=subsection reason=relative-heading-rank; base=h2

\subsection{十二、算账的时候}

但宜考察比喻中的那人——君王——愿与自己的仆人算账是在什么时候，以及所讲的事应归属哪一时期。如果是在终结之后，或在所期待的审判之时，那么我们如何解释那欠一百银钱的人被那蒙赦免了许多塔冷通的人扼喉的事呢？但如果是在审判之前，我们又如何解释君王在此之前与自己的仆人算账呢？我们应当一般性地思考每一个比喻，即那些没有被福音作者记录解释的比喻，尽管耶稣私下给祂的门徒解释了一切\BibleRef{谷 4:34}；为此缘故，福音书的作者隐藏了比喻的清楚解释，因为比喻所指的事物超出了言辞本性的表达能力，而这类比喻的每一种解决和解释都如此深奥，以致整个世界本身也容纳不下为这些比喻所应写的书\BibleRef{若 21:25}。然而，或许会有一颗适宜的心被找到，因其纯洁，能够接受比喻解释的文字，使它们被永生天主的圣神写在其中。但有人会说，我们或许行事不虔敬，因为对某些出自天主的圣经的隐秘和奥秘含义，我们希望它们是象征性的，并努力解释它们，即使从\OriginalTerm{假设}{ex hypothesi}看来我们似乎对它们的含义有准确的认识。但对此我们必须说，如果有人获得了准确理解这些事物的恩赐，他们知道应当做什么；至于我们，承认自己缺乏洞察此处所象征事物深度的能力，即使我们对经文略有所悟，我们也要说：某些我们经过多次考察探究似乎发现的道理，无论是因天主的恩宠还是凭我们心智的力量，我们不敢写入著作；但有些道理，为了我们自己心智的修炼以及偶然读到它们之人的益处，我们将在一定程度上加以陈述。那么，这些就当作为辩解之言，因为比喻的深奥；但关于比喻中的那人——君王——愿与自己的仆人算账是在什么时候的问题，我们要说：这似乎发生在所宣告的审判时期。这一点由两个比喻所证实，一个在我们面前的福音结尾处\BibleRef{玛 25:14}，一个出自路加福音\BibleRef{路 19:12}。为了不因逐字引用而延长讨论——因为任何愿意的人都可以亲自从圣经中取用——我们要说，按玛窦的比喻记载：\Quote{因为就像一个人要往外国去，叫了自己的仆人，把自己的家业交给他们；给了一个五个塔冷通，另一个两个，另一个一个}\BibleRef{玛 25:14}，然后他们对所托付的采取了行动，过了许久，那些仆人的主人来了，并记载说，他也与他们算账\BibleRef{玛 25:19}。又比较这话：\Quote{他开始算账的时候}\BibleRef{玛 18:24}，并思考：他将家主往外国去称为\Quote{我们住在肉身内，但与主远离}\BibleRef{格后 5:6}的时候，但他的来临——\Quote{过了许久，那些仆人的主人来了}\BibleRef{玛 25:19}——是终结审判的时期；因为过了许久，那些仆人的主人来了，与他们算账，随后的事就发生了。但路加福音中的比喻更清楚地表明：\Quote{有一个贵人往远方去，为自己取得王位而后回来}\BibleRef{路 19:12}，去的时候\Quote{他叫了十个仆人，给了他们十个米纳，对他们说：你们去做生意，直到我回来}\BibleRef{路 19:13}。但这贵人被自己的公民所恨，他们在他后面派了使节，因为他们不愿他统治他们；他得了王位回来后，叫了那些得银钱的仆人到他面前来，要知道他们做生意赚了多少。看到他们所行的，他对那将一个米纳赚成十个的人以这话加以赞扬：\Quote{好，善仆，你既然在小事上忠诚}\BibleRef{路 19:17}，就给了他权柄管理十座城，即那些属于他王国的城。对另一个将米纳赚了五倍的人，他没有给予给第一个人的赞扬，也没有如第一个那样具体说\Quote{权柄}一词，而是对他说：\Quote{你也要管理五座城}。但对那将米纳包在手帕中的人，他说：\Quote{凭你自己的口，我要定你的罪，恶仆}\BibleRef{路 19:22}；并对旁边站着的人说：\Quote{从他手里拿过那米纳来，给那有十个米纳的}\BibleRef{路 19:24}。那么，对于这比喻，谁不会说那往远方去为自己取得王位而后回来的贵人就是基督，祂仿佛往别国去，要取得这世界及其中的万国？那些领受了十个塔冷通的人，就是受托管理所交付给他们的圣言职务的人。而那些当祂按降生成人曾是这世界的公民时，不愿祂统治他们的公民\BibleRef{路 19:14}，或许是不信祂的以色列人，或许还有不信祂的外邦人。\par

% structure: p0026 source=h2 target=subsection reason=relative-heading-rank; base=h2

\subsection{十三、不宽容者不得宽恕。}

只是，我说这些话，为的是将祂带着祂的国来临时的归来，指向终结，那时祂命令那些领了银子的仆人被召到祂面前，为要知道他们经商的收获；并且希望借此，并从塔冷通的比喻，表明\BibleQuote{他愿意同自己的仆人算账}\BibleRef{玛 18:23} 这段话应指向终结，那时祂已是君王，接受了王国，为此，按照另一个比喻，\BibleRef{路 19:12} 祂往远方去，为自己取得王国并返回。因此，当祂接受了王国返回后，祂愿意同自己的仆人算账。而\BibleQuote{祂刚开始算账时，有人带了一个欠祂许多塔冷通的人来}\BibleRef{玛 18:24}，他被带到王面前，如同被那些被指派为祂臣仆的人带来——我想，是天使。或许他是王国领导之一，曾被委托重大管理职责，却没有善加管理，反而浪费了所托付的，以致欠下了他失去的许多塔冷通。这个人，也许无力偿还，就被王下令连同他的妻子一起卖掉，他与妻子交合成为某些孩子的父亲。但要看出父亲、母亲和孩子们在理智上的含义并非易事。这究竟在真理上意味着什么，只有天主知道；祂是否赐予了我们洞察力，有能者自可判断。不过，这只是我们对这段经文的看法：正如\BibleQuote{天上的耶路撒冷}是\BibleQuote{母亲}\BibleRef{迦 4:26}，是\Person{保禄}及类似他的人的母，同样，按照耶路撒冷之母的类比，也可能有其他人的母，例如埃及的\Place{息安}，或\Place{漆冬}，或圣经中所称的许多城市之母。继而，正如耶路撒冷是\BibleQuote{盛装打扮的新娘，等候她的丈夫}\BibleRef{默 21:2}——基督，同样，也可能有某些权能的母，她们被分配给他们作为妻子或新娘。正如耶路撒冷作为母，基督作为父，有某些儿女，同样，\Place{息安}、\Place{曼菲斯}、\Place{提洛}或\Place{漆冬}以及管辖它们的长官，也有某些儿女。那么，也许这个欠许多塔冷通被带到王面前的人，如我们所说，也有妻子和儿女，起初王下令将他们连同他一切所有的都卖掉；但后来，王动了怜悯之心，释放了他，免除了一切债务；这并非好像王不知道将来之事，而是为了我们能够理解所发生的事，才如此记载。因此，每一个如我们所说有妻子和儿女的人，都将在王接受王国并返回后来算账时，交账；而他们每个人，作为某个\Place{息安}、\Place{曼菲斯}、\Place{提洛}、\Place{漆冬}或类似之地的地方长官，也有欠债者。于是，这个人被释放并免除了全部债务后，\BibleQuote{从王那里出来，遇见了一个与他同作仆人的}\BibleRef{玛 18:28}等等；为此，我想，他掐住那人的喉咙，是在他离开王之后，因为若他没有离开，他不会掐住自己同伴的喉咙。然后请留意圣经的精确性：一人俯伏\BibleQuote{朝拜}，另一人也俯伏，却不是朝拜，而是\BibleQuote{恳求}；王动了怜悯之心，释放了他，免除了全部债务，但那个仆人甚至连同情自己的同伴都不愿；而王在释放他之前，曾下令把他和他的一切卖掉，然而被赦免的人却把自己的同伴投入监牢。还请留意，他的同伴们并没有指控或\BibleQuote{说}，而是\BibleQuote{告诉}\BibleRef{玛 18:31}；并且他（王）起初并没有在关于损失钱财的事上用\BibleQuote{邪恶}这一称呼，而是后来针对他对同伴的行为才使用。也要留意王的节制：他没有说\BibleQuote{你朝拜了我}，而是说\BibleQuote{你恳求了我}；他不再下令把他和他的一切卖掉，而是更糟，把他交给刑役，因为他恶毒。\BibleRef{玛 18:34} 但这些刑役若不是那些被指派掌管惩罚的人，又会是谁呢？但同时，要留意异端信徒使用这个比喻的方式：如果他们指责造物主动怒，因为经文宣告了天主的愤怒，那么他们也应当指责这个王，因为\BibleQuote{他发怒}而把欠债者交给刑役。但还必须对那些持\BibleQuote{没有人被耶稣交给刑役}观点的人说：请解释给我们，善人们，这位把恶仆交给刑役的王是谁？也请他们注意这话：\BibleQuote{因此，我天上的父也要这样对待你们}\BibleRef{玛 18:35}；而且对同一群人，或许更应引用善天主子在十个米纳的比喻中所说的话：\BibleQuote{至于那些反对我作他们君王的人}\BibleRef{路 19:27}等等。然而，这个比喻的结论也适合较单纯的人；因为我们所有获得自身罪过赦免，却没有宽恕自己弟兄的人，立刻被教导：我们将会遭受那个被赦免却不宽恕同伴之人的命运。\par

% structure: p0028 source=h2 target=subsection reason=relative-heading-rank; base=h2

\subsection{十四、耶稣如何结束祂的话。}

\BibleQuote{\Emph{当耶稣说完了这些话}} \BibleRef{玛 19:1} 凡详细而完整地叙述自己面前的每一个问题、以致毫无遗漏的人，就结束了自己的话。但那以极大勤勉致力于通读旧约和新约的人，将更有信心对此点作出说明；因为如果\Quote{他说完了这些话}这一说法不能用于其他人，既不能用于\Person{梅瑟}，也不能用于任何先知，而只能用于\Person{耶稣}，那么人将敢说，只有\Person{耶稣}独自终了祂的话，祂来是为了终结事物，并完成律法中缺陷的部分，藉着说：\BibleQuote{你们一向听对古人说} \BibleRef{玛 5:33}等等，以及：\Quote{为使先知们所说的话得以应验。}然而，如果这些话语也在他们中某处被记载，那么你便可以比较和对照他们所终了的话语与\Person{救主}所终了的话语，以找出它们之间的差异。然而在此点上，也值得探究：在藉着神谕所说的事物中，\Quote{他终了}这一说法究竟是用于\Person{梅瑟}所说的话，还是用于任何先知，或是用于二者合起来所说的话；因为仔细观察会向那些知道如何\BibleQuote{以属神的事比较属神的事}，并因此\BibleQuote{不是用人的智慧所教的言语，而是用\Person{圣神}所教的言语} \BibleRef{格前 2:13}说话的人提出非常重要的思想。但或许另外有人，以过于好奇的精神关注\Quote{终了}一词，这词被指派给更奥秘的事物，正如我们说某人将神秘和\OriginalTerm{终成}{perfecting}之礼交付给那些在他控制之下的人，但并非以值得赞许的方式，而另一个人则将\Person{天主}的奥秘和与这些奥秘相称的终成之礼交付给那些配得的人，他可能说，在入门之后，他施行了一种终成之礼，藉此终成，话语被证明是有能力的，以致\Person{耶稣}的福音被传遍天下，并藉着神圣的终成之德能，获得了\Person{圣父}吸引到\Person{圣子}那里的每一个灵魂的掌控，按照\Person{救主}所说的话：\BibleQuote{除非派遣我的父吸引他，没有人能到我这里来} \BibleRef{若 6:44}。因此，\Quote{言语}——那些藉着\Person{天主}的恩宠作福音使者的人——和\Quote{他们的宣讲，不是用智慧的动人言词，而是用\Person{圣神}和德能的明证}，对于他们，\Person{耶稣}教义的言语已经终了。你会因此观察到\Quote{祂终了}一词被说了多少次，关于哪些事物被说，并且你可以以\WorkTitle{真福八端}以及随后附有的\BibleQuote{当\Person{耶稣}说完了这些话，众群众都希奇祂的教训} \BibleRef{玛 7:28}的全部言论作为例证。但现在，\Quote{\Person{耶稣}说完了这些话}这句话也直接指那非常奥秘的比喻，根据这比喻，天国好比一个国王，而且超出这比喻，更指在它之前所写的段落。\par

% structure: p0030 source=h2 target=subsection reason=relative-heading-rank; base=h2

\subsection{十五、人们如何跟随\Person{耶稣}。}

只是，当\Person{耶稣}说完了这些话，在\Place{加里肋亚}靠近\Place{葛法翁}说了这些话之后，然后\BibleQuote{祂就从那里起身，来到\Place{犹太}边界} \BibleRef{玛 19:1}，这不同于\Place{加里肋亚}。但祂来到\Place{犹太}边界，并非来到中部，而是仿佛来到最外部的地方，在那里有许多群众跟随祂 \BibleRef{玛 19:2}，祂在\BibleQuote{\Place{约但河}对岸的\Place{犹太}边界}——那里曾施行洗礼——医治了他们。\BibleRef{若 1:28} 但你会注意到那些仅仅跟随的群众与\Person{伯多禄}和其他舍弃一切而跟随的人，以及\Person{玛窦}——他起来跟随了祂——之间的区别；\BibleRef{玛 9:9} 他并非仅仅跟随，而是\Quote{起来}；因为\Quote{起来}是一个重要的补充。因此，总有许多人像那些群众一样跟随，他们没有起来以便跟随，也没有舍弃以前所有的一切，但只有少数人起来跟随了，他们也要在重生时坐在十二个宝座上。\BibleRef{玛 19:28} 只是，若有人渴望被治愈，就让他跟随\Person{耶稣}。\par

% structure: p0032 source=h2 target=subsection reason=relative-heading-rank; base=h2

\subsection{十六、论\Person{法利塞人}和\Person{经师}试探\Person{耶稣}（藉着询问）人是否可以因任何缘故休妻。}

此后经上记载说：\Quote{\Emph{有法利塞人前来试探祂说：人无论因什么缘故，休妻可以吗？}}\BibleRef{玛 19:3} 马尔谷也记载了类似的话。\BibleRef{谷 10:2} 因此，那些前来问耶稣的人中，有些是发问来试探祂的。如果我们的救主如此卓越尚且被试探，那么祂门徒中被立为教导者的人，若因某些人发问——并非为求知，而是为试探——而感到烦恼，这又何必呢？你若收集起来，会发现许多段落记载法利塞人试探我们的耶稣，还有别的人，例如一位律法师，\BibleRef{玛 22:35} 或许还有一位经师，\BibleRef{谷 12:28} 这样，通过收集关于试探祂之人的记载，你探究时便可发现对此类提问有用的教训。只是，救主对那些试探祂的人，却制定了道理。因为他们说：\Quote{人无论因什么缘故，休妻可以吗？}祂回答说：\Quote{你们没有念过，那起初造人的，是造男造女？}\BibleRef{玛 19:4} 等等。我认为，法利塞人提出这个问题，是为了无论祂说什么，他们都能攻击祂；例如，若祂说\Quote{可以}，他们就会控告祂因小事而拆散婚姻；若祂说\Quote{不可以}，他们就会控告祂允许人与妻子同住，甚至带着罪。同样，在纳税的问题上，\BibleRef{玛 22:17} 若祂叫他们给，他们就会控告祂使百姓服从罗马人，而不服从天主的法律；若祂叫他们不给，他们就会控告祂挑起战争和叛乱，煽动那些无力对抗如此强大军队的人。但他们没有看出祂如何无可指责且明智地回答：首先，祂否定了可以因任何缘故休妻的观点；其次，祂回答了关于休书的问题。因为祂看到并非任何缘故都是解除婚姻的合理理由，并且丈夫必须与妻子同住，视她为较软弱的器皿，分给尊荣，\BibleRef{伯前 3:7} 并分担她罪过的重担。\BibleRef{迦 6:2} 祂引用\BibleBook{创世}{纪}的记载，质问那些以\Person{梅瑟}著作自夸的法利塞人说：\Quote{你们没有念过，那起初造人的，是造男造女？}等等，并因着\Quote{二人成为一体}这句话，祂接着教导与一体相合的道理，即\Quote{这样，他们不再是两个人，而是一体了。}\BibleRef{玛 19:4} 为了说服他们不应因任何缘故休妻，又说：\Quote{天主所结合的，人不可拆散。}\BibleRef{玛 19:6} 不过，在解释福音所引用\BibleBook{创世}{纪}的话时，要注意这些话并非如福音所写的那样是连续说的；我认为甚至不是关于同一批人，即那些按天主肖像受造的人和那些从地上的尘土及\Person{亚当}一条肋骨所造的人。因为那里说\Quote{祂造男造女}\BibleRef{创 1:27}，指的是按\Quote{肖像}受造的人；但那里又说\Quote{为此，人要离开自己的父母}\BibleRef{创 2:24}等等，则不是指按肖像受造的人；因为此后不久，上主天主用尘土造了人，并从他的肋骨造了助手。同时请注意，在按\Quote{肖像}受造的人那里，用的词不是\Quote{丈夫和妻子}，而是\Quote{男和女}。我们也在希伯来文中注意到这一点：男人用词\OriginalTerm{伊施}{is}表示，而男性用词\OriginalTerm{匝加尔}{zachar}表示；同样，女人用词\OriginalTerm{依莎}{essa}表示，而女性用词\OriginalTerm{纳克巴}{agkeba}表示。因为按\Quote{肖像}受造的人决不是\Quote{女人}或\Quote{男人}，而是更高级的性别——男性，和次级的性别——女性。而且，若有人离开母亲和父亲，他不是依附于女性，而是依附于自己的妻子，并且\Quote{成为一体}，因为男人和女人在肉体上是同一的。然后，救主描述那些由天主结合的人应当怎样，以便他们以合乎天主的方式结合，祂补充说：\Quote{他们不再是两个人；}\BibleRef{玛 19:6} 并且，凡夫妻之间确实有和谐、同心与和睦之处——当丈夫如同统治者，妻子服从\Quote{他要管辖你}\BibleRef{创 3:16}这句话时——对这样的人，我们真可以说：\Quote{他们不再是两个人。}既然为着\Quote{那与主联合的}，应当保留\Quote{成为一灵与祂联合}\BibleRef{格前 6:17}这一点，所以对于由天主结合的人，在\Quote{他们不再是两个人}之后，接着说：\Quote{而是一体。}是天主将二人合而为一，使他们不再是两个人，从女人嫁给男人之时起。既然天主结合了他们，因此在由天主结合的人身上，有一种\Quote{恩赐}。\Person{保禄}知道按照天主之道的婚姻是一种\Quote{恩赐}，正如神圣的独身也是一种恩赐，他说：\Quote{我愿众人像我一样，但每人都有自己从天主得来的恩赐：有人这样，有人那样。}\BibleRef{格前 7:7} 那些由天主结合的人既思念又遵守这诫命：\Quote{你们作丈夫的，要爱你们的妻子，如同基督爱了教会。}\BibleRef{弗 5:25} 于是救主命令：\Quote{天主所结合的，人不可拆散，}\BibleRef{玛 19:6} 但人却想拆散天主所结合的，当他\Quote{离弃正道，听从欺诈的邪神和魔鬼的教义，这教义是因说谎者的伪善，他们自己的良心像被烙铁烙过，禁止……}不仅禁止淫乱，而且\Quote{禁止嫁娶}\BibleRef{弟前 4:1}时，他甚至拆散了那些先前已由天主眷顾结合的人。这些话应参照救主在回答法利塞人时明确讲述的关于男女、男人和女人的话来讲。\par

% structure: p0034 source=h2 target=subsection reason=relative-heading-rank; base=h2

\subsection{十七、基督与教会的结合。}

但既然宗徒理解经文\BibleQuote{\Emph{二人成为一体}}\BibleRef{玛 19:5}是论及基督与教会\BibleRef{弗 5:31}，我们必须说：基督持守这句话\BibleQuote{天主所结合的，人不可拆散}\BibleRef{玛 19:6}，并没有休弃祂的前妻——也就是说，从前的会堂——除非是因为那妻子犯了奸淫，被那恶者变成了淫妇，并与恶者同谋背叛了她的丈夫，并杀死了祂，说\BibleQuote{从地上除掉这样的人，钉祂在十字架上，钉祂在十字架上}。所以是她自己离弃了丈夫，而不是丈夫休弃了她。因此，\BibleBook{依撒意亚}{书}责斥她离弃丈夫，说\BibleQuote{你们母亲的休书在哪里？我写了休书休了她吗？}\BibleRef{依 50:1}那在起初按肖像创造了\BibleQuote{具有天主的形象}的那一位的，造了祂为男，造了教会为女，使双方在肖像上合而为一。为了教会的缘故，主——丈夫——离开了祂原所见到的\BibleQuote{具有天主的形象}\BibleRef{斐 2:6}的圣父，也离开了祂的母亲，即那属于天上的耶路撒冷的儿子，而结合于那堕落于此的妻子，这两人在此成为一体。因为为了她，祂自己也成了血肉，当\BibleQuote{圣言成了血肉，居住在我们中间}\BibleRef{若 1:14}，他们不再是两个，而是已成为一体，因为对妻子说\BibleQuote{你们就是基督的身体，各自都是肢体}\BibleRef{格前 12:27}；因为基督的身体并非与教会——祂的身体——以及各自为肢体的成员分离的别的什么。天主结合了这不再为两个、已成为一体的，命令人不可将教会与主分离。那谨慎自守以免分离的人，确信自己不可能被分离，说\BibleQuote{谁能使我们与基督的爱隔绝？}\BibleRef{罗 8:35}因此，这里所说的\BibleQuote{天主所结合的，人不可拆散}\BibleRef{玛 19:6}虽是针对法利塞人而写，但对于那些超越法利塞人的人，却可以说\BibleQuote{所以，天主所结合的，没有什么能拆散}——无论是首领、权能者都不能；因为那结合的天主比任何人所能想象和称呼的这一切都更强大。\par

% structure: p0036 source=h2 target=subsection reason=relative-heading-rank; base=h2

\subsection{十八、休书。}

此后，我们将讨论法利塞人对耶稣所说的话\BibleQuote{\Emph{那么，为什么梅瑟命令给休书休她呢？}}\BibleRef{玛 19:7}。为此，我们有理提出\BibleBook{申命}{纪}中关于休书的段落，如下：\BibleQuote{如果一人娶了妻，与她同住，若她在丈夫眼中不蒙喜爱，因为他在她身上发现了什么不体面的事}等等，直到\BibleQuote{你们不可玷污上主你们的天主赐给你们为产业的土地}\BibleRef{申 24:1}。现在，我查问：在这些按此法律的事上，我们是否除了字面之外不应寻求任何东西，因为天主并未赐下它？或者，对引用\BibleQuote{梅瑟命令给休书休她}的法利塞人，有必要说：\BibleQuote{梅瑟为了你们的心硬，准许你们休妻，但起初并不是这样}\BibleRef{玛 19:8}？但如果有人上升到耶稣基督的福音——这福音教导法律是属神的——他也会寻求这法律属神的意义。那愿意用比喻解释这些事的人会说：正如保禄因所蒙的恩宠而自信地说：\BibleQuote{丈夫活着的时候，妻子是被束缚的；丈夫若死了，她就自由了，可以随意再嫁，只是要嫁在主里的人。然而按我的意见，如果她守节，就更有福；我想我也有天主的圣神}\BibleRef{格前 7:39-40}（因为在这里，为了不让\InnerQuote{按我的意见}这句话因缺乏圣神而被轻视，他很好地加上了\InnerQuote{我想我也有天主的圣神}），同样，梅瑟也可以因他立法的权能——他因百姓的心硬容忍了一些事，其中包括休妻——而凭自己的判断坚信他所颁布的法律也带有天主的圣神。他将会说：除非一部法律是属神的，另一部不是，但这部法律是属神的，它的属神意义应当被探究。\par

% structure: p0038 source=h2 target=subsection reason=relative-heading-rank; base=h2

\subsection{十九、以色列的休弃。}

现在，我们铭记以上关于依撒意亚论休书的那段话，我们要说：人民的母亲在没有得到休书的情况下，离开了她的丈夫基督；但后来在她身上发现了一件不体面的事，她在祂面前不再蒙恩宠，休书便为她写了出来。因为当新的盟约将外邦人召到那休弃了前妻者的家中时，它实际上是将休书给了那先前离开她丈夫——法律和圣言——的妻子。因此，祂（基督）也离开了她，可以说是另娶了一位，将休书交给了前妻；因此，他们不能再做法律所命令的事，因为有了休书。一个表明她已领受休书的迹象是：耶路撒冷被毁灭，连同他们称为圣所的地方（那里有被视为圣物之物），以及全燔祭台和与之相关的一切敬礼。并且休书的另一迹象是：他们不能守他们的庆节，即使法律明文命令他们在上主天主所指定的地方守节；还有此：整个会堂已不能以石头砸死犯了这样或那样的罪的人；此外，成千上万的诫命也是休书的迹象；还有事实：\Quote{再也没有先知了}，并且他们说：\Quote{我们不再看见异兆了}；因为主说：\Quote{祂从犹大和耶路撒冷除去了}，按照依撒意亚的话，\Quote{有能力的男子和有能力的女子，一个强壮的巨人}等等，直到\Quote{明理的听者}。\BibleRef{依 3:1} 现在，作为基督的那一位可能曾娶了会堂为妻并与之同居，但后来她在祂眼中不再蒙恩宠；她不再蒙恩宠的原因，是在她身上发现了一件不体面的事；因为有什么比以下情况更不体面呢？当提议在庆节释放一人时，他们要求释放强盗\Person{巴辣巴}，而定耶稣的罪？\BibleRef{玛 27:21} 还有什么比以下事实更不体面呢？他们全都论及祂说：\Quote{钉死祂，钉死祂}和\Quote{除掉这个人，从地上除去}？\BibleRef{若 19:15} 这句话能摆脱不体面之责吗？\Quote{祂的血归在我们和我们子孙的身上}？\BibleRef{玛 27:25} 因此，当祂被报复时，耶路撒冷被军队包围，它的荒凉近了，\BibleRef{路 21:20}，他们的房屋也被夺去，\BibleQuote{熙雍的女子被留下，像葡萄园中的茅屋，像黄瓜田中的草棚，像被围攻的城邑。}\BibleRef{依 1:8} 而且，大约在同一时间，我想，丈夫给前妻写了休书，交在她手中，从祂家中将她逐出；那从外邦人而来的妻子的债据已被取消，关于此，宗徒说：\BibleQuote{涂抹了那在规条上所写的、反对我们的债据，并将它除去，钉在十字架上；}\BibleRef{哥 2:14} 因为保禄和其他人也成了以色列的归化者，为了那从外邦人而来的她。因此，第一任妻子在她丈夫面前未蒙恩宠，因为在她身上发现了一件不体面的事，便从丈夫的居所出来，离开后，与另一个男人结合，她使自己服从于他；这丈夫我们应称之为强盗\Person{巴辣巴}（他象征魔鬼），或某种邪恶势力。并且在那会堂的一些人身上，发生了法律上所写的头一件事；而在另一些人身上，发生了第二件事。因为最后的丈夫，\BibleRef{申 24:3}，恨了他的妻子，并将在事物终局时为她写休书，当天主如此安排时，就将休书交在她手中，将她从居所逐出；因为正如好天主将在蛇与女人之间、他的后裔与她的后裔之间放置敌对\BibleRef{创 3:15}，祂也将安排最后的丈夫恨她。\par

% structure: p0040 source=h2 target=subsection reason=relative-heading-rank; base=h2

\subsection{二十、基督与外邦人。}

如今，有些人的情况是，那个男人与他们同住，却没有憎恨他们，因为他们住在最后一位丈夫的家里，这位丈夫娶了他们的会堂为妻。但在这类情况下，后夫也死了，\BibleRef{申 24:3} 也许是在基督最后的敌人——死亡被消灭的时候。但无论这些事中哪一件发生在妻子身上，经上说，前夫——就是休了她的那位——不能回转，在她被玷污之后，再娶她为妻，因为经上说：\BibleQuote{它是可憎恶的}，\BibleQuote{在上主你的天主面前}。\BibleRef{申 24:4} 但这些话似乎与这句不一致：\BibleQuote{如果外邦人的数目满了，以色列全家都要得救}。\BibleRef{罗 11:25} 不过请思量，是否可以这样说：如果她因前夫回转并娶她为妻而得救，那么她必然是在被玷污之后才得救。那么，司铎不可娶曾为妓女或被遗弃的女子为妻，\BibleRef{肋 21:14} 但其他地位低于司铎的人并不禁止这样做。但若你从外邦人的蒙召角度寻找妓女，你可以使用这段经文：\BibleQuote{你去娶一个淫妇为妻，生淫乱的子女}，\BibleRef{欧 1:2} 等等；因为，正如\BibleQuote{在圣殿里的司铎违犯了安息日，也不算为罪}，\BibleRef{玛 12:5} 同样，那休弃前妻、在适当季节娶了\BibleQuote{一个淫妇}的人，是依照那位的命令而行；那位在必要之时且只要必要，就说：\Quote{他不可娶妓女为妻}；而当合理之时，祂又说：\BibleQuote{你去娶一个淫妇为妻}。因为正如人子是安息日的主，\BibleRef{玛 12:8} 不像百姓那样是安息日的奴隶，同样，颁赐法律的那位有权柄赐下法律\BibleQuote{直到改革的时期}，\BibleRef{希 9:10} 并更改法律；当改革的时期临近时，祂也能在先前的方式和先前的心志之后，赐予另一种方式和另一种心志，\BibleQuote{在悦纳的时候，在救恩的日子}。\BibleRef{格后 6:2} 就让这些话按照我们对休书律法的解释来说吧。\par

% structure: p0042 source=h2 target=subsection reason=relative-heading-rank; base=h2

\subsection{二十一、天使与人类灵魂的结合}

但有人可能会询问，人的灵魂可否比喻为妻子，而那位管理她、作她主宰的天使——她与这位至高者交谈——可否称为她的丈夫？按照这种说法，每个人与配得上神圣天使守护的灵魂合法同居；但有时，在长久居留与交往之后，灵魂中可能产生一个原因，使她在她的主上、统治者的天使眼中失宠，因为在她身上发现了一件不体面的事；于是可能写出一份条约——正如通常所写的那样——并将休书写好，放在被逐者手中，这样当她被逐出他的住所后，就不再有与她先前守护者的亲密关系。即使那离开前住所的人嫁给了另一个丈夫，与他在一起也不幸，不仅像先前一样，因在她身上发现了不体面的事而在他眼中失宠，甚至还被他憎恨。是的，甚至可能从第二个丈夫那里写出一份休书，并由这最后的丈夫交到她手中，将她逐出他的住所。但关于天使与人的生命是否可能发生这样的变化，以至于就他们与我们的关系而言，相当于他们的死亡，有人可能会提出这个冒昧的问题；但无论如何，那一旦离开前夫的女子不会再回到他那里，因为休了她的前夫不能回转，在她被玷污之后，再娶她为妻。\BibleRef{申 24:4} 如果有人胆敢使用教会内流传、但并非所有人都承认为神圣的经文来缓和这一诫命，则可以引用\WorkTitle{牧者}中的一段话，论及某些人一信就服从\Person{弥额尔}，但因贪图享乐而背离他，然后服从享乐的天使，接着服从惩罚的天使，之后服从悔改的天使；因为你会注意到，那曾一度委身于享乐的妻子或灵魂不再回到第一个统治者那里，而且除了受惩罚外，她还服从一个比\Person{弥额尔}低下的天使；因为悔改的天使比他低微。因此，我们必须警惕，免得我们身上有任何不体面的事，以致在我们的丈夫基督或我们之上所派的天使眼中失宠。因为如果我们不警惕，或许我们也会收到休书，要么失去我们的守护者，要么归向另一个男人。但我认为，将天使与我们灵魂的婚姻想象成吉兆是不妥的。\par

% structure: p0044 source=h2 target=subsection reason=relative-heading-rank; base=h2

\subsection{二十二、教会显要人物的婚姻}

但是，在处理这段经文时，我想说，现在我们或许能够理解并清楚地阐述一个难以把握和洞悉的问题，即关于宗徒有关教会事务的立法；因为保禄不希望教会中那些在圣事管理上超越众人而获得某种尊荣的人尝试第二次婚姻。因为他在弟茂德前书中关于主教制定法律时说：\BibleQuote{如果有人渴望主教的职务，他向往的是善工。所以，主教必须无可指责，只作一个妻子的丈夫，有节制，能持重，} \BibleRef{弟前 3:1}等等；关于执事，他说：\BibleQuote{执事应当只作一个妻子的丈夫，善于管理自己的儿女和家庭，} \BibleRef{弟前 3:12}等等。是的，并且在选拔寡妇时，他说：\BibleQuote{寡妇要登记，年纪必须六十岁以下，且只作过一个丈夫的妻子；} \BibleRef{弟前 5:9}在这之后，他又说了些附加的话，其重要性次于此事。并且在弟铎书中，他说：\BibleQuote{我留你在克里特，是要你整顿那些尚未办妥的事，并照我所吩咐的，在各城设立长老。如果有人无可指责，只作一个妻子的丈夫，有信主的儿女} \BibleRef{铎 1:5}等等。现在，当我们看到有些结过两次婚的人可能比只结过一次婚的人好得多时，我们感到困惑：为什么保禄完全不允许那些结过两次婚的人被任命担任教会尊职？因为在我看来，这件事也值得探讨：一个人可能在两次婚姻中都不幸，在年轻时就失去了第二任妻子，此后在余生直至老年都生活在极大的自制和贞洁中。那么，谁会不自然感到困惑：当我们在寻找教会统治者时，为什么我们完全不任命这样一个虽然结过两次婚的人，因为关于婚姻的表述，却抓住那个只结过一次婚的人作为我们的统治者，即使他碰巧与妻子同住到老，有时可能并未在贞洁和节制方面受过锻炼？但是，从律法关于休书的话中，我反思：既然主教、司铎和执事是按这些名称实际存在之事的象征，他是否希望任命那些象征性地只结过一次婚的人，以便那能留心此事的人，能从属灵法律中找出不配担任教会统治的人，其灵魂未得丈夫欢心，因为在她身上发现了不体面的事，因而她已配得休书；因为这样一个灵魂，与第二个丈夫同居后，又被这样的丈夫憎恨，在第二次休书之后，不能再回到前夫那里。因此，很可能那些比我们更智慧、更有洞察力的人，也将在关于休书的律法中，或在禁止结过两次婚的人统治教会或被优先推举主持教会的宗徒著作中，找到其他论据。但是，在找到更好的、能以知识之光的极大辉煌将我们所说的话置于阴影中的东西之前，我们已经说出了在这些经文上临到我们的事。\par

% structure: p0046 source=h2 target=subsection reason=relative-heading-rank; base=h2

\subsection{二十三、一些因让步于人性软弱而赐下的法律。}

但是，即使我们似乎在经文上触及了超出我们能力范围的深奥之事，然而，由于字面表达，还必须说这些事：有些法律写下来，并非因其卓越，而是作为对领受法律者软弱的迁就；因为类似的事在以下话中表明：\BibleQuote{梅瑟因你们的心硬，准许你们休妻} \BibleRef{玛 19:8}；但那超越并高于为他们的心硬而写下的法律的，在以下话中表明：\BibleQuote{但从起初并不是这样}。但在新约中，也有一些法律禁令与\BibleQuote{梅瑟因你们的心硬，准许你们休妻}属于同一层次；例如，因我们的心硬，因我们的软弱而写下：\BibleQuote{但为免淫乱的事，男人当各有自己的妻子，女人也当各有自己的丈夫} \BibleRef{格前 7:2}；以及\BibleQuote{丈夫当还妻子所欠的债，妻子也当照样还丈夫所欠的债} \BibleRef{格前 7:3}。在这些话之后接着又说：\BibleQuote{我说这话是出于宽容，并不是出于命令} \BibleRef{格前 7:6}。此外，保禄因我们的心硬和软弱，对那些不愿意切慕更大的恩赐、成为更有福的人说了这话：\BibleQuote{妻子在丈夫活着的时候是被束缚的；但如果丈夫死了，她就可以随意嫁人，只要在主内} \BibleRef{格前 7:39} \BibleRef{格前 12:31}。但现在，有些教会统治者甚至违背了所写的，允许女人在丈夫还活着时嫁人，这违反了所写的：\BibleQuote{妻子在丈夫活着的时候是被束缚的}，以及\BibleQuote{所以，丈夫活着的时候，她若与别的男人结合，便被称为淫妇} \BibleRef{罗 7:3}。这并非完全没有理由，因为这种让步可能是相对于更坏的事而被允许的，违反起初由法律所规定和所写的。\par

% structure: p0048 source=h2 target=subsection reason=relative-heading-rank; base=h2

\subsection{二十四、犹太人对基督法律的批评。}

但也许有些敢于反对我们救主教诲的犹太人会说，当耶稣说：\BibleQuote{凡休自己妻子的，除非为了奸淫的缘故，就是使她犯奸淫，}\BibleRef{玛5:32}祂也像梅瑟一样允许休妻，而梅瑟被祂说成是因人民心硬而赐下法律；并且他们会认为\BibleQuote{因为他在她身上发现了不雅之事}\BibleRef{申24:1}这句话应等同于奸淫，因这缘故，妻子有充分理由可以被丈夫休弃。但必须对他说的是：如果犯奸淫的妇女按法律应被石头砸死，那么显然不雅之事不应作此理解。因为奸淫或如此严重的不雅行为，无需写休书交到妻子手中；但也许梅瑟把每一样罪都称为不雅之事，丈夫在妻子身上发现这罪，因妻子在他眼中不蒙喜悦，就写下休书，将妻子从丈夫家中打发走；\BibleQuote{但从起初并不是这样。}\BibleRef{玛19:8}此后，我们的救主说，绝不允许因任何其他罪而解除婚姻，除了唯独妻子的奸淫：\BibleQuote{凡休自己妻子的，除非为了奸淫的缘故，就是使她犯奸淫。}\BibleRef{玛5:32}但值得探讨的是，是否因此祂禁止任何人因任何其他原因休妻，除非她犯奸淫，例如下毒，或趁丈夫不在家时杀害他们所生的婴儿，或任何一种谋杀。再者，如果妻子被发现掠夺和洗劫丈夫的家宅，尽管她没有犯奸淫，有人可能会问，他是否有理由休弃这样的妻子，因为救主禁止任何人休妻，除非为了奸淫的缘故。这两种情况似乎都有些骇人听闻——是否真的骇人听闻，我不知道；因为忍受如此严重、似乎比奸淫更坏的罪，看来是不合理的；但另一方面，违背救主教导的意图，每个人都会承认是不虔敬的。因此我奇怪为什么祂没有说：除了奸淫的缘故，谁也不要休妻；却说：\BibleQuote{凡休自己妻子的，除非为了奸淫的缘故，就是使她犯奸淫。}\BibleRef{玛5:32}因为显然，休弃并非犯奸淫的妻子的人，就他而言，是使她成了奸妇；因为\BibleQuote{丈夫活着的时候，她若归于别的男人，便称为奸妇。}\BibleRef{罗7:3}当他休妻时，给了她再婚的借口，很明显的，他这样做就是使她成了奸妇。但是否她在下毒或杀人的现场被抓住，能为他的休妻提供任何辩护，你们可以自行探究；因为丈夫除了休妻外，还可以其他方式使自己的妻子犯奸淫；例如，允许她做超出本分的事，纵容她与她所意愿的男人交往，因为常常由于丈夫的单纯，妻子会犯这样的过失；但这样的丈夫在这些过失中是否有辩护的理由，你们要仔细探究，并对我们就这段经文提出的难题也给出你们的意见。甚至那克制自己不与妻子同房的人，当他不满足妻子的欲望时，也常常使她成为奸妇，即使他表面上是为了更大的庄重和节制。也许这个人更有罪，因为他不满足妻子的欲望，就他而言，是使她成为奸妇，比那因奸淫以外的原因——如下毒、谋杀或任何最严重的罪——而休妻的人更有罪。但是，正如一个女人即使似乎嫁给了另一个男人，只要前夫还活着，她就是奸妇；同样，那个似乎娶了被休之妻的男人，按照我们救主的宣告，与其说是娶她，不如说是与她犯奸淫。\par

% structure: p0050 source=h2 target=subsection reason=relative-heading-rank; base=h2

\subsection{二十五、贞洁与祈祷。}

此后，门徒们思量婚姻中可能发生多少意外，丈夫必须忍受这些，并因此遭受极大的艰辛；若他不忍受，就会违犯基督的话。于是他们对祂说，视独身生活为较婚姻更容易、更有益的途径，尽管后者看似有益：\BibleQuote{\Emph{若丈夫与妻子的情形是这样，倒不如不娶。}}\BibleRef{玛 19:10}对此，救主教导我们说，绝对的贞洁是天主所赐的恩赐，不仅是修炼的成果，而是藉着祈祷由天主赐予：\BibleQuote{这话不是人人都能领悟的，只有那些得了恩赐的人才能领悟。}\BibleRef{玛 19:11}然而，有些人诡辩地攻击\BibleQuote{得了恩赐的人}这句话，仿佛那些愿意在独身中保持纯洁、却被情欲所胜的人有了借口。我们必须说：如果我们相信圣经，为何我们只抓住\BibleQuote{只有那些得了恩赐的人}这句话，而不留意\BibleQuote{你们求，必要给你们}\BibleRef{玛 7:7}，以及随后所说的\BibleQuote{因为凡求的，就必得到}\BibleRef{玛 7:8}？因为如果他们——\BibleQuote{得了恩赐的人}——能够接受这关于绝对纯洁的话，那么愿意的人就当祈求，顺服并相信那说\BibleQuote{你们求，必要给你们}\BibleRef{玛 7:7}的主，且不怀疑\BibleQuote{凡求的，就必得到}\BibleRef{玛 7:8}这句话。但你要查问谁是那祈求的，因为凡未得着的人，没有人曾祈求，即使他看似祈求了，因为不能说\BibleQuote{凡求的，就必得到}\BibleRef{玛 7:8}是谎言。那么，谁是祈求的人呢？岂不是那听从耶稣说\BibleQuote{你们祈祷时，要相信必得着，就必给你们成就}\BibleRef{谷 11:24}的人？但祈求的人必须尽力而为，好能\BibleQuote{用神魂}祈祷，也\BibleQuote{用理智}\BibleRef{格前 14:15}祈祷，并\BibleQuote{不断祈祷}\BibleRef{得前 5:17}，也当记住这话：\BibleQuote{耶稣给他们设了一个比喻，论及人应当时常祈祷，不要灰心，说：某城中有一个判官}\BibleRef{路 18:1}等等。知道什么是求，什么是得，以及\BibleQuote{凡求的，就必得到}\BibleRef{玛 7:8}和\BibleQuote{我告诉你们：他纵然不为了他是朋友而起来给他，但因他恬不知耻地求，就必起来照他所需用的给他}\BibleRef{路 11:8}的意思，是有益的。因此又加上：\BibleQuote{所以我告诉你们：你们求，必要给你们}等等。更进一步，让这话\BibleQuote{这话不是人人都能领悟的，只有那些得了恩赐的人才能领悟}\BibleRef{玛 19:11}激励我们相称地祈求领受；以及\BibleQuote{你们中间哪有为父亲的，儿子向他求鱼，反将蛇当鱼给他呢？}\BibleRef{路 11:11}等等。因此，天主必将美好的恩赐，即独身和贞洁中的完全纯洁，赐给那些全心、以信心、在不间断的祈祷中祈求祂的人。\par

\SourceRecord{Translated by John Patrick.；Ante-Nicene Fathers；Vol. 9.；Edited by Allan Menzies.；Buffalo, NY: Christian Literature Publishing Co.,；1896.；<http://www.newadvent.org/fathers/101614.htm>.}
